% 2026-11-21 / Q-v1. Independent quantization paper; compile from repo root.
% !TeX root = cubic-quantization.tex
\RequirePackage{fix-cm}
\documentclass[a4paper,11pt]{ltjsarticle}
\usepackage[margin=25mm,headheight=16pt]{geometry}
\usepackage{amsmath,amssymb,amsthm,mathtools,bm,mathrsfs}
\DeclareMathSizes{9}{9}{6}{5}
\DeclareMathSizes{9.85501}{9}{6}{5}
\usepackage{booktabs,longtable,array,enumitem}
\usepackage{xcolor}
\usepackage[unicode,colorlinks=true,linkcolor=blue!40!black,citecolor=blue!40!black,urlcolor=blue!45!black]{hyperref}
\usepackage{fancyhdr}
\usepackage{microtype}
\setlength{\parindent}{1em}
\setlength{\parskip}{0.25em}
\setlist{itemsep=0.2em,topsep=0.4em}
\allowdisplaybreaks[2]
\pagestyle{fancy}
\fancyhf{}
\fancyhead[L]{\small 三次超曲面の代数的量子化}
\fancyhead[R]{\small 2026-11-21 / Q-v1}
\fancyfoot[C]{\thepage}
\newcommand{\CC}{\mathbb C}
\newcommand{\ZZ}{\mathbb Z}
\newcommand{\PP}{\mathbb P}
\newcommand{\OO}{\mathcal O}
\newcommand{\cK}{\mathcal K}
\newcommand{\cP}{\mathcal P}
\newcommand{\cD}{\mathcal D}
\newcommand{\cC}{\mathcal C}
\newcommand{\cE}{\mathcal E}
\newcommand{\cF}{\mathcal F}
\newcommand{\cM}{\mathcal M}
\newcommand{\fm}{\mathfrak m}
\newcommand{\Perf}{\operatorname{Perf}}
\newcommand{\Ku}{\operatorname{Ku}}
\newcommand{\MF}{\operatorname{MF}}
\newcommand{\HH}{\operatorname{HH}}
\newcommand{\QGr}{\operatorname{QGr}}
\newcommand{\qgr}{\operatorname{qgr}}
\newcommand{\Dsg}{D_{\mathrm{sg}}}
\newcommand{\RHom}{\operatorname{RHom}}
\newcommand{\REnd}{\operatorname{REnd}}
\newcommand{\Hom}{\operatorname{Hom}}
\newcommand{\Ext}{\operatorname{Ext}}
\newcommand{\Ind}{\operatorname{Ind}}
\newcommand{\Def}{\operatorname{Def}}
\newcommand{\thick}{\operatorname{thick}}
\newcommand{\Spec}{\operatorname{Spec}}
\newcommand{\Spf}{\operatorname{Spf}}
\newcommand{\Id}{\operatorname{Id}}
\newcommand{\sym}{\operatorname{sym}}
\newcommand{\Cone}{\operatorname{Cone}}
\newcommand{\rank}{\operatorname{rank}}
\newcommand{\Pic}{\operatorname{Pic}}
\newcommand{\otL}{\mathbin{\otimes^{\mathbf L}}}
\newcommand{\DocDate}{2026年11月21日}
\newcommand{\DocVersion}{量子化論文 v1}
\theoremstyle{plain}
\newtheorem{theorem}{定理}[section]
\newtheorem{proposition}[theorem]{命題}
\newtheorem{lemma}[theorem]{補題}
\newtheorem{corollary}[theorem]{系}
\theoremstyle{definition}
\newtheorem{definition}[theorem]{定義}
\newtheorem{problem}[theorem]{問題}
\newtheorem{example}[theorem]{例}
\theoremstyle{remark}
\newtheorem{remark}[theorem]{注意}
\renewcommand{\proofname}{証明}
\hypersetup{pdfauthor={},pdfsubject={三次超曲面の代数的量子化・中心元・同変降下}}

\numberwithin{equation}{section}
\hypersetup{pdftitle={滑らかな三次超曲面の代数的量子化}}
\title{滑らかな三次超曲面の代数的量子化\\
\large 中心三次元・多パラメータ族・同変降下}
\author{}
\date{2026年11月21日・量子化論文（公表準備稿 v1）}
\fancyhead[L]{\small 三次超曲面の代数的量子化}
\fancyhead[R]{\small 2026-11-21 / Q-v1}
\begin{document}
\maketitle
\begin{abstract}
複素数体上の滑らかな射影三次超曲面
$X\subset\PP^{n+1}$、$n\ge3$の正則Poisson構造の代数的量子化を研究する。
Poisson正規形の構造定理を入力に、三変数Jacobianブロックごとの
中心元を構成し、全次元・全ブロックへ組み立てる。
各Poisson構造を、原点を含む一変数Zariski開集合上の
Koszul Artin--Schelter正則代数と、その正則中心三次商に量子化する。
中心元は、滑らかな三変数制限に対する普遍Hesse公式と、
古典Casimir核および普遍交換子行列の階数を用いて代数的に構成する。
この構成は三変数制限が特異または零の場合も含み、
固定した変数分解の全滑らかな係数空間上に中心元を選べる。
階数二の場合には、supportの三平面を全て動かし、三変数制限の型によらず、
全座標群に共変な中心元を、普遍Hessianの二次補正まで指定して零切断近傍に選ぶ。
証明にはGrassmann多様体上の純中心束のコホモロジー消滅と、
有限中心核のNoetherian完備化を用いる。
従って$\gcd(n+2,3)=1$なら、被約な階数二のPoisson対の全モジュライと
変形直線の積の零切断近傍へ量子化族を降ろせる。
特に三次三様体では、全ての非零Poisson対を同時に扱う。
一般次元でも三次Veroneseは全階数二層へ降下する。
全ての非零階数の被約層で、中心スカラーgerbeの位数は$\gcd(n+2,3)$であり、
これが次数一を保つ線束補正による降下の障害を測る。
中心変数が残らない満support層でも、ブロックの順序を底に加えずに、
指定二次jetを持つ共変な中心元を零切断近傍に選べる。
この場合は係数空間のアフィン性により有限jetを直接持ち上げ、
三次Veroneseを被約Poissonモジュライへ降ろす。
特に三次四様体の階数四の全Poisson対を同時に扱う。
制限が恒等的零の層では、PBW対称化した三次式自体が
全パラメータで全ブロック群に共変な正則中心元となる。
ただし$n\ge4$では、この標準中心直線を全support上の選択と一致させることは
二次jetの相対コホモロジー障害によって妨げられる。
節点・尖点層では特異点直線によって次数重みを除き、
中心変数の個数によらず、Poisson対の各層と変形直線の積の零切断近傍へ量子化族を降ろす。
また、固定した正規形のPoisson線形族を同時に量子化する。

付録では、既知のSklyanin中心公式との対応、および偏極付き量子超曲面からの
二次ambientと中心直線の復元を扱う。
\end{abstract}
\tableofcontents

\section{序論}\label{sec:cq-introduction}
本稿の目的は、Poisson括弧を所定の古典極限に持つ次数付き代数族を
有限な代数的構成によって与え、その中心元の規格化、座標変換への共変性、
Poissonモジュライへの降下を区別して制御することである。
$X=V(F)\subset\PP(V)$、$\dim V=N=n+2\ge5$とし、基礎体は$\CC$とする。
古典的な正規形・退化軌跡・Poissonモジュライは\cite{PoissonCubic}、
非可換三次超曲面の可換復元・構造対象付きモジュライ・Kuznetsov成分との比較は
\cite{NCCubicMain}で扱う。
本稿はその間の代数的量子化を独立に扱い、一般の複数ブロックの結果を
良い一ブロック成分の仮定へ狭めない。
周囲の二次代数を$A$、中心商を$B=A/(\Omega)$と書く。

\subsection*{代数的量子化と中心元}
定理\ref{thm:pg-quantization}は、固定した任意のPoisson構造に対し、
$0$を含むZariski開集合$U\subset\mathbb A^1_\hbar$上に
二次代数$A_U$と正則中心三次元$\Omega_U$を構成する。
各次数成分は底上で平坦であり、
$A_U$の閉ファイバーは$N$生成元、Hilbert級数$(1-t)^{-N}$、
大域次元$N$のKoszul AS正則代数となる。
中心商$B_U=A_U/(\Omega_U)$は$\operatorname{Sym}(V^*)/(F)$の平坦量子化で、
各閉ファイバーはAS Gorenstein、Hilbert級数$(1-t^3)/(1-t)^N$を持つ。
関係式は、各ブロックについて巡回的に
\begin{equation}\label{eq:intro-quantum-relations}
 [x_{s2},x_{s3}]
 =\hbar\,\operatorname{sym}(\partial_{x_{s1}}f_s)
\end{equation}
であり、異なるブロックと中心変数の間の交換子は零である。

中心三次元の構成は、以下の有限次元の議論による。
まず滑らかな三変数制限を持つブロックについて、
Hesse座標で全混合係数を許す厳密な多項式中心公式を与える
（定理\ref{thm:pg-fermat-mixed-universal}、
\ref{thm:pg-universal-hesse-mixed}）。
次に、\emph{全体の}三次式$f$が滑らかであれば、
三変数制限の特異性によらず古典Casimir核は
\begin{equation}\label{eq:intro-classical-kernel-all}
 \ker\left(S_3\xrightarrow{p\mapsto(\{p,x_i\}_f)_{i=1}^3}
 (S_4)^3\right)=\mathbb C f\oplus\mathbb C[u]_3
\end{equation}
となる。ここで$S=\mathbb C[x_1,x_2,x_3,u]$である。
滑らかな三変数制限を持つ係数の稠密開集合での中心公式から、
普遍量子交換子行列の所定の小行列式が恒等的に消える。
これと古典核の一定次元性を合わせ、Cramer公式で
$\hbar=0$において正則な有理式の中心元を得る
（定理\ref{thm:pg-universal-determinant}）。
中心項を配分して各ブロック全体を滑らかにできることは
補題\ref{lem:pg-central-splitting-smooth}で示す。
このため量子化定理には形式的中心元の存在を仮定する必要がない。

固定した変数分解の全滑らかな係数空間についても、
零切断全体を含む一つの開集合上に偶な中心元を選べる
（定理\ref{thm:pg-det-global-coefficient-lift}）。
ここでは$\hbar$で割った交換子の核を定階数のベクトル束とし、
純中心部分による商線束をfactorialな係数環上で自明化する。
異なる規格化の間の遷移は一般に
$\Omega'=a\Omega+c$（$a$は単元、$c$は純中心三次式）である。
同じ古典極限を持つ中心元でも、自然な周囲代数の同一視で
同じ中心イデアルを定めるとは限らない
（例\ref{ex:pg-det-distinct-normalized-ideals}）。
従って、係数空間上の存在とモジュライへの同変な降下は区別する。

同変な選択については、まず
$G_0=(\mathrm{GL}_3\times\mathrm{GL}_m)\times\mathbb G_m$
という変数分解を保つ群に対し、零切断全体を含む不変開上で
中心元を選べることを証明する
（定理\ref{thm:pg-equivariant-levi-center}）。
証明には既知のReynolds平均を用い、
中心元だけでなく、その定義域と古典極限も同変に保つ。
さらに$g=f(x,0)$が滑らかな部分、および節点・尖点の各軌道上では、
$g$の方向微分が張る二次式の部分束を同変に分裂し、
混合二次項を用いてactive座標の中心を選ぶ。
この操作により$X=Mx+Cu,\ U=Nu$という
中心変数による平行移動を含む群に対しても同変な中心元を得る
（定理\ref{thm:pg-equivariant-parabolic-center}）。
特異な二軌道では、基準形の安定化群と同変な左逆を明示し、
軌道上へ正則に移す（補題\ref{lem:pg-equivariant-singular-retractions}）。
Poisson係数と変形変数を分けると、変形変数を固定したまま
対応する商スタック上に中心商の族が降りる
（系\ref{cor:pg-equivariant-central-quotient}）。
また$g=0$ならJacobian括弧は中心係数環上でactive変数についてaffine-linearとなり、
PBW対称化の随伴同変性から
\begin{equation}\label{eq:intro-zero-restriction-center}
 \Omega_{f,t}=\sym(f),\qquad
 [x_i,\sym(p)]=t\,\sym(\{x_i,p\}_f)
\end{equation}
を得る。この場合は全$t\in\mathbb A^1$でPBW性と全ブロック群への共変性が成り立ち、
例外パラメータの除去も純中心係数の変更も不要である
（定理\ref{thm:pg-equivariant-zero-center}）。
これらの層ごとの構成とは別に、階数二の全supportにわたる共変な選択を
定理\ref{thm:pg-intrinsic-rank-two-center}で証明する。
その証明は中心調整を経由せず、可約・非被約・零の三変数制限も含む。
$g$が滑らかな場合、全ブロック群$G$による商は、
階数二のPoisson対の被約層内の対応する開部分上の
$\mathbb G_m$-gerbeである
（命題\ref{prop:pg-equivariant-poisson-gerbe}）。
$\gcd(m,3)=1$なら指標を用いて次数線形化を補正でき、
次数付き代数の族がこの被約層へ降りる
（系\ref{cor:pg-equivariant-character-descent}）。
特に三次三様体では$m=2$なので、十七次元の被約Poissonモジュライの稠密開上に
この族を得る。
ここで十七は底の次元であり、固定した非零変形パラメータでの
非可換代数の同型類の像の次元を主張するものではない。
零制限の層はsupport平面が$X$に含まれる条件であり、$m\ge3$のときに限り非空となる。
この層でも同じgerbeと指標補正を用い、$\gcd(m,3)=1$なら
全変形直線上の代数族をPoissonモジュライへ降ろせる
（命題\ref{prop:pg-equivariant-zero-stratum}）。

節点・尖点の各層は階数二層内でそれぞれ余次元$1,2$である。
ここでは一意な特異点の直線$\ell_g\subset W$が正則に動く。
引戻しの線形化で$K_{\mathrm{sc}}$が$\ell_g$に重み$-1$で作用するため、
\begin{equation}\label{eq:intro-singular-point-descent}
 \mathcal A_d^{\mathrm{sp}}=\mathcal A_d\otimes\ell_g^{\otimes d},
 \qquad
 \mathcal B_d^{\mathrm{sp}}=\mathcal B_d\otimes\ell_g^{\otimes d}
\end{equation}
は重み零になる。
従ってこの二層では$\gcd(m,3)=1$という条件を要せず、
各層と変形直線の積の零切断近傍へ実際の次数付き代数族が降りる
（命題\ref{prop:pg-equivariant-singular-strata}）。
三次三様体ではその底の次元はそれぞれ$16,15$である。
これらの明示的な層ごとの選択が互いに一致するとは限らない。
型の境界をまたぐ一つの族は、次の内在的な構成によって選ぶ。

$C_3^\times(V)\subset\bigwedge^3V$を非零分解可能三重ベクトルの錐とし、
$\mathcal T_N$を全三次式が滑らかな係数空間とする。
$\mathcal S_N=\mathcal T_N\times C_3^\times(V)$では
support $W_P$が三平面の全体を動く。
古典Casimir核は自然に
\begin{equation}\label{eq:intro-intrinsic-classical-kernel}
 \mathcal K_0=\OO_{\mathcal S_N}F\oplus
 \operatorname{Sym}^3(V/W_P)^*
\end{equation}
と分裂する。Grassmann多様体上のBorel--Weil--Bott計算により
$H^1(\mathcal S_N,\mathcal K_0)=0$であるため、
普遍二次Hessian補正を初期値として$\hbar$進の中心核へ順次持ち上げられる。
Stiefel開上の正則関数の延長によって係数を有限生成の不変式環に置き、
有限加群の完備化を$\hbar^3$で還元し、指定二次jetを持つ\emph{代数的な}中心元を得る。
最後に還元的群$\mathrm{GL}(V)\times\mathbb G_m$で平均する。
この方法で全$\mathcal S_N\times\{0\}$を含む不変開上の共変な中心元と、
次数別局所自由な正則中心商族を構成する
（定理\ref{thm:pg-intrinsic-rank-two-center}）。
形式的級数の全体を代数的な式と同一視するのではなく、
中心性と指定された有限jetを持つ有限核の元を選ぶ点が要となる。

さらに$\gcd(N,3)=1$なら、$Nr-3s=1$を満たす整数を使った
指標$\theta(L,\lambda)=(\det L)^r\lambda^s$で次数重みを補正し、
被約な階数二のPoisson対の\emph{全}モジュライと変形直線の積の
零切断近傍へ代数族を降ろす
（系\ref{cor:pg-intrinsic-rank-two-descent}）。
三次三様体では$N=5$なので、十七次元の非零Poisson対の全被約モジュライが底となる。
これは変形原点の近傍における存在であり、全変形直線への延長や一意性は主張しない。
一般の$N$でも、次数$3d$部分を$\lambda^d$で捩る三次Veronese
$\mathcal A^{[3]}_d=\mathcal A_{3d}\otimes\lambda^d$とその中心商は、
全階数二層の零切断近傍へ降下する
（系\ref{cor:pg-intrinsic-third-veronese-descent}）。
その古典偏極は各ファイバー上で$\OO(3)$である。
一方、$\operatorname{Pic}(\mathcal S_N)=0$と同変可逆層の重みを計算すると、
中心スカラーgerbeの位数は正確に$\gcd(N,3)$となる
（命題\ref{prop:pg-intrinsic-gerbe-period}）。
従って$3\mid N$では、指標に限らず、可逆層の次数冪による補正によって
元の次数一の族を降ろすことはできない。

零制限層の標準中心直線との一致には、これとは別の障害がある。
$N\ge6$では、この層上で$\mathscr H(P,F)$は非零の純中心切断になり得るが、
全support上の純中心束の$P$重み二の大域切断はBott計算によって零である。
従って$\Omega^{\mathrm z}=\sym(F)$への全層での一致は、
任意の大域的な偶中心元について、同変性を要求しなくても二次jetで不可能となる
（命題\ref{prop:pg-intrinsic-zero-matching-obstruction}）。
これは所定の周囲代数における中心直線の比較であり、中心商の抽象的な非同型は意味しない。
全係数での補空間の同変選択は実際には不可能である
（命題\ref{prop:pg-equivariant-global-centering-obstruction}）が、
この障害は中心核そのものの持上げを妨げない。

固定した正規形では、独立なパラメータ$\eta_1,\ldots,\eta_k$を導入し、
各ブロックの中心元を加えることで
\begin{equation}\label{eq:intro-independent-center}
 \Omega_{\boldsymbol\eta}=\sum_{s=1}^k\Omega_{f_s,\eta_s}
\end{equation}
を得る。ブロック間の追加補正は不要である。
$\eta_s=\lambda_s\hbar$により、$\sum_s\lambda_s\pi_s$という
Poisson線形族の全方向を一つの相対代数族で量子化できる
（定理\ref{thm:pg-all-dimensions}、\ref{thm:pg-ad-all-poisson}）。
滑らかな三変数制限では高々$\eta_s^4$の多項式公式を用い、
零制限では上のPBW対称化を、非零の特異な制限を含む一般の場合には
原点で正則な有理式を用いる。
また、定数三重ベクトルとHessianの縮約
\begin{equation}\label{eq:intro-intrinsic-hessian}
 \mathscr H(P,F)=\frac16P^{ijk}P^{abc}F_{ia}F_{jb}F_{kc}
\end{equation}
により、中心元の二次補正を内在的に表す
（命題\ref{prop:pg-universal-casimir-jet}、
\ref{prop:pg-ad-intrinsic-hessian}）。
二次jetの普遍恒等式とHesse表示での全次数中心公式は、異なる主張である。

中心変数が残らない$N=3k$の満support層では、固定した正規形を離れても
一つの共変な中心元を選べる
（定理\ref{thm:pg-intrinsic-full-support-center}）。
まず、三重ベクトルから復元される順序のない$k$個のブロックを
有限\'etale被覆として相対化する
（補題\ref{lem:pg-intrinsic-block-cover}）。
満supportの三重ベクトルの軌道は
$\mathrm{GL}_N/(\mathrm{SL}_3^k\rtimes\mathfrak S_k)$であり、
既知のアフィン商の定理によってアフィンである。
適合する滑らかな三次式の係数空間も、ベクトル束の判別式主開としてアフィンになる。
この底上で割った交換子の核を階数$k$の束とし、
$\sym(F-\hbar^2\mathscr H(P,F)/12)$を法$\hbar^3$の切断として直接持ち上げる。
還元群と偶性による平均はこのjetを保ち、
全零切断を含む不変開上に代数的な中心元を与える。
ブロックごとの次数によるfiltrationから正則性を示すので、
テンソル積代数全体のdomain性を仮定しない。
得られる三次Veroneseと中心商は、実際の被約Poisson層へ降下する
（系\ref{cor:pg-full-support-veronese-descent}）。
満supportの底の次元は$2k$であり、三次四様体の階数四では四次元となる。
これはPoisson対の底の次元であり、固定した非零$\hbar$での
非可換同型類の像の次元ではない。

また、中心スカラーgerbeの位数$\gcd(N,3)$は階数二に限らない。
命題\ref{prop:pg-intrinsic-all-rank-gerbe}では、各非零階数$2k$の全被約層で
これを証明する。$3\mid N$の場合、Fermat点の安定化群内の二つの元の交換子が
中心の三乗根となることから、全ての同変可逆層の中心重みが三の倍数だと分かる。
この証明は底のPicard群の消滅を必要としない。
従って満support層の三次Veronese降下は、
元の次数一の族の可逆層による補正とは区別される。

この量子化は各構造、固定した正規形の線形族、および固定した変数分解の
係数空間上での構成を含み、上記の範囲では座標変換との同変性も備える。
次数付き代数を載せる商スタックでは、次数$d$に重み$d$で作用する
スカラーの$\mathbb G_m$をまず保持し、指標補正・特異点直線による捻り、または三次Veroneseによって
上記の範囲においてこれを除く。
中心変数を残す一般の複数ブロック層での同時選択や、
全閉点での$\operatorname{qgr}B_\hbar$のsmoothnessは結論に含めない。

\subsection*{ambientの復元}
三次以上の一つの斉次関係で二次代数を割っても、次数一・二の積は変わらない。
この事実から量子超曲面の二次包絡としてambientを復元し、
商写像の三次核から中心方程式の直線を得る
（定理\ref{thm:pg-ambient-recovery-general}）。
左右noetherianなAS Gorenstein性と所定のinjective dimensionの下では、
偏極付き三つ組からsection algebraを復元して同じ構成を適用できる
（定理\ref{thm:pg-ambient-recovery-polarized}）。
三次三様体については、Serre双対性・二変数中心・中心上の自由性などを仮定して、
構造対象付き導来圏から偶数Veroneseを経由する復元も示す
（系\ref{cor:pg-ambient-derived-threefold}）。
一般次元でSerre関手から得られるVeroneseが元の代数を決定するかは別の問題である。


\subsection*{先行研究との関係と各定理の位置づけ}
以下では証明に用いる既知の入力と、本稿で証明する追加部分を区別する。
古典的なPoisson正規形は\cite{PoissonCubic}からの入力であり、
本稿では第\ref{sec:cq-classical-input}節にその仮定と必要な結論を明記する。

\paragraph{一ブロックの量子化：既知の存在定理と相対的規格化。}
三生成元同次代数のAS正則性・3-Calabi--Yau性の判定には
Mori--Smith\cite[Theorems~1.4, 4.2]{PMoriSmith}を用いる。
このうちTheorem~1.4はArtin--Tate--Van den Berghの判定の再掲である。
低次potentialによるPBW変形はBerger--Taillefer
\cite[Theorem~1.1(i), p.~242]{PBergerTaillefer}の結果であり、
命題\ref{prop:pg-PBW}では中心係数環とactive次数の濾過に即して適用する。

さらに、三変数Calabi--Yau最高次potentialに低次項を加えた代数に
非スカラー中心元が存在することは、Etingof--Ginzburg
\cite[Theorem~3.3.2(ii), p.~1385]{EG}が既に証明している。
その証明は同論文\S8.3, pp.~1410--1411で、形式的量子化、Rees構成と特殊化を用いる。
またProposition~8.1.1(i), pp.~1407--1408は、正の重み$(a,b,c)$、$d=a+b+c$と
$\deg\Phi_j\le d$を満たす$\Phi=xyz-yxz+\sum_{j\ge1}\hbar^j\Phi_j$に対し、
古典剰余$(\Phi_1)^{\mathrm{ab}}$を持つ次数$d$以下の形式中心元を与える。
この項には孤立特異点の仮定はなく、その仮定は同命題(ii), (iii)で追加される。
証明には相対係数環と還元群の作用に対応した形式性も既に現れる。
従って、指定した古典剰余の形式中心の存在や、相対的・同変な形式構成自体も
本稿の新規性の根拠とはしない。
一ブロックの中心元存在や一般係数から特殊係数への延長の原理自体を、本稿独自の結果とはしない。
本稿の定理\ref{thm:pg-universal-determinant}と
定理\ref{thm:pg-det-global-coefficient-lift}が与えるのは、
\emph{指定した古典Casimir $f$}を極限に持つ中心元の、
独立な中心混合係数を含む係数環上の選択と、$\hbar=0$における正則性である。
その証明では、古典核の一定次元性と普遍交換子行列の小行列式を用い、
三変数制限が特異・零の場合も同じ有限線形代数で扱う。
中心元の存在、左右正則性、および中心商の平坦性は別々に証明する。
零制限については、定理\ref{thm:pg-equivariant-zero-center}で
PBW対称化が随伴作用と可換する既知の原理
\cite[Corollaries~13.7--13.8, p.~75]{PEtingofLie}を用いて、
全パラメータの中心元を直接構成する。
この原理自体を新規とせず、中心混合係数を動かす全ブロック群への共変性と
support平面を含む層上の族への適用を明示する。

\paragraph{明示中心公式と全ブロックへの組立て。}
純Sklyanin中心元の基本公式は既知である
\cite{SmithTate,DeLaetLeBruyn}。
Etingof--Ginzburg\cite[\S9, p.~1414]{EG}には、
一変数低次項を許す中心元の明示公式も収められている。
一般Hesse表示の混合二次項は、Hesse係数を可逆にすると
中心平行移動でこの表示へ移せる
（注意\ref{rem:pg-hesse-eg-translation}）。
従って、混合係数を含むという事実だけを公式の新規性の根拠とはしない。
定理\ref{thm:pg-universal-hesse-mixed}では固定した係数環のまま、
この平行移動による極を導入しない規格化と普遍恒等式を与える。
有限の次数四線形代数が証明する中心性と、
全次数のPBW性を明確に区別する。

節点型についても、Chekhov--Mazzocco--Rubtsov
\cite[Lemma~4.6, 式(4.39), p.~22]{PCMR}が、
一般化Etingof--Ginzburg代数の$\gamma=0$の場合の中心三次元を明示している。
同じaccepted versionのTable~5, p.~41には、$xyz+x^3+y^3$に
対角二次項・一般一次項を加えた古典式との対応がある。
注意\ref{rem:pg-cmr-nodal-comparison}では、一般の二次混合項をこの表示へ移す
安定化群同変な左逆、対称Jacobian関係とのパラメータ対応、
古典極限を$f$にする規格化と純中心補正を明示する。
節点型の中心元の存在と基準座標での公式を本稿の新規成果とはしない。

全次元の構成は、この一ブロックの入力を正規形定理と組み合わせるものである。
定理\ref{thm:pg-quantization}では中心項の配分、ブロックの組立て、
最高active記号による中心正則性を証明し、
定理\ref{thm:pg-all-dimensions}と\ref{thm:pg-ad-all-poisson}では
独立な多パラメータを保持する。
同変性については定理\ref{thm:pg-equivariant-levi-center}、
\ref{thm:pg-equivariant-parabolic-center}、\ref{thm:pg-equivariant-zero-center}で
所定の群への同変な選択を、系\ref{cor:pg-equivariant-character-descent}と
命題\ref{prop:pg-equivariant-zero-stratum}、\ref{prop:pg-equivariant-singular-strata}で
所定の条件下の降下を証明する。
節点・尖点軌道への移送には既知の軌道商の理論
\cite[Propositions~7.12, 7.17]{PMilneGroups}を、
特異点直線の捻りにはgerbe上の重み零の層の既知の降下
\cite[Proposition~4.10, p.~1487]{PBerghSchnurerGerbes}を用いる。
これらの一般原理自体を新規性の根拠とはしない。
定理\ref{thm:pg-intrinsic-rank-two-center}では、既知のBorel--Weil--Bott消滅、
不変式環の有限生成、Noetherian完備化の完全性を、
全support上のCasimir核と有限量子交換子行列に適用する。
追加されるのは、三変数制限の型を限定しない共変な中心元の選択と、
系\ref{cor:pg-intrinsic-rank-two-descent}の範囲での全階数二層への降下である。
ここでは、点ごとの非スカラー中心元や所定の形式中心の存在に加えて、
全supportを動かす底の上で古典剰余を固定した代数的な大域切断を選び、
全$\mathrm{GL}_N\times\mathbf G_m$への共変性を課している。
ここに挙げた文献の箇所から、この同時選択と降下が直接与えられるわけではない。
また、重み零の降下、Veronese同値、および底の可逆層による次数別の捩りは
既知の一般原理である\cite{PBerghSchnurerGerbes,PSmithVeronese,PVdBNoncommP1Bundles}。
本稿ではそれらを全supportの族へ適用し、さらに階数二の同変Picard群と、
全非零階数層の安定化群を用いてgerbeの位数を決定する。
零制限の標準中心直線への非延長も別に計算する。
満supportの構成では、既知の還元群によるアフィン商の定理
\cite[\S2, Proposition~2.2, Theorem~2.1]{PArzhantsevMatsushima}と
有限射影加群の切断の持上げを用いる。
ブロックの点ごとの内在性は\cite[縮約二形式による復元補題]{PoissonCubic}に既に含まれ、
補題\ref{lem:pg-intrinsic-block-cover}はその族としての表示である。
定理\ref{thm:pg-intrinsic-full-support-center}で加わる内容は、
順序を選ばない全満support層における指定二次jet付きの代数的な同時選択と、
系\ref{cor:pg-full-support-veronese-descent}の中心商族である。
ただし、関連文献全体に対する新規性の最終確認は残る。
一般の形式的量子化とunimodularityに関する
Kontsevich・Dolgushevの理論\cite{PKontsevich,PDolgushev}は背景であり、
本稿の有限Cramer構成の証明入力ではない。


\paragraph{ambient復元の位置づけ。}
二次包絡によるambient復元は低次数の核の議論であり、
偏極付き圏のsection algebra復元はArtin--Zhangの理論
\cite{PAZAmbient}の適用として記す。
Serre関手からの復元には中心・自由性・滑らかさ等の追加仮定が必要である。
これらは中心商から保たれるデータを明確にするものであり、
全ての非可換三次超曲面の分類や、全閉変形パラメータでの
残余圏のsmoothnessを与えるものではない。

\subsection*{証明の依存関係と本文の構成}
第\ref{sec:cq-classical-input}節で古典的入力と記法を固定する。
第\ref{sec:pg-quantization-foundations}節の一ブロックのPBW性から始め、
第\ref{sec:pg-fermat-mixed-universal}--\ref{sec:pg-universal-determinant}節で
具体的・普遍的な中心元を構成する。
第\ref{sec:pg-quantization}--\ref{sec:pg-all-dimensions}節で全ブロックへ組み立て、
第\ref{sec:pg-equivariant-central-lifts}節で層ごとの同変化と降下を扱う。
第\ref{sec:pg-intrinsic-rank-two-center}節は階数二の全support、
第\ref{sec:pg-full-support-quantization}節は全階数層のgerbeと
満support層の同時量子化を扱う。
その後にFermat型と非楕円型の具体例を置き、付録に既知中心公式とambient復元を収める。

本稿における新規性の網羅的確認と外部査読は完了していない。
中心元の有限次数の厳密計算、全次数のPBW性・中心正則性の証明、
$qgr$のsmoothness、および非可換同型類の比較は異なる主張として扱う。
中心変数を残す一般の複数ブロックでの全座標群への同変な同時選択、
全変形直線への延長、規格化の異なる中心商の同型・Morita比較は残る問題である。

% !TeX root = ../../cubic-quantization.tex
\section{古典的入力と記法}\label{sec:cq-classical-input}
本稿では$V$を$N=n+2\ge5$次元の複素ベクトル空間、
$F\in\operatorname{Sym}^3(V^*)$を滑らかな射影超曲面
$X=V(F)\subset\PP(V)$を定める三次式とする。
\cite[大域双ベクトル場の表示およびPoisson正規形定理]{PoissonCubic}により、
$X$上の正則Poisson構造は定数三重ベクトル$P\in\Lambda^3V$から
$\pi_X=(\iota_{dF}P)|_X$と表され、対応する錐上の双ベクトル$\pi=\iota_{dF}P$もPoissonとなる。
以下ではその構造定理を入力として用いる。

\begin{theorem}[古典的入力：Poisson正規形 \cite{PoissonCubic}]\label{thm:pg-normal}
任意の非零 Poisson 構造に対して、$N=3k+m$ となる $k\geq1$ と線形座標
$x_s=(x_{s1},x_{s2},x_{s3})$、$u=(u_1,\ldots,u_m)$ が存在し、
\begin{equation}
P=\sum_{s=1}^k\partial_{x_{s1}}\wedge\partial_{x_{s2}}\wedge\partial_{x_{s3}},
\qquad F=\sum_{s=1}^k f_s(x_s,u)\label{eq:pg-normal}
\end{equation}
となる。括弧は
\begin{equation}
\{x_{s1},x_{s2}\}=(f_s)_{x_{s3}},\quad
\{x_{s2},x_{s3}\}=(f_s)_{x_{s1}},\quad
\{x_{s3},x_{s1}\}=(f_s)_{x_{s2}}\label{eq:pg-bracket}
\end{equation}
と、異なるブロック間および $\{u_\alpha,-\}$ が零であることで定まる。
逆に\eqref{eq:pg-normal}を満たす滑らかな $F$ はこの括弧によって Poisson 構造を定める。
$\pi$ と $\pi_X$ の一般階数はともに $2k$ である。
\end{theorem}

この定理のEuler--Koszul表示・Hessian線形化・Lie代数の根基排除を用いた
証明は\cite{PoissonCubic}にある。本稿で証明する量子化の存在をその証明に用いない。
三変数制限$f_s(x_s,0)$は滑らかとは限らず、特異または恒等的に零でもよい。
また$f_s$への純中心三次項の配分は一意でない。
量子化に必要な各$f_s$全体の滑らかさを実現する配分は
補題\ref{lem:pg-central-splitting-smooth}で証明する。

\subsection{相対階数層を扱う際の入力}
$P$のsupportは$\bigoplus_sW_s$であり、順序を忘れた各$W_s$は
縮約二形式の平方零スキームから復元される
\cite[縮約二形式による復元補題]{PoissonCubic}。
固定した$P$と適合する三次式は、異なるブロック方向の混合二階偏微分が
零となるものに一致する。このHessian条件と、滑らかな三次式の開条件を
用いた相対Poisson階数層の記述は
\cite[Jacobi条件の線形化および相対Poisson成分定理]{PoissonCubic}による。
本稿の第\ref{sec:pg-full-support-quantization}節では、この入力を有限\'etaleな
ブロック被覆として相対化し、量子化の同変構成と降下に用いる。
固定した超曲面上のPoissonスキーム、その自己同型群による商、
三次式自体を動かす相対モジュライは区別する。
階数層への降下では、明記しない限りその被約構造を取る。

\subsection{代数と中心商の規約}
周囲の二次代数を$A$、中心三次元を$\Omega$、中心商を$B=A/(\Omega)$と書く。
族には$\mathcal A,\mathcal B$を使い、$u_1,\ldots,u_m$は中心変数である。
量子化パラメータ$\hbar$、または各ブロックの独立パラメータは次数零、
$x_{si},u_\alpha$は次数一とする。
純中心三次項を加えた$\Omega'=a\Omega+c$は、同じ古典極限を持っても
一般には同じ中心イデアルを定めない。この比較は、商の抽象同型や
$\operatorname{qgr}$同値の判定とは別である。
\cite{NCCubicMain}のモジュライ部分ではambientと中心商の文字の使い方が
逆になる箇所があるため、引用時には代数の役割によって対応させる。
\par

% !TeX root = ../../cubic-quantization.tex
\section{一ブロックの相対量子化とPBW性}
\label{sec:pg-quantization-foundations}
ここから代数的量子化を構成する。周囲の二次代数を$A$、
その正則中心三次商を$B=A/(\Omega)$と書く。
まず一ブロックの全次数PBW性と中心正則拡大の入力を準備する。
中心公式の有限次数の検算、全次数のPBW性、および
中心元の左右正則性は、それぞれ別の主張として証明する。

\subsection{次数付き量子化とAS条件}

\begin{definition}
連結次数付き可換代数 $C$ に次数零の Poisson 括弧があるとする。
$0\in U\subset\mathbb A^1_\hbar$、$K=\OO(U)$ 上の次数付き量子化とは、
$\deg\hbar=0$、$\mathcal C_0=K$、各 $\mathcal C_d$ が有限自由な
次数付き代数 $\mathcal C$ と $\mathcal C/(\hbar)\simeq C$ であって、
\begin{equation}
\frac{[\widetilde a,\widetilde b]}\hbar\bmod\hbar=\{a,b\}
\end{equation}
を満たすものをいう。特殊 fiber の可換性と平坦性により持上げの選択に依存しない。
\end{definition}
連結次数付き代数 $A$ が Koszul であるとは、自明加群 $\CC=A/A_{\geq1}$ が線形自由分解をもつことをいう。
左右の injective dimension が $d$ で、左右それぞれ
\begin{equation}
\Ext_A^i(\CC,A)=
\begin{cases}\CC(\ell)&i=d,\\0&i\neq d\end{cases}
\end{equation}
ならば、次元 $d$、AS parameter $\ell$ の AS-Gorenstein 代数という。
さらに大域次元 $d$ と有限 GK 次元をもてば AS-regular という。
次数移動は $M(\ell)_j=M_{j+\ell}$ とする。
以下の Ext は全内部次数を含み、injective/global dimension は通常の左右加群圏で取る。

\subsection{ブロック代数と PBW 性}
$R=\CC[u_1,\ldots,u_m]$ とし、$\sym_R$ は中心係数を外へ出して active 変数の語を平均する対称化である。
一ブロック $f(x_1,x_2,x_3,u)$ に対し
\begin{equation}
A_f=R[\hbar]\langle x_1,x_2,x_3\rangle/
\bigl([x_2,x_3]-\hbar\sym_R(f_{x_1}),\ [x_3,x_1]-\hbar\sym_R(f_{x_2}),\
[x_1,x_2]-\hbar\sym_R(f_{x_3})\bigr).\label{eq:pg-block-algebra}
\end{equation}
これは相対 cyclic potential
\begin{equation}
W=x_1x_2x_3-x_2x_1x_3-\hbar\sym_R(f)\label{eq:pg-potential}
\end{equation}
の active cyclic derivatives で定まる。
cyclic word は自由代数を交換子の線形包で割った類であり、
$\partial_i^{\mathrm{cyc}}(a_1\cdots a_r)
=\sum_{a_j=x_i}a_{j+1}\cdots a_ra_1\cdots a_{j-1}$ とする。

\begin{proposition}\label{prop:pg-PBW}
有限個のブロックに共通な $0\in U_{\mathrm{reg}}\subset\mathbb A^1$ を選び、以下を満たせる。
$g_s(x_s)=f_s(x_s,0)$ とおくと、全 $\lambda\in U_{\mathrm{reg}}$ で
\begin{equation}
S_{s,\lambda}=\CC\langle x_{s1},x_{s2},x_{s3}\rangle/
([x_{s2},x_{s3}]-\lambda\sym((g_s)_{x_{s1}}),\text{循環置換})
\end{equation}
は三次元 Koszul AS-regular domain であり、3-Calabi--Yau である。
その Hilbert 級数は $(1-t)^{-3}$。
active 次数で filter し、$u_\alpha$ の filtration degree を零にすると
\begin{equation}
\operatorname{gr}_xA_{f_s,\lambda}\simeq S_{s,\lambda}\otimes_\CC R.
\end{equation}
族全体にも対応する同型があり、$A_{f_s}|_{U_{\mathrm{reg}}}$ は
$R_{\mathrm{reg}}=\OO(U_{\mathrm{reg}})[u]$ 上自由である。
\end{proposition}
\begin{proof}
一つの $g$ について関係式の列は $M_\lambda x$ で表される。ただし
\begin{equation}
M_\lambda=
\begin{pmatrix}
-\frac\lambda2g_{11}&-x_3-\frac\lambda2g_{12}&x_2-\frac\lambda2g_{13}\\
x_3-\frac\lambda2g_{21}&-\frac\lambda2g_{22}&-x_1-\frac\lambda2g_{23}\\
-x_2-\frac\lambda2g_{31}&x_1-\frac\lambda2g_{32}&-\frac\lambda2g_{33}
\end{pmatrix}.
\end{equation}
反対称成分が三つの交換子を与えるので関係式は独立である。
potential 由来の表示は\cite[Corollary 2.7]{PMoriSmith}の意味で standard であり、
同論文 Theorem 1.4に再掲されたArtin--Tate--Van den Berghの判定により、AS-regular 性は全 $\PP^2$ 上で $\rank M_\lambda\geq2$ と同値である。
階数一以下の incidence は $\mathbb A^1\times\PP^2$ の閉集合で、その底への像も閉じている。
$M_0$ は全 $\PP^2$ 上で階数二の交代行列なので、像を除けば原点を含む開集合を得る。
有限個のブロックについて交わりを取る。
3-CY 性と domain 性は\cite[Corollary 2.9]{PMoriSmith}による。
三生成元三二次関係式の三次元 AS-regular 代数では、最小分解を双対化すると最後の shift は三となり、
\begin{equation}
0\longrightarrow S_{s,\lambda}(-3)\longrightarrow S_{s,\lambda}(-2)^3
\longrightarrow S_{s,\lambda}(-1)^3\longrightarrow S_{s,\lambda}
\longrightarrow\CC\longrightarrow0.
\end{equation}
これが Koszul 性と Hilbert 級数を与える。

$\lambda$ を固定する。$u=a$ への特殊化では、\eqref{eq:pg-potential}は
$g$ の三次 potential に active 次数二以下の項を加えたものである。
\cite[Theorem 1.1(i)]{PBergerTaillefer}により、その代数は $S_{s,\lambda}$ の PBW 変形である。
$S_{s,\lambda}$ の斉次基底を語の線形結合として持ち上げる。
最高 active 次数の関係式から次数帰納法で、これらは $A_{f_s,\lambda}$ を $R$ 上生成する。
線形関係があれば全 $u=a$ で PBW 基底間の関係となり、係数多項式が全点で零だから元の関係も零となる。
従って所要の PBW 性と自由性を得る。

族の斉次 $S_{s,U_{\mathrm{reg}}}$ の各次数部分は底上有限生成で、全 fiber の次元が一定である。
底環は $\CC[\hbar]$ の局所化で PID なので、各部分は有限自由である。
torsion があればその台で fiber 次元が増えるからである。
その各次数の基底を自由代数へ持ち上げ、同じ帰納法と全 $(\lambda,a)$ への特殊化を使えば、族上の自由性も従う。
\end{proof}
\begin{remark}\label{rem:pg-regular-all}
$g_s$ が滑らかな平面三次曲線を定めるとき、
\cite[Theorem 4.2(1)]{PMoriSmith}により $S_{s,\lambda}$ の AS-regular 性のために除く $\lambda$ はない。
$g_s=0$ の場合も常に多項式環である。
中心元の代数化・正則性のためには、以下でさらに開集合を縮める。
ambient の AS-regular 性は、その中心超曲面商の圏論的滑らかさとは別条件である。
\end{remark}
以下
$A_{\mathrm{reg}}=A_{f_1}\otimes_{R_{\mathrm{reg}}}\cdots\otimes_{R_{\mathrm{reg}}}A_{f_k}$
とする。異なるブロックは可換で、各全次数部分は底上有限自由であり、
$A_{\mathrm{reg}}/(\hbar)\simeq S$。関係式の一次項は\eqref{eq:pg-bracket}である。

\subsection{代数的中心元の規格化と正則中心拡大}
任意の滑らかな三次ブロックにおける中心元の存在とCramer表示は、
後の定理\ref{thm:pg-universal-determinant}で証明する。
同定理が使う古典的Casimir核の計算は、
三変数制限が特異・零の場合も成立する。
さらに複数ブロックを持つ滑らかな$F$では
補題\ref{lem:pg-central-splitting-smooth}によって
中心変数だけの三次項を配分し、各ブロックを
\emph{全変数に関して滑らかな三次式}に選べる。
これらを合わせれば全Poisson構造に対して
有限次数の行列式演算で代数的中心三次元を得る。

\begin{lemma}[中心三次元の偶性]\label{lem:pg-center-parity}
命題\ref{prop:pg-PBW}の代数的族$A_f$について、
$\hbar\mapsto-\hbar$で不変な$0$を含む開集合へ必要なら縮める。
生成元に
\begin{equation}
 x_i^\dagger=x_i,\qquad u_\alpha^\dagger=u_\alpha,\qquad
 \hbar^\dagger=-\hbar,\qquad(ab)^\dagger=b^\dagger a^\dagger
\end{equation}
を指定すると、$\CC$-線形の反自己同型$\dagger$が定まる。
中心元$\Omega\in Z(A_f)_3$、$\Omega\bmod\hbar=f$があれば、
同じ剰余をもつ中心元を$\Omega^\dagger=\Omega$となるように選べる。
特に形式的中心元も同様に選べ、Weyl対称化表示では
\begin{equation}
 \widehat\Omega
 =\sym_R\bigl(f+\hbar^2 f_2+\hbar^4 f_4+\cdots\bigr),
 \qquad f_{2r}\in R[x_1,x_2,x_3]_3
\end{equation}
である。ここで添字$3$は中心変数$u_\alpha$も次数一と
数えた全次数を意味する。
\end{lemma}
\begin{proof}
定義関係式
$r_i=[x_j,x_k]-\hbar\sym_R(f_{x_i})$に対し、
$r_i^\dagger=-r_i$である。したがって反自己同型が商へ降り、
中心を保つ。$\dagger^2=1$なので
$(\Omega+\Omega^\dagger)/2$は中心であり、$\hbar=0$では$f$に一致する。
完備化でも同じ操作が可能である。

全次数三のWeyl対称化
$\sym_R:R[x_1,x_2,x_3]_3[[\hbar]]\to\widehat A_{f,3}$
は$\hbar=0$で恒等写像となり、
PBW性と完備Nakayamaにより同型である。
代数的族についてもこの有限自由加群間の写像は
原点で可逆なので、さらに開集合を縮めれば同型となる。
対称化した単語は反転で不変であるから、
$\dagger$はその係数の$\hbar\mapsto-\hbar$に対応する。
したがって不変な中心元の対称化係数は$\hbar$の偶冪だけを含む。
\end{proof}

\begin{proposition}[普遍的なCasimirの二次補正]\label{prop:pg-universal-casimir-jet}
$f\in R[x_1,x_2,x_3]$を全次数三の任意の多項式とし、
\begin{equation}
 H_f=\det\left(\frac{\partial^2f}
 {\partial x_i\partial x_j}\right)_{i,j=1}^3,\qquad
 \Theta_f=\sym_R\left(f-\frac{\hbar^2}{12}H_f\right)\in A_f
\end{equation}
とおく。命題\ref{prop:pg-PBW}の開集合上で
\begin{equation}
 [\Theta_f,x_i]\in\hbar^5 A_f\qquad(i=1,2,3),
 \qquad [\Theta_f,u_\alpha]=0.
 \label{eq:pg-casimir-jet}
\end{equation}
従って$\Theta_f$は$\hbar^5$を法として中心元であり、
法$\hbar$で$f$となる。
さらに補題\ref{lem:pg-center-parity}で正規化した
任意の形式的中心元について、$\hbar^2$係数$f_2$は
\begin{equation}
 \left\{f_2+\frac1{12}H_f,x_i\right\}_f=0\quad(i=1,2,3)
 \label{eq:pg-jet-casimir-condition}
\end{equation}
を満たす。
最後の条件は$f_2=-H_f/12$が
\emph{全次数で厳密な中心元に延長する}ことを主張しない。
\end{proposition}
\begin{proof}
まず$\hbar^4$を法とした恒等式を示す。
係数が代数的に独立な二十個の不定元である
\begin{equation}
 f_{\mathrm{univ}}=\sum_{a+b+c\le3}
 t_{abc}x_1^a x_2^b x_3^c
\end{equation}
を考える。これは中心係数環を任意に特殊化できる普遍式である。
自由語では、隣り合った逆順の文字を
\begin{equation}
 x_2x_1\rightsquigarrow x_1x_2-\hbar\sym(f_{x_3}),\quad
 x_3x_1\rightsquigarrow x_1x_3+\hbar\sym(f_{x_2}),\quad
 x_3x_2\rightsquigarrow x_2x_3-\hbar\sym(f_{x_1})
 \label{eq:pg-jet-reductions}
\end{equation}
で置き換える。各置換は代数の定義関係式そのものである。
$\hbar$次数が三以下の範囲では、
逆転数が減るか$\hbar$次数が増えるため、還元は必ず停止する。
$\sym(f_{\mathrm{univ}}-\hbar^2H_{f_{\mathrm{univ}}}/12)$と
各$x_i$との交換子を展開して\eqref{eq:pg-jet-reductions}で還元すると、
次数三以下の各係数が$\mathbb Q[t_{abc}]$内で零になる。
これを古典Jacobian括弧$\{-,-\}_f$で書き直すと、
より明示的な普遍恒等式
\begin{equation}
 [\sym_R(f),x_i]\equiv
 \frac{\hbar^3}{12}\sym_R(\{H_f,x_i\}_f)
 \pmod{\hbar^4 A_f}
 \label{eq:pg-jet-hessian-identity}
\end{equation}
を得る。実際、$\sym_R(H_f)$の$x_i$との交換子の一次項は
$\hbar\sym_R(\{H_f,x_i\}_f)$なので、
この恒等式は$\Theta_f$の可換性の計算と同値である。
二十個の独立係数を保持して全係数を厳密に展開する
再現可能な検算コードは
\texttt{computations/quantization/verify\_cubic\_casimir\_jet.py}である。
数値の特殊化や標本点での検査ではない。
よって任意の中心係数環への特殊化で
$[\Theta_f,x_i]\in\hbar^4A_f$となる。

一方、$\Theta_f^\dagger=\Theta_f$であり、
$[\Theta_f,x_i]^\dagger=-[\Theta_f,x_i]$である。
PBW性により$\hbar$は非零因子なので、
$[\Theta_f,x_i]=\hbar^4b_i$と書くと
$b_i^\dagger=-b_i$となる。
法$\hbar$では反転は可換多項式環の恒等写像だから
$b_i\bmod\hbar=0$、従って$b_i\in\hbar A_f$。
これで法$\hbar^5$の結論を得る。
中心変数との可換性は定義から従う。

最後に$\widehat\Omega=\sym_R(f+\hbar^2f_2+O(\hbar^4))$
との交換子を比較する。
前二段の恒等式および量子化の一次Poisson条件から
$\hbar^3\{f_2+H_f/12,x_i\}_f=0\bmod\hbar^4$
となり、PBW性から\eqref{eq:pg-jet-casimir-condition}が従う。
\end{proof}
\begin{remark}[Hesse公式との照合]
$f=(x_1^3+x_2^3+x_3^3)/3+t x_1x_2x_3$なら
\begin{equation}
 H_f=2(t^3+4)x_1x_2x_3-2t^2(x_1^3+x_2^3+x_3^3).
\end{equation}
したがって$-\hbar^2H_f/12$は
定理\ref{thm:explicit-hesse-center}の
中心式の$\hbar^2$係数と一致する。
これは既知のSklyanin公式との整合性チェックであって、
一般$f$の全高次中心公式が同じ形になるという主張ではない。
\end{remark}

\begin{lemma}\label{lem:pg-central-linear}
$D$ を連結次数付き代数、$z\in D_1$ を中心正則元とする。
$C=D/(z)$ が次元 $d$、AS parameter $d$ の Koszul AS-regular なら、
$D$ は次元 $d+1$、AS parameter $d+1$ の Koszul AS-regular である。
\end{lemma}
\begin{proof}
$0\to D(-1)\xrightarrow{z}D\to C\to0$ より、
$\operatorname{Tor}^D_q(C,\CC)$ は $q=0,1$ でそれぞれ $\CC,\CC(-1)$、他は零である。
change-of-rings spectral sequence
\begin{equation}
E^2_{p,q}=\operatorname{Tor}^C_p(\CC,\operatorname{Tor}^D_q(C,\CC))
\Longrightarrow\operatorname{Tor}^D_{p+q}(\CC,\CC)
\end{equation}
の各項は内部次数 $p+q$ に集中する。
微分は内部次数を保ち全ホモロジー次数を一減らすため全て零であり、
$D$ は Koszul、$j>d+1$ で $\operatorname{Tor}^D_j(\CC,\CC)=0$、$j=d+1$ では非零である。
$D^!$ を二次双対とすると、Koszul 双加群複体の項は $D\otimes(D_i^!)^*\otimes D$。
全内部次数を固定し右の $D$ 次数で filter すれば有限 filtration となり、
随伴複体の正ホモロジーは左 Koszul 分解で消える。
よって増大双加群複体は長さ $d+1$ の完全列であり、右加群として分裂完全なので
任意の左加群とテンソルしても完全である。従って通常の大域次元も $d+1$。
さらに $\RHom_D(C,D)\simeq C(1)[-1]$ と導来随伴から
\begin{equation}
\Ext_D^{i+1}(\CC,D)\simeq\Ext_C^i(\CC,C)(1).
\end{equation}
右加群でも同様で、左右 AS 条件と parameter を得る。
$\operatorname{Hilb}_D=\operatorname{Hilb}_C/(1-t)$ なので GK 次元も有限である。
\end{proof}
\begin{lemma}\label{lem:pg-central-cubic}
$A$ を次元 $N$、AS parameter $N$ の AS-regular 代数、$\eta\in A_3$ を中心正則元とする。
$B=A/(\eta)$ は左右 injective dimension $N-1$、AS parameter $N-3$ の AS-Gorenstein 代数で、
$\operatorname{Hilb}_B=(1-t^3)\operatorname{Hilb}_A$ である。
\end{lemma}
\begin{proof}
$0\to A(-3)\xrightarrow{\eta}A\to B\to0$ が Hilbert 級数を与える。
導来随伴により任意の左 $B$-加群 $M$ について
$\Ext_B^i(M,B)\simeq\Ext_A^{i+1}(M,A)(-3)$。
右辺は $i\geq N$ で零、$M=\CC$、$i=N-1$ で $\CC(N-3)$ となる。
よって injective dimension はちょうど $N-1$ で AS 条件も成立する。右側も同様。
\end{proof}


% !TeX root = ../../cubic-quantization.tex
\section{Fermat型の普遍混合中心三次公式}
\label{sec:pg-fermat-mixed-universal}
本節では代数的な普遍中心元の公式を直接構成するため、
Fermat型の三変数Jacobianブロックの\emph{全混合係数}に対する
有限多項式の中心三次元を構成する。
六つの係数を互いに独立な記号として保ち、
$\hbar$も不定元とした全係数恒等式の証明を与える。

\begin{theorem}[普遍Fermat混合中心三次元]
\label{thm:pg-fermat-mixed-universal}
$R$を標数零の可換次数付き$\CC$-代数とし、
$A,B,C\in R_1$、$D,E,H\in R_2$、$G\in R_3$を任意に取る。
\begin{equation}
 F=\frac{x^3+y^3+z^3}{3}
   +Axy+Byz+Czx+Dx+Ey+Hz+G
 \label{eq:pg-fermat-mixed-F}
\end{equation}
とおき、
\begin{equation}
 [x,y]=\hbar(z^2+By+Cx+H),\quad
 [y,z]=\hbar(x^2+Ay+Cz+D),\quad
 [z,x]=\hbar(y^2+Ax+Bz+E)
 \label{eq:pg-fermat-mixed-rel}
\end{equation}
で定まる$R[\hbar]$-代数を$\mathcal A_F$とする。
$T=xyz+xzy+yxz+yzx+zxy+zyx$、
$\sym(xy)=(xy+yx)/2$などと書く。
このとき
\begin{align}
\Omega_{F,\hbar}={}&
 \frac{x^3+y^3+z^3}{3}-\frac{\hbar^2}{9}T
 +A\sym(xy)+B\sym(yz)+C\sym(zx)\notag\\
 &+\left(D+\frac{\hbar^2B^2}{6}-\frac{\hbar^4AC}{18}\right)x
 +\left(E+\frac{\hbar^2C^2}{6}-\frac{\hbar^4AB}{18}\right)y\notag\\
 &+\left(H+\frac{\hbar^2A^2}{6}-\frac{\hbar^4BC}{18}\right)z
 +\frac{\hbar^4}{18}(Az^2+Bx^2+Cy^2)+G
 \label{eq:pg-fermat-mixed-Omega}
\end{align}
は、少なくとも$\hbar=0$を含む共通のPBW開部分上で
$\mathcal A_F$の\emph{厳密な中心元}となり、
$\Omega_{F,0}=F$を満たす。
これは有限次ジェットの主張ではなく、全ての交換子が
厳密に零であるという代数的恒等式である。
特に式は$R$や中心変数の個数によらず有効である。
\end{theorem}
\begin{proof}
証明は有限次元の非可換自由語の\emph{普遍的な整数・有理関数恒等式}
で行う。係数は最初から独立な中心不定元としてよいので、
まず
\begin{equation}
 R_0=\mathbb Q[A,B,C,D,E,H],\qquad
 \deg(A,B,C)=1,\quad\deg(D,E,H)=2
\end{equation}
を考える。$x,y,z$も次数一、$\hbar$は次数零とする。
自由な語の各中央係数は可換に並べ、active変数の語順だけを保持する。

全次数四の自由語は、中央係数の重み$0,1,2,3,4$と
active語の長さ$4,3,2,1,0$ごとに列挙すると339個ある。
三つの関係式\eqref{eq:pg-fermat-mixed-rel}を
左右から全次数合計二の中央係数・active語で掛けると
162個の関係ベクトルを得る。
ここでactive語を長さ$a,b$に分けて
中央重み$c$を選び、$a+b+c=2$とする。
\begin{equation}
 \mathcal V_4=\mathrm{span}_{\mathbb Q(\hbar)}
  \{\text{全次数四の339語}\},\quad
 I_4=\mathrm{span}_{\mathbb Q(\hbar)}
  \{\text{上記162関係ベクトル}\}.
\end{equation}
これらは全次数四で必要な\emph{全ての}二次関係式の倍数である。

有理関数体$\mathbb Q(\hbar)$上で
160個余りの有限行列を正確に行簡約すると
$\dim I_4=144$であり、
\begin{equation}
 [\Omega_{F,\hbar},x],\quad
 [\Omega_{F,\hbar},y],\quad
 [\Omega_{F,\hbar},z]
 \quad\in I_4
 \label{eq:pg-fermat-mixed-certificate}
\end{equation}
が三つとも成立する。
全次数・全中央係数単項式・全active語を保持する
独立な計算実装
\texttt{computations/quantization/verify\_universal\_fermat\_mixed\_center.py}
は$\mathbb Q(\hbar)$の分数体演算だけを用いる。
乱数、有限体への特殊化、混合係数への代入は一切行わない。
スクリプトは339個の列、162個の関係ベクトル、
階数144と三つの零余剰を全て検査する。
したがってこれは任意の$A,B,C,D,E,H$への特殊化に先立つ
\emph{普遍的な恒等式}である。

分数体上での所属から$\hbar=0$を含む開部分へ
戻る操作にも、同じ有限次元行列で証明を与える。
$\hbar=0$では三関係式が通常の可換化関係式となり、
次数四の商は
$\mathbb Q[A,B,C,D,E,H,x,y,z]_4$である
（中央変数の重みは上記のとおり）。
中央重み$0,1,2,3,4$の単項式数は順に
$1,3,9,19,39$であり、
従って商の次元は
\begin{equation}
 1\cdot15+3\cdot10+9\cdot6+19\cdot3+39\cdot1
 =195,\qquad
 \operatorname{rank}I_4|_{\hbar=0}=339-195=144.
\end{equation}
この階数144は$\mathbb Q(\hbar)$上の階数144と等しく、
付属スクリプトでは双方を\emph{別々の行簡約}として検算する。
ゆえに離散付値環$\mathbb Q[\hbar]_{(\hbar)}$上の
この有限行列をSmith標準形にすると、
144個の非零不変因子は全て原点で単元である。
次数四の商は$\hbar$-torsion-freeとなり、
$I_4$は原点で$\hbar$-飽和している。
従って\eqref{eq:pg-fermat-mixed-certificate}は
$\hbar=0$で正則な係数の\emph{多項式関係式の恒等式}として成立する。
その後の任意の可換中心係数環$R$への特殊化は、
分母が原点で非零のスカラー多項式しか含まないため許される。
この最後の降下には全次数のPBW定理を仮定する必要すらない。
ただし周囲の代数のAS正則性や中心元の正則性を
結論したい場合には、別途命題\ref{prop:pg-PBW}などが必要である。
式の剰余は$\hbar=0$を代入すれば直ちに$F$となる。
\end{proof}

\begin{corollary}[対角二次混合を含む全Fermat型]
\label{cor:pg-fermat-all-mixed}
$u$を任意個の中心次数一変数とし、
\begin{equation}
 F=\frac{x^3+y^3+z^3}{3}
 +p\,x^2+q\,y^2+r\,z^2+
 Axy+Byz+Czx+Dx+Ey+Hz+G
 \label{eq:pg-fermat-all-mixed}
\end{equation}
とする。ただし$p,q,r,A,B,C$は中心一次式、
$D,E,H$は中心二次式、$G$は中心三次式である。
$X=x+p$, $Y=y+q$, $Z=z+r$を置き、
\begin{equation}
 \begin{aligned}
 D'&=D-p^2-Aq-Cr,\\
 E'&=E-q^2-Ap-Br,\\
 H'&=H-r^2-Bq-Cp,\\
 G'&=G+\tfrac23(p^3+q^3+r^3)
     +Apq+Bqr+Crp-Dp-Eq-Hr
 \end{aligned}
 \label{eq:pg-fermat-translation}
\end{equation}
とすると、
$F$は$(X,Y,Z)$に関して
\eqref{eq:pg-fermat-mixed-F}の形となり、
その係数は$A,B,C,D',E',H',G'$である。
従って定理\ref{thm:pg-fermat-mixed-universal}の公式へ
これらを代入し、$X,Y,Z$を$x+p,y+q,z+r$に戻すことで、
\emph{任意の}中心変数混合項を含むFermat型の
明示的な有限多項式中心三次元を得る。
\end{corollary}
\begin{proof}
$F(x,y,z,u)$に$x=X-p$, $y=Y-q$, $z=Z-r$を代入して
三次式を展開すると\eqref{eq:pg-fermat-translation}を得る。
$p,q,r$は中心なので、線形平行移動は交換子を保ち、
Weyl対称化と可換である。
ゆえに定理の中心元をこの座標同型で引き戻せる。
\end{proof}

\begin{remark}[Hessianジェットとの一致と新規性]
\label{rem:pg-fermat-mixed-hessian}
\eqref{eq:pg-fermat-mixed-F}のactive Hessian行列式は
\begin{equation}
 \det\operatorname{Hess}_{x,y,z}(F)
 =2(4xyz-A^2z-B^2x-C^2y+ABC).
\end{equation}
従って公式\eqref{eq:pg-fermat-mixed-Omega}は、
Hessianを用いて次の短い\emph{厳密な}表示にも書き直せる：
\begin{equation}
\boxed{\begin{aligned}
\Omega_{F,\hbar}
 =\sym\!\Bigl(&F-\frac{\hbar^2}{12}\det\operatorname{Hess}_{x,y,z}(F)
 +\frac{\hbar^2}{6}ABC\\
 &+\frac{\hbar^4}{18}
 (Az^2+Bx^2+Cy^2-ACx-ABy-BCz)\Bigr).
\end{aligned}}
\label{eq:pg-fermat-mixed-hessian}
\end{equation}
ここで$\sym$は全active単項式へWeyl対称化を施し、
中央係数はそのままとする。
この式は命題\ref{prop:pg-universal-casimir-jet}の二次Hessian補正と整合し、
残りの純中心項$\hbar^2ABC/6$はCasimir規格化の不定性を表す。
重要なのは、$\hbar^4$補正も含む本式が
有限次ジェットでなく\emph{厳密な中心元}であることである。

本節では、全ての中心係数を独立な記号として保持して
式\eqref{eq:pg-fermat-mixed-Omega}の中心性を厳密に証明し、
\eqref{eq:pg-fermat-all-mixed}への明示的拡張を与えた。
先行文献に同じ公式が含まれるかは未確認であり、
無条件な新規性は主張しない。
一般のHesseパラメータ$t$については
定理\ref{thm:pg-universal-hesse-mixed}で
全中心混合方向の短い$\hbar$の多項式公式を既に証明した。
さらに三変数制限が特異または零の場合も、
全体の三次式が滑らかなら
定理\ref{thm:pg-universal-determinant}の
代数的Cramer中心元が存在する。
本節の式はその$t=0$の場合の
簡明な閉公式として位置づける。
\end{remark}

% !TeX root = ../../cubic-quantization.tex
\section{一般Hesse型の全混合方向に対する普遍中心三次式}
\label{sec:pg-hesse-mixed-universal}
本節ではFermat型に限らず、Hesse三次形式
\begin{equation}
 g_t(x,y,z)=\frac{x^3+y^3+z^3}{3}+txyz
\end{equation}
の任意のパラメータ$t$に対して、中心変数を含む全ての三次混合項を
\emph{有限個の多項式係数}で量子化した中心三次元を与える。
この構成は形式的中心元の存在定理を使用しない。
非零の混合係数を独立な記号のまま保持し、
$\hbar,t$の有理関数体上で自由語の関係式への所属を
厳密に検証することで証明する。
得られる式はFermat型の
定理\ref{thm:pg-fermat-mixed-universal}を$t=0$として含む。

\subsection{短い多項式による普遍公式}
$R$を標数零の可換次数付き$\CC$-代数とする。
$A,B,C\in R_1$、$D,E,H\in R_2$、$G\in R_3$、
$t\in\CC$を任意に取って
\begin{equation}
 F=g_t(x,y,z)+Axy+Byz+Czx+Dx+Ey+Hz+G
 \label{eq:pg-hm-F}
\end{equation}
とおく。$R[\hbar]$上の次数付き代数
$\mathcal A_{F,\hbar}$を
\begin{equation}
\begin{aligned}
 [x,y]&=\hbar\bigl(z^2+t\sym(xy)+By+Cx+H\bigr),\\
 [y,z]&=\hbar\bigl(x^2+t\sym(yz)+Ay+Cz+D\bigr),\\
 [z,x]&=\hbar\bigl(y^2+t\sym(zx)+Ax+Bz+E\bigr)
\end{aligned}
\label{eq:pg-hm-rel}
\end{equation}
で定める。$\deg\hbar=0$であり、全ての$R$の元は中心とする。
記法
\begin{equation}
\begin{aligned}
 S={}&x^3+y^3+z^3,&
 T={}&\sym(xyz),\\
 Q={}&A\sym(xy)+B\sym(yz)+C\sym(zx),&
 Z_2={}&Az^2+Bx^2+Cy^2,\\
 L={}&Dx+Ey+Hz,&
 R_2={}&B^2x+C^2y+A^2z,\\
 R_{11}={}&ACx+ABy+BCz
\end{aligned}
\end{equation}
を用いる。ここで$\sym(xyz)$は六語の平均である。

\begin{theorem}[Hesse型の普遍混合中心三次式]
\label{thm:pg-universal-hesse-mixed}
$q=\hbar^2$とし、次の七つの多項式を定義する。
\begin{equation}
\begin{aligned}
 U={}&1+\frac{qt^2}{2}+\frac{q^2t}{6}-\frac{q^2t^4}{48},\\
 W={}&\frac{(4-qt^2)(12t+q(t^3-8))}{48},\\
 c={}&\frac{(4-qt^2)(12+qt^2)}{48},&
 d={}&\frac{q(12t+q(2-t^3))}{36},\\
 e={}&1+\frac{q^2t(8-t^3)}{144},&
 f={}&\frac{q(4-qt^2)}{24},\\
 g={}&-\frac{q(12t+q(t^3+4))}{72}.
\end{aligned}
\label{eq:pg-hm-coefficients}
\end{equation}
このとき
\begin{equation}
 \boxed{\quad
 \Omega_{F,\hbar}
 =\frac{U}{3}S+WT+cQ+dZ_2+eL+fR_2+gR_{11}+G
 \quad}
 \label{eq:pg-hm-center}
\end{equation}
は$\hbar=0$を含む適切なZariski開部分上で
$\mathcal A_{F,\hbar}$の\emph{厳密な中心元}となる。
また$\Omega_{F,0}=F$である。
これは$\hbar$について有限な多項式公式であり、
$n$個の中心変数に係数$A,\ldots,G$を特殊化する場合でも
中心変数の個数に依存せず同じ表示が有効である。
$t^3=-1$を排除する条件はこの\emph{中心性の恒等式そのもの}には不要だが、
$g_t=0$が滑らかな平面三次であるためには$t^3\ne-1$が必要である。
\end{theorem}
\begin{proof}
基底係数を
$R_0=\mathbb Q[A,B,C,D,E,H]$、
$\deg A=\deg B=\deg C=1$、
$\deg D=\deg E=\deg H=2$とする。
$t,\hbar$は独立な次数零の不定元とする。
$G$は中心なので零としてよい。
非可換自由語では中央係数を可換な単項式としてまとめ、
$x,y,z$の順序付き語をそのまま保持する。
以下、全次数四の自由語空間を$V_4$、式
\eqref{eq:pg-hm-rel}の三つの二次関係式の
左右からの全次数二の倍数が張る空間を$I_4$とする。

中央係数の重み$0,1,2,3,4$の単項式の個数は
順に$1,3,9,19,39$である。各重み$c$について
$x,y,z$からなる長さ$4-c$の語を列挙すると
\begin{equation}
 \dim_{\mathbb Q(\hbar,t)}V_4
 =1\cdot81+3\cdot27+9\cdot9+19\cdot3+39=339.
\end{equation}
全次数二の左右の倍数を列挙すると
162個の関係ベクトルが得られる。
\emph{これらは全次数四の二次関係式のイデアル成分を生成する。}

式\eqref{eq:pg-hm-coefficients}を
\eqref{eq:pg-hm-center}に代入し、三交換子
$[\Omega_{F,\hbar},x]$、$[\Omega_{F,\hbar},y]$、
$[\Omega_{F,\hbar},z]$を自由語で展開する。
162関係ベクトルの行列を$\mathbb Q(\hbar,t)$上で
厳密に行簡約すると
\begin{equation}
 \operatorname{rank}I_4=144,\qquad
 [\Omega_{F,\hbar},x],\
 [\Omega_{F,\hbar},y],\
 [\Omega_{F,\hbar},z]\in I_4
 \label{eq:pg-hm-certificate}
\end{equation}
となる。この恒等式は六つの中心係数を独立不定元として扱い、
$\hbar,t$を特殊化せずに得られた。
検証コード
\texttt{computations/quantization/verify\_universal\_hesse\_mixed\_center.py}
が全339基底、162関係と三交換子を再生成して
有理関数体上の行簡約を実行する。
有限体や有理数標本での補間は証明に用いない。

分数体上での関係式所属を$\hbar=0$まで降下させるには
階数を比較する。$\hbar=0$で二次関係は通常の可換化関係なので
\begin{equation}
 \dim(V_4/I_4)|_{\hbar=0}
 =15+3\cdot10+9\cdot6+19\cdot3+39=195.
\end{equation}
従って$\operatorname{rank}I_4|_{\hbar=0}=339-195=144$であり、
上の一般階数と一致する。
ここで最後の等式は任意の$t$について成立する。
従って$(\hbar,t)$-平面の直線$\hbar=0$を含む
あるZariski開近傍で、関係行列の階数は144に一定であり、
階数144の小行列は単元として取れる。
行空間とその商はこの近傍上局所自由であるから、
分数体上で零の三交換子はその局所自由商でも零となる。
ゆえに任意の$t$に特殊化しても、
$\hbar=0$の適切な近傍で厳密な中心性が成立する。
この降下に、全次数のPBW性を別途仮定する必要はない。

最後に$q=0$で$U=c=e=1$、$W=t$、
$d=f=g=0$なので$\Omega_{F,0}=F$。
\end{proof}

\begin{corollary}[Hessianによるコンパクトな普遍式]
\label{cor:pg-hesse-hessian-form}
定理\ref{thm:pg-universal-hesse-mixed}の$F$に対し
\begin{equation}
 H_F=\det\operatorname{Hess}_{x,y,z}(F)
\end{equation}
とおく。また上記の$Q,Z_2,L,R_2,R_{11}$を使い
\begin{equation}
\begin{aligned}
 \Psi_t(F)={}&
 \frac{t(8-t^3)}{48}g_t
 -\frac{t^4}{48}Q
 +\frac{2-t^3}{36}Z_2\\
 &+\frac{t(8-t^3)}{144}L
 -\frac{t^2}{24}R_2
 -\frac{t^3+4}{72}R_{11}.
\end{aligned}
\end{equation}
すると普遍中心三次元は
\begin{equation}
 \boxed{\quad
 \Omega_{F,\hbar}
 =\sym\left(
 F-\frac{\hbar^2}{12}H_F
 +\frac{\hbar^2}{6}ABC
 +\hbar^4\Psi_t(F)
 \right).
 \quad}
 \label{eq:pg-hesse-hessian-closed}
\end{equation}
特に$\hbar^2$項は普遍Hessian補正
（命題\ref{prop:pg-universal-casimir-jet}）であり、
中心係数$ABC/6$は中心元の規格化の自由度を反映する。
本式は$\hbar^4$項まで含めることで
\emph{全ての$\hbar$次数について厳密に中心}となる。
\end{corollary}
\begin{proof}
$H_F$をactive変数だけで微分して展開すると
\begin{equation}
 H_F=-2t^2S+2(t^3+4)xyz+
 2t^2Q-4tZ_2-2R_2+2tR_{11}+2ABC
\end{equation}
である。
これを右辺に代入して$\hbar^2$と$\hbar^4$の係数を比較すると、
式\eqref{eq:pg-hm-center}と恒等的に一致する。
従って定理の中心性から従う。
\end{proof}

\begin{remark}[既存の中心元との整合性]
$t=0$では
$U=1$、$W=-2\hbar^2/3$、$c=e=1$、
$d=\hbar^4/18$、$f=\hbar^2/6$、
$g=-\hbar^4/18$となる。
ゆえに定理\ref{thm:pg-fermat-mixed-universal}の
全混合Fermat式を正確に回復する。
$A=B=C=D=E=H=G=0$では
定理\ref{thm:explicit-hesse-center}の
Sklyanin中心元を回復する。
それらは既知の入力・先行研究上の貢献と区別する。
\end{remark}

\begin{remark}[Etingof--Ginzburgの中心公式との関係]
\label{rem:pg-hesse-eg-translation}
Etingof--Ginzburg\cite[\S9, p.~1414]{EG}には、
一変数ごとの二次・一次項を含む三生成元代数に対する
Rainsの中心元公式が記載されている。
式\eqref{eq:pg-hm-F}でも$t\ne0$なら、中心係数による平行移動
\begin{equation}
 \begin{aligned}
 X&=x+\frac Bt,\qquad Y=y+\frac Ct,\qquad Z=z+\frac At,\\
 F^{\flat}(X,Y,Z)
 &=F\left(X-\frac Bt,Y-\frac Ct,Z-\frac At\right).
 \end{aligned}
 \label{eq:pg-hm-eg-coordinates}
\end{equation}
により混合二次項を除ける。具体的には
\begin{equation}
 \begin{aligned}
 F^{\flat}={}&g_t(X,Y,Z)
 -\frac BtX^2-\frac CtY^2-\frac AtZ^2\\
 &+\widetilde D X+\widetilde E Y+\widetilde H Z+\widetilde G,
 \end{aligned}
 \label{eq:pg-hm-eg-polynomial}
\end{equation}
ただし
\begin{equation}
 \begin{aligned}
 \widetilde D&=D+\frac{B^2}{t^2}-\frac{AC}{t},\\
 \widetilde E&=E+\frac{C^2}{t^2}-\frac{AB}{t},\\
 \widetilde H&=H+\frac{A^2}{t^2}-\frac{BC}{t},\\
 \widetilde G&=G-\frac{BD+CE+AH}{t}
 +\frac{2ABC}{t^2}-\frac{A^3+B^3+C^3}{3t^3}.
 \end{aligned}
 \label{eq:pg-hm-eg-coefficients}
\end{equation}
これらは三次式の展開による恒等式である。
平行移動の係数は中心かつ次数一なので、
Jacobian交換子を保つ次数付き代数の同型となる。
さらに$1-\hbar t/2$を可逆にして
\begin{equation}
 q_{\mathrm{EG}}=\frac{1+\hbar t/2}{1-\hbar t/2},\qquad
 \mu_{\mathrm{EG}}=\frac{\hbar}{1-\hbar t/2}
 \label{eq:pg-hm-eg-parameters}
\end{equation}
とおけば、交換関係は
\begin{equation}
 \begin{aligned}
 XY-q_{\mathrm{EG}}YX
 &=\mu_{\mathrm{EG}}\left(Z^2-\frac{2A}{t}Z+\widetilde H\right),\\
 YZ-q_{\mathrm{EG}}ZY
 &=\mu_{\mathrm{EG}}\left(X^2-\frac{2B}{t}X+\widetilde D\right),\\
 ZX-q_{\mathrm{EG}}XZ
 &=\mu_{\mathrm{EG}}\left(Y^2-\frac{2C}{t}Y+\widetilde E\right)
 \end{aligned}
 \label{eq:pg-hm-eg-relations}
\end{equation}
となり、同文献の表示に一致する。
ここで$q_{\mathrm{EG}}$は定理中の$q=\hbar^2$とは異なる。
従って$t$が可逆な範囲では、混合項を含むこと自体を
既知中心公式から独立な新規性の根拠にすることはできない。
本節の公式は、固定したHesse座標と独立な中心係数を保ったまま、
$t=0$も含む多項式係数で$\Omega_{F,0}=F$という規格化を実現する。
上の座標変換が$t=0$で定義されないことは、
別の座標変換による既存表示への帰着を排除するものではない。
式\eqref{eq:pg-hm-eg-polynomial}--\eqref{eq:pg-hm-eg-relations}の
係数恒等式は\texttt{computations/quantization/verify\_hesse\_translation.py}で
独立係数のまま再現できる。
\end{remark}

\subsection{対角二次混合と任意の滑らかな三変数三次式}

\begin{corollary}[全ての中心混合項を含むHesse型]
\label{cor:pg-universal-hesse-all-mixed}
中心次数一の$\ell_x,\ell_y,\ell_z$に対し
\begin{equation}
 \widetilde F=g_t(x,y,z)
 +\ell_xx^2+\ell_yy^2+\ell_zz^2
 +Axy+Byz+Czx+Dx+Ey+Hz+G
\end{equation}
を考える。$X=x+\ell_x$、$Y=y+\ell_y$、
$Z=z+\ell_z$と置く。係数を
\begin{equation}
\begin{aligned}
 A'&=A-t\ell_z,& B'&=B-t\ell_x,& C'&=C-t\ell_y,\\
 D'&=D-\ell_x^2-A\ell_y-C\ell_z+t\ell_y\ell_z,\\
 E'&=E-\ell_y^2-A\ell_x-B\ell_z+t\ell_x\ell_z,\\
 H'&=H-\ell_z^2-B\ell_y-C\ell_x+t\ell_x\ell_y,\\
 G'&=G+\tfrac23(\ell_x^3+\ell_y^3+\ell_z^3)
 +A\ell_x\ell_y+B\ell_y\ell_z+C\ell_z\ell_x\\
 &\qquad-D\ell_x-E\ell_y-H\ell_z
 -t\ell_x\ell_y\ell_z
\end{aligned}
\label{eq:pg-hm-translation}
\end{equation}
と定めると
\begin{equation}
 \widetilde F
 =g_t(X,Y,Z)+A'XY+B'YZ+C'ZX+D'X+E'Y+H'Z+G'.
\end{equation}
したがって定理\ref{thm:pg-universal-hesse-mixed}の公式に
$A',B',C',D',E',H',G'$を代入して
$X=x+\ell_x$などへ引き戻すと、
\emph{任意の}Hesse型三次式と中心変数の混合項に対する
有限多項式の厳密な中心三次元が得られる。
\end{corollary}
\begin{proof}
表示式は三次式の展開による恒等式である。
$\ell_x,\ell_y,\ell_z$は中心元なので、
平行移動はJacobian量子化の交換子を保ち、
次数付き代数の同型を定める。
従って中心元を引き戻せる。
\end{proof}

\begin{corollary}[任意の滑らかな平面三次制限]\label{cor:pg-smooth-ternary-explicit}
$F(x,y,z,u)$を斉次三次式とし、
$g(x,y,z)=F(x,y,z,0)$が滑らかな平面三次曲線を定めるとする。
Hesse標準形\cite[Lemma~1]{ArtebaniDolgachev}により、
$M\in\mathrm{GL}_3(\CC)$、$\lambda\ne0$、
$t^3\ne-1$を選んで
$g(x)=\lambda\,g_t(Mx)$と書ける。
\begin{equation}
 F'(X,u)=\lambda^{-1}F(M^{-1}X,u),
 \qquad \hbar'=\lambda\det(M)\,\hbar
\end{equation}
とおく。$F'$はその混合係数を
系\ref{cor:pg-universal-hesse-all-mixed}の形に一意に書ける。
同系によって構成した中心元$\Omega_{F',\hbar'}$から
\begin{equation}
 \boxed{\quad
 \Omega_{F,\hbar}(x,u)
 =\lambda\Omega_{F',\hbar'}(Mx,u)
 \quad}
 \label{eq:pg-hm-arbitrary-ternary}
\end{equation}
を得る。これは$F$のJacobian量子化において
$\hbar=0$近傍の\emph{厳密な}中心三次元であり、
$\Omega_{F,0}=F$である。
形式的中心元の存在補題を入力せず、
Hesse標準形の選択と有限個の多項式演算だけで構成できる。
\end{corollary}
\begin{proof}
定数三重ベクトルの変換則から、
$X=Mx$の座標ではJacobian交換子の量子化パラメータが
$\hbar'=\lambda\det(M)\hbar$になる。
中心元は代数同型で中心元に移り、
$\lambda$倍によって零ファイバーの剰余が$F$となる。
\end{proof}

\begin{corollary}[滑らかな制限を持つ全Poissonブロックの明示的代数的量子化]
\label{cor:pg-smooth-block-algebraic-quant}
定理\ref{thm:pg-normal}の正規形
$F=\sum_{s=1}^k f_s(x_s,u)$について、
各$g_s=f_s(x_s,0)$が滑らかな平面三次曲線を定めるとする。
各ブロックで系\ref{cor:pg-smooth-ternary-explicit}の
有限多項式中心元$\Omega_{s,\hbar}$を構成すれば
\begin{equation}
 \Omega_\hbar=\sum_{s=1}^k\Omega_{s,\hbar},
 \qquad \Omega_0=F
\end{equation}
は代数的な厳密中心三次元となる。
命題\ref{prop:pg-PBW}の開集合を必要に応じて縮めると
$\Omega_\hbar$は各ファイバーで左右正則であり、
周囲の量子化代数はKoszul AS-regular、
中心商は定理\ref{thm:pg-quantization}の
Hilbert級数とAS-Gorenstein条件を満たす。
この部分族に対する量子化の構成は
形式的中心元存在補題を全く使用しない。
\end{corollary}
\begin{proof}
各ブロックに中心変数を共有させ、
異なるブロックのactive生成元を可換と取る。
個別の中心元は他ブロックと可換なので、その和は中心である。
各$g_s$が滑らかであるから、命題\ref{prop:pg-PBW}は
各ブロックのPBW自由性と$u=0$でのAS正則性を与える。
定理\ref{thm:pg-quantization}の証明のうち
中心元存在の補題の利用を
系\ref{cor:pg-smooth-ternary-explicit}に置き換える。
残る中央正則拡大、active次数filtrationによる
$\Omega_\hbar$の左右正則性、
中心三次商のAS-Gorenstein性の証明はそのまま適用できる。
\end{proof}

\begin{remark}[中心性・平坦性・座標依存性の範囲]
本定理は、非可換量子化の中心性を\emph{全係数の恒等式}として
証明するものであって、各閉ファイバーのAS正則性・domain性や
中心商の平坦性をこの有限次数の検算のみから結論しない。
それらは命題\ref{prop:pg-PBW}、定理\ref{thm:pg-quantization}で
別途扱う。
一つの固定した元の座標に対するGL$_3$同変な短い不変式を
得たとは主張しない。座標選択でHesse標準形を求めた後に
適用できる\emph{有限多項式の統一的公式}である。
三変数制限が節点型・尖点型・零の場合の
\emph{短い多項式閉公式}は本節からは従わないが、
全体の三次式が滑らかなときは
定理\ref{thm:pg-universal-determinant}が
有限Cramer有理式で中心三次元を与える。
従ってそれらの特異制限についても
代数的中心元の\emph{存在}は証明済みである。
最高active次数のpotentialが三次元Calabi--Yau代数を定める範囲では、
一ブロックの非スカラー中心元の存在は既に
Etingof--Ginzburg\cite[Theorem~3.3.2(ii), p.~1385]{EG}の範囲に含まれる。
本節は、注意\ref{rem:pg-hesse-eg-translation}の既存公式との関係の下で、
指定された古典元への規格化と独立な中心係数上の恒等式を明示する。
全係数空間での正則な選択、複数ブロックへの組立て、
座標変換についての同変降下は、それぞれ後続の定理で扱う。
\end{remark}

% !TeX root = ../../cubic-quantization.tex
\section{任意の滑らかなJacobianブロックの代数的中心三次元}
\label{sec:pg-universal-determinant}
本節では、$\hbar$進の形式的中心元、
相対formalityおよびDolgushevのunimodularity定理を
\emph{使わずに}、任意の三変数Jacobianブロックの中心三次元を
代数的に構成する。
重要なのは、三変数制限
$g=f(x,y,z,0)$の滑らかさではなく、
\emph{全変数に関する三次超曲面$f=0$の滑らかさ}である。
$g$が節点型・尖点型・可約、あるいは恒等的に零の場合も含む。
証明は（i）古典的Casimir核の一定次元性、
（ii）Hesse型で既知の明示中心元、
（iii）普遍交換子行列の小行列式の恒等的消滅、
（iv）Cramer公式からなる。
これはEtingof--Ginzburg\cite[\S8.3]{EG}が
potentialの中心元を特殊化する際に用いた
rank一定性・射影的閉性の考え方と近縁だが、
ここでは固定次数の行列式によって
$\hbar=0$の規格化も直接制御する。

$R=\CC[u_1,\ldots,u_m]$、
$S=R[x,y,z]$、$f\in S_3$とし、
\begin{equation}
 \begin{aligned}
 A_{f,\hbar}
 =R[\hbar]\langle x,y,z\rangle/
 \bigl(&[y,z]-\hbar\sym_R(f_x),\\
       &[z,x]-\hbar\sym_R(f_y),\
       [x,y]-\hbar\sym_R(f_z)\bigr)
 \end{aligned}
 \label{eq:pg-det-quantum}
\end{equation}
とおく。$\sym_R$はactive変数のWeyl対称化である。
$f_{\mathrm{cent}}=f(0,0,0,u)$と書く。

\begin{lemma}[滑らかさとactive偏微分の最大公約因子]
\label{lem:pg-det-gradient-coprime}
$V(f)\subset\PP^{m+2}$が滑らかであるとき、
\begin{equation}
 \gcd(f_x,f_y,f_z)=1
 \quad\text{in }\CC[x,y,z,u_1,\ldots,u_m].
 \label{eq:pg-det-gradient-coprime}
\end{equation}
特に$f$には少なくとも一つactive変数を含む三次単項式が現れる。
\end{lemma}
\begin{proof}
$f$がactive変数に依存しなければ、
$\PP(x,y,z)$上で$f$と全偏微分が零となり、
滑らかさに反する。
$f_x,f_y,f_z$が正次数の共通因子を持つと仮定し、
その既約斉次因子を$h$とする。
$h=0$は$\PP^{m+2}$の次元$m+1$の射影超曲面である。
そこへ残る$m$個の斉次方程式
$f_{u_1}=\cdots=f_{u_m}=0$を課しても、
Krullの高さ定理によりその共通零点の
アフィン錐の次元は少なくとも
\begin{equation}
 (m+3)-(m+1)=2.
 \label{eq:pg-det-gradient-dimension}
\end{equation}
従って零ベクトル以外の共通零点が存在する。
そこでは全ての偏微分が零であり、
Euler等式から$f=0$でもある。
これは$V(f)$の滑らかさに矛盾する。
ゆえに共通非定数因子はない。
\end{proof}

\begin{lemma}[全三次式の古典的Casimir核]
\label{lem:pg-det-classical-kernel}
$V(f)\subset\PP^{m+2}$が滑らかなら、
三変数Jacobian Poisson括弧に対し
\begin{equation}
 \boxed{\quad
 \ker\left(S_3\xrightarrow{p\mapsto
 (\{p,x\}_f,\{p,y\}_f,\{p,z\}_f)}
 S_4^{\oplus3}\right)
 =\CC f\oplus R_3.
 \quad}
 \label{eq:pg-det-classical-kernel}
\end{equation}
$m=0$では$R_3=0$と解釈する。
$f(x,y,z,0)$は特異でも零でもよい。
\end{lemma}
\begin{proof}
$S$はUFDであり、補題
\ref{lem:pg-det-gradient-coprime}により
$(f_x,f_y,f_z)$の最大公約因子は1である。
$p\in S_3$がCasimirなら
\begin{equation}
 d_xp\wedge d_xf=0,
 \qquad
 (p_x,p_y,p_z)=a(f_x,f_y,f_z)
 \quad\text{in }\operatorname{Frac}(S)^3
 \label{eq:pg-det-casimir-proportional}
\end{equation}
となる。
既約因子の付値を比較すると、$p_x,p_y,p_z\in S$かつ
$\gcd(f_x,f_y,f_z)=1$より$a\in S$。
両辺は斉次次数2なので、$d_xp\ne0$なら$a$は
次数零の多項式、すなわち定数$c\in\CC$である。
$d_xp=0$なら$p\in R_3$である。
従っていずれの場合も$p-cf\in R_3$。
逆に$f$と$R$はJacobian括弧のCasimirなので
等号を得る。
$f$はactive変数に依存するから直和である。
\end{proof}

\subsection{普遍交換子行列の階数が跳ばないこと}
$\mathcal T_m\subset S_3$を、
$V(f)\subset\PP^{m+2}$が滑らかな三次式の
非空Zariski開集合とする。
$\mathcal T_m$はアフィン空間の開部分なので既約である。
$\mathcal T_m^\circ$を$f(x,y,z,0)$が滑らかな
平面三次を定める非空開部分とする。
$\mathcal T_m^\circ$も稠密である。

命題\ref{prop:pg-PBW}のrank条件を係数変数も含めて考える。
$S_{g,\hbar}$を定める$3\times3$の線形行列
$M_{g,\hbar}$は係数について多項式的であり、
$\hbar=0$では全非零$x\in\PP^2$でrank2である。
従ってrank1以下のincidenceを射影$\PP^2$から
パラメータ空間へproperに押し出すと閉集合が得られ、
その補集合は
\begin{equation}
 \mathcal U_m\subset\mathcal T_m\times\mathbb A^1_\hbar,
 \qquad
 \mathcal T_m\times\{0\}\subset\mathcal U_m
 \label{eq:pg-det-universal-open}
\end{equation}
というZariski開集合となる。
各ファイバーで三次元AS-regular性と
Berger--TailleferのPBW定理
\cite{PMoriSmith,PBergerTaillefer}が成立する。
固定内部次数$d$の商は有限表示であり、
全ファイバー次元が
$\binom{m+d+2}{d}$に等しい。
ここから各次数部分の局所自由性が従うことを明示する。
底のアフィン開部分上でこの商の有限表示行列を取り、
その行数を$q$とすると、全ての閉点で行列の階数は
$r=q-\binom{m+d+2}{d}$である。
底は被約な有限型$\CC$-スキームなので、
全ての$(r+1)$次小行列式は恒等的に零である。
各点の近傍で非零の$r$次小行列式を可逆にし、
行と列の基本変形を行うと行列は
$\operatorname{diag}(I_r,0)$の形となる。
従ってその余核、すなわち次数$d$部分は局所自由である。
$d=3,4$のWeyl対称化基底は$\hbar=0$で
通常の単項式基底なので、さらに
$\mathcal U_m$を零切断の近傍へ縮めれば、
その基底が全てのファイバーで有効となる。

$N_d=\binom{m+d+2}{d}$、
$d_m=\binom{m+2}{3}$、
$r_m=N_3-d_m-1$とおく。
Weyl座標表示で次の有限行列を得る：
\begin{equation}
 \mathcal D_{f,\hbar}:S_3\longrightarrow S_4^{\oplus3},
 \qquad
 p\longmapsto
 \frac1\hbar\left(
 [\sym_R(p),x],[\sym_R(p),y],[\sym_R(p),z]\right).
 \label{eq:pg-det-commutator-matrix}
\end{equation}
$\hbar=0$で交換子が消え、商が$\hbar$-flatなので
右辺の$1/\hbar$は正則であり、
$\mathcal D_{f,0}$は古典Jacobian括弧を与える。

\begin{lemma}[量子交換子行列の普遍階数恒等式]
\label{lem:pg-det-universal-rank}
$\mathcal U_m$の上で
$\mathcal D_{f,\hbar}$の全ての$(r_m+1)$次小行列式は
恒等的に零である。
特に任意の滑らかな$f$について、
$\hbar=0$を含む開集合上で
\begin{equation}
 \operatorname{rank}\mathcal D_{f,\hbar}\le r_m.
 \label{eq:pg-det-universal-rank}
\end{equation}
\end{lemma}
\begin{proof}
まず任意の閉点$f\in\mathcal T_m^\circ$を固定する。
系\ref{cor:pg-smooth-ternary-explicit}が
その$\hbar=0$近傍で$f$を持ち上げる
\emph{代数的な厳密中心三次元}を与える。
その中心元と$d_m$個の純中心三次単項式は
法$\hbar$で線形独立なので、
必要ならこの近傍を縮めれば、
$\mathcal D_{f,\hbar}$の核の次元は
少なくとも$d_m+1$である。
従って全$(r_m+1)$次小行列式は
既約な開曲線$\mathcal U_m\cap(\{f\}\times\mathbb A^1)$の
非空開部分で零であり、その全体で零となる。
これを全ての$f\in\mathcal T_m^\circ$に適用すると、
小行列式は$\mathcal U_m$の稠密開部分
$\mathcal U_m\cap(\mathcal T_m^\circ\times\mathbb A^1)$で零となる。
この議論にはHesse座標を$f$について代数的に同時選択する必要はない。
これらの小行列式は
$\mathcal U_m$上の\emph{正則関数}であり、
$\mathcal U_m$は既約かつ被約である。
したがって小行列式は$\mathcal U_m$上恒等的に零。
これは特異な三変数制限への
\emph{係数についての恒等式の特殊化}であり、
稠密開上の中心元をそのまま特殊化して
分母が消えないと仮定する議論ではない。
\end{proof}

\begin{theorem}[全滑らかな三次ブロックに対する普遍Cramer中心元]
\label{thm:pg-universal-determinant}
$V(f)\subset\PP^{m+2}$が滑らかな任意の$f\in S_3$に対し、
三変数制限$f(x,y,z,0)$が節点型、尖点型、可約、
または恒等的零であっても、
$0$を含むZariski開集合$U_f\subset\mathbb A^1_\hbar$上に
\begin{equation}
 \boxed{\quad
 \Omega_{f,\hbar}\in
 Z\!\left(A_{f,\hbar}|_{U_f}\right)_3,\qquad
 \Omega_{f,0}=f
 \quad}
 \label{eq:pg-det-all-centers}
\end{equation}
が存在する。
しかも$\Omega_{f,\hbar}$は
\emph{有限個の非可換三次単項式の線形結合}であり、
その係数は$\hbar^2$の有理関数で、
全ての分母は$\hbar=0$で非零となる。
以下のCramer公式が、与えられた$f$から中心元を
有限回の四則演算と行列式計算で構成するアルゴリズムである。

$f$のactive変数を含む三次単項式のうち、
係数$a\ne0$を持つ一つを$w_0$とする。
純中心三次単項式と$w_0$を除いた
$r_m$個の三次単項式を
$w_1,\ldots,w_{r_m}$とする。
$\mathcal D_{f,0}|_{\langle w_1,\ldots,w_{r_m}\rangle}$
は単射なので、$r_m$個の行を選んで
\begin{equation}
 M_{ij}(\hbar)=\ell_i(\mathcal D_{f,\hbar}(w_j)),
 \qquad
 b_i(\hbar)=\ell_i(\mathcal D_{f,\hbar}(w_0)),
 \qquad
 \det M(0)\ne0
 \label{eq:pg-det-minor}
\end{equation}
となるようにする。
$M_j$を$M$の第$j$列を$b$に換えた行列とすれば、
\begin{equation}
 \boxed{\quad
 \Omega_{f,\hbar}
 =f_{\mathrm{cent}}+
 a\left(
 \sym_R(w_0)-\sum_{j=1}^{r_m}
 \frac{\det M_j(\hbar)}{\det M(\hbar)}\sym_R(w_j)
 \right).
 \quad}
 \label{eq:pg-det-universal-center}
\end{equation}
が所要の中心元である。
\end{theorem}
\begin{proof}
補題\ref{lem:pg-det-classical-kernel}によって
$\operatorname{rank}\mathcal D_{f,0}=r_m$。
また、$f$の$w_0$係数は非零なので、
古典Casimir核$\CC f\oplus R_3$と
$W=\langle w_1,\ldots,w_{r_m}\rangle$の交叉は零。
従って$\mathcal D_{f,0}|_W$は単射であり、
指定した$r_m\times r_m$小行列$M$が選べる。
$\det M(0)\ne0$なので、$\hbar=0$を含む開集合へ
縮めれば$M$は可逆となる。

一方、補題\ref{lem:pg-det-universal-rank}により
$\mathcal D_{f,\hbar}$のrankは高々$r_m$。
$M$が可逆なので少なくとも$r_m$でもあり、
従って$M$を構成する$r_m$行は
交換子行列の\emph{全ての行}の行空間を張る。
Cramer公式は
\begin{equation}
 \mathcal D_{f,\hbar}
 \left(w_0-\sum_j\frac{\det M_j(\hbar)}{\det M(\hbar)}w_j\right)=0
 \label{eq:pg-det-cramer-kernel}
\end{equation}
を満たし、これは三つの生成元$x,y,z$との
\emph{全交換子が厳密に零}であることを意味する。
$u_\alpha$は定義から中心なので
$\Omega_{f,\hbar}$は代数全体の中心元である。

$\hbar=0$では括弧の核は$\CC f\oplus R_3$であり、
括弧内の元は$w_0$係数が1、
純中心三次係数が0である。
したがってその特殊化は
$(f-f_{\mathrm{cent}})/a$である。
これが$\Omega_{f,0}=f$を示す。

最後に$\hbar\mapsto-\hbar$と語の反転による
反自己同型（補題\ref{lem:pg-center-parity}）を考える。
Weyl基底の各元は不変であり、
$[\sym_R(p),x_i]/\hbar$のWeyl係数は
$\hbar\mapsto-\hbar$に関して不変。
従って$\mathcal D_{f,\hbar}$、$M,b$と
Cramer公式の係数は$\CC(\hbar^2)$に属する。
全ての操作は次数三・四の有限自由加群で完結し、
形式冪級数も無限次数の中心補正も必要ない。
\end{proof}

\subsection{中心元の規格化と係数空間上の構成}
\label{subsec:pg-det-center-family}
普遍交換子行列の階数がちょうど$r_m$となる開部分を
$\mathcal V_m\subset\mathcal U_m$とする。
補題\ref{lem:pg-det-classical-kernel}により、
$\mathcal V_m$は零切断$\mathcal T_m\times\{0\}$全体を含む。
この上で
\begin{equation}
 \mathcal E_d=S_d\otimes_\CC\OO_{\mathcal V_m},\qquad
 \mathcal C_m=R_3\otimes_\CC\OO_{\mathcal V_m},\qquad
 \rho:\mathcal E_3\longrightarrow\mathcal C_m,
 \quad p\longmapsto p(0,0,0,u)
 \label{eq:pg-det-degree-three-bundles}
\end{equation}
とおき、$\mathcal E_3^{\mathrm{act}}=\ker\rho$とする。
従って$\mathcal E_3=\mathcal C_m\oplus\mathcal E_3^{\mathrm{act}}$である。
$\mathcal A$をこの上の普遍量子代数とし、
以下、Weyl同型によって$\mathcal E_3$と$\mathcal A_3$を同一視する。

\begin{lemma}[割った交換子の核と底変更]
\label{lem:pg-det-divided-kernel}
\begin{equation}
 \mathcal K=\ker\mathcal D,\qquad
 \mathcal L=\ker(\mathcal D|_{\mathcal E_3^{\mathrm{act}}})
 \label{eq:pg-det-kernel-line}
\end{equation}
とおくと、$\mathcal K$は階数$d_m+1$、$\mathcal L$は階数一の
局所自由加群であり、
\begin{equation}
 \mathcal K=\mathcal C_m\oplus\mathcal L,
 \qquad \mathcal K/\mathcal C_m\simeq\mathcal L.
 \label{eq:pg-det-kernel-splitting}
\end{equation}
任意の底変更$\alpha:T\to\mathcal V_m$について、
核・像・余核および\eqref{eq:pg-det-kernel-splitting}の形成は底変更と両立する。
特に
\begin{equation}
 \alpha^*\mathcal K\simeq\ker(\alpha^*\mathcal D),
 \qquad
 \alpha^*\mathcal L\simeq
 \ker(\alpha^*\mathcal D|_{\alpha^*\mathcal E_3^{\mathrm{act}}}).
 \label{eq:pg-det-kernel-base-change}
\end{equation}
\end{lemma}
\begin{proof}
$\mathcal D$は純中心三次部分を消す。
各点の近傍で$r_m$次小行列式を可逆にし、行と列の基本変形を行う。
全ての$(r_m+1)$次小行列式は恒等的に零なので、
行列は単位行列$I_{r_m}$と零行列からなる形となる。
従って核・像・余核は局所自由で、対応する完全列は局所的に分裂する。
この形は非被約な底への変更でも保たれる。
$\mathcal E_3^{\mathrm{act}}$の階数は$r_m+1$であるから、その核は階数一である。
純中心部分への射影$\rho$が分裂を与える。
\end{proof}

\begin{remark}[通常の中心の形成との区別]
\label{rem:pg-det-center-base-change}
この補題が扱うのは$\hbar$で\emph{割った}交換子の核である。
$h_T=\alpha^*\hbar$、$\mathcal D_T=\alpha^*\mathcal D$と書くと、
通常の中心については
\begin{equation}
 \begin{aligned}
 Z(\mathcal A_T)_3&=\ker(h_T\mathcal D_T),\\
 Z(\mathcal A_T)_3/\alpha^*\mathcal K
 &\simeq
 \operatorname{Ann}_{\OO_T}(h_T)
 \otimes_{\OO_T}\operatorname{im}\mathcal D_T.
 \end{aligned}
 \label{eq:pg-det-ordinary-center-base-change}
\end{equation}
実際、$\mathcal D_T$の像への全射は局所的に分裂しており、
中心元の像はちょうど$h_T$で消える部分である。
$h_T$が非零因子なら両方の核は一致する。
しかし$h_T=0$では代数は可換なので通常の中心は$\mathcal E_{3,T}$全体であり、
割った交換子の核は古典Casimir部分
$\mathcal C_{m,T}+\OO_T f$である。
従って\eqref{eq:pg-det-kernel-base-change}を、通常の中心が
任意の底変更と両立するという主張に読み替えてはならない。
\end{remark}

\begin{proposition}[Cramer規格化の一意性と遷移]
\label{prop:pg-det-normalization-transition}
全次数三でactive変数を含む単項式$w_i$を固定し、
その係数写像を$\varepsilon_i$、$f$の$w_i$係数を$\alpha_i$とする。
$\alpha_i$が可逆で、$\varepsilon_i|_{\mathcal L}$が自明化を与える開部分では、
Cramer中心元$\Omega_i$は次の三条件によって一意に定まる：
\begin{equation}
 \mathcal D\Omega_i=0,\qquad
 \rho(\Omega_i)=f_{\mathrm{cent}},\qquad
 \varepsilon_i(\Omega_i-f_{\mathrm{cent}})=\alpha_i.
 \label{eq:pg-det-normalization-conditions}
\end{equation}
特に、同じ$w_i$を用いる限り、可逆な小行列を作る\emph{行の選択}に依存しない。
二つの規格化が有効な共通開部分では
\begin{equation}
 \begin{aligned}
 \Omega_j&=a_{ij}\Omega_i+c_{ij},\\
 a_{ij}&=
 \frac{\alpha_j}{\varepsilon_j(\Omega_i-f_{\mathrm{cent}})}
 \in\OO^\times,\qquad
 c_{ij}=(1-a_{ij})f_{\mathrm{cent}}\in\mathcal C_m.
 \end{aligned}
 \label{eq:pg-det-affine-transition}
\end{equation}
一般に、$\mathcal K/\mathcal C_m$の基底へ写る二つの局所切断
$\Omega,\Omega'$の間には、一意な単元$a$と純中心三次式$c$による
$\Omega'=a\Omega+c$が成り立つ。
両方が法$\hbar$で$f$となるなら、$a|_{\hbar=0}=1$、$c|_{\hbar=0}=0$である。
\end{proposition}
\begin{proof}
$\Omega_i-f_{\mathrm{cent}}$は$\mathcal L$の切断であり、
同線束の自明化$\varepsilon_i$における値を指定すれば一意に決まる。
定理\ref{thm:pg-universal-determinant}のCramer式はこの条件を満たすので、
行の選択を変えても同じ元を与える。
共通開上では$\Omega_i-f_{\mathrm{cent}}$と
$\Omega_j-f_{\mathrm{cent}}$がともに$\mathcal L$の基底である。
その比を$\varepsilon_j$で計算すると\eqref{eq:pg-det-affine-transition}となる。
一般の場合も商線束での比を$a$とすれば、差は$\mathcal C_m$に属する。
法$\hbar$で$\CC f\oplus R_3$が直和であることから最後の主張を得る。
従って同じ古典極限を持つという条件だけでは、$a=1$とは結論できない。
\end{proof}

\begin{theorem}[固定した変数分解の全係数空間における中心元]
\label{thm:pg-det-global-coefficient-lift}
全三次式の係数を$\boldsymbol c$、滑らかさの判別式を$\Delta$とし、
$B_0=\CC[\boldsymbol c,\Delta^{-1}]$とおく。
従って$\mathcal T_m=\operatorname{Spec}B_0$である。
$\tau$を新しい不定元とすると、ある$Q\in B_0[\tau]$、$Q(0)=1$が存在し、
\begin{equation}
 B=B_0[\tau,Q(\tau)^{-1}],\qquad
 B_h=B_0[\hbar,Q(\hbar^2)^{-1}]
 \simeq B[\hbar]/(\hbar^2-\tau)
 \label{eq:pg-det-even-coefficient-rings}
\end{equation}
上で次が成り立つ。
固定したactive変数$x,y,z$と中心変数$u$に関する普遍代数には、
\begin{equation}
 \Omega_{\mathrm{univ}}\in Z(\mathcal A_{B_h})_3,
 \qquad
 \Omega_{\mathrm{univ}}\bmod\hbar=f_{\mathrm{univ}},
 \qquad
 \rho(\Omega_{\mathrm{univ}})=f_{\mathrm{cent}}
 \label{eq:pg-det-global-center}
\end{equation}
を満たす中心元が存在する。
そのWeyl係数は$B$に属し、$\hbar$について偶である。
さらに$\Omega_{\mathrm{univ}}-f_{\mathrm{cent}}$は商線束の大域的な基底となる。
従って全ての滑らかな$f$について、零切断全体を含む\emph{一つの}係数開集合上に、
係数が正則な中心三次元を同時に選べる。
この選択の正準性や座標変換に関する同変性は主張しない。
同変な選択の存在については、第\ref{sec:pg-equivariant-central-lifts}節の
定理\ref{thm:pg-equivariant-levi-center}と
定理\ref{thm:pg-equivariant-parabolic-center}で扱う。
\end{theorem}
\begin{proof}
$\mathcal V_m$の補集合を閉集合として定めるイデアルを
$I\subset B_0[\hbar]$とする。
補集合は零切断と交わらないので$I+(\hbar)=B_0[\hbar]$であり、
$p\in I$、$p(0)=1$を選べる。
$p(\hbar)p(-\hbar)$は偶多項式であるから、
これを$Q(\hbar^2)$と書く。
そのprincipal開は$\mathcal V_m$に含まれ、零切断全体を含む。
必要なら最初に$\hbar\mapsto-\hbar$で移した開部分との共通部分を取っておく。

Weyl基底では$\mathcal D$の全成分が$\hbar$について偶であるから、
$\mathcal D$は$B$上の行列へ降りる。
$B_h=B\oplus\hbar B$は$B$上忠実平坦である。
従って$r_m$次小行列式が生成するイデアルは$B$上でも単位イデアルであり、
全ての$(r_m+1)$次小行列式は零である。
補題\ref{lem:pg-det-divided-kernel}の行列操作を$B$上で行うと、
active部分の核$L_B$は階数一の有限射影$B$-加群となる。
$B$は多項式環の局所化なのでUFDであり、
$\operatorname{Pic}(B)=0$である\cite[Tag~0BCH]{PStacks}。
従って$L_B$の基底$\ell$を選べる。

$Q(0)=1$より$B/(\tau)=B_0$である。
古典Casimir核の計算から、$L_B/(\tau)$の基底は
$f_{\mathrm{univ}}-f_{\mathrm{cent}}$である。
実際、この切断は全ての滑らかな$f$で非零なので線束の基底となる。
従ってある$b\in B_0^\times$について
\begin{equation}
 \ell\bmod\tau=b(f_{\mathrm{univ}}-f_{\mathrm{cent}}).
 \label{eq:pg-det-global-generator-normalization}
\end{equation}
$b^{-1}$を$B$の単元とみて$\ell$を$b^{-1}\ell$へ換え、
$\Omega_{\mathrm{univ}}=f_{\mathrm{cent}}+\ell$とすればよい。
$\mathcal D\ell=0$は全交換子の消滅を意味し、
$B$に属する係数は$\hbar$について偶である。
\end{proof}

\begin{example}[規格化の変更と中心三次イデアル]
\label{ex:pg-det-distinct-normalized-ideals}
$f=(x^3+y^3+z^3)/3+xyz+u^3$を取る。
これは滑らかなHesse三次と$u^3$の和なので全三次超曲面は滑らかである。
$S=x^3+y^3+z^3$、$T=\sym(xyz)$、$\tau=\hbar^2$とし、
\begin{equation}
 U=1+\frac{\tau}{2}+\frac{7\tau^2}{48},\qquad
 W=1-\frac{5\tau}{6}+\frac{7\tau^2}{48},\qquad
 U-W=\frac{4\tau}{3}
 \label{eq:pg-det-hesse-transition-coefficients}
\end{equation}
とおく。定理\ref{thm:pg-universal-hesse-mixed}の$t=1$の式により、
$US/3+WT+u^3$は中心元である。
$U,W$が可逆な原点近傍で、
\begin{equation}
 \Omega_x=u^3+\frac S3+\frac WU T,\qquad
 \Omega_T=u^3+\frac U{3W}S+T
 \label{eq:pg-det-hesse-normalized-centers}
\end{equation}
も中心元となる。
命題\ref{prop:pg-det-normalization-transition}の一意性により、
これらはそれぞれ$x^3$の係数を$1/3$、$xyz$の係数を$1$に固定したCramer規格化である。
両方の古典極限と純中心係数は同じだが、
\begin{equation}
 \Omega_T=\frac UW\Omega_x+
 \left(1-\frac UW\right)u^3
 \label{eq:pg-det-hesse-affine-transition}
\end{equation}
であり、その差は純中心項だけではない。
通常のWeyl基底も有効な共通開へ縮めると、
$\hbar\ne0$では、自然な周囲代数の同一視の下で
$(\Omega_x)$と$(\Omega_T)$は異なるイデアルである。
実際、両イデアルが一致すれば次数三部分を比較して
$\Omega_T=b\Omega_x$となる。
$u^3$係数より$b=1$だが、$S$の係数は$U/W\ne1$なので矛盾する。
これは中心商の抽象的な代数同型の不存在を主張する例ではなく、
affineな遷移だけでは同じ中心イデアルを得ないことを示す。
\end{example}

\begin{remark}[ブロックを保つ座標変換]
\label{rem:pg-det-block-covariance}
以下の比較は、双方の構成が定義される共通開部分上で行う。
定理で選んだprincipal開そのものの座標変換に関する不変性は主張しない。
$M\in\mathrm{GL}_3(\CC)$、$N\in\mathrm{GL}_m(\CC)$、
$C\in\operatorname{Mat}_{3\times m}(\CC)$、$\lambda\in\CC^\times$に対し、
\begin{equation}
 \begin{aligned}
 X&=Mx+Cu,\qquad U=Nu,\\
 f'(X,U)&=\lambda^{-1}
 f\bigl(M^{-1}(X-CN^{-1}U),N^{-1}U\bigr),\qquad
 \hbar'=\lambda\det(M)\hbar
 \end{aligned}
 \label{eq:pg-det-block-coordinate-change}
\end{equation}
とおくと、代入は$A_{f',\hbar'}\simeq A_{f,\hbar}$を与える。
中心変数は可換であり、Weyl対称化はこの線形代入と可換する。
三重ベクトルの変換則から、代入を$\Phi$と書けば
\begin{equation}
 \Phi^{\oplus3}(\mathcal D_{f',\hbar'}p')
 =\frac1{\lambda\det M}\,
 M\mathcal D_{f,\hbar}(\Phi p')
 \label{eq:pg-det-divided-matrix-covariance}
\end{equation}
となる。従って$\mathcal K$、純中心部分、およびその商線束はこの変換と両立する。
Poisson係数を$\eta=\nu\hbar$と分けて記す場合は、
$\nu'=\lambda\det(M)\nu$とすることで共通の変形変数$\hbar$を固定できる。
一方、$f'$へ規格化した中心元を$\Omega'$とすると、
$\lambda\Phi(\Omega')$と$f$へ規格化した$\Omega$の間に保証されるのは
$\lambda\Phi(\Omega')=a\Omega+c$という遷移である。
特に中心平行移動は純中心部分への射影$\rho$を保たないため、
この代数の変換則から中心元の選択の同変性までを結論することはできない。
\end{remark}

\begin{remark}[係数空間上の存在とモジュライへの降下]
\label{rem:pg-det-global-descent-scope}
定理\ref{thm:pg-det-global-coefficient-lift}は、固定した変数分解の係数空間上で
中心元を一括して選べることを示す。
この底では商線束の非自明性を障害とする必要はない。
しかし例\ref{ex:pg-det-distinct-normalized-ideals}の通り、
別々に規格化した中心商が自然な周囲代数の同一視で貼り合わさるとは限らない。

純中心部分との交わりが零で、商線束へ同型に写る中心三次直線を選ぶことは、
完全列$0\to\mathcal C_m\to\mathcal K\to\mathcal L\to0$の分裂を選ぶことと同値である。
その選択空間は$\mathcal Hom(\mathcal L,\mathcal C_m)$の作用するaffine束であり、
中心三次イデアルの選択を記述する。
モジュライへ降下するには、この選択を座標変換に関して同変にし、
零切断で$\CC f$を選ぶ条件および遷移のcocycleと両立させる必要がある。
複数ブロックでは、ブロック分解と純中心三次の配分の変化との両立も必要である。
これらは本定理のfactorialな係数環上での基底選択とは別の条件であり、
第\ref{sec:pg-equivariant-central-lifts}節では変数分解を保つreductive群について、また三変数制限が滑らかな場合には
中心平行移動を含む群について、これらの条件を満たす選択を構成する。
特異な三変数制限を含む全座標群への拡張と、複数ブロックの配分との整合性を含めた
全Poissonモジュライへの大域量子化の降下は主張しない。
\end{remark}

\begin{lemma}[中心三次式の配分による全ブロックの滑らか化]
\label{lem:pg-central-splitting-smooth}
$X=V(F)\subset\PP^{N-1}$が滑らかで、
正規形定理\ref{thm:pg-normal}によって
\begin{equation}
 F=\sum_{s=1}^k\widetilde f_s(x_s,u)+c(u),
 \qquad
 \widetilde f_s(x_s,0)\in\CC[x_s]_3,\quad
 c(u)\in\CC[u]_3
 \label{eq:pg-det-central-splitting}
\end{equation}
と書けるとする。各$\widetilde f_s$は
純中心三次部分を含まないように選ぶ。
このとき中心三次部分を
\begin{equation}
 c(u)=c_1(u)+\cdots+c_k(u),\qquad
 f_s(x_s,u)=\widetilde f_s(x_s,u)+c_s(u)
 \label{eq:pg-det-center-distribution}
\end{equation}
と配分して、\emph{全ての}
$V(f_s)\subset\PP(W_s\oplus U)$が
滑らかになるようにできる。
この配分は$\partial_{x_s}f_s$を変えないため、
元のPoisson構造を一切変えない。
\end{lemma}
\begin{proof}
まず$u=0$の座標平面$\PP(W_s)$を考える。
そこでは$f_s=\widetilde f_s$、
その$x_s$偏微分は$F$の対応偏微分と一致する。
他のactive変数$x_t$の偏微分は、
$x_t=0,u=0$では斉次次数三のため零である。
また$f_{s,u}$はこの平面上で$F_u$と一致し、
$c_s$の偏微分はそこで零である。
従って$\PP(W_s)$上で$f_s$が特異なら、
同じ点で$F$も特異となり矛盾する。
ゆえに$f_s$は全ての$c_s$に対し、
$u=0$上では滑らかである。

次に$m>0$とする。
$u\ne0$の開集合では、中心三次式$\CC[u]_3$の
線形系は基点自由である。
従って$\langle\widetilde f_s,\CC[u]_3\rangle$が定める
射影線形系に標数零のBertini定理を適用すると、
一般の元はこの開集合上で滑らかである。
その非空開条件と$\widetilde f_s$の係数が非零という
非空開条件を同時に満たし、その係数を1に規格化する。
これにより、固定した$\widetilde f_s$へ一般の$c_s(u)$を足す
アフィン線形系についても、
$u\ne0$でも滑らかな$f_s$を与える$c_s$の集合は
非空Zariski開となる。
$k=1$なら$f_1=F$なので既に滑らか。
$k\ge2$では
\begin{equation}
 \mathcal C=
 \{(c_1,\ldots,c_k)\in(\CC[u]_3)^k:
 \sum_sc_s=c\}\simeq(\CC[u]_3)^{k-1}
 \label{eq:pg-det-distribution-space}
\end{equation}
は既約なアフィン空間であり、
各$c_s$への射影は全射である。
従って各$f_s$の滑らかさを課す非空開条件を
$\mathcal C$上で同時に満たせる。
$m=0$では中心三次式はなく、
$F$の滑らかさから各純active$f_s$が既に滑らかである。
最後の主張は$\partial_{x_s}c_s=0$から明らかである。
\end{proof}

\begin{example}[三変数制限が恒等的零でも中心元を得る]
\label{ex:pg-det-zero-ternary}
$m=3$、$u=(u,v,w)$とし
\begin{equation}
 f=u x^2+v y^2+w z^2+u^3+v^3+w^3
 \label{eq:pg-det-zero-ternary-example}
\end{equation}
を考える。$f(x,y,z,0,0,0)=0$である。
しかし$2ux=x^2+3u^2=0$から$x=u=0$
（$y,v$、$z,w$についても同様）なので、
$V(f)\subset\PP^5$は滑らか。
古典交換子行列$S_3\to S_4^{\oplus3}$は
\begin{equation}
 \dim S_3=56,\qquad
 \dim R_3=10,\qquad
 \operatorname{rank}\mathcal D_{f,0}=45
 \label{eq:pg-det-zero-ternary-rank}
\end{equation}
を満たす。これは補題
\ref{lem:pg-det-classical-kernel}の帰結であり、
\path{computations/quantization/verify_singular_cubic_classical_kernel.py}
でも有理数演算で厳密検算できる。
しかもこの例ではCramer式を計算するまでもなく、
\begin{equation}
 [y,z]=2\hbar ux,\quad
 [z,x]=2\hbar vy,\quad
 [x,y]=2\hbar wz,
 \qquad
 \Omega_{f,\hbar}=f
 \label{eq:pg-det-zero-ternary-explicit}
\end{equation}
が\emph{全ての$\hbar$で厳密に中心}である。
実際
$[vy^2+wz^2,x]
=-2\hbar vw(yz+zy)+2\hbar vw(zy+yz)=0$で、
他の生成元との可換性も循環置換で従う。
これは$f(x,y,z,0)=0$であっても
代数的中心元が実際に存在することの直接的な例であり、
零制限の全係数についての同じ現象は
定理\ref{thm:pg-equivariant-zero-center}でPBW対称化から導く。
任意の三変数制限に対する存在は定理\ref{thm:pg-universal-determinant}による。
\end{example}

\begin{remark}[先行研究と確実性]
\label{rem:pg-det-scope}
Hesse型・Fermat型では
定理\ref{thm:pg-universal-hesse-mixed}と
定理\ref{thm:pg-fermat-mixed-universal}の
短い多項式閉公式があり、これらを
普遍交換子行列のgeneric rank恒等式の入力とした。
本定理は、それらの中心性を代数的恒等式として用い、
正規形の各ブロックの全体の滑らかさだけから
\emph{特異制限も含む}代数的中心元を構成する。
中心式は通常は有限行列式の\emph{有理式}であって、
特異な制限の全てに共通な短い多項式閉式が存在するとは
主張しない。全体を滑らかにする配分の選択も一般に非一意。
新規性の全面的な先行研究調査は未完了である。
\end{remark}

% !TeX root = ../../cubic-quantization.tex
\section{全Poisson構造の代数的量子化}
\label{sec:pg-quantization}
ここでは全体の記法に戻り、$X=V(F)\subset\PP(V)$を滑らかな三次超曲面、
$S=\operatorname{Sym}(V^*)$、$N=\dim V=n+2\ge5$とする。
非零構造では正規形\eqref{eq:pg-normal}を固定し、$N=3k+m$とする。
前節までに、一ブロックのPBW性、普遍Hesse中心公式、
および特異・零の三変数制限を含む有限Cramer構成を準備した。
本節では正規形の全ブロックを組み立て、
中心三次元の左右正則性と中心商の平坦性を証明する。
三変数potential代数の非スカラー中心元の存在は
Etingof--Ginzburg\cite[Theorem~3.3.2(ii), p.~1385]{EG}に含まれる。
ここで用いるのは、指定した$f_s$への古典極限と原点での正則性を備えた、
定理\ref{thm:pg-universal-determinant}の相対的な中心元である。

\begin{theorem}[代数的量子化]\label{thm:pg-quantization}\label{pthm:quantization}
任意の Poisson 構造 $\pi_X$ に対し、$0\in U\subset\mathbb A^1_\hbar$、
$K=\OO(U)$ 上の二次関係式で定まる次数付き代数 $A$、および三次の中心元 $\Omega$ が存在して次を満たす。
\begin{enumerate}
\item $A$ は $S$ 上の $q$ の量子化であり、全ての閉点 $\lambda\in U$ で
$A_\lambda$ は Koszul AS-regular、
$\operatorname{gldim}A_\lambda=N$、$\operatorname{Hilb}_{A_\lambda}(t)=(1-t)^{-N}$、AS parameter $N$ である。
\item $\Omega\bmod\hbar=F$ で、各 $\Omega_\lambda$ は左右正則元である。
\item $B=A/(\Omega)$ は $S/(F)$ の量子化であり、各 $B_\lambda$ は AS-Gorenstein で、
\begin{equation}
\operatorname{Hilb}_{B_\lambda}(t)=\frac{1-t^3}{(1-t)^N},\qquad
\operatorname{injdim}B_\lambda=N-1,\qquad
\text{AS parameter}=N-3=n-1.
\end{equation}
\item $U$ を$\hbar\mapsto-\hbar$で不変に取れ、
語の反転と$\hbar$の符号反転による反自己同型に関して
$\Omega^\dagger=\Omega$と選べる。
そのWeyl対称化表示には$\hbar$の奇数次の項が現れない。
\end{enumerate}
\end{theorem}

\subsection{代数的量子化定理の証明}
\begin{proof}[定理\ref{thm:pg-quantization}の証明]
零構造は $A=S[\hbar]$、$\Omega=F$ でよい。
非零構造に対して $A_{\mathrm{reg}}$ を構成する。
PBW 基底により各 fiber で $u_1,\ldots,u_m$ は中心正則列で、
\begin{equation}
A_\lambda/(u_1,\ldots,u_m)\simeq\bigotimes_{s=1}^kS_{s,\lambda}.
\end{equation}
右辺の自明加群の分解、双加群分解およびその双対は各因子のもののテンソル積で得られる。
従って次元 $3k$、AS parameter $3k$ の Koszul AS-regular 代数である。
補題\ref{lem:pg-central-linear}を $m$ 回適用して $A_\lambda$ の次元 $N=3k+m$ と AS parameter $N$ を得る。
Hilbert 級数は PBW 基底から $(1-t)^{-N}$ である。

最初に$F$の純中心三次部分を
補題\ref{lem:pg-central-splitting-smooth}に従って各ブロックへ分配する。
これにより各$V(f_s)\subset\PP(x_s,u)$は滑らかになり、
active偏微分とPoisson構造は変化しない。
定理\ref{thm:pg-universal-determinant}を各$f_s$に適用すると、
\emph{三変数制限$f_s(x_s,0)$が特異または恒等的零でも}、
ある原点を含むZariski開集合上に
有限Cramer式で表せる厳密中心三次元
$\Omega_s\in Z(A_{f_s}|_U)_3$、
$\Omega_s\bmod\hbar=f_s$を得る。
有限個のブロックについて開集合を共通に縮める。
異なるブロックは可換なので $\Omega=\sum_s\Omega_s$ は $A$ の中心元で、residue は $F$。
さらに$U$を$\hbar\mapsto-\hbar$で不変な開集合へ縮め、
補題\ref{lem:pg-center-parity}の反自己同型で各$\Omega_s$を平均する。
これにより$\Omega^\dagger=\Omega$と選べ、
次数三のWeyl対称化係数は$\hbar$に関し偶関数となる。
この正規化は中心性・剰余・正則性の証明を損なわない。

少なくとも一つの $f_s$ は active 変数に依存する。
そうでなければ $F$ は中心変数だけに依存し、active 座標点に特異点をもつ。
その $s$ を固定し、PBW 基底での $\Omega_s$ の active 成分のうち原点で非零の係数が消えないようさらに縮める。
各 $\lambda$ で $\Omega_{s,\lambda}$ の active 次数は正となる。
最高次数が $\lambda$ によらず一定である必要はない。
$\lambda$ を固定し $x_s$ 次数だけで filter すると
\begin{equation}
\operatorname{gr}_sA_\lambda\simeq
(S_{s,\lambda}\otimes_\CC R)\otimes_R A_{\mathrm{rest},\lambda},
\qquad A_{\mathrm{rest},\lambda}=\bigotimes_{t\neq s,R}A_{f_t,\lambda}.
\end{equation}
他の $\Omega_t$ は filtration degree 零なので、$\Omega_\lambda$ の最高記号は
$\Omega_{s,\lambda}$ の非零最高記号 $a_s\in S_{s,\lambda}\otimes R$ である。
$S_{s,\lambda}$ は domain だから多項式拡大も domain で、$a_s$ の左右乗法は単射。
$R$-自由な $A_{\mathrm{rest},\lambda}$ とテンソルしても単射なので、最高次数を比較して
$\Omega_\lambda$ の左右正則性が従う。非可換 domain のテンソル積が domain であるとは仮定していない。

従って各 fiber に
$0\to A_\lambda(-3)\xrightarrow{\Omega_\lambda}A_\lambda\to B_\lambda\to0$
があり、所要の Hilbert 級数を得る。
各次数の余核は PID $\OO(U)$ 上で全 fiber 次元が一定なので有限自由である。
法 $\hbar$ での商は $S\to S/(F)$ で、交換子の一次項もその商の括弧と一致する。
補題\ref{lem:pg-central-cubic}が残りの AS-Gorenstein の主張を与える。
\end{proof}
二次関係式は各生成元対に一つで、反対称成分により独立なので個数は $\binom N2$ である。
\begin{remark}\label{rem:pg-quantization-scope}
定理\ref{thm:pg-quantization}は、個々の Poisson 構造に対する量子化の存在定理である。
代数の同型分類、中心全体の記述、$\qgr B_\lambda$ の滑らかさは含まない。
とくに本定理だけを根拠に全パラメータの Kuznetsov 成分の smooth proper 性や可換復元を結論することはできない。
\end{remark}


% !TeX root = ../../cubic-quantization.tex
\section{任意次元のJacobianブロックに対する多変数量子化}
\label{sec:pg-all-dimensions}
本節は三変数一ブロックの普遍Hesse中心公式を、
\emph{任意の次元}の滑らかな三次超曲面に対する
\emph{複数の量子化パラメータを持つ同時量子化}へ拡張する。
論文全体の正規形定理\ref{thm:pg-normal}を入力とするが、
本節での各中心三次元の構成には
形式的中心元の存在定理を使わない。
最初の多項式公式では各三変数制限の滑らかさを追加仮定する。
後半の定理\ref{thm:pg-ad-all-poisson}ではこれを除き、
全ての滑らかなcubic上のPoisson構造に対して
原点で正則な有理式の代数的中心元を構成する。

\subsection{Jacobianブロックの可換Poisson線形族}
$n\ge3$、$N=n+2=3k+m$（$k\ge1,m\ge0$）とし、
\begin{equation}
 V=\bigoplus_{s=1}^{k}W_s\oplus U,\quad
 \dim W_s=3,\quad\dim U=m,\quad
 F=\sum_{s=1}^k f_s(x_s,u)\in\CC[V]_3.
 \label{eq:pg-ad-normal}
\end{equation}
ここで$F$は滑らかな射影三次超曲面
$X=V(F)\subset\PP^{n+1}$を定め、
\begin{equation}
 g_s(x_s)=f_s(x_s,0)
 \label{eq:pg-ad-smooth}
\end{equation}
は各$s$について滑らかな平面三次曲線を定めると仮定する。
中心三次部分の$f_s$間での分配は任意に固定する。

$P_s=\partial_{x_{s1}}\wedge\partial_{x_{s2}}
\wedge\partial_{x_{s3}}$、
$q_s=\iota_{df_s}P_s$とおく。
$q_s$は$x_s$方向のみにベクトル添字を持ち、
その係数は$x_s,u$だけに依存する。

\begin{proposition}[一般次元の可換Poisson pencil]
\label{prop:pg-ad-compatible}
各$q_s$は$F$をCasimirとするPoisson双ベクトルであり、
\begin{equation}
 [q_s,q_t]_{\mathrm{Sch}}=0\qquad(s,t=1,\dots,k).
\end{equation}
よって任意の$\boldsymbol\lambda=(\lambda_1,\ldots,\lambda_k)$に対し
\begin{equation}
 q_{\boldsymbol\lambda}=\sum_s\lambda_s q_s,\qquad
 \pi_{\boldsymbol\lambda}=q_{\boldsymbol\lambda}|_X
\end{equation}
はPoisson構造である。
各$q_s|_X$は一次独立であり、
$\pi_{\boldsymbol\lambda}$の一般Poisson階数は
$2\,\#\{s:\lambda_s\ne0\}$となる。
したがって$\boldsymbol\lambda$は$k$次元の線形Poisson族を与える。
\end{proposition}
\begin{proof}
三変数Jacobian括弧のJacobi恒等式は直接成立する。
$F$の$x_s$偏微分は$f_s$のそれに等しいので
$q_s(F,-)=0$。
$s\ne t$なら、$q_s$のベクトル添字は$x_s$方向、
$q_t$の係数は$x_t,u$だけに依存する。
したがって$q_s$の導分は$q_t$の係数を消し、
逆向きも同様なのでSchouten括弧は零である。

$g_s$の滑らかさにより各$q_s$は非零であり、
異なるブロックの$\bigwedge^2W_s$成分に属するから一次独立である。
一般点で各非零係数のブロックは階数二を持ち、
直和なのでアフィン錐上の一般階数は$2r$である
（$r=\#\{s:\lambda_s\ne0\}$）。
射影化で階数が落ちるのはEuler方向が像に入るときだけである。
$m>0$なら$u\ne0$の稠密開集合上でEuler方向は中心成分を持ち、
像はactive方向のみなので階数は落ちない。
$m=0$の場合、$r<k$なら非零のinactive座標成分から同様に落ちず、
$r=k$ならEulerが像に入るためには
$3g_s(x_s)=\sum_i x_{si}g_{s,x_{si}}=0$が全$s$で必要となる。
$X$が$\bigcap_sV(g_s)$に含まれることはない
（$k\ge2$のとき各ブロックは非零三次式である）。
よって一般点では階数は落ちない。
\end{proof}

\subsection{任意階数に対する独立パラメータ量子化}
$\boldsymbol\eta=(\eta_1,\ldots,\eta_k)$を次数零の独立な変数とし、
$R=\CC[u_1,\ldots,u_m]$とおく。
$R[\boldsymbol\eta]$上の代数$\mathcal A_{\boldsymbol\eta}$を
全active生成元$x_{si}$、中心生成元$u_\alpha$、
および関係式
\begin{equation}
\begin{aligned}
 [x_{s2},x_{s3}]
 &=\eta_s\operatorname{sym}_R\bigl((f_s)_{x_{s1}}\bigr),\\
 [x_{s3},x_{s1}]
 &=\eta_s\operatorname{sym}_R\bigl((f_s)_{x_{s2}}\bigr),\\
 [x_{s1},x_{s2}]
 &=\eta_s\operatorname{sym}_R\bigl((f_s)_{x_{s3}}\bigr),\\
 [x_{si},x_{tj}]&=0\ (s\ne t),\qquad [u_\alpha,x_{si}]=0
\end{aligned}
\label{eq:pg-ad-relations}
\end{equation}
で定める。すなわち
\begin{equation}
 \mathcal A_{\boldsymbol\eta}
 =\mathcal A_{f_1,\eta_1}
 \otimes_{R[\boldsymbol\eta]}\cdots
 \otimes_{R[\boldsymbol\eta]}\mathcal A_{f_k,\eta_k}.
 \label{eq:pg-ad-tensor}
\end{equation}
ここでは各ブロックの係数に\emph{同じ中心変数$u$}を使う。
積のテンソルは中心係数環上で取り、
中心変数をブロックごとに複製しない。

\begin{theorem}[任意次元の明示的な同時代数量子化]
\label{thm:pg-all-dimensions}
\eqref{eq:pg-ad-normal}--\eqref{eq:pg-ad-smooth}を仮定する。
各$f_s$に対し、
系\ref{cor:pg-smooth-ternary-explicit}の
有限多項式の中心三次元$\Omega_{s,\eta_s}$を構成する。
すると
\begin{equation}
 \boxed{\qquad
 \Omega_{\boldsymbol\eta}=\sum_{s=1}^k\Omega_{s,\eta_s},
 \qquad
 \Omega_{\boldsymbol0}=F
 \qquad}
 \label{eq:pg-ad-universal-center}
\end{equation}
は$\mathcal A_{\boldsymbol\eta}$の\emph{厳密な中心元}である。
また原点を含むZariski開集合
$\boldsymbol0\in\mathcal U\subset\mathbb A^k_{\boldsymbol\eta}$を選び、
以下を全て満たせる。
\begin{enumerate}[label=(\arabic*)]
\item 各$\boldsymbol a\in\mathcal U$について
$A_{\boldsymbol a}=\mathcal A_{\boldsymbol\eta}
 \otimes_{\CC[\boldsymbol\eta]}\CC_{\boldsymbol a}$
は$N$生成元、$\binom N2$個の二次関係式による
Koszul AS-regular代数であり、
\begin{equation}
 \operatorname{gldim}A_{\boldsymbol a}=N,\quad
 \operatorname{Hilb}_{A_{\boldsymbol a}}(z)=(1-z)^{-N},
 \quad\text{AS parameter}=N.
\end{equation}
各内部次数成分は$\OO(\mathcal U)$上局所自由である。
\item $\Omega_{\boldsymbol a}$は全てのファイバーで
左右正則元となる。中心三次商
$B_{\boldsymbol a}=A_{\boldsymbol a}/(\Omega_{\boldsymbol a})$
はAS-Gorensteinで
\begin{equation}
 \operatorname{injdim}B_{\boldsymbol a}=N-1,\quad
 \text{AS parameter}=N-3=n-1,\quad
 \operatorname{Hilb}_{B_{\boldsymbol a}}(z)
 =\frac{1-z^3}{(1-z)^N}.
\end{equation}
$\mathcal B=\mathcal A/(\Omega)$の各内部次数成分も
$\OO(\mathcal U)$上局所自由であり、
$\mathcal B_{\boldsymbol0}=\CC[V]/(F)$である。
\item 原点での独立な一次変形方向は
命題\ref{prop:pg-ad-compatible}の
$k$個のPoisson構造$q_s|_X$に対応する。
任意の$\boldsymbol\lambda\in\CC^k$について
$\eta_s=\lambda_s\hbar$と特殊化すれば、
$\pi_{\boldsymbol\lambda}$の
\emph{一変数代数的量子化}が得られる。
その中心三次元も式\eqref{eq:pg-ad-universal-center}で明示される。
\end{enumerate}
Hesse座標を各ブロックごとに選ぶと、
$\Omega_{\boldsymbol\eta}$は各$\eta_s$に関して
次数高々四の\emph{有限多項式}であり、
異なるパラメータ$\eta_s,\eta_t$を同時に含む
新たな補正項は現れない。
したがって$3k+m=n+2$が任意であっても、
この構成は形式的中心元の存在定理を必要としない。
\end{theorem}
\begin{proof}
各$f_s$の純active三次$g_s$が滑らかなので、
Hesse標準形への線形変換と中心変数の平行移動により、
系\ref{cor:pg-smooth-ternary-explicit}を適用できる。
その公式は係数を数値特殊化する必要のない
$\eta_s^2$の高々二次の多項式であり、
$\Omega_{s,0}=f_s$、各active生成元との交換子は
まず$\eta_s=0$の開近傍で零となる。
ここでは$g_s$が滑らかなので、
注意\ref{rem:pg-regular-all}により全ての$\eta_s$で
$S_{s,\eta_s}$はAS-regularである。
従って命題\ref{prop:pg-PBW}の族に関する証明を
$U_{\mathrm{reg}}=\mathbb A^1$に適用でき、
$\mathcal A_{f_s,\eta_s}$は$R[\eta_s]$上自由となる。
その局所化への写像は単射なので、
多項式である$\Omega_{s,\eta_s}$の交換子は
全ての$\eta_s$について恒等的に零である。
他ブロックの生成元は$x_s$と可換であるから
$\Omega_{s,\eta_s}$は$\mathcal A_{\boldsymbol\eta}$全体で中心となり、
和も中心である。
これで式\eqref{eq:pg-ad-universal-center}を得る。
特に\emph{異なるブロック間の量子的混合補正は必要ない}。

命題\ref{prop:pg-PBW}を各$s$と各$\eta_s$について適用する。
独立なパラメータを持つので、原点を含む各ブロックの
Zariski開集合の積を取ればよい。
$R[\boldsymbol\eta]$-自由なPBW基底を各因子で選び、
\eqref{eq:pg-ad-tensor}でテンソルすると、
全体の各内部次数は底上局所自由で、
$u_1,\ldots,u_m$は各ファイバーで中心正則列である。
その商は
\begin{equation}
 A_{\boldsymbol a}/(u_1,\ldots,u_m)
 \simeq\bigotimes_{s=1}^k
  S_{s,a_s},\qquad
 S_{s,a_s}=\mathcal A_{f_s,a_s}/(u).
\end{equation}
各$S_{s,a_s}$は三次元Koszul AS-regular domain。
それらのKoszul分解と双対のテンソル積から、
右辺は次元$3k$、AS parameter $3k$の
Koszul AS-regular代数となる。
補題\ref{lem:pg-central-linear}を
$m$回適用し、$A_{\boldsymbol a}$の
次元$N=3k+m$のKoszul AS-regular性を得る。
PBW基底がHilbert級数と二次関係数を与える。

$\Omega$の左右正則性には
定理\ref{thm:pg-quantization}の証明と同じ
\emph{一ブロックactive次数}のfiltrationを使う。
第一ブロック$x_1$の次数を数え、
他のactive変数と$u$のfiltration次数を零とする。
$\operatorname{gr}_{x_1}A_{\boldsymbol a}
\simeq
 (S_{1,a_1}\otimes_\CC R)
 \otimes_R A_{\mathrm{rest},\boldsymbol a}
 \simeq S_{1,a_1}\otimes_\CC A_{\mathrm{rest},\boldsymbol a}$
が得られる。
$\Omega_{1,a_1}$のactive次数三の最高記号は、
$\boldsymbol a=\boldsymbol0$で$g_1\ne0$なので、
必要なら開集合をさらに縮めることにより
$S_{1,a_1}$の非零元となる。
他の$\Omega_{s,a_s}$はactive次数零なので、
$\Omega_{\boldsymbol a}$の最高記号は
$\Omega_{1,a_1}$の非零active三次記号である。
$S_{1,a_1}$はdomainであり、
$A_{\mathrm{rest},\boldsymbol a}$は$R$-自由なので、
最高記号の左右乗法は単射。
ゆえに$\Omega_{\boldsymbol a}$は左右正則元である。
補題\ref{lem:pg-central-cubic}で商のAS-Gorenstein性を得る。

各内部次数の$\mathcal A$は底上有限局所自由で、
$\Omega_{\boldsymbol a}$の正則性により
\begin{equation}
 0\longrightarrow(A_{\boldsymbol a})_{d-3}
 \xrightarrow{\cdot\Omega_{\boldsymbol a}}
 (A_{\boldsymbol a})_d
 \longrightarrow(B_{\boldsymbol a})_d
 \longrightarrow0
\end{equation}
のrankは全$\boldsymbol a\in\mathcal U$で一定である。
有限表示行列の定階数条件より、
$\mathcal B_d$も$\OO(\mathcal U)$上局所自由である。
よって多変数底上の平坦性を得る。

最後に定義関係式の$\eta_s$線形項と
$\Omega_{\boldsymbol\eta}\bmod\boldsymbol\eta=F$から、
原点での一次Poisson括弧が$q_s$であることが分かる。
$\eta_s=\lambda_s\hbar$による底変更が
全ての方向を同時に実現する。
\end{proof}

\begin{theorem}[全てのPoissonブロックの多変数同時量子化]
\label{thm:pg-ad-all-poisson}
$n\ge3$、$X=V(F)\subset\PP^{n+1}$を
\emph{滑らかな}三次超曲面とし、
定理\ref{thm:pg-normal}の正規形
$F=\sum_{s=1}^k f_s(x_s,u)$を取る。
ここでは$f_s(x_s,0)$の滑らかさを\emph{仮定しない}。
純中心三次の配分を
補題\ref{lem:pg-central-splitting-smooth}により選び直し、
各$V(f_s)\subset\PP(x_s,u)$を滑らかにする。
定理\ref{thm:pg-universal-determinant}のCramer式により
\begin{equation}
 \Omega_{s,\eta_s}\in
 Z(A_{f_s,\eta_s})_3,\qquad
 \Omega_{s,0}=f_s
 \label{eq:pg-ad-general-block-center}
\end{equation}
を、それぞれ$\eta_s=0$を含む適切な開集合上に構成できる。

従って$\eta_1,\ldots,\eta_k$を独立に動かして
\begin{equation}
 \boxed{\qquad
 \Omega_{\boldsymbol\eta}=\sum_{s=1}^k\Omega_{s,\eta_s},
 \qquad
 \Omega_{\boldsymbol0}=F
 \qquad}
 \label{eq:pg-ad-general-center}
\end{equation}
を厳密な代数的中心三次元とする
$k$パラメータの次数付き量子化が得られる。
原点を含むZariski開集合上で、
周囲の各ファイバーは次元$N=n+2$のKoszul AS-regular、
中心三次商はAS-Gorensteinかつflatで、
Hilbert級数はそれぞれ
\begin{equation}
 \operatorname{Hilb}_{A_{\boldsymbol\eta}}(z)=(1-z)^{-N},
 \qquad
 \operatorname{Hilb}_{B_{\boldsymbol\eta}}(z)
 =(1-z^3)(1-z)^{-N}
 \label{eq:pg-ad-general-hilbert}
\end{equation}
となる。
$\eta_s=\lambda_s\hbar$とする底変更によって、
固定した$X$の当該$k$次元Poisson線形族の
\emph{全方向}を零切断の近傍で同時に量子化できる。

一般の特異三変数制限に対する中心元は、
$\eta_s^2$の\emph{有理式}である。
Hesse型の滑らかな制限の場合のように、
すべての係数が$\eta_s^4$で打ち切られる
多項式表示を主張するものではない。
この定理の証明に形式的量子化は必要ない。
\end{theorem}
\begin{proof}
各ブロックの全三次式を滑らかにする
中心三次項の配分が存在することは
補題\ref{lem:pg-central-splitting-smooth}により証明済み。
配分はactive偏微分を変えないので、
対応する三つの二次関係式もPoisson一次括弧も変化しない。
各$f_s$について定理
\ref{thm:pg-universal-determinant}を適用し、
原点で正則な有限行列式の比で中心三次元を得る。
異なるブロックのactive生成元は互いに可換、
中心変数$u$は共通であるので、和は厳密に中心である。

命題\ref{prop:pg-PBW}は三変数制限の滑らかさを
仮定せずに原点近傍で成立する。
各ブロックのPBW自由性を共通中心環上でテンソルし、
中央一次正則列を使えば、
全体のKoszul AS正則性と次元$N$を得る。
中心三次元の左右正則性は、
定理\ref{thm:pg-quantization}の証明と同様、
active次数のfiltrationの最高記号における
三次元AS-regular domainの乗法単射性による。
各ブロックの$f_s$はactive変数に依存するから
その最高active記号は原点で非零で、
必要に応じて開を縮めても非零である。
中心元による完全列と中央超曲面の
change of ringsからAS-Gorenstein性を得る。

各内部次数に関するPBW基底と、
中心元による乗法行列のfiberwise rank一定性から
多変数底上の局所自由性・flatnessも得られる。
最後の主張は$\boldsymbol\eta=\hbar\boldsymbol\lambda$
での底変更と一次交換子の計算による。
\end{proof}

\begin{corollary}[Poisson線形族の零切断近傍での一様量子化]
\label{cor:pg-ad-whole-pencil}
定理\ref{thm:pg-all-dimensions}の$k$次元Poisson線形族
$L=\langle\pi_1,\ldots,\pi_k\rangle$に対して、
\begin{equation}
 \mathscr U=
 \{(\boldsymbol\lambda,\hbar)\in
 \mathbb A^k_{\boldsymbol\lambda}\times\mathbb A^1_\hbar:
 (\lambda_1\hbar,\ldots,\lambda_k\hbar)\in\mathcal U\}
\end{equation}
とおく。
$\mathscr U$は\emph{全ての}
$(\boldsymbol\lambda,0)$を含むZariski開集合であり、
$\mathcal A_{\boldsymbol\eta},\Omega_{\boldsymbol\eta}$
を$\eta_s=\lambda_s\hbar$で引き戻した族は、
全線形空間$L$の各Poisson構造に対する
一変数量子化を一つの代数的族として与える。
量子化の中心三次元は
\begin{equation}
 \Omega_{\boldsymbol\lambda,\hbar}
 =\sum_{s=1}^k\Omega_{f_s,\lambda_s\hbar}
\end{equation}
であり、全$\boldsymbol\lambda$について
$\hbar=0$で$F$となる。
$\lambda_s=0$を含む全ての部分線形空間にも
同じ式が制限される。
\end{corollary}
\begin{proof}
写像
$(\boldsymbol\lambda,\hbar)\mapsto
(\lambda_1\hbar,\ldots,\lambda_k\hbar)$は代数的で、
その$0$切断は全て$\boldsymbol\eta=0$へ写る。
従って$\mathcal U$の逆像は指定の開集合であり、
定理\ref{thm:pg-all-dimensions}の局所自由性・正則性・
中心性を底変更によって保つ。
原点での一次交換子は$\sum_s\lambda_s q_s$なので
$\pi_{\boldsymbol\lambda}$の量子化となる。
\end{proof}

\begin{proposition}[正規形の係数も動く相対普遍公式]
\label{prop:pg-ad-relative-chart}
$k,m$とブロック座標を固定する。
各$f_s$の純active三次式をHesse形
$g_{t_s}(x_s)$に正規化した上で、
中心変数を含む全ての斉次三次混合係数と
$t_1,\ldots,t_k$を独立に動かす
\emph{正規形の表示パラメータ空間}$\mathcal T_{k,m}$を考える。
このとき\eqref{eq:pg-ad-universal-center}は
\begin{equation}
 \Omega\in
 \CC[\mathcal T_{k,m}][\boldsymbol\eta,u,x_1,\ldots,x_k]_3
\end{equation}
の\emph{非可換Weyl対称化として読める有限多項式}であり、
$\eta_s$の次数は高々四である。
この表示パラメータ空間上の、$X$と全$g_{t_s}$が滑らかな開部分
$\mathcal T_{k,m}^{\mathrm{sm}}$について、
その零切断を含むZariski開近傍では
$\Omega$は中心となり、
定理\ref{thm:pg-all-dimensions}の
PBW/AS正則性と中心商の平坦性を同時に満たす。

これは固定した$F$に対する量子化を超え、
正規形の係数$F$と独立量子化方向$\boldsymbol\eta$を
\emph{同時に動かす相対的な代数的構成}である。
ただし$\mathcal T_{k,m}$は
座標とブロック分解を指定した\emph{表示空間}であり、
一般の非可換三次超曲面のモジュライ商ではない。
またHesse座標の全係数空間上の正準選択や
異なる表示チャートの大域的な中心元の接着は主張しない。
\end{proposition}
\begin{proof}
定理\ref{thm:pg-universal-hesse-mixed}の中心元は
$t_s$、混合係数、$\eta_s$についての
普遍的な多項式恒等式である。
中心平行移動による混合二次項の消去も
式\eqref{eq:pg-hm-translation}の多項式演算なので、
全係数を変数としたまま
$\Omega=\sum_s\Omega_{s,\eta_s}$を作れる。
中心性の多項式恒等式は
各表示チャートの$\boldsymbol\eta=0$近傍で成立する。

AS正則性は、定理\ref{thm:pg-all-dimensions}の
各係数値に対して$\boldsymbol\eta=0$を含む
開条件として得られる。
その開条件を表示係数との積空間で課せば、
$\mathcal T_{k,m}^{\mathrm{sm}}$の零切断を含む開近傍を得る。
共通中心変数上のPBW基底と
中心正則性のfiltration証明は係数の底変更と両立し、
各内部次数の商の有限定階数条件から
相対的な局所自由性・平坦性が従う。
\end{proof}

\begin{corollary}[一般次元での全方向の有限Hessian公式]
\label{cor:pg-ad-hessian}
各ブロックのHesse正規化と中心平行移動を済ませ、
その際の量子化変数を$\eta_s$と書く。
$f_s$の混合係数を$A_s,B_s,C_s,D_s,E_s,H_s,G_s$、
Hesse係数を$t_s$とする。
\begin{equation}
 {\mathscr H}_s=
 \det\left(\frac{\partial^2 f_s}
 {\partial x_{si}\partial x_{sj}}\right)_{i,j=1}^3.
\end{equation}
系\ref{cor:pg-hesse-hessian-form}の
有限多項式$\Psi_{t_s}(f_s)$を用いると、
\begin{equation}
 \boxed{\quad
 \Omega_{\boldsymbol\eta}=
 \operatorname{sym}\left(
 F-\sum_{s=1}^k\frac{\eta_s^2}{12}{\mathscr H}_s
 +\sum_{s=1}^k\frac{\eta_s^2}{6}A_sB_sC_s
 +\sum_{s=1}^k\eta_s^4\Psi_{t_s}(f_s)
 \right).
 \quad}
 \label{eq:pg-ad-hessian-center}
\end{equation}
これは$\boldsymbol\eta$の多項式であり、
全次数で厳密な中心元である。
特に$\eta_s=\lambda_s\hbar$としたとき
$\hbar^2$のHessian補正は$\lambda_s^2$、
$\hbar^4$補正は$\lambda_s^4$を係数とする。
異なるブロックの混合Hessian項
$\eta_s^2\eta_t^2$は生じない。
一般の$g_s$へはHesse正規化時の
$\eta_s$の行列式スケール変換を施してから適用する。
\end{corollary}
\begin{proof}
系\ref{cor:pg-hesse-hessian-form}の式を各ブロックに用いて
和を取れば得られる。
各補正は$x_s,u$にのみ依存し、
他ブロックの生成元と可換なため、
高次の交叉補正を必要としない。
\end{proof}

\begin{proposition}[次元に依存しない内在的Hessian二次補正]
\label{prop:pg-ad-intrinsic-hessian}
一般の$N=n+2$次元ベクトル空間$V$、
定数三重ベクトル$P\in\bigwedge^3V$、
斉次三次式$F\in\operatorname{Sym}^3V^*$に対し、
座標を一時的に選んで
\begin{equation}
 \mathscr H(P,F)=
 \frac16\sum_{i,j,k,a,b,c}
 P^{ijk}P^{abc}\,
 F_{ia}F_{jb}F_{kc}\in\operatorname{Sym}^3V^*
 \label{eq:pg-ad-intrinsic-hessian}
\end{equation}
と定義する。これは座標に依存しない多項式である。
\eqref{eq:pg-ad-normal}の正規形と
$P_{\boldsymbol\lambda}=\sum_s\lambda_sP_s$については
\begin{equation}
 \mathscr H(P_{\boldsymbol\lambda},F)
 =\sum_{s=1}^k\lambda_s^2
  \det\operatorname{Hess}_{x_s}(f_s).
 \label{eq:pg-ad-intrinsic-blocks}
\end{equation}
各$g_s=f_s(x_s,0)$が滑らかなら、
定理\ref{thm:pg-all-dimensions}の
$\eta_s=\lambda_s\hbar$による厳密な中心三次元を
純中心三次元$\hbar^2c_{\boldsymbol\lambda}(u)$だけ
調整して
\begin{equation}
 \boxed{\quad
 \widetilde\Omega_{\boldsymbol\lambda,\hbar}
 =\operatorname{sym}\left(
 F-\frac{\hbar^2}{12}\mathscr H(P_{\boldsymbol\lambda},F)
 +\hbar^4K_{\boldsymbol\lambda,F}
 \right)
 \quad}
 \label{eq:pg-ad-intrinsic-center}
\end{equation}
と書ける。$K_{\boldsymbol\lambda,F}$は
全次数三の明示可能な多項式であり、
Hesse座標と中心平行移動の選択後には
$\sum_s\lambda_s^4\Psi_{t_s}(f_s)$の座標変換で与えられる。
従って$\hbar^2$補正は任意の次元に対し
\emph{ブロック分解によらない不変な式}で指定できる。
高次の$K_{\boldsymbol\lambda,F}$まで
座標選択に依存しない単一の不変多項式として表したとは主張しない。
\end{proposition}
\begin{proof}
式\eqref{eq:pg-ad-intrinsic-hessian}は
$P\otimes P$と三つのHessianテンソルの縮約なので
座標変換に不変である。
三次元で$P=\partial_x\wedge\partial_y\wedge\partial_z$なら
Levi--Civita記号の恒等式から
$\mathscr H(P,F)=\det\operatorname{Hess}_{x,y,z}(F)$。
正規形では異なるactiveブロック間の
Hessian成分$F_{x_{si}x_{tj}}$（$s\ne t$）は零なので、
$P_{\boldsymbol\lambda}\otimes P_{\boldsymbol\lambda}$の交叉項は
縮約で消え、式\eqref{eq:pg-ad-intrinsic-blocks}を得る。

系\ref{cor:pg-hesse-hessian-form}の一ブロック中心元は
$\hbar^2$係数が
$-\det\operatorname{Hess}(f_s)/12$と
純中心三次$A_sB_sC_s/6$の和であり、
$\hbar^4$項までで厳密に中心となる。
中心平行移動および線形Hesse変換では
Hessian縮約の変換則が
量子化パラメータの行列式スケール変換と一致する。
純中心項は変換後も$\CC[u]_3$に属するので、
全ブロックの和から$\hbar^2$倍の純中心項を引けばよい。
残る$\hbar^4$項を$K_{\boldsymbol\lambda,F}$とする。
$\CC[u]$は量子化代数の中心にあるので、この調整は
厳密な中心性を変えない。
\end{proof}

\begin{example}[共通中心変数を持つ六次元cubicと独立な二方向]
\label{ex:pg-ad-sixfold}
$n=6$、$N=8$、$k=2$、$m=2$とし、
ブロック座標を$(x,y,z)$、$(a,b,c)$、
中心座標を$(u,v)$とする。
\begin{equation}
 F=
 g_{1/2}(x+u,y+v,z+u+v)
 +g_{-1/3}(a+v,b+u,c)
 +\frac{u^3+v^3}{3}.
\end{equation}
$\eta_1,\eta_2$を独立な量子化パラメータとして、
一ブロックHesse中心元を
$\Omega_{g_t,\eta}=\frac{U(\eta,t)}3(X^3+Y^3+Z^3)
 +W(\eta,t)\sym(XYZ)$
（定理\ref{thm:explicit-hesse-center}）と書く。
このとき
\begin{equation}
\begin{split}
\Omega_{\eta_1,\eta_2}={}&
 \Omega_{g_{1/2},\eta_1}(x+u,y+v,z+u+v)\\
 &+\Omega_{g_{-1/3},\eta_2}(a+v,b+u,c)
 +\frac{u^3+v^3}{3}
\end{split}
\end{equation}
は厳密な中心元である。
各$\eta_s$について高々四次の多項式であり、
$\eta_1^i\eta_2^j$（$i,j>0$）の項は現れない。

$g_{1/2}$と$g_{-1/3}$はともに滑らかなHesse三次である。
与えた座標の中心平行移動は可逆なので、
$F$は二つの滑らかなHesse三次と
$u^3+v^3$の直和に変換できる。
全偏微分の共通零点がアフィン原点だけなので$X$は滑らか。
$\eta_s=\lambda_s\hbar$の下で
$(\lambda_1,\lambda_2)$の全線形族を同時に量子化し、
両係数が非零なら一般Poisson階数は四である。
\end{example}

\begin{remark}[次元と独立パラメータ数]
\label{rem:pg-ad-dimension-table}
許されるブロック数は$k\le\lfloor(n+2)/3\rfloor$。
例えば
\begin{equation}
\begin{array}{c|ccccc}
 n&3&4&6&7&10\\ \hline
 \max k&1&2&2&3&4
\end{array}
\end{equation}
となる。
同じ$n$でも$k$の値によって異なるPoisson階数が現れ、
その正規形ごとに独立な$k$個の量子化パラメータを取れる。
この表はPoissonモジュライの既約成分の個数を表すものではない。
\end{remark}

\begin{remark}[全次元・全成分についての正確な範囲]
\label{rem:pg-ad-scope}
$k$は$1\le k\le\lfloor(n+2)/3\rfloor$を動ける。
この上限は正規形における\emph{許される最大ブロック数}であり、
全Poissonモジュライの成分数や族の次元と同一視しない。
固定したブロック分解における三次式$F$の線形パラメータ空間では、
滑らかな$X$と全ての滑らかな$g_s$という条件は
非空なZariski開条件である。
実際$\sum_{s,i}x_{si}^3+\sum_\alpha u_\alpha^3$
はその一例である。
従って短い多項式公式を与える定理\ref{thm:pg-all-dimensions}は、
任意次元の各正規形型の
\emph{非空な一般部分}に適用される。
滑らかな$X$であっても一部の$g_s$が特異または零となる場合は、
定理\ref{thm:pg-ad-all-poisson}が扱う。
その場合にも全てのPoisson構造の同時代数量子化が得られるが、
中心元は一般に原点で正則な有理式のCramer表示となり、
各パラメータについて次数高々四の多項式とは主張しない。

固定した$F$と固定したブロック分解に対する
$k$変数量子化を構成したのであり、
相対Poissonモジュライの\emph{全既約成分}の
普遍的量子化空間を構成したという意味ではない。
またこの量子化から$\operatorname{qgr}B$の
幾何学的smoothnessや可換Kuznetsov復元は従わない。
Hesse標準形を全係数空間にわたって
正準的に同時選択できるとも主張しない。
\end{remark}

% !TeX root = ../../cubic-quantization.tex
\section{中心三次元の同変選択と中心商族}
\label{sec:pg-equivariant-central-lifts}
本節では単一のJacobianブロックを扱う。
定階数核の構成を同変に行うことで、固定したactive・中心変数の分解を保つ
還元的群について、古典極限を指定した中心元を同時に選べる。
さらに三変数制限が滑らかな部分、および節点・尖点の各軌道上では、
二次混合項の規格化によって中心変数による平行移動も処理できる。
三変数制限が恒等的零の場合には、PBW対称化そのものが
全パラメータ上で全ブロック群に共変な中心元を与える。
平均化そのものは既知のReynolds作用素の性質であり、
ここで用いるのはその定階数核への適用である
\cite[\S2, p.~369]{PDerksenInvariants}。

$W=\CC^3$、$Z=\CC^m$とし、$x=(x,y,z)$と$u$をそれぞれの座標とする。
変形パラメータを一旦$t$と書き、
\eqref{eq:pg-det-quantum}の$\hbar$を$t$に置き換えた代数を$A_{f,t}$とする。
群とその作用を
\begin{equation}
 \begin{aligned}
 G_0&=(\mathrm{GL}(W)\times\mathrm{GL}(Z))\times\mathbf G_m,\\
 G&=\left\{
 \begin{pmatrix}M&C\\0&N\end{pmatrix}
 :M\in\mathrm{GL}(W),\ N\in\mathrm{GL}(Z),\ C\in\operatorname{Hom}(Z,W)
 \right\}\times\mathbf G_m,\\
 X&=Mx+Cu,\qquad U=Nu,\\
 f'(X,U)&=\lambda^{-1}
 f\bigl(M^{-1}(X-CN^{-1}U),N^{-1}U\bigr),\qquad
 t'=\lambda\det(M)t
 \end{aligned}
 \label{eq:pg-equivariant-groups}
\end{equation}
で定める。$G_0\subset G$は$C=0$の部分群である。
代入による同型を$\Phi:A_{f',t'}\to A_{f,t}$とする。
注意\ref{rem:pg-det-block-covariance}により、Weyl同型、割った交換子の核
$\mathcal K$、純中心部分$\mathcal C_m$および商線束
$\mathcal K/\mathcal C_m$は$G$と両立する。
純中心部分への射影$\rho(p)=p(0,u)$は$G_0$と両立する。

\subsection{不変開と中心元の同変持上げ}

\begin{lemma}[アフィン底上の同変持上げ]
\label{lem:pg-equivariant-affine-lifting}
$H$を複素還元的線形代数群、$Y$をアフィン$H$-スキームとする。
次が成り立つ。
\begin{enumerate}
\item $Z_0\subset Y$を$H$-安定閉部分、$V\subset Y$を$Z_0$を含む
$H$-安定開部分とする。
ある$H$-不変関数$q\in\Gamma(Y,\OO_Y)$について、
$q|_{Z_0}=1$かつ$D(q)\subset V$とできる。
\item $\mathcal E$を$Y$上の同変ベクトル束とする。
$\mathcal E|_{Z_0}$の不変切断は$\mathcal E$の不変切断へ持ち上がる。
任意の指標で線形化をtwistした場合も同じである。
\item $Y$上の同変ベクトル束の短完全列は同変に分裂する。
\end{enumerate}
\end{lemma}
\begin{proof}
アフィン底上の同変準連接層の大域切断は有理$H$-加群であり、
各元の軌道の線形包は有限次元である。
その有限次元部分において、Zariski稠密な最大compact部分群で平均すれば
Reynolds射影が得られる。
これは$H$-安定部分加群を保ち、同変線形写像と可換する。

第一の主張では、$Y\setminus V$のイデアルを$I$、$Z_0$のイデアルを$J$とすると
$I+J=\Gamma(Y,\OO_Y)$である。
$p\in I$、$p\equiv1\pmod J$を選び、平均して$q$とすればよい。
第二の主張では、アフィン性による全射
\begin{equation}
 \Gamma(Y,\mathcal E)\longrightarrow
 \Gamma(Z_0,\mathcal E|_{Z_0})
 \label{eq:pg-equivariant-restriction-surjection}
\end{equation}
において任意の持上げを平均する。
第三の主張では、商ベクトル束が有限射影加群に対応するので、
まず$\OO_Y$-線形な分裂を選べる。
Hom加群におけるその平均も$\OO_Y$-線形であり、
合成が恒等写像であるという条件を保つ。
\end{proof}

\begin{theorem}[零制限における全パラメータの同変中心元]
\label{thm:pg-equivariant-zero-center}
\begin{equation}
 \mathcal T_m^{\mathrm z}
 =\{f\in\mathcal T_m:f(x,0)=0\}
 \label{eq:pg-equivariant-zero-base}
\end{equation}
とおく。この係数空間と\emph{全ての}$t\in\mathbb A^1$の積上で、
$A_{f,t}$は$R[t]$、$R=\CC[u]$上の順序付き単項式
$x_1^{a_1}x_2^{a_2}x_3^{a_3}$をPBW基底に持つ。
係数も動かした場合には$\Gamma(\mathcal T_m^{\mathrm z},\OO)[u,t]$上で同じ基底が成立する。
特にactive次数の濾過について
\begin{equation}
 \operatorname{gr}_x A_{f,t}\simeq R[t][x_1,x_2,x_3]
 \label{eq:pg-equivariant-zero-pbw}
\end{equation}
であり、Weyl対称化$\sym_R$は全次数で加群の同型となる。

中心三次元を
\begin{equation}
 \Omega^{\mathrm z}_{f,t}=\sym_R(f)
 \label{eq:pg-equivariant-zero-center}
\end{equation}
と取れる。この元のWeyl係数は$t$に依存せず、
古典極限は$f$、純中心部分は$f_{\mathrm{cent}}$である。
全$G$に関して
\begin{equation}
 \Phi(\Omega^{\mathrm z}_{f',t'})=\lambda^{-1}\Omega^{\mathrm z}_{f,t}
 \label{eq:pg-equivariant-zero-covariance}
\end{equation}
が成立する。割った交換子の次数三の核をこの底でも$\mathcal K$と書けば、
\begin{equation}
 \mathcal K=\OO\Omega^{\mathrm z}\oplus\mathcal C_m
 \label{eq:pg-equivariant-zero-kernel}
\end{equation}
である。いずれにも零切断近傍への縮小は不要である。
\end{theorem}
\begin{proof}
$g=0$という条件から、$x=(x_1,x_2,x_3)^{\mathsf t}$を用いて
\begin{equation}
 f=\tfrac12x^{\mathsf t}H(u)x+b(u)^{\mathsf t}x+c(u),
 \quad H^{\mathsf t}=H,\quad
 \deg_u H_{ij}=1,\quad \deg_u b_i=2,\quad\deg_u c=3
 \label{eq:pg-equivariant-zero-quadratic-form}
\end{equation}
と書ける（各成分は零でもよい）。
従ってJacobian括弧は$R$上で$x$についてaffine-linearである。
中心記号$e_0$を加え、$R[t]$上の自由加群
$\mathfrak l=R[t]e_0\oplus\bigoplus_{i=1}^3R[t]e_i$に
\begin{equation}
 [e_i,e_j]_{\mathfrak l}
 =t\sum_{k=1}^3\varepsilon_{ijk}
       \left(\sum_{\ell=1}^3 H_{k\ell}(u)e_\ell+b_k(u)e_0\right),
 \qquad [e_0,e_i]_{\mathfrak l}=0
 \label{eq:pg-equivariant-zero-lie-algebra}
\end{equation}
を定める。ここで$\varepsilon_{123}=1$である。
Jacobian括弧のJacobi恒等式を一次項と定数項に分ければ、
これがLie括弧であることが分かる。
その包絡代数について
\begin{equation}
 A_{f,t}\simeq U_{R[t]}(\mathfrak l)/(e_0-1)
 \label{eq:pg-equivariant-zero-enveloping}
\end{equation}
となる。
自由基底を持つLie代数のPBW論証はこの係数環上でも成立する。
実際、逆順の二文字を交換子で置き換えるmonicな規則は、
語長と逆転数について停止し、三文字の重なりの差はJacobi恒等式で消える。
$e_0$を含む重なりはその中心性で消える。
従ってDiamond lemma\cite[Theorem~1.2]{BergNE12}により、
$U(\mathfrak l)$は$R[t,e_0]$上で順序付き$e_i$単項式の自由加群である。
$e_0=1$と置いてもこの基底は保たれる。
この議論は三次式の係数を記号として扱う場合にも同じであり、
PBW性と\eqref{eq:pg-equivariant-zero-pbw}を得る。

対称化が随伴作用と可換するという既知のPBWの原理
\cite[Corollaries~13.7--13.8, p.~75]{PEtingofLie}を、ここでは直接計算できる。
任意の$p\in R[t][x_1,x_2,x_3]$について
\begin{equation}
 [x_i,\sym_R(p)]=t\,\sym_R(\{x_i,p\}_f)
 \label{eq:pg-equivariant-zero-exact-commutator}
\end{equation}
である。実際、単項式の全順序の平均を取り、交換子のLeibniz則を適用する。
各因子を$\{x_i,x_j\}_f$で置き換えた和が現れ、この式が$x$の一次式と
中心定数項だけからなるので、その和は右辺の対称化に一致する。
$p=f$に適用すると$\{x_i,f\}_f=0$なので、中心性が従う。
Weyl係数に補正を加えていないため、古典極限、偶性、純中心係数も所要の通りである。
ここまでのPBW性と中心性には$f$の滑らかさを用いておらず、
任意の可換$\CC$-係数環上でactive次数が二以下の式に同じ証明が適用できる。

\eqref{eq:pg-equivariant-zero-exact-commutator}は、Weyl表示での
割った交換子が全$t$で古典作用素と等しいことも示す。
補題\ref{lem:pg-det-classical-kernel}により各$f$での三次核は
$\CC f\oplus\CC[u]_3$である。
係数底$\mathcal T_m^{\mathrm z}$は被約であり、この古典行列は定階数であるから、
可逆小行列を用いた局所分裂により\eqref{eq:pg-equivariant-zero-kernel}を得る。
この核の形成は任意の底変更と両立する。
ただし$t=0$の通常の中心は次数三全体へ増えるので、
通常の中心と割った交換子の核は区別する。

最後に$g=0$の条件は$G$-安定である。
中心変数を含む線形代入はWeyl対称化と可換するため、
\eqref{eq:pg-equivariant-groups}から
\eqref{eq:pg-equivariant-zero-covariance}が直接従う。
\end{proof}

\begin{theorem}[還元的部分群に関する中心元の同変選択]
\label{thm:pg-equivariant-levi-center}
$\mathcal T_m$を全三次式$f$が滑らかな射影超曲面を定める係数空間とする。
$G_0$-不変な偶関数$Q(f,t)$、$Q(f,0)=1$を用いたprincipal開
\begin{equation}
 \mathcal W_m^0=D(Q)\subset\mathcal T_m\times\mathbb A^1_t
 \label{eq:pg-equivariant-levi-open}
\end{equation}
上で、Weyl係数が正則な中心元$\Omega^0_{f,t}$を、
\begin{equation}
 \begin{aligned}
 \mathcal D\Omega^0_{f,t}&=0,&
 \Omega^0_{f,0}&=f,&
 \rho(\Omega^0_{f,t})&=f_{\mathrm{cent}},\\
 \Omega^0_{f,-t}&=\Omega^0_{f,t},&
 \Phi(\Omega^0_{f',t'})&=\lambda^{-1}\Omega^0_{f,t}\qquad(C=0)
 \end{aligned}
 \label{eq:pg-equivariant-levi-center}
\end{equation}
となるように選べる。
ここで偶性はWeyl係数についての等式である。
さらに$\Omega^0_{f,t}$の$\mathcal K/\mathcal C_m$への像は
この線束の基底となる。
零制限の全パラメータの族とも
\begin{equation}
 Q|_{\mathcal T_m^{\mathrm z}\times\mathbb A^1}=1,
 \qquad
 \Omega^0|_{\mathcal T_m^{\mathrm z}\times\mathbb A^1}
 =\Omega^{\mathrm z}=\sym(f)
 \label{eq:pg-equivariant-levi-zero-compatibility}
\end{equation}
となるように整合させられる。
この選択は一般に一意ではない。
\end{theorem}
\begin{proof}
命題\ref{prop:pg-PBW}と
補題\ref{lem:pg-det-divided-kernel}で用いた開集合を、
三変数斉次代数のAS-regular条件、次数三・四のWeyl同型、
および割った交換子の階数$r_m$という条件で定める。
これらの条件は座標変換と両立するので、得られる開集合$\mathcal V_m$は
$G$-安定であり、$\mathcal T_m\times\{0\}$を含む。
$g=0$では定理\ref{thm:pg-equivariant-zero-center}により
PBW性、Weyl同型と核の定階数性が全$t$で成立するので、
$\mathcal T_m^{\mathrm z}\times\mathbb A^1$も$\mathcal V_m$に含まれる。
$t\mapsto-t$は代数を反対代数へ移し、これらの条件も保つ。

$\mathcal T_m$は判別式を可逆にしたアフィン空間である。
補題\ref{lem:pg-equivariant-affine-lifting}(1)を
$H=G_0\times\boldsymbol\mu_2$と、その安定な被約閉部分
\begin{equation}
 Z_0=(\mathcal T_m\times\{0\})\cup
     (\mathcal T_m^{\mathrm z}\times\mathbb A^1)
 \label{eq:pg-equivariant-levi-prescribed-locus}
\end{equation}
に適用すると、
$\mathcal V_m$内に、不変かつ偶な$Q$による
\eqref{eq:pg-equivariant-levi-open}を$Q|_{Z_0}=1$として取れる。
$t$は$G_0$の指標$\lambda\det(M)$で変換するが、
$I_g$を$g$の十係数のイデアルとすると、
$I(Z_0)=(t)\cap I_g=tI_g$は$H$-安定である。
従って$Q|_{Z_0}=1$という両方の規格化が平均で保たれる。

このアフィン開上で$\mathcal K$は同変ベクトル束である。
Weyl表示における$\mathcal D$は$t$について偶なので、
$\boldsymbol\mu_2$はWeyl記号を変えずに$\mathcal K$へ作用する。
補題\ref{lem:pg-det-classical-kernel}と
定理\ref{thm:pg-equivariant-zero-center}より、
$Z_0$上でWeyl記号$f$は$\mathcal K|_{Z_0}$の切断である。
実際、被約な$Z_0$上で$\mathcal Df$の各成分は全閉点で消える。
引戻しの線形化を$\lambda\Phi$へtwistしてこの切断を不変とみなし、
補題\ref{lem:pg-equivariant-affine-lifting}(2)で持ち上げる。
得られた切断を$\widetilde\Omega$とし、
\begin{equation}
 \Omega^0=\widetilde\Omega-\rho(\widetilde\Omega)+f_{\mathrm{cent}}
 \label{eq:pg-equivariant-pure-center-adjustment}
\end{equation}
と置く。
$\rho$は$G_0$および$\boldsymbol\mu_2$と両立し、
純中心三次式は$\mathcal D$で消えるので、
\eqref{eq:pg-equivariant-levi-center}を得る。
この補正は$Z_0$上のWeyl記号$f$を変えない。

商線束における$\Omega^0$の像は$Z_0$全体で消えない。
その零点集合を除いた安定開に同補題(1)を再適用する。
必要なら不変な分母を払えば、元のアフィン底でも
$Q|_{Z_0}=1$を満たす不変な偶関数のprincipal開として取れる。
この開では商線束への像が基底となる。
同時に\eqref{eq:pg-equivariant-levi-zero-compatibility}が保たれる。
\end{proof}

\begin{lemma}[節点・尖点軌道上の方向微分の同変分裂]
\label{lem:pg-equivariant-singular-retractions}
$L=\mathrm{GL}(W)\times\mathbf G_m$を、
$(M,\lambda)\cdot g=\lambda^{-1}g(M^{-1}x)$によって作用させる。
次の二つの被約軌道を考える：
\begin{equation}
 \begin{aligned}
 g_{\mathrm{nod}}&=x^3+y^3+xyz,&
 \mathcal O_{\mathrm{nod}}&=L\cdot g_{\mathrm{nod}},\\
 g_{\mathrm{cusp}}&=y^2z-x^3,&
 \mathcal O_{\mathrm{cusp}}&=L\cdot g_{\mathrm{cusp}}.
 \end{aligned}
 \label{eq:pg-equivariant-singular-orbits}
\end{equation}
それぞれは既約な節点平面三次式、尖点平面三次式の軌道で、
滑らかかつ既約、次元はそれぞれ$9,8$である。
各軌道上で$\delta_g(v)=\partial_vg$は単射であり、
正則な$L$-同変左逆$\sigma_g:\operatorname{Sym}^2W^*\to W$を持つ。
二次式への作用は$\rho_2(M,\lambda)q=\lambda^{-1}q(M^{-1}x)$とする。
\end{lemma}
\begin{proof}
$S_{xy}$を$x,y$の交換行列、$\zeta$を原始三乗根とし、
$\Gamma=\langle\operatorname{diag}(\zeta,\zeta^{-1},1),S_{xy}\rangle$とおく。
基準形の安定化群は
\begin{equation}
 \begin{aligned}
 H_{\mathrm{nod}}&=\{(aT,a^{-3}):a\in\mathbf G_m,\ T\in\Gamma\},\\
 H_{\mathrm{cusp}}&=\{(\operatorname{diag}(a,b,a^3b^{-2}),a^{-3}):
                      a,b\in\mathbf G_m\}
 \end{aligned}
 \label{eq:pg-equivariant-singular-stabilizers}
\end{equation}
である。実際、両方の唯一の特異点は$[0:0:1]$である。
節点ではその二つの接線を保つので、必要なら$x,y$を交換した後、
$M(x,y,z)=(ax,by,ex+fy+cz)$と書ける。
$g_{\mathrm{nod}}(Mx)$が$g_{\mathrm{nod}}(x)$のスカラー倍となる条件は、
$e=f=0$、$a^3=b^3=abc$であり、上の群を得る。
尖点では接線$y=0$も保つので、
$M(x,y,z)=(ax+by,dy,ex+fy+hz)$と書ける。
$y^2z-x^3$の係数を比較すると$b=e=f=0$、$d^2h=a^3$となる。
これが第二の群である。

二次単項式$r$の係数を$[r]q$と書けば、基準形での左逆は
\begin{equation}
 \begin{aligned}
 \sigma_{\mathrm{nod}}(q)
   &=\bigl([x^2]q/3,\ [y^2]q/3,\ [xy]q\bigr),\\
 \sigma_{\mathrm{cusp}}(q)
   &=\bigl(-[x^2]q/3,\ [yz]q/2,\ [y^2]q\bigr)
 \end{aligned}
 \label{eq:pg-equivariant-singular-retractions}
\end{equation}
で与えられる。
それぞれの勾配$(3x^2+yz,3y^2+xz,xy)$、$(-3x^2,2yz,y^2)$に
適用すると$\sigma_\gamma\delta_{g_\gamma}=\mathrm{id}_W$である。
\eqref{eq:pg-equivariant-singular-stabilizers}の対角行列の重みと
$x,y$の交換を調べれば、$\sigma_\gamma$は$H_\gamma$-同変である。

従って$\ell=(M,\lambda)\in L$に対して
\begin{equation}
 \sigma_{\ell\cdot g_\gamma}
   =M\sigma_\gamma\rho_2(\ell)^{-1}
 \label{eq:pg-equivariant-singular-transport}
\end{equation}
と置くと、右$H_\gamma$作用に不変な正則写像となる。
軌道写像$L\to L/H_\gamma\simeq\mathcal O_\gamma$は忠実平坦であり、
この写像は軌道上の正則な左逆へ降りる
\cite[Propositions~7.12, 7.17, pp.~142--143]{PMilneGroups}。
これは各軌道上の構成であり、軌道のアフィン性を仮定しない。
安定化群の次元が$1,2$なので、軌道の次元は$10-1,10-2$である。

標準形の範囲も直接確認できる。
既約な特異平面三次の特異点を$[0:0:1]$へ移すと、式は
$zq_2(x,y)+q_3(x,y)$となる。
$q_2=0$なら二変数の三次式となって可約なので、$q_2\ne0$である。
節点では$q_2=xy$とでき、$z$の線形置換で$x^2y,xy^2$を消せる。
残る$x^3,y^3$の係数は既約性により非零なので、スカラー変換で節点の基準形を得る。
尖点では$q_2=y^2$とでき、$x^3$の非零係数を使った$x$の置換で$x^2y$を消し、
$z$の置換で残る$xy^2,y^3$を消せば尖点の基準形となる。
\end{proof}

\begin{theorem}[各三変数制限層上の全ブロック群に関する選択]
\label{thm:pg-equivariant-parabolic-center}
\begin{equation}
 \begin{aligned}
 \mathcal T_m^{\mathrm{ell}}
 &=\{f\in\mathcal T_m:g(x)=f(x,0)\text{ が滑らかな平面三次を定める}\},\\
 \mathcal T_m^\gamma
 &=\{f\in\mathcal T_m:f(x,0)\in\mathcal O_\gamma\}
       \qquad(\gamma=\mathrm{nod},\mathrm{cusp})
 \end{aligned}
 \label{eq:pg-equivariant-elliptic-base}
\end{equation}
とおく。
各$\gamma=\mathrm{ell},\mathrm{nod},\mathrm{cusp}$について、
$\mathcal T_m^\gamma\times\mathbb A^1_t$には、零切断全体を含む
$G$-不変なprincipal開$\mathcal W_m^G$と、そこに正則な中心元
$\Omega^G_{f,t}$が存在し、
\begin{equation}
 \mathcal D\Omega^G_{f,t}=0,\qquad
 \Omega^G_{f,0}=f,\qquad
 \Omega^G_{f,-t}=\Omega^G_{f,t},\qquad
 \Phi(\Omega^G_{f',t'})=\lambda^{-1}\Omega^G_{f,t}
 \label{eq:pg-equivariant-parabolic-center}
\end{equation}
を満たす。
商線束$\mathcal K/\mathcal C_m$への像も基底となる。
純中心係数を$f_{\mathrm{cent}}$に固定する条件は課さない。
選択は三つの層でそれぞれ行い、型の境界をまたぐ整合は結論しない。
\end{theorem}
\begin{proof}
滑らかな三変数三次式$g$のアフィン係数空間を$\mathcal T_0$とする。
方向微分の写像
\begin{equation}
 \delta_g:W\longrightarrow\operatorname{Sym}^2W^*,
 \qquad v\longmapsto\partial_vg
 \label{eq:pg-equivariant-gradient-embedding}
\end{equation}
は各点で単射である。
実際、非零方向で$\partial_vg=0$なら、$g$はその方向の変数に依存せず、
その座標点が平面三次の特異点となる。
従って余核もベクトル束である。
二次形式側の作用に$\lambda^{-1}$を含めれば
\eqref{eq:pg-equivariant-gradient-embedding}は
$\mathrm{GL}(W)\times\mathbf G_m$-同変である。
滑らかな制限上では補題\ref{lem:pg-equivariant-affine-lifting}(3)によって、
正則な同変左逆$\sigma_g$を一つ選ぶ：
\begin{equation}
 \sigma_g:\operatorname{Sym}^2W^*\longrightarrow W,
 \qquad \sigma_g\delta_g=\mathrm{id}_W.
 \label{eq:pg-equivariant-gradient-retraction}
\end{equation}
節点・尖点の各軌道では、補題\ref{lem:pg-equivariant-singular-retractions}の
左逆を用いる。以下の構成はこの三場合で共通である。

$f$のactive次数二の部分を$q_f(x,u)=\sum_\alpha q_\alpha(x)u_\alpha$と書き、
\begin{equation}
 B_f(u)=\sum_\alpha\sigma_g(q_\alpha)u_\alpha,
 \qquad
 f^0(X,u)=f(X-B_f(u),u)
 \label{eq:pg-equivariant-centering}
\end{equation}
と定める。
全ての係数は対応する$\mathcal T_m^\gamma$上で正則であり、
$f^0$の二次混合部分には$\sigma_g$を適用すると零になる。
$G$の作用下での変換を計算すると、
\begin{equation}
 B_{f'}=(MB_f-C)N^{-1},\qquad
 X+B_{f'}U=M(x+B_fu),\qquad
 (f')^0(X,U)=\lambda^{-1}f^0(M^{-1}X,N^{-1}U)
 \label{eq:pg-equivariant-centering-covariance}
\end{equation}
となる。
第一の等式では、$g$を$x\mapsto x-Cu$で移すと
二次混合項が$-\delta_g(Cu)$だけ変わることと
$\sigma_g\delta_g=\mathrm{id}$を用いた。
従って$(f,t)\mapsto(f^0,t)$は中心平行移動に関して不変で、
残る群$G_0$について同変である。

定理\ref{thm:pg-equivariant-levi-center}の開と中心元を用いて
\begin{equation}
 \begin{aligned}
 \mathcal W_m^G
 &=\{(f,t):(f^0,t)\in\mathcal W_m^0\},\\
 \Omega^G_{f,t}(x,u)
 &=\Omega^0_{f^0,t}(x+B_f(u),u)
 \end{aligned}
 \label{eq:pg-equivariant-centered-quantization}
\end{equation}
と置く。
開集合は不変関数$Q(f^0,t)$で定まり、$Q(f^0,0)=1$である。
中心変数による線形代入はWeyl対称化と可換し、代数の同型を与えるので、
中心性、偶性および商線束の生成性が保たれる。
\eqref{eq:pg-equivariant-centering-covariance}と
\eqref{eq:pg-equivariant-levi-center}から全$G$に関する共変性を得る。
古典極限は$f^0(x+B_fu,u)=f(x,u)$である。
\end{proof}

\begin{remark}[節点型の既知の明示中心元との関係]
\label{rem:pg-cmr-nodal-comparison}
Chekhov--Mazzocco--Rubtsov~\cite{PCMR}は、一般化Etingof--Ginzburg代数の
$\gamma=0$の場合に中心三次元の明示式を与えている。
以下で用いるのは、著者がaccepted versionと明記した
arXiv:1905.02772v2のLemma~4.6、式~(4.39)、p.~22であり、
対応する関係式は式~(4.38)、p.~21にある。
同版Table~5、p.~41には、古典式$xyz+x^3+y^3$に
対角二次項・一般一次項・定数項を加えた場合との対応も記されている。
従って節点型における中心三次元の存在と、この基準座標での明示式は既知結果に帰属する。

まず、この既知の式は節点型の一般の下位係数にも適用できる。
$g=x^3+y^3+xyz$に対して
\begin{equation}
 \widetilde\sigma_g(q)
   =\bigl([yz]q,[xz]q,[xy]q\bigr),
 \qquad
 \delta_g(v)=v_1(3x^2+yz)+v_2(3y^2+xz)+v_3xy
 \label{eq:pg-cmr-alternative-centering}
\end{equation}
と置けば$\widetilde\sigma_g\delta_g=\mathrm{id}$である。
この左逆も節点型の安定化群に同変であり、
核$\langle x^2,y^2,z^2\rangle$は安定化群で保たれる。
一般の二次項を
\begin{equation}
 q=q_{xx}x^2+q_{xy}xy+q_{xz}xz+q_{yy}y^2+q_{yz}yz+q_{zz}z^2
 \label{eq:pg-cmr-general-quadratic}
\end{equation}
と書くと、$B=\widetilde\sigma_g(q)$による$f^0(X)=f(X-B)$の二次項は
\begin{equation}
 (q_{xx}-3q_{yz})X^2+(q_{yy}-3q_{xz})Y^2+q_{zz}Z^2
 \label{eq:pg-cmr-diagonal-quadratic}
\end{equation}
となる。一次項・定数項もこの同じ三次式全体の並進で変換する。
従って、以下では独立な中心係数を用いて
\begin{equation}
 f^0=x^3+y^3+xyz+A_1x^2+A_2y^2+A_3z^2
       +\ell_1x+\ell_2y+\ell_3z+\kappa
 \label{eq:pg-cmr-diagonal-potential}
\end{equation}
を扱えばよい。斉次三次式$f(x,u)$への適用では、$A_i$を$u$の一次式、
$\ell_i$を二次式、$\kappa$を三次式とする。

$p=1+t/2$、$s=1-t/2$と置き、$sp$を可逆にする。
本稿の対称Jacobian関係式
\begin{equation}
 [x,y]=t\,\operatorname{sym}(f^0_z),\qquad
 [y,z]=t\,\operatorname{sym}(f^0_x),\qquad
 [z,x]=t\,\operatorname{sym}(f^0_y)
 \label{eq:pg-cmr-jacobian-relations}
\end{equation}
は、CMRの式~(4.38)に
\begin{equation}
 \begin{gathered}
 q_{\mathrm{CMR}}=p/s,\qquad
 \alpha=\beta=3t/s,\qquad \gamma=0,\\
 (a_1,b_1,c_1)=-2t(A_1,A_2,A_3)/s,\qquad
 (a_2,b_2,c_2)=-t(\ell_1,\ell_2,\ell_3)/s
 \end{gathered}
 \label{eq:pg-cmr-parameter-comparison}
\end{equation}
を代入したものと一致する。
実際、CMRの三関係式をそれぞれ$s$倍すると
\begin{equation}
 \begin{gathered}
 sxy-pyx-2tA_3z-t\ell_3=0,\\
 syz-pzy-3tx^2-2tA_1x-t\ell_1=0,\\
 szx-pxz-3ty^2-2tA_2y-t\ell_2=0
 \end{gathered}
 \label{eq:pg-cmr-comparison-relations}
\end{equation}
を得る。

CMRの式~(4.39)の中心元を$\Omega_{\mathrm{CMR}}$とする。
上の代入後、その全係数は$t$で割れ、
$\Omega_{\mathrm{CMR}}/(6t)$は$t=0$で正則である。
また$s^5$を掛ければ係数の分母を除ける。従って
\begin{equation}
 Q=s^5\frac{\Omega_{\mathrm{CMR}}}{6t}+\kappa
 \label{eq:pg-cmr-normalized-center}
\end{equation}
は$t,A_i,\ell_i,\kappa$の多項式を係数とする語の有限和であり、
$Q\bmod t=f^0$である。
ここで$t$で割る操作は係数ごとの可除性によるものであり、
最終的な係数環で$t$を可逆にしているわけではない。
中心性はCMRの公式から従い、$t=0$への延長はPBWによる$t$-torsionの消滅から従う。

この規格化を安定化群に正確に共変にする純中心補正も明示できる。
$x,y$の交換に伴って
\begin{equation}
 A_1\leftrightarrow A_2,\qquad
 \ell_1\leftrightarrow\ell_2,\qquad t\longmapsto-t
 \label{eq:pg-cmr-swap-action}
\end{equation}
と変換する。最後の符号は$t'=\lambda\det(M)t$による。
変換した式を元の代数に引き戻したものを$Q^{\mathrm{sw}}$と書く。
CMRの公式に上のパラメータを代入し、$\nu=1+t^2/12$と置くと
\begin{equation}
 \begin{aligned}
 Q^{\mathrm{sw}}-Q
 &=\nu z(sxy-pyx)+A_3\nu(s^2-p^2)z^2
      +6A_3t^2(pyx-sxy)\\
 &\quad+\bigl((s-p)\nu\ell_3+6(s+p)A_3^2t^3\bigr)z.
 \end{aligned}
 \label{eq:pg-cmr-swap-computation}
\end{equation}
\eqref{eq:pg-cmr-comparison-relations}の第一式と$s+p=2$、$s-p=-t$を用いて
\begin{equation}
 Q^{\mathrm{sw}}-Q=-6A_3\ell_3t^3
 \label{eq:pg-cmr-swap-discrepancy}
\end{equation}
を得る。従って
\begin{equation}
 \widehat Q=Q-3A_3\ell_3t^3
 \label{eq:pg-cmr-equivariant-center}
\end{equation}
はこの交換に正確に共変である。
他の生成子$rI$と$\operatorname{diag}(\zeta,\zeta^{-1},1)$、$\zeta^3=1$については、
各単項式がそれぞれ次数重み$3$、$\mu_3$重み$0$を持つことから共変性が従う。
こうしてCMRの既知の公式から、節点基準形上の安定化群に共変な中心元を得る。

一般の下位係数への復帰は$\widehat Q_{f^0,t}(x+B)$で与えられる。
中心化の変換則と軌道上への正則な移送を組み合わせれば、本稿で用いる全ブロック群との
整合性へ接続する。この接続、相対的な係数選択、およびPoisson層への降下の主張は、
CMRの点ごとの明示式と区別して述べる。
多項式表示\eqref{eq:pg-cmr-equivariant-center}だけから、$t=\pm2$でのPBW性・平坦性や
型の境界を越える同変な貼合せは結論しない。
\end{remark}

\begin{remark}[層ごとの選択と貼合せの限界]
\label{rem:pg-equivariant-stratumwise-centering}
上の三つの層を合わせた不可約三次式の開全体には、
$g$だけに依存する一つの正則な左逆$\sigma_g$を選べない。
実際、可約三次式の射影的な閉集合は、一次式と二次式の積による
$\PP^2\times\PP^5$の像であり、一般の分解は一意なので次元は$7$である。
従って$\mathbb A^{10}$内の可約式の錐の余次元は$2$である。
その補集合上の正則関数は$\mathbb A^{10}$へ延長するので、
左逆の行列成分が存在すれば$\sigma_g\delta_g=\mathrm{id}_W$が
全空間上の多項式恒等式となり、$g=0$で矛盾する。
これはこのcenteringの方法の限界であり、混合係数も用いる別の中心元選択や、
中心直線そのものの貼合せの不存在を意味しない。
\end{remark}

\subsection{中心商の族と規格化の条件}

\begin{corollary}[同変な正則中心商族]
\label{cor:pg-equivariant-central-quotient}
上の三定理のそれぞれの底と群について、
$\mathcal J=\OO\Omega\subset\mathcal A_3$は同変な中心直線部分束である。
各幾何学的ファイバーで$\Omega$は左右正則であり、
中心商$\mathcal B=\mathcal A/(\mathcal J)$の各次数部分は底上局所自由である。
そのファイバーは
\begin{equation}
 \operatorname{Hilb}_{\mathcal B}(z)
 =\frac{1-z^3}{(1-z)^{m+3}}
 \label{eq:pg-equivariant-quotient-hilbert}
\end{equation}
を持ち、左右injective dimension $m+2$、AS parameter $m$の
AS-Gorenstein代数である。

変形パラメータを群作用から独立に固定するには、
\begin{equation}
 t=\eta\hbar,\qquad
 \eta'=\lambda\det(M)\eta,\qquad \hbar'=\hbar
 \label{eq:pg-equivariant-fixed-parameter}
\end{equation}
とおく。
定理\ref{thm:pg-equivariant-levi-center}、\ref{thm:pg-equivariant-parabolic-center}から得る引戻しは、それぞれ
\begin{equation}
 \left[(\mathcal T_m\times\mathbb A^1_\eta)/G_0\right]
       \times\mathbb A^1_\hbar,
 \qquad
 \left[(\mathcal T_m^\gamma\times\mathbb A^1_\eta)/G\right]
       \times\mathbb A^1_\hbar
 \label{eq:pg-equivariant-quotient-stacks}
\end{equation}
の、$\hbar=0$全体を含む開部分上に次数付き代数の族として降りる。
第二の式では$\gamma=\mathrm{ell},\mathrm{nod},\mathrm{cusp}$をそれぞれ取る。
零制限の場合は縮小を要せず、
\begin{equation}
 \left[(\mathcal T_m^{\mathrm z}\times\mathbb A^1_\eta)/G\right]
       \times\mathbb A^1_\hbar
 \label{eq:pg-equivariant-zero-quotient-stack}
\end{equation}
全体に同じ結論が成立する。
特殊ファイバーは$\CC[x,u]/(f)$であり、一次括弧は
$\eta\{-,-\}_f$である。
\end{corollary}
\begin{proof}
共変性は$\mathcal J$およびその生成する両側イデアルを保つ。
命題\ref{prop:pg-PBW}により、各ファイバーで
\begin{equation}
 \operatorname{gr}_x A_{f,t}
 \simeq S_{g,t}\otimes_\CC\CC[u_1,\ldots,u_m]
 \label{eq:pg-equivariant-fiber-domain}
\end{equation}
はdomainであるから、$A_{f,t}$もdomainである。
商線束を生成する$\Omega$は零ではなく、従って左右正則である。
各次数の乗法写像
\begin{equation}
 \mathcal J\otimes\mathcal A_{d-3}\longrightarrow\mathcal A_d
 \label{eq:pg-equivariant-regular-multiplication}
\end{equation}
は全ファイバーで単射なので、ベクトル束の写像として局所的に分裂する。
その余核$\mathcal B_d$は局所自由であり、Hilbert級数も従う。
周囲代数$A_{f,t}$において$u_1,\ldots,u_m$はPBW性による中心正則列で、
その商$S_{g,t}$は三次元Koszul AS-regularである。
零制限の場合には$S_{0,t}=\CC[x_1,x_2,x_3]$であり、
定理\ref{thm:pg-equivariant-zero-center}により全$t$で同じ議論が成立する。
補題\ref{lem:pg-central-linear}を$m$回適用し、
さらに補題\ref{lem:pg-central-cubic}を適用するとAS条件を得る。
中心商$\mathcal B$における$u_1,\ldots,u_m$の正則性を用いているのではない。
特に零制限の場合には$\Omega^{\mathrm z}\bmod(u)=0$である。

\eqref{eq:pg-equivariant-fixed-parameter}は群作用と両立する。
元の開は$t=0$全体を含むので、引戻した開は$\hbar=0$全体を含む。
同変代数と同変イデアルの降下により
\eqref{eq:pg-equivariant-quotient-stacks}の開部分、または
\eqref{eq:pg-equivariant-zero-quotient-stack}全体の上の族が得られる。
その平坦性はアフィンatlas上で示した次数別局所自由性から従う。
定義関係式を$\hbar$で割って特殊化すれば一次括弧は
$\eta\{-,-\}_f$となる。
\end{proof}

\begin{remark}[純中心部分がない単一ブロック]
\label{rem:pg-equivariant-line-and-gerbe}
$m=0$では$\mathcal C_m=0$なので、定階数核$\mathcal K$自身が
正準な同変中心直線である。
従って同変な中心イデアル$(\mathcal K)$とその商には、
中心元の基底やReynolds平均の選択は不要である。
これは単一ブロックについての主張であり、複数ブロックの中心元の和の
相対規格化を定めるものではない。

座標スカラーの作用と射影座標のモジュライへの降下は、後の命題
\ref{prop:pg-equivariant-poisson-gerbe}と系
\ref{cor:pg-equivariant-character-descent}で区別して扱う。
\end{remark}

\begin{remark}[純中心係数の固定と平行移動]
\label{rem:pg-equivariant-pure-center-obstruction}
$m\geq1$では、零切断全体を含む$G$-不変開上で、全$G$に関する共変性と
$\rho(\Omega_f)=f_{\mathrm{cent}}$を同時に要求することは一般にできない。
実際、中心平行移動$X=x+Cu$とこの規格化を全ての$C$に適用すると、
$\Omega_f$のWeyl記号$P_f(x,u)$について
\begin{equation}
 P_f(-Cu,u)=f(-Cu,u)\qquad\text{for all }C
 \label{eq:pg-equivariant-pure-center-evaluation}
\end{equation}
を得る。
$u\ne0$なら$Cu$は任意の$x\in W$を動くので、$P_f=f$が強制される。
しかし例えば
\begin{equation}
 f=\frac{x^3+y^3+z^3}{3}+\sum_{\alpha=1}^m u_\alpha^3
 \label{eq:pg-equivariant-pure-center-fermat}
\end{equation}
について、$t\ne0$の適切な開上では$\sym(f)$は中心元ではない。
定理\ref{thm:explicit-hesse-center}のFermat公式が与える商線束の基底は
$(x^3+y^3+z^3)/3-2t^2\sym(xyz)/3$であり、
Weyl基底の独立性から$(x^3+y^3+z^3)/3$はこの直線に入らない。
従って定理\ref{thm:pg-equivariant-parabolic-center}で
純中心係数の固定を課さないことは必要である。
同定理は滑らかな三変数制限に加え、節点・尖点の各層でも全$G$を扱う。
一方、三変数制限が零の場合には定理\ref{thm:pg-equivariant-zero-center}により
$P_f=f$そのものが中心であり、全$G$への共変性と純中心係数の固定は両立する。
節点・尖点の各軌道についても、定理\ref{thm:pg-equivariant-parabolic-center}の
全$G$同変選択を得る。
可約・非被約制限も含めた型をまたぐ一つの選択は、
定理\ref{thm:pg-intrinsic-rank-two-center}で中心核を直接扱って構成する。
この選択と本節の各明示式との一致は課していない。
\end{remark}

\subsection{階数二のPoissonモジュライへの降下}

\begin{proposition}[階数二のPoissonスタックと座標スカラーのgerbe]
\label{prop:pg-equivariant-poisson-gerbe}
$m\ge2$、$n=m+1$、$V=W\oplus Z$、$\dim W=3$、$\dim Z=m$とする。
\eqref{eq:pg-equivariant-elliptic-base}の係数空間と
\eqref{eq:pg-equivariant-groups}の群$G$について、
\begin{equation}
 \begin{aligned}
 \mathscr S_m&=\mathcal T_m^{\mathrm{ell}}\times\mathbf G_{m,\eta},\\
 K_{\mathrm{sc}}&=
 \{(aI_3,aI_m,0,a^{-3}):a\in\mathbf G_m\}\subset G,
 \qquad \overline G=G/K_{\mathrm{sc}}
 \end{aligned}
 \label{eq:pg-equivariant-gerbe-parameter}
\end{equation}
とおく。$G$は$f$に\eqref{eq:pg-equivariant-groups}で、
$\eta$に$\eta'=\lambda\det(M)\eta$で作用する。

$\mathfrak P_{n,1}^{\mathrm{ell}}$を、相対Poissonスタック
$\mathfrak{Pois}^{\mathrm{rel}}_n$の一般階数二の
\emph{被約な局所閉層}のうち、表示$\pi=\iota_{dF}P$の
support三次元空間$W_P$への制限$F|_{W_P}$が
滑らかな平面三次を定める非空稠密開部分とする。
このとき
\begin{equation}
 [\mathscr S_m/\overline G]\simeq
 \mathfrak P_{n,1}^{\mathrm{ell}},\qquad
 [\mathscr S_m/G]\longrightarrow
 \mathfrak P_{n,1}^{\mathrm{ell}}
 \quad\text{は}\ \mathbf G_m\text{-gerbe}
 \label{eq:pg-equivariant-poisson-gerbe}
\end{equation}
である。ここで開部分という条件は階数二の層の中での条件であり、
一般には相対Poissonスタック全体の開部分ではない。
特に$n=3$では底スタックの次元は
$\dim\mathfrak P_{3,1}^{\mathrm{ell}}=17$であり、
$[\mathscr S_2/G]$の次元は$16$である。
\end{proposition}
\begin{proof}
$\mathsf P_W\subset\mathrm{GL}(V)$を$W$を保つ放物部分群とする。
従って$G=\mathsf P_W\times\mathbf G_{m,\lambda}$である。
$P_0=e_1\wedge e_2\wedge e_3\in\bigwedge^3W$を固定し、
\begin{equation}
 \widetilde{\mathscr R}^{\mathrm{ell}}
 =\{(F,P):V(F)\text{は滑らか},\quad
 P\ne0\text{は分解可能},\quad
 F|_{\operatorname{supp}P}\text{は滑らかな平面三次を定める}\}
 \label{eq:pg-equivariant-decomposable-incidence}
\end{equation}
とおく。これは滑らかな三次式のアフィン係数空間と
非零Pl\"ucker錐の積の開部分である。
非零の分解可能三重ベクトルはsupportの三次元空間と
その最高外積の非零元を一意に定める。
$\mathrm{GL}(V)$は三次元部分空間全体に推移的に作用するので、
\begin{equation}
 \mathrm{GL}(V)\times^{\mathsf P_W}\mathscr S_m
 \xrightarrow{\ \sim\ }\widetilde{\mathscr R}^{\mathrm{ell}},\qquad
 [g,(f,\eta)]\longmapsto
 \bigl(f\circ g^{-1},\eta(\bigwedge^3g)P_0\bigr)
 \label{eq:pg-equivariant-associated-incidence}
\end{equation}
を得る。ここでは$\mathsf P_W$の作用は$\lambda=1$の作用である。
この同型は任意の底の上でも成立する。
実際、supportの部分束に適合する枠を局所的に選べば逆写像を構成でき、
枠の違いはちょうど$\mathsf P_W$の作用となる。

階数層のスキーム構造についても確認する。
縮約写像$\bigwedge^2V^*\to V$が階数三である局所閉条件を課すと、
その像は階数三の部分束$\mathcal W\subset V\otimes\OO$となる。
全ての二重縮約が$\mathcal W$に入ることから、
$P\in\bigwedge^3\mathcal W$である。
像の階数が三なので、この最高外積の係数は局所的に単元である。
従ってこの行列式的な階数三の層は、非零Pl\"ucker錐と
スキームとして同型である。
その上では$\iota_{dF}P$は三変数Jacobian括弧であり、
任意の$F$に対してJacobi恒等式を満たす。
よってJacobi方程式は追加条件を課さない。
正規形定理による一般階数二との対応と合わせると、
\eqref{eq:pg-equivariant-decomposable-incidence}は
本稿の被約な階数二の層のアフィン三次式による表示である。
全Jacobiスキームの他の層やその厚みについての同定は用いない。
滑らかな三次超曲面の一般の平面切断はBertiniの定理により滑らかなので、
この開条件は各滑らかな三次式の上で非空である。
非零Pl\"ucker錐と滑らかな三次式の係数空間は既約であり、
その開部分の商である上記の層でも稠密となる。

三次式の代表を忘れる自由な作用は
\begin{equation}
 \lambda\cdot(F,P)=(\lambda^{-1}F,\lambda P)
 \label{eq:pg-equivariant-equation-scalar}
\end{equation}
である。その商を$\mathscr R^{\mathrm{ell}}$と書けば、
$\mathscr R^{\mathrm{ell}}$は
$\bigwedge^3V\otimes\OO_{U_N}(-1)$内の上記の局所閉部分であり、
$\mathfrak P_{n,1}^{\mathrm{ell}}
 =[\mathscr R^{\mathrm{ell}}/\mathrm{PGL}(V)]$である。
一方、\eqref{eq:pg-equivariant-associated-incidence}から
\begin{equation}
 [\mathscr S_m/G]
 \simeq
 [\widetilde{\mathscr R}^{\mathrm{ell}}/
       (\mathrm{GL}(V)\times\mathbf G_{m,\lambda})]
 \label{eq:pg-equivariant-two-quotients}
\end{equation}
を得る。右辺の作用では
$(aI_V,a^{-3})$が$(F,P)$を固定する。
この二つのスカラーを区別するため、
\begin{equation}
 \widehat G=(\mathrm{GL}(V)\times\mathbf G_{m,\lambda})/K_{\mathrm{sc}},
 \qquad
 1\longrightarrow\mathbf G_{m,\lambda}
 \longrightarrow\widehat G
 \longrightarrow\mathrm{PGL}(V)\longrightarrow1
 \label{eq:pg-equivariant-effective-group-sequence}
\end{equation}
と書く。左の単射は$\lambda\mapsto[(I_V,\lambda)]$である。
実際、fppf局所的に$\mathrm{PGL}(V)$への射の核の元を$[(aI_V,\lambda)]$と表せば、
$[(I_V,a^3\lambda)]$と一意に書ける。
従ってまず$K_{\mathrm{sc}}$で割り、次に自由な
\eqref{eq:pg-equivariant-equation-scalar}の作用で割ると、
$\mathrm{PGL}(V)$による商を得る。
これが\eqref{eq:pg-equivariant-poisson-gerbe}の第一の同型を与える。
$K_{\mathrm{sc}}$は$G$の中心にあり$\mathscr S_m$に自明に作用するので、
第二の射はこの$K_{\mathrm{sc}}\simeq\mathbf G_m$をbandとするgerbeである
（商のgerbe判定について\cite[Lemma~78.27.1, Tag~\texttt{06PD}]{PStacks}も参照）。
特に幾何学的点$(f,\eta)$における安定化群の関係は、fppf層としての完全列
\begin{equation}
 1\longrightarrow K_{\mathrm{sc}}
 \longrightarrow G_{(f,\eta)}
 \longrightarrow\operatorname{Aut}(X,\pi)\longrightarrow1
 \label{eq:pg-equivariant-stabilizer-sequence}
\end{equation}
である。右端は射影的なPoisson対の自己同型群であり、
方程式の代表を変える自由なスカラーを安定化群に余分に含めていない。

$n=3$、$m=2$では、三次式の係数空間の次元は$35$なので、
\begin{equation}
 \dim\mathfrak P_{3,1}^{\mathrm{ell}}
 =35+1-(9+4+6+1)+1=17.
 \label{eq:pg-equivariant-threefold-stack-dimension}
\end{equation}
$\overline G$の代わりに$G$で割るgerbeの次元は一小さい。
\end{proof}

\begin{corollary}[指標による次数線形化の降下]
\label{cor:pg-equivariant-character-descent}
$\gcd(m,3)=1$とする。整数$r,s$を$mr-3s=1$となるように選び、
\begin{equation}
 \theta:G\longrightarrow\mathbf G_m,
 \qquad \theta(M,N,C,\lambda)=(\det N)^r\lambda^s,
 \qquad H=\ker\theta
 \label{eq:pg-equivariant-scalar-character}
\end{equation}
とおく。この選択により$G\simeq K_{\mathrm{sc}}\times H$、
$H\simeq\overline G$となり、
\eqref{eq:pg-equivariant-poisson-gerbe}のgerbeは中立になる。

系\ref{cor:pg-equivariant-central-quotient}の族を
$\eta\ne0$へ制限する。各次数$d$の線形化を指標$\theta^{-d}$で
twistすると、周囲代数、中心三次イデアルおよびその商は、
\begin{equation}
 \mathcal U\subset
 \mathfrak P_{n,1}^{\mathrm{ell}}\times\mathbb A^1_\hbar,
 \qquad
 \mathfrak P_{n,1}^{\mathrm{ell}}\times\{0\}\subset\mathcal U
 \label{eq:pg-equivariant-rigidified-family-base}
\end{equation}
上の次数付き$\OO$-代数の族へ降りる。
その各次数部分は局所自由で、幾何学的ファイバーのHilbert級数、
AS-Gorenstein性、古典極限および一次Poisson括弧は
系\ref{cor:pg-equivariant-central-quotient}のものと同じである。
この族の選択は、先の中心元の選択と指標$\theta$に依存してよい。
特に$m=2$、$n=3$では
\begin{equation}
 \theta=(\det N)^2\lambda
 \label{eq:pg-equivariant-threefold-character}
\end{equation}
を取れるため、上の十七次元の底スタック上でこの代数族を得る。
\end{corollary}
\begin{proof}
$k(a)=(aI_3,aI_m,0,a^{-3})$と書くと、
$\theta(k(a))=a^{mr-3s}=a$である。
$K_{\mathrm{sc}}$は中心部分群なので、
\begin{equation}
 K_{\mathrm{sc}}\times H\xrightarrow{\ \sim\ }G,
 \qquad (k,h)\longmapsto kh,
 \qquad
 g\longmapsto
 \bigl(k(\theta(g)),k(\theta(g))^{-1}g\bigr)
 \label{eq:pg-equivariant-character-splitting}
\end{equation}
は群スキームの同型である。
$K_{\mathrm{sc}}$が底に自明に作用することから、
$[\mathscr S_m/G]\simeq[\mathscr S_m/H]\times B\mathbf G_m$であり、
$[\mathscr S_m/H]\simeq\mathfrak P_{n,1}^{\mathrm{ell}}$を得る。

元の次数$d$の線形化を$\rho_d(g)$と書き、
\begin{equation}
 \rho_d^{\theta}(g)=\theta(g)^{-d}\rho_d(g)
 \label{eq:pg-equivariant-character-twist}
\end{equation}
へ換える。
$K_{\mathrm{sc}}$の元$k(a)$は元の代数の次数$d$部分に$a^d$で作用するので、
新しい線形化では自明に作用する。
従って$\rho_d^{\theta}$は$\overline G$を通じて作用する。
積は次数を加えるため
$\theta^{-d}\theta^{-e}=\theta^{-(d+e)}$となり、積と単位の同変性は保たれる。
中心三次直線の線形化も$\theta^{-3}$でtwistすれば、
その包含および生成するイデアルの乗法写像は同変である。
従って代数と商は$\overline G$による商へ降りる。
ここでは商スタック上の準連接層と同変準連接層の同値
\cite[Proposition~96.14.3, Tag~\texttt{06WT}]{PStacks}を
次数別に用いており、積とイデアルも同変写像として降りる。

底の開集合は$G$-不変であり$K_{\mathrm{sc}}$が底に自明に作用するので、
\eqref{eq:pg-equivariant-rigidified-family-base}へ降りる。
線形化の変更はアフィンatlas上の代数、中心イデアル、乗法を変えない。
従って次数別局所自由性と各ファイバーの主張は前の系から従う。
$m=2$では$r=2,s=1$が所要の等式を満たす。
\end{proof}

\begin{remark}[降下した古典極限と中心三次直線]
\label{rem:pg-equivariant-classical-family}
前の系の降下した代数を$\mathcal A^\theta$、$\mathcal B^\theta$とし、
$E_\theta=\mathcal A^\theta_1|_{\hbar=0}$を次数一のベクトル束とする。
古典極限は底スタック上で
\begin{equation}
 \mathcal A^\theta|_{\hbar=0}\simeq\operatorname{Sym}(E_\theta),
 \qquad
 \mathcal B^\theta|_{\hbar=0}
 \simeq\operatorname{Sym}(E_\theta)/(\mathcal J_0)
 \label{eq:pg-equivariant-classical-family}
\end{equation}
と書ける。ここで$\mathcal J_0\subset\operatorname{Sym}^3(E_\theta)$は、
局所的に古典方程式$f$が張る中心三次直線であり、
右辺ではそれが生成する斉次イデアルを取る。
この表示は各atlasでの多項式環と方程式の表示を降ろしたものである。
指標による一様な次数twistは相対$\operatorname{Proj}$を変えないので、
その射影化は底Poissonスタック上の元の普遍三次超曲面族となる。
降下するのは中心三次直線とその生成イデアルであり、
底スタック上でこの直線を自明化する一つの大域生成元$\Omega$を
選べることまでは主張しない。
\end{remark}

\begin{proposition}[support平面が超曲面に含まれる層の量子化]
\label{prop:pg-equivariant-zero-stratum}
$\mathcal T_m^{\mathrm z}$が非空であることと$m\ge3$は同値である。
$m\ge3$、$n=m+1$とし、
$\mathfrak P_{n,1}^{\mathrm z}$を、一般階数二の被約Poisson層のうち
\begin{equation}
 \PP(W_P)\subset X=V(F),\qquad W_P=\operatorname{supp}P
 \label{eq:pg-equivariant-contained-support-plane}
\end{equation}
を満たす被約閉部分とする。
これは滑らかかつ既約であり、階数二の層の中で余次元十を持つ。
$D=\binom{m+5}{3}$とおけば
\begin{equation}
 \dim\mathfrak P_{n,1}^{\mathrm z}=D-m^2-3m-18
 \label{eq:pg-equivariant-zero-stratum-dimension}
\end{equation}
である。

$\mathscr S_m^{\mathrm z}=\mathcal T_m^{\mathrm z}\times\mathbf G_{m,\eta}$とし、
群$G$、$K_{\mathrm{sc}}$、$\overline G$を
\eqref{eq:pg-equivariant-gerbe-parameter}のものとすると、
\begin{equation}
 [\mathscr S_m^{\mathrm z}/\overline G]\simeq\mathfrak P_{n,1}^{\mathrm z},
 \qquad
 [\mathscr S_m^{\mathrm z}/G]\longrightarrow\mathfrak P_{n,1}^{\mathrm z}
 \quad\text{は}\ \mathbf G_m\text{-gerbe}
 \label{eq:pg-equivariant-zero-gerbe}
\end{equation}
である。系\ref{cor:pg-equivariant-central-quotient}の零制限の族は
$[\mathscr S_m^{\mathrm z}/G]\times\mathbb A^1_\hbar$全体に定義される。
$\gcd(m,3)=1$なら、\eqref{eq:pg-equivariant-scalar-character}の
指標$\theta$によって次数線形化を補正すると、周囲代数、中心三次イデアルと商が
\begin{equation}
 \mathfrak P_{n,1}^{\mathrm z}\times\mathbb A^1_\hbar
 \label{eq:pg-equivariant-zero-descended-base}
\end{equation}
全体の上の代数族へ降りる。
各次数部分は局所自由であり、全幾何学的ファイバーで同系のHilbert級数とAS条件を満たす。
\end{proposition}
\begin{proof}
$g=0$なら$\PP(W)$上でactive偏微分は全て零である。
残りの$m$個の偏微分$f_{u_\alpha}|_{u=0}$は$\PP^2$上の二次式であり、
$m\le2$なら共通零点を持つ。
その点は$X$の特異点となるため、$\mathcal T_m^{\mathrm z}$は空である。
逆に$m\ge3$なら
\begin{equation}
 f=\sum_{i=1}^3(u_ix_i^2+u_i^3)+\sum_{\alpha=4}^m u_\alpha^3
 \label{eq:pg-equivariant-zero-smooth-example}
\end{equation}
は滑らかな例を与える。
$u_ix_i=x_i^2+3u_i^2=0$から$x_i=u_i=0$となり、
$\alpha>3$では$f_{u_\alpha}=3u_\alpha^2$なので、全偏微分の共通零点は錐の原点だけである。
$m=3$は例\ref{ex:pg-det-zero-ternary}の式である。

$f|_W=0$は、三変数三次式の十個の係数が零という独立な線形条件である。
従って$\mathcal T_m^{\mathrm z}$は次元$D-10$の線形空間の非空開部分で、
滑らかかつ既約である。
共通のincidence
\begin{equation}
 \widetilde{\mathscr R}
 =\{(F,P):V(F)\text{は滑らか},\ P\ne0\text{は分解可能}\}
 \label{eq:pg-equivariant-common-incidence}
\end{equation}
に$F|_{W_P}=0$を課した閉部分を
$\widetilde{\mathscr R}^{\mathrm z}$と書く。
support部分束への三次式の制限は階数十のベクトル束への全射なので、
この部分は$\widetilde{\mathscr R}$の滑らかな余次元十の部分となる。
命題\ref{prop:pg-equivariant-poisson-gerbe}の証明と同じく、
supportに適合する枠を選ぶと
\begin{equation}
 \mathrm{GL}(V)\times^{\mathsf P_W}\mathscr S_m^{\mathrm z}
 \simeq\widetilde{\mathscr R}^{\mathrm z}
 \label{eq:pg-equivariant-zero-associated-incidence}
\end{equation}
を得る。
ここでは滑らかな平面三次の開部分ではなく、
\eqref{eq:pg-equivariant-common-incidence}の上記の閉条件を用いている。
方程式の自由なスカラーと$K_{\mathrm{sc}}$については同命題の群の完全列がそのまま適用でき、
\eqref{eq:pg-equivariant-zero-gerbe}を得る。
これにより滑らかさ、既約性、余次元も商へ降りる。
\begin{equation}
 \dim[\mathscr S_m^{\mathrm z}/\overline G]
 =(D-10)+1-(9+m^2+3m)=D-m^2-3m-18
 \label{eq:pg-equivariant-zero-dimension-count}
\end{equation}
が次元の式である。

定理\ref{thm:pg-equivariant-zero-center}と
系\ref{cor:pg-equivariant-central-quotient}は全$t$にわたるので、
$t=\eta\hbar$によって全$\hbar$の族を得る。
$\gcd(m,3)=1$なら、系\ref{cor:pg-equivariant-character-descent}の証明の
指標分裂と次数ごとの$\theta^{-d}$による補正を同じ群$G$に適用する。
この操作は底の開への縮小を加えないから、
\eqref{eq:pg-equivariant-zero-descended-base}全体への降下となる。
古典極限は注意\ref{rem:pg-equivariant-classical-family}と同様に
次数一束の対称代数を中心三次直線で割ったもので、その相対$\operatorname{Proj}$は元の三次族である。
例えば$m=3$、$n=4$では底の次元は二十であり、全パラメータの族はgerbe上にある。
この場合には上記の指標条件を満たさないので、その方法による代数族の降下は結論しない。
\end{proof}

\begin{proposition}[節点・尖点層の中心商と特異点直線による降下]
\label{prop:pg-equivariant-singular-strata}
$m\ge2$、$n=m+1$、$D=\binom{m+5}{3}$とし、
$\gamma\in\{\mathrm{nod},\mathrm{cusp}\}$、
$c_{\mathrm{nod}}=1$、$c_{\mathrm{cusp}}=2$とする。
$\mathfrak P_{n,1}^{\gamma}$を、一般階数二の被約Poisson層のうち、
support平面上の切断$V(F|_{W_P})$が$\gamma$型の既約平面三次曲線となる
局所閉部分に被約構造を入れたものとする。
これは非空、滑らかかつ既約であり、一般階数二の層の中で余次元$c_\gamma$を持つ。
その次元は
\begin{equation}
 \dim\mathfrak P_{n,1}^{\gamma}
 =D-m^2-3m-8-c_\gamma
 \label{eq:pg-equivariant-singular-stratum-dimension}
\end{equation}
である。特に三次三様体では節点層と尖点層の次元はそれぞれ$16$、$15$である。

$\mathscr S_m^\gamma=\mathcal T_m^\gamma\times\mathbf G_{m,\eta}$とし、
$G$、$K_{\mathrm{sc}}$、$\overline G$を
\eqref{eq:pg-equivariant-gerbe-parameter}のものとする。このとき
\begin{equation}
 [\mathscr S_m^\gamma/\overline G]\simeq\mathfrak P_{n,1}^{\gamma},
 \qquad
 \mathcal G_\gamma:=[\mathscr S_m^\gamma/G]
 \longrightarrow\mathfrak P_{n,1}^{\gamma}
 \quad\text{は中立な}\ \mathbf G_m\text{-gerbe}
 \label{eq:pg-equivariant-singular-gerbe}
\end{equation}
である。この中立性には$\gcd(m,3)=1$を仮定しない。

具体的には、$g=f(x,0)$の唯一の特異点を表す直線
$\ell_g\subset W$が$\mathscr S_m^\gamma$上の正則な$G$-同変直線部分束をなす。
本文の引戻し規約ではその同変写像は
$M^{-1}:\ell_{g'}\to\ell_g$であり、$K_{\mathrm{sc}}$の元$k(a)$は
$\ell$に$a^{-1}$で作用する。
定理\ref{thm:pg-equivariant-parabolic-center}の$\gamma$型上の中心元を用い、
$t=\eta\hbar$とおいて得る代数、中心三次直線、商を
$\mathcal A$、$\mathcal J\subset\mathcal A_3$、$\mathcal B$とする。
次数ごとに
\begin{equation}
 \widetilde{\mathcal A}_d=\mathcal A_d\otimes\ell^{\otimes d},
 \qquad
 \widetilde{\mathcal J}=\mathcal J\otimes\ell^{\otimes3},
 \qquad
 \widetilde{\mathcal B}_d=\mathcal B_d\otimes\ell^{\otimes d}
 \label{eq:pg-equivariant-singular-line-twist}
\end{equation}
と捻ると、これらは
\begin{equation}
 \mathfrak P_{n,1}^{\gamma}\times\{0\}
 \ \subset\ \mathcal U_\gamma
 \ \subset\ \mathfrak P_{n,1}^{\gamma}\times\mathbb A^1_\hbar
 \label{eq:pg-equivariant-singular-descended-base}
\end{equation}
という開部分上の次数付き代数、中心三次イデアルと商へ降りる。
各次数部分は局所自由で、幾何学的ファイバーのHilbert級数とAS条件は
系\ref{cor:pg-equivariant-central-quotient}のものと同じである。
古典極限の相対$\operatorname{Proj}$は元の普遍三次超曲面族であり、
一次Poisson括弧も元のPoisson対のものとなる。
これらは型$\gamma$ごとの主張であり、全$\hbar$直線上への延長や、
節点型から尖点型への退化をまたぐ選択の整合性は結論しない。
\end{proposition}
\begin{proof}
補題\ref{lem:pg-equivariant-singular-retractions}の軌道
$\mathcal O_{\mathrm{nod}}$、$\mathcal O_{\mathrm{cusp}}$は、
三変数三次式の十次元空間内の滑らかな既約局所閉部分で、
次元はそれぞれ$9$、$8$である。
制限$\operatorname{Sym}^3V^*\to\operatorname{Sym}^3W^*$は全射なので、
その$\mathcal O_\gamma$の逆像は
$\mathcal O_\gamma\times\mathbb A^{D-10}$と同型である。
$\mathcal T_m^\gamma$はこの逆像のうち全三次式が滑らかである開部分である。

この開部分の非空性は、$m\ge1$について次の例から分かる：
\begin{equation}
 \begin{aligned}
  g_{\mathrm{nod}}&=x^3+y^3+xyz,&
  g_{\mathrm{cusp}}&=y^2z-x^3,\\
  f_\gamma&=g_\gamma+u_1z^2+\sum_{\alpha=1}^m u_\alpha^3.
 \end{aligned}
 \label{eq:pg-equivariant-singular-smooth-examples}
\end{equation}
全偏微分の共通零点では$\alpha>1$について$u_\alpha=0$である。
$u=u_1$と書く。節点型の残る四つの勾配方程式は
\begin{equation}
 3x^2+yz=3y^2+xz=xy+2uz=z^2+3u^2=0
 \label{eq:pg-equivariant-nodal-example-gradient}
\end{equation}
である。$u=0$なら直ちに$x=y=z=0$となる。
$u\ne0$なら$z\ne0$かつ$xy\ne0$であり、最初の二式を掛けると
$z^2=9xy=-18uz$となる。
最後の式と合わせると$z=u/6$を得るが、これは$z^2=-3u^2$に反する。
尖点型では
\begin{equation}
 -3x^2=2yz=y^2+2uz=z^2+3u^2=0
 \label{eq:pg-equivariant-cuspidal-example-gradient}
\end{equation}
であり、$x=0$、$yz=0$から$y=z=u=0$を得る。
従っていずれも錐の原点以外に勾配の共通零点を持たず、滑らかである。
なお$m=0$なら全三次式が$g$自身となるので、この係数空間は空である。
以上から$\mathcal T_m^\gamma$は非空、滑らかかつ既約で、次元$D-c_\gamma$を持つ。

共通のincidenceに制限の型を課した部分を
\begin{equation}
 \widetilde{\mathscr R}^{\gamma}
 =\left\{(F,P)\ \middle|\,
 \begin{gathered}
 V(F)\text{は滑らか},\quad P\ne0\text{は分解可能},\\
 F|_{W_P}\in\mathcal O_\gamma
 \end{gathered}\right\}
 \label{eq:pg-equivariant-singular-incidence}
\end{equation}
と書く。ここで最後の条件はsupportの枠に依存しない。
非零Pl\"ucker錐上のsupport部分束への制限は、階数十のベクトル束への全射である。
supportの枠を局所的に取ると上記の軌道への制限条件となるので、
$\widetilde{\mathscr R}^{\gamma}$は滑らかな余次元$c_\gamma$の局所閉部分であり、
適合枠による表示は
\begin{equation}
 \mathrm{GL}(V)\times^{\mathsf P_W}\mathscr S_m^\gamma
 \simeq\widetilde{\mathscr R}^{\gamma}
 \label{eq:pg-equivariant-singular-associated-incidence}
\end{equation}
である。これは既約性も示す。
命題\ref{prop:pg-equivariant-poisson-gerbe}の証明で用いた
方程式代表の自由なスカラー作用、中心部分群$K_{\mathrm{sc}}$、
および二つの商の同定はこの不変局所閉部分へ制限できる。
従って\eqref{eq:pg-equivariant-singular-gerbe}の商表示とgerbeの主張、
ならびに底の滑らかさ、既約性、余次元を得る。
次元は
\begin{equation}
 \begin{aligned}
 \dim[\mathscr S_m^\gamma/\overline G]
 &=(D-c_\gamma)+1-(9+m^2+3m)\\
 &=D-m^2-3m-8-c_\gamma,\\
 \dim\mathcal G_\gamma
 &=D-m^2-3m-9-c_\gamma
 \end{aligned}
 \label{eq:pg-equivariant-singular-dimension-count}
\end{equation}
となる。

次に、各基準形$g_\gamma$の唯一の特異点は$[0:0:1]$である。
その安定化群はこの点を保つので、軌道表示から正則な同変射
\begin{equation}
 \nu_\gamma:\mathcal O_\gamma\longrightarrow\PP(W),
 \qquad g\longmapsto\operatorname{Sing}(V(g))_{\mathrm{red}},
 \qquad
 \ell=\nu_\gamma^*\OO_{\PP(W)}(-1)\subset W\otimes\OO
 \label{eq:pg-equivariant-singular-point-line}
\end{equation}
が得られる。正則性は軌道へのfppf射からの降下で従い、
特異点スキームの長さを一定にする主張を用いていない。
全$G$の作用でも$u=0$上では$X=Mx$であるから、
点の順方向の写像は$M:\ell_g\to\ell_{g'}$となる。
代数の$\Phi$と同じ向きの引戻しデータはその逆であり、
\begin{equation}
 \ell_{g'}\xrightarrow{\ M^{-1}\ }\ell_g,
 \qquad k(a)|_\ell=a^{-1},
 \qquad k(a)|_{\mathcal A_d\otimes\ell^{\otimes d}}
 =a^d a^{-d}=1
 \label{eq:pg-equivariant-singular-line-weights}
\end{equation}
となる。

$L=\ell^\vee$は$\mathcal G_\gamma$上の重み一の可逆層である。
$L$のframeを付けた対象のスタックを考える。
gerbeをfppf局所的に$B\mathbf G_m$と同定すれば、
重み一の線束のframeを指定することで相対自己同型群が消え、
このスタックは底と同型になる。
この局所同型は降下するので、frame付きスタックは
$\mathfrak P_{n,1}^\gamma$そのものである。
frameを忘れる射がgerbeの切断を与え、中立性が従う。
ここでは$G$を一定の直積$K_{\mathrm{sc}}\times H$へ分裂させていない。

補題\ref{lem:pg-equivariant-singular-retractions}と
定理\ref{thm:pg-equivariant-parabolic-center}により、各型の全零切断を含む
$G$-不変開上で中心元と中心商族を得る。
中心調整写像で既存の不変関数を引き戻した$Q(f^0,t)$は
$t=0$で常に$1$なので、この段階で係数底の一部を除く必要はない。
$t=\eta\hbar$と引き戻せば、定義開は$\hbar=0$全体を含む。
\eqref{eq:pg-equivariant-singular-line-weights}によって
\eqref{eq:pg-equivariant-singular-line-twist}の各次数は$K_{\mathrm{sc}}$に自明であり、
$\overline G$の同変層として降下する。
これは商スタック上の準連接層の降下
\cite[Proposition~96.14.3, Tag~\texttt{06WT}]{PStacks}であり、
gerbe上の重み零の層による記述では
\cite[Proposition~4.10, p.~1487]{PBerghSchnurerGerbes}の特別な場合である。
積に対しては$\ell^{\otimes d}\otimes\ell^{\otimes e}
\simeq\ell^{\otimes(d+e)}$を用いるので、単位、積、
中心三次直線の包含と生成イデアルも同時に降りる。
従って実際の次数付き代数とその中心商を
\eqref{eq:pg-equivariant-singular-descended-base}上に得る。
可逆層による次数ごとの捻りは局所的に元の代数と同型であるため、
系\ref{cor:pg-equivariant-central-quotient}の左右正則性、次数別局所自由性、
Hilbert級数およびAS条件が保たれる。

降下した古典次数一束を$E_\gamma$とすると、そのatlasへの引戻しは
$V^*\otimes\ell$であり、古典中心三次直線を$\mathcal J_{\gamma,0}$と書けば
\begin{equation}
 \widetilde{\mathcal A}|_{\hbar=0}
 \simeq\operatorname{Sym}(E_\gamma),\qquad
 \widetilde{\mathcal B}|_{\hbar=0}
 \simeq\operatorname{Sym}(E_\gamma)/(\mathcal J_{\gamma,0})
 \label{eq:pg-equivariant-singular-classical-family}
\end{equation}
である。可逆層による一様な次数捻りは相対$\operatorname{Proj}$を変えない。
従ってこの射影族は元の普遍三次族であり、局所表示$t=\eta\hbar$から
一次括弧も元の$\eta\{-,-\}_f$に戻る。
底上の中心三次直線の大域生成元の存在は主張しない。
また、構成に用いた左逆と特異点直線は型ごとの軌道上で定義されているので、
異なる型の境界をまたぐ貼合せはこの証明には含まれない。
\end{proof}

\begin{remark}[次数重みと十七次元の意味]
\label{rem:pg-equivariant-gerbe-degree-weights}
\eqref{eq:pg-equivariant-fixed-parameter}によって$\hbar$を固定したときも、
$K_{\mathrm{sc}}$は$(f,\eta,\hbar)$に自明に作用する。
しかし座標代入は各次数一生成元を$a$倍するので、
周囲代数$\mathcal A_d$および中心商$\mathcal B_d$には$a^d$で作用する。
中心三次直線には$a^3$で作用するため、そのイデアルは同変であるが、
次数付き代数のこの非自明な重みは消えない。
従って元の線形化のままの代数はgerbe上の族である。
系\ref{cor:pg-equivariant-character-descent}と
命題\ref{prop:pg-equivariant-zero-stratum}は$\gcd(m,3)=1$のとき、
この次数重みを指標によって取り除いている。
$3\mid m$なら、$G$の任意の指標の$K_{\mathrm{sc}}$上の重みは$3$の倍数である。
実際、指標は$\det M,\det N,\lambda$の整数冪の積であり、
それらの重みはそれぞれ$3,m,-3$である。
この場合は上の指標による分裂は得られないが、
gerbeが非中立であることや別の降下法の不存在までは結論しない。
実際、命題\ref{prop:pg-equivariant-singular-strata}の節点・尖点層では、
特異点直線による捻りが全$m$で重みを除き、gerbeの中立性も与える。

十七という数は$n=3$での底Poissonスタックの次元である。
固定した$\hbar\ne0$における抽象的な非可換代数のモジュライへの写像が
generically finiteであることや、その像の次元が十七であることは主張しない。
一方、$\hbar$を保つ次数付き変形族全体の同型は、
特殊ファイバーの可換代数と$\hbar^{-1}[-,-]\bmod\hbar$を保つ。
従ってその射影化から元のPoisson対$(X,\pi)$が回復される。
これは一般ファイバーだけからの回復とは異なる。
\end{remark}

% !TeX root = ../../cubic-quantization.tex
\section{階数二の全係数空間における中心直線}
\label{sec:pg-intrinsic-rank-two-center}
前節の中心調整は、三変数制限の型ごとに適した座標を選ぶ構成であった。
本節ではsupportの三次元部分空間を動かし、中心核そのものを扱う。
Grassmann多様体上の一次コホモロジーの消滅と有限加群の完備化を用いることで、
三変数制限の型を区別せず、全線形群に共変な中心三次元を同時に選ぶ。
この中心元は普遍的な二次Hessian補正まで指定できる。
三次Veroneseでは次元の合同条件なしに全階数二層へ降下し、
次数一の降下を妨げるgerbeの位数も計算する。
一方、零制限層の既存中心直線との大域的一致には二次jetの障害がある。
形式的中心元の存在自体は、所定の古典剰余を持つ場合を含め、
Etingof--Ginzburg\cite[Proposition~8.1.1(i), pp.~1407--1408]{EG}の既知の結果である。
ここで証明するのは、動く全support上での代数的な同時選択である。
形式的持上げには補題\ref{lem:pg-det-divided-kernel}の局所核を用い、
その大域化に必要なコホモロジーを直接計算する。

$N\ge5$、$m=N-3$、$V=\CC^N$とし、
$\mathcal T_N\subset\operatorname{Sym}^3V^*$を、
滑らかな射影超曲面を定める三次式の係数空間とする。
これは判別式を可逆にしたアフィンスキームである。
次の底と群を用いる：
\begin{equation}
 \begin{aligned}
 C_3^\times(V)&=
 \{0\ne P\in\Lambda^3V:P\text{ は分解可能}\},\\
 \mathcal S_N&=\mathcal T_N\times C_3^\times(V),
 &H_N&=\mathrm{GL}(V)\times\mathbf G_{m,\lambda},\\
 F'(X)&=\lambda^{-1}F(L^{-1}X),
 &P'&=\lambda(\Lambda^3L)P,\qquad \hbar'=\hbar.
 \end{aligned}
 \label{eq:pg-intrinsic-rank-two-parameters}
\end{equation}
ここで$(L,\lambda)\in H_N$、$X=Lv$である。
座標$v_1,\ldots,v_N$で$P$の交代成分を$P^{ijk}$と書き、
\begin{equation}
 \mathcal A_{F,P,\hbar}
 =\CC\langle v_1,\ldots,v_N\rangle\Big/
 \left(
 [v_i,v_j]-\hbar\sum_{k=1}^N
 P^{ijk}\sym\left(\frac{\partial F}{\partial v_k}\right)
 \right)_{i,j}
 \label{eq:pg-intrinsic-rank-two-algebra}
\end{equation}
とする。係数を動かす場合には同じ表示を構造層上で読む。
代入$X=Lv$は同型
$\Phi:\mathcal A_{F',P',\hbar}\to\mathcal A_{F,P,\hbar}$を与える。
$P$のsupportに適合する基底では、この代数は中心変数$m$個を持つ
一ブロック代数である。

\begin{theorem}[全階数二係数空間上の同変中心元]
\label{thm:pg-intrinsic-rank-two-center}
$H_N$および$\hbar\mapsto-\hbar$で不変なZariski開
\begin{equation}
 \mathcal W_N\subset\mathcal S_N\times\mathbb A^1_\hbar,
 \qquad \mathcal S_N\times\{0\}\subset\mathcal W_N
 \label{eq:pg-intrinsic-rank-two-open}
\end{equation}
と、その上でWeyl係数が正則な中心三次元$\Omega_{F,P,\hbar}$が存在し、
次を満たす。
\begin{enumerate}
\item 古典極限、偶性および全線形群に関する共変性は
\begin{equation}
 \Omega_{F,P,0}=F,\qquad
 \Omega_{F,P,-\hbar}=\Omega_{F,P,\hbar},\qquad
 \Phi(\Omega_{F',P',\hbar})
       =\lambda^{-1}\Omega_{F,P,\hbar}
 \label{eq:pg-intrinsic-rank-two-covariance}
\end{equation}
である。偶性の等号はWeyl表示での等号を表す。
さらに、命題\ref{prop:pg-ad-intrinsic-hessian}の内在的Hessian縮約
\begin{equation}
 \mathscr H(P,F)=\frac16\sum_{i,j,k,a,b,c}
 P^{ijk}P^{abc}F_{ia}F_{jb}F_{kc}
 \label{eq:pg-intrinsic-prescribed-hessian}
\end{equation}
を用いて、二次jetを
\begin{equation}
 \Omega_{F,P,\hbar}
   =\sym\left(F-\frac{\hbar^2}{12}\mathscr H(P,F)+O(\hbar^4)\right)
 \label{eq:pg-intrinsic-prescribed-second-jet}
\end{equation}
と指定できる。ここで$O(\hbar^4)$は全零切断に沿うWeyl係数の表示である。
\item $W_P\subset V$を$P$のsupportとする。
割った交換子の三次核$\mathcal K$と、
純中心三次式の部分束$\mathcal C$について
\begin{equation}
 \mathcal C_{(F,P,\hbar)}
   =\operatorname{Sym}^3(V/W_P)^*,\qquad
 \mathcal K=\mathcal C\oplus\OO_{\mathcal W_N}\Omega
 \label{eq:pg-intrinsic-rank-two-kernel-splitting}
\end{equation}
が成り立つ。従って$\mathcal J=\OO_{\mathcal W_N}\Omega$は
$H_N$-同変な中心直線である。
\item $\mathcal A$と中心商$\mathcal B=\mathcal A/(\mathcal J)$の各次数部分は
$\mathcal W_N$上で有限局所自由である。
各閉ファイバーの$\mathcal A$は次元$N$、AS parameter $N$の
Koszul AS-regular domainであり、$\Omega$は左右正則である。
各閉ファイバーの$\mathcal B$は左右injective dimension $N-1$、
AS parameter $N-3$のAS-Gorenstein代数で、
\begin{equation}
 \operatorname{Hilb}_{\mathcal A}(z)=(1-z)^{-N},\qquad
 \operatorname{Hilb}_{\mathcal B}(z)=\frac{1-z^3}{(1-z)^N},\qquad
 \mathcal B/(\hbar)\simeq
 \operatorname{Sym}(V^*)\otimes\OO_{\mathcal S_N}/(F)
 \label{eq:pg-intrinsic-rank-two-quotient}
\end{equation}
である。特殊ファイバーの一次Poisson括弧は
$\{v_i,v_j\}=\sum_kP^{ijk}\partial_kF$である。
\end{enumerate}
三変数制限が可約、非被約または恒等的零となる点も含め、
$\mathcal S_N$の全零切断を一つの開集合で扱う。
全$\hbar$直線上での存在や、既に選んだ型別の中心式との一致は主張しない。
\end{theorem}

証明で使うコホモロジーの消滅を先に示す。
$\mathbf G=\operatorname{Gr}(3,V)$上のtautological完全列を
\begin{equation}
 0\longrightarrow\mathcal W\longrightarrow V\otimes\OO_{\mathbf G}
 \longrightarrow\mathcal Q\longrightarrow0,\qquad
 \OO_{\mathbf G}(1)=\det\mathcal W^*
 \label{eq:pg-rank-two-tautological-sequence}
\end{equation}
と書く。

\begin{lemma}[中心三次束の一次コホモロジー]
\label{lem:pg-rank-two-bott-vanishing}
任意の整数$j$について
\begin{equation}
 H^1(\mathbf G,\OO_{\mathbf G}(j))=0,\qquad
 H^1\left(\mathbf G,
 \operatorname{Sym}^3\mathcal Q^*\otimes\OO_{\mathbf G}(j)\right)=0.
 \label{eq:pg-rank-two-bott-vanishing}
\end{equation}
従ってsupportを取る写像
$\pi:C_3^\times(V)\to\mathbf G$について
\begin{equation}
 H^1(C_3^\times(V),\OO)=0,\qquad
 H^1(C_3^\times(V),\pi^*\operatorname{Sym}^3\mathcal Q^*)=0.
 \label{eq:pg-rank-two-cone-cohomology}
\end{equation}
\end{lemma}
\begin{proof}
Borel--Weil--Bottの計算規則を用いる
\cite[Theorem~5, p.~3]{PLurieBWB}。
Grassmann多様体上での重みの規約は
\cite[\S2.2, Theorem~2.3, pp.~5--6]{PGuoLuSharpe}のものとする。
$\operatorname{Sym}^3\mathcal Q^*\otimes\OO(j)$に対応する重みと、
$\rho=(N-1,N-2,\ldots,0)$を加えた列は
\begin{equation}
 \begin{aligned}
 \alpha_j&=(j,j,j,3,0,\ldots,0),\\
 \alpha_j+\rho
   &=(j+N-1,j+N-2,j+N-3,N-1,N-5,N-6,\ldots,0).
 \end{aligned}
 \label{eq:pg-rank-two-bott-weights}
\end{equation}
$j=0,1,2$では$N-1$が重複するため全コホモロジーが消える。
$j\ge3$では既に狭義単調減少なので、非零となり得るのは$H^0$だけである。
$j<0$では第四成分$N-1$が最初の三成分のいずれよりも大きい。
従って重複がなければ、狭義単調減少に並べ替える置換の長さは少なくとも三である。
重複がある場合にもコホモロジーは零なので、いずれも$H^1$は消える。

$\OO(j)$では重みは$(j,j,j,0,\ldots,0)$である。
$j\ge0$では$H^0$だけが可能であり、$-N<j<0$では
$\rho$を加えた最初の三成分の少なくとも一つが後続成分と重複する。
$j\le-N$では、重複のない列を並べ替える長さは$3(N-3)>1$である。
これが\eqref{eq:pg-rank-two-bott-vanishing}を示す。

$C_3^\times(V)$は$\det\mathcal W=\OO_{\mathbf G}(-1)$の
零切断を除いた全空間であり、$\pi$はアフィンな$\mathbf G_m$-torsorである。
従って
\begin{equation}
 \pi_*\OO_{C_3^\times(V)}
    =\bigoplus_{j\in\mathbb Z}\OO_{\mathbf G}(j),\qquad
 \pi_*\pi^*\operatorname{Sym}^3\mathcal Q^*
    =\bigoplus_{j\in\mathbb Z}
      \operatorname{Sym}^3\mathcal Q^*\otimes\OO_{\mathbf G}(j).
 \label{eq:pg-rank-two-cone-direct-images}
\end{equation}
有限アフィン被覆の\v{C}ech複体を用いれば、ここでのコホモロジーは
直和と可換する。よって\eqref{eq:pg-rank-two-cone-cohomology}を得る。
\end{proof}

\begin{proof}[定理\ref{thm:pg-intrinsic-rank-two-center}の証明]
\emph{局所核と形式的持上げ。}
$C_3^\times(V)$上で局所的に$W_P$の基底と補空間を選び、
$P=\eta\,\partial_{x_1}\wedge\partial_{x_2}\wedge\partial_{x_3}$と書く。
ここで$\eta$は単元である。
\eqref{eq:pg-intrinsic-rank-two-algebra}は
$t=\eta\hbar$とした一ブロック代数となる。
命題\ref{prop:pg-PBW}の係数も動かした開
\eqref{eq:pg-det-universal-open}と、
補題\ref{lem:pg-det-divided-kernel}を適用する。
これらの条件は代数、Weyl写像、割った交換子および$W_P$から定まるので、
基底の変更と両立する。
従って全零切断を含む$H_N$-不変開上で、
各次数の代数は局所自由であり、次数三、四のWeyl写像は同型である。
また割った交換子
\begin{equation}
 \mathcal D(p)=
 \left(\hbar^{-1}[\sym(p),v_i]\right)_{i=1}^N,
 \qquad \mathcal K=\ker\mathcal D
 \label{eq:pg-intrinsic-rank-two-divided-map}
\end{equation}
の核・像・余核は局所自由で、任意の底変更と両立する。
この開を$\mathcal U_{\mathrm{loc}}$と書く。
$\hbar\mapsto-\hbar$と反対代数の同型を併用すれば、
$\mathcal U_{\mathrm{loc}}$をこの対称性でも不変に取れる。

以下、$\pi^*\operatorname{Sym}^3\mathcal Q^*$を
$\mathcal S_N$へ引き戻した束も$\mathcal C$と書く。
補題\ref{lem:pg-det-classical-kernel}により、零切断上では
\begin{equation}
 \mathcal K_0:=\mathcal K/(\hbar)
     =\OO_{\mathcal S_N}F\oplus\mathcal C,
 \qquad H^1(\mathcal S_N,\mathcal K_0)=0.
 \label{eq:pg-intrinsic-rank-two-classical-splitting}
\end{equation}
後半は補題\ref{lem:pg-rank-two-bott-vanishing}と
$\mathcal T_N$のアフィン性から従う。
ここでは$F$は普遍三次式であり、$F$が純中心式となることは滑らかさに反する。

各$\hbar$進厚みは$\mathcal U_{\mathrm{loc}}$に含まれる。
$\mathcal K_r=\mathcal K/\hbar^{r+1}\mathcal K$とおけば、
局所自由性により$\mathcal S_N$上の完全列
\begin{equation}
 0\longrightarrow\mathcal K_0
 \xrightarrow{\ \hbar^{r+1}\ }\mathcal K_{r+1}
 \longrightarrow\mathcal K_r\longrightarrow0
 \qquad(r\ge0)
 \label{eq:pg-intrinsic-rank-two-formal-sequences}
\end{equation}
を得る。\eqref{eq:pg-intrinsic-rank-two-classical-splitting}から
大域切断上の右の写像は全射である。
命題\ref{prop:pg-universal-casimir-jet}を$t=\eta\hbar$として適用すると、
\begin{equation}
 \Theta=F-\frac{\hbar^2}{12}\mathscr H(P,F),\qquad
 \mathcal D\Theta\in\hbar^4
       \bigl(\operatorname{Sym}^4V^*\bigr)^{\oplus N}[[\hbar]]
 \label{eq:pg-intrinsic-initial-second-jet}
\end{equation}
である。Hessian縮約は基底に依存せず、$F$と同じ$\lambda^{-1}$の共変性を持つ。
局所核の底変更との両立から、$\Theta\bmod\hbar^3$は
$\mathcal K_2=\mathcal K/\hbar^3\mathcal K$の大域切断となる。
この切断を初期値として\eqref{eq:pg-intrinsic-rank-two-formal-sequences}を
$r\ge2$について順次用いると、Weyl表示で
\begin{equation}
 \widehat\Omega\in
 \Gamma(\mathcal S_N,\OO)[[\hbar]]
       \otimes_\CC\operatorname{Sym}^3V^*,
 \qquad \widehat\Omega\equiv\Theta\pmod{\hbar^3},
 \qquad \mathcal D\widehat\Omega=0
 \label{eq:pg-intrinsic-rank-two-formal-center}
\end{equation}
を得る。この段階では同変性を課さない。

\emph{アフィン係数環への延長。}
Stiefel開部分$\operatorname{St}(3,V)\subset\operatorname{Mat}_{N\times3}$から
列の外積を取る写像は、$C_3^\times(V)$上の
Zariski局所自明な$\mathrm{SL}_3$-torsorである。
実際、非零の三次小行列を固定した開では、指定された外積を持つ基底を正則に選べる。
階数二以下の行列の閉集合は余次元$N-2\ge2$なので、
\begin{equation}
 \Gamma(\operatorname{St}(3,V),\OO)
      =\CC[\operatorname{Mat}_{N\times3}],\qquad
 \Gamma(C_3^\times(V),\OO)
      =\CC[\operatorname{Mat}_{N\times3}]^{\mathrm{SL}_3}=:B_N.
 \label{eq:pg-intrinsic-rank-two-stiefel-functions}
\end{equation}
最初の等号は次の初等的な延長による。
正則関数を多項式環の分数体の既約分数として書いたとき、
分母の既約因子があればその零点因子の一般点はStiefel開に属し、
そこで極を持つことになる。従って分母は定数である。
第二の等号はtorsorに沿う正則関数の降下である。
Hilbertの有限生成定理
\cite[Theorem~2.1, pp.~368--369]{PDerksenInvariants}により$B_N$は有限生成であり、
\begin{equation}
 A_N=\Gamma(\mathcal T_N,\OO)\otimes_\CC B_N
     =\Gamma(\mathcal S_N,\OO),\qquad
 \overline{\mathcal S}_N=\operatorname{Spec}A_N
 \label{eq:pg-intrinsic-rank-two-affine-base}
\end{equation}
は有限型の整域で定義される。
群作用はこのアフィン化へ延長する。
外積の成分$P^{ijk}$は三次小行列式で与えられる不変式なので、
\eqref{eq:pg-intrinsic-rank-two-algebra}は$A_N[\hbar]$上でも定義される。
ここでは不変式環とPl\"ucker方程式の商環の同定や、錐の正規性を必要としない。

\emph{有限次数のWeyl開。}
$d=3,4$について、語長$d$の自由加群と二次関係式の倍数からなる
有限表示行列$R_d$を取る。
$\hbar=0$でその余核は
$\operatorname{Sym}^dV^*\otimes A_N$である。
また$\mathcal U_{\mathrm{loc}}$上のPBW性から、一般点の余核の次元も
$a_d=\binom{N+d-1}{d}$である。
従って$r_d=N^d-a_d$とおくと、$R_d$の全ての$(r_d+1)$次小行列式は
整域$A_N[\hbar]$において恒等的に零である。

$\hbar=0$で独立な$r_d$本の関係式の列を固定する。
これらと標準単項式のWeyl対称化$a_d$本とを並べた正方行列の行列式は、
$\hbar=0$で非零定数である。
これを規格化した関数$q_d$は$q_d\bmod\hbar=1$を満たす。
$q_d$を可逆にすると、選んだ関係式の列とWeyl列が基底になる。
上の高次小行列式の消滅により、残りの関係式も選んだ関係式の列で生成される。
従って次数$d$のWeyl写像は同型である。

次数三、四のWeyl写像が同型となる開は
$H_N\times\boldsymbol\mu_2$-不変で、
$\overline{\mathcal S}_N\times\{0\}$全体を含む。
ここで$\boldsymbol\mu_2$は$\hbar\mapsto-\hbar$と語の順序反転で作用する。
補題\ref{lem:pg-equivariant-affine-lifting}(1)により、
この開に含まれるprincipal開を
\begin{equation}
 D(q),\qquad q\in A_N[\hbar]^{H_N\times\boldsymbol\mu_2},
 \qquad q\bmod\hbar=1
 \label{eq:pg-intrinsic-rank-two-uniform-weyl-open}
\end{equation}
と取れる。
これは中心性を有限行列で扱うための二つの次数だけの主張であり、
全次数のPBW性は別に$\mathcal U_{\mathrm{loc}}$から用いる。

\emph{有限核の完備化による代数的持上げ。}
\begin{equation}
 R=A_N[\hbar,q^{-1}],\qquad
 \widehat R=A_N[[\hbar]],\qquad
 D:R\otimes\operatorname{Sym}^3V^*
     \longrightarrow R\otimes(\operatorname{Sym}^4V^*)^{\oplus N},
 \qquad K=\ker D
 \label{eq:pg-intrinsic-rank-two-finite-kernel}
\end{equation}
とする。$D$は\eqref{eq:pg-intrinsic-rank-two-divided-map}のWeyl行列である。
次数四部分が自由で、交換子は$\hbar=0$で零となるため、
$\hbar$で割った行列も$R$上で定義される。
$R$はNoether環なので$K$は有限加群である。
\eqref{eq:pg-intrinsic-rank-two-formal-center}の各係数は
\eqref{eq:pg-intrinsic-rank-two-affine-base}により$A_N$に属する。
中心性の各係数等式もこの環で成立するから、
$\widehat\Omega\in\ker\widehat D$である。

有限加群の完備化の完全性
\cite[Lemmas~10.97.1--10.97.2, Tags~00MA--00MB]{PStacks}から
\begin{equation}
 \widehat K=\ker\widehat D,\qquad
 \widehat K/\hbar^\ell\widehat K\simeq K/\hbar^\ell K\quad(\ell\ge1)
 \label{eq:pg-intrinsic-rank-two-completion}
\end{equation}
となる。従って$\widehat\Omega$の$\hbar^3$を法とする剰余類の代表を
$\Omega\in K$として選べる。
有限自由加群への包含と剰余写像の可換性により
\begin{equation}
 D\Omega=0,\qquad \Omega\equiv\Theta\pmod{\hbar^3}
 \label{eq:pg-intrinsic-rank-two-algebraic-lift}
\end{equation}
である。指定した形式級数全体の代数化を要求しているのではなく、
指定した有限二次jetを保つ一つの代数的中心元を得ている。
特に$P=0$で$K$の階数が跳ぶことは、有限加群に対する
\eqref{eq:pg-intrinsic-rank-two-completion}の適用を妨げない。

アフィン底$D(q)$上の$K$は有理$H_N$-加群であり、
その引戻しdatumを$\lambda\Phi$でtwistすれば$\Theta\bmod\hbar^3$は不変切断になる。
$\Omega$をReynolds作用素で平均する。
平均は$K$と剰余条件\eqref{eq:pg-intrinsic-rank-two-algebraic-lift}を保ち、
\eqref{eq:pg-intrinsic-rank-two-covariance}の共変性を与える。
さらに$\boldsymbol\mu_2$について平均すればWeyl係数は偶になる。
$\Theta$も偶なので指定jetは保たれ、差の三次係数も零となるため
\eqref{eq:pg-intrinsic-prescribed-second-jet}を得る。
この平均化には$K$の局所自由性は不要である。
各切断の軌道の線形包が有限次元であることと、
Reynolds作用素が同変線形写像と可換することだけを用いる。

\emph{全次数の性質と中心商。}
$D(q)$を$\mathcal S_N\times\mathbb A^1_\hbar$へ引き戻し、
$\mathcal U_{\mathrm{loc}}$と交わらせる。
この上で$\mathcal K/\mathcal C$は線束であり、
$\Omega$の像は零切断の全点で生成元である。
その像が零となる閉集合を除いて$\mathcal W_N$を得る。
これは$H_N$と$\hbar\mapsto-\hbar$で不変で、全零切断を含む。
$\Omega$の像が商線束を自明化することから
\eqref{eq:pg-intrinsic-rank-two-kernel-splitting}が従う。

各閉点でsupportに適合する基底を選ぶと、
命題\ref{prop:pg-PBW}により
\begin{equation}
 \operatorname{gr}_x\mathcal A_{F,P,\hbar}
   \simeq S_{g,\eta\hbar}\otimes_\CC\CC[u_1,\ldots,u_m],
 \qquad g=F|_{W_P}
 \label{eq:pg-intrinsic-rank-two-fiber-domain}
\end{equation}
はdomainである。従って$\mathcal A_{F,P,\hbar}$もdomainであり、
非零な中心元$\Omega$は左右正則である。
同じPBW性により、周囲代数の$u_1,\ldots,u_m$は中心正則列となる。
その商$S_{g,\eta\hbar}$は三次元Koszul AS-regularなので、
補題\ref{lem:pg-central-linear}を$m$回適用すると
周囲代数のAS条件と次元を得る。
$g$の滑らかさは使わず、$g=0$もこの議論に含まれる。

各次数$d$での乗法写像は完全列
\begin{equation}
 0\longrightarrow\mathcal J\otimes\mathcal A_{d-3}
 \xrightarrow{\ \mathrm{mult}\ }\mathcal A_d
 \longrightarrow\mathcal B_d\longrightarrow0
 \label{eq:pg-intrinsic-rank-two-flat-quotient-sequence}
\end{equation}
を与える。左の写像は全閉ファイバーで単射なので、
有限局所自由加群の写像として局所的に分裂する。
従って$\mathcal B_d$は有限局所自由である。
Hilbert級数、特殊化および一次括弧は
\eqref{eq:pg-intrinsic-rank-two-algebra}と
\eqref{eq:pg-intrinsic-rank-two-covariance}から従い、
AS-Gorenstein条件は補題\ref{lem:pg-central-cubic}による。
\end{proof}

\begin{remark}[型別の選択と大域的選択]
\label{rem:pg-intrinsic-rank-two-choice}
定理\ref{thm:pg-intrinsic-rank-two-center}は、supportの補空間を
同変に選ぶことなく、一つの中心直線を全ての三変数制限の型にわたって構成する。
純中心係数を固定する規格化や、前節の型別の中心式との一致は課していない。
零制限上での全$\hbar$直線にわたる式
$\Omega^{\mathrm z}=\sym(F)$は、その強い範囲を持つ別の選択である。
$N\ge6$では、この標準中心直線への全層での一致は、
同変性を要求しなくても二次jetで妨げられる
（命題\ref{prop:pg-intrinsic-zero-matching-obstruction}）。
命題\ref{prop:pg-equivariant-global-centering-obstruction}の中心調整の非存在は、
この中心直線の存在と両立する。
\end{remark}

\subsection{全階数二層への降下}

\begin{corollary}[階数二の被約Poisson層全体への次数付き代数の降下]
\label{cor:pg-intrinsic-rank-two-descent}
$N\ge5$、$n=N-2$、$V=\CC^N$とする。
滑らかな全三次式の係数空間を$\mathcal T_N$、
非零の分解可能三重ベクトルの空間を$C_3^\times(V)$とし、
\begin{equation}
 \begin{aligned}
 \mathcal S_N&=\mathcal T_N\times C_3^\times(V),\\
 H_N&=\mathrm{GL}(V)\times\mathbf G_{m,\lambda},\\
 (L,\lambda)\cdot(F,P)
   &=\bigl(\lambda^{-1}F\circ L^{-1},
            \lambda(\bigwedge^3L)P\bigr),\\
 K_{\mathrm{sc}}&=\{(cI_V,c^{-3}):c\in\mathbf G_m\},
 \qquad \overline H_N=H_N/K_{\mathrm{sc}}
 \end{aligned}
 \label{eq:pg-intrinsic-descent-base}
\end{equation}
とおく。$\mathfrak P_{n,1}$を、相対Poissonスタックのうち
一般階数が二である局所閉層に被約構造を入れたものとする。
このとき
\begin{equation}
 [\mathcal S_N/\overline H_N]\simeq\mathfrak P_{n,1},
 \qquad
 [\mathcal S_N/H_N]\longrightarrow\mathfrak P_{n,1}
 \quad\text{は}\ \mathbf G_m\text{-gerbe}
 \label{eq:pg-intrinsic-descent-gerbe}
\end{equation}
である。

定理\ref{thm:pg-intrinsic-rank-two-center}の族は、
このgerbeと$\mathbb A^1_\hbar$の積の零切断近傍に定義される。
さらに$\gcd(N,3)=1$なら、$Nr-3s=1$となる整数$r,s$を選び、
\begin{equation}
 \theta(L,\lambda)=(\det L)^r\lambda^s
 \label{eq:pg-intrinsic-descent-character}
\end{equation}
とおく。次数$d$の線形化を$\theta^{-d}$でtwistすることにより、
周囲代数、中心三次直線とその生成イデアル、および中心商は、
\begin{equation}
 \mathcal U_N\subset\mathfrak P_{n,1}\times\mathbb A^1_\hbar,
 \qquad
 \mathfrak P_{n,1}\times\{0\}\subset\mathcal U_N
 \label{eq:pg-intrinsic-descended-open}
\end{equation}
上の実際の次数付き$\OO$-代数の族へ降りる。
各次数部分の局所自由性、左右正則な中心三次元による商、
Hilbert級数とAS条件、古典極限および一次Poisson括弧は、
同定理のatlas上の族から保たれる。
上記のgerbeもこの指標の選択により中立になる。

特に$N=5$では$\mathfrak P_{3,1}$は非零Poisson対の被約層全体であり、
\begin{equation}
 \theta(L,\lambda)=(\det L)^2\lambda^3,
 \qquad \dim\mathfrak P_{3,1}=17
 \label{eq:pg-intrinsic-threefold-descent}
\end{equation}
を用いることができる。
\end{corollary}
\begin{proof}
非零の分解可能な$P$に対して$\pi=\iota_{dF}P$は
三変数Jacobian型のPoisson構造であり、
正規形定理により一般階数二のPoisson構造はちょうどこの形である。
階数層の被約構造と非零Pl\"ucker錐との対応は、
命題\ref{prop:pg-equivariant-poisson-gerbe}の証明と同じである。
そこでは$F|_{\operatorname{supp}P}$の滑らかさを
対応する開部分を指定するためだけに用いたので、
その条件を外すとここでの全階数二の層を得る。

方程式の代表を忘れる作用
$(F,P)\mapsto(\lambda^{-1}F,\lambda P)$は自由である。
また$K_{\mathrm{sc}}$は$\mathcal S_N$に自明に作用し、
\begin{equation}
 1\longrightarrow\mathbf G_{m,\lambda}
 \longrightarrow\overline H_N
 \longrightarrow\mathrm{PGL}(V)\longrightarrow1
 \label{eq:pg-intrinsic-effective-group}
\end{equation}
は完全である。左の射は$\lambda\mapsto[(I_V,\lambda)]$である。
従って方程式の自由なスカラーを先に除き、次に
$\mathrm{PGL}(V)$で割ることで第一の商同型を得る。
$K_{\mathrm{sc}}$が中心にあり底に自明に作用することから、
第二の射はこの群をbandとするgerbeである
\cite[Lemma~78.27.1, Tag~\texttt{06PD}]{PStacks}。
定理\ref{thm:pg-intrinsic-rank-two-center}の
$H_N$-不変開と共変な代数・中心三次直線を商に移せば、
まずこのgerbe上の族を得る。

$k(c)=(cI_V,c^{-3})$と書けば$\theta(k(c))=c$である。
従って
\begin{equation}
 H_N\simeq K_{\mathrm{sc}}\times\ker\theta,
 \qquad
 h\longmapsto
 \bigl(k(\theta(h)),k(\theta(h))^{-1}h\bigr),
 \qquad
 \ker\theta\simeq\overline H_N
 \label{eq:pg-intrinsic-character-splitting}
\end{equation}
となり、gerbeは中立である。
元の代数の次数$d$部分には$k(c)$が$c^d$で作用するため、
$\theta^{-d}$によるtwist後にはこの作用が自明となる。
積の次数加法性により、これらのtwistは乗法と単位を保つ。
中心三次直線にも$\theta^{-3}$を用いると、
その包含と生成イデアルの乗法写像も降下する。
したがって商スタック上の同変準連接層の降下
\cite[Proposition~96.14.3, Tag~\texttt{06WT}]{PStacks}
を次数別に適用すれば、所要の代数族を得る。
この操作はatlas上の代数を変えないので、
局所自由性およびファイバーごとの主張は元の族から従う。

$N=5$では正規形に現れる互いに独立な三次元ブロックは一つだけであり、
非零Poisson構造の一般階数は必ず二である。
係数空間の次元は$35$、$C_3^\times(\CC^5)$の次元は$7$、
$\overline H_5$の次元は$25$なので、底スタックの次元は$17$である。
\end{proof}

\begin{remark}[降下の範囲と中心三次直線]
\label{rem:pg-intrinsic-descent-scope}
この系の底は$P\ne0$の階数二の層であり、零Poisson構造は含めていない。
零切断は変形パラメータ$\hbar=0$を意味する。
降下した次数一の束を$E_\theta$、中心三次直線の古典極限を
$\mathcal J_0\subset\operatorname{Sym}^3(E_\theta)$とすれば、
\begin{equation}
 \mathcal A^\theta|_{\hbar=0}\simeq\operatorname{Sym}(E_\theta),
 \qquad
 \mathcal B^\theta|_{\hbar=0}
 \simeq\operatorname{Sym}(E_\theta)/(\mathcal J_0)
 \label{eq:pg-intrinsic-descended-classical-fibre}
\end{equation}
であり、その相対$\operatorname{Proj}$は元の普遍三次超曲面族となる。
中心三次直線の一つの大域生成元を底スタック上で選ぶことは、
この降下の主張には含まれない。
また、各制限型で先に構成した明示中心元との一致、
零制限の全$\hbar$にわたる明示族との一致、
あるいは全$\mathbb A^1_\hbar$への延長も主張していない。
$3\mid N$の場合、$H_N$の指標の$K_{\mathrm{sc}}$上の重みは
$Nr-3s\in3\mathbb Z$なので、上記の重み一の指標は存在しない。
命題\ref{prop:pg-intrinsic-gerbe-period}では、指標以外の線束による補正も含めて
このgerbeの非中立性を示す。一方、系\ref{cor:pg-intrinsic-third-veronese-descent}の
三次Veroneseの降下は全ての$N$で成立する。
\end{remark}

\begin{proposition}[全階数二層のgerbeの位数]
\label{prop:pg-intrinsic-gerbe-period}
$N\ge5$とし、\eqref{eq:pg-intrinsic-descent-base}の記法を用いる。
$K_{\mathrm{sc}}\simeq\mathbf G_m$の同一視は
$c\mapsto(cI_V,c^{-3})$で固定する。
\begin{equation}
 \mathfrak G_N=[\mathcal S_N/H_N],\qquad
 \alpha_N=[\mathfrak G_N\to\mathfrak P_{n,1}]
 \in H^2_{\mathrm{fppf}}(\mathfrak P_{n,1},\mathbf G_m)
 \label{eq:pg-intrinsic-gerbe-class}
\end{equation}
とおく。このとき
\begin{equation}
 \begin{aligned}
 \operatorname{Pic}(\mathcal S_N)&=0,\\
 \operatorname{im}\left(
 \operatorname{wt}:\operatorname{Pic}(\mathfrak G_N)\to\mathbb Z
 \right)&=N\mathbb Z+3\mathbb Z=\gcd(N,3)\mathbb Z,\\
 \operatorname{ord}(\alpha_N)&=\gcd(N,3).
 \end{aligned}
 \label{eq:pg-intrinsic-gerbe-period}
\end{equation}
ここで$\operatorname{wt}(\mathcal L)=r$は、$k(c)\in K_{\mathrm{sc}}$が
$\mathcal L$に$c^r$で作用することを意味する。
特に$3\mid N$ならgerbeは非中立で、位数は正確に三である。

さらに$3\mid N$とし、
$\mathfrak P_{n,1}\times\{0\}$全体を含む任意の開
$\mathcal U\subset\mathfrak P_{n,1}\times\mathbb A^1_\hbar$を取る。
この開上の引戻しgerbe
\begin{equation}
 \mathfrak G_{N,\mathcal U}
 = (\mathfrak G_N\times\mathbb A^1_\hbar)
     \times_{\mathfrak P_{n,1}\times\mathbb A^1_\hbar}\mathcal U
 \label{eq:pg-intrinsic-gerbe-neighborhood}
\end{equation}
の位数も三であり、重み$1$または$-1$の可逆層は存在しない。
従って定理\ref{thm:pg-intrinsic-rank-two-center}の族を、
各次数$d$で一つの可逆層の$d$乗をtensorする方法によって、
元の次数一の構造を保った代数族へ降ろすことはできない。
\end{proposition}

\begin{proof}
\emph{係数空間のPicard群。}
三次式の係数を$\boldsymbol a$、判別式を$\Delta$と書き、
\begin{equation}
 R_N=\CC[\boldsymbol a,\Delta^{-1}],\qquad
 X=\mathcal T_N\times\operatorname{Mat}_{N\times3},\qquad
 Y=\mathcal T_N\times\operatorname{St}(3,V)
 \label{eq:pg-intrinsic-gerbe-stiefel-cover}
\end{equation}
とおく。$X$の座標環$R_N[A_{i\beta}]$はUFDであり、
$X\setminus Y$の余次元は$N-2\ge2$である。
従って
\begin{equation}
 \Gamma(Y,\OO)=R_N[A_{i\beta}],\qquad
 \Gamma(Y,\OO)^\times=R_N^\times,\qquad
 \operatorname{Pic}(Y)=0.
 \label{eq:pg-intrinsic-gerbe-stiefel-picard}
\end{equation}
関数の等号は\eqref{eq:pg-intrinsic-rank-two-stiefel-functions}と同じ
余次元二以上の延長による。
Picard群については、$Y$上の可逆層の非零有理切断を一つ選ぶ。
その因子の各余次元一成分を$X$へ閉包するとWeil因子となり、
$X$がfactorialなので主因子である。
$Y$への制限から元の可逆層が自明であることが従う。

列の外積を取る写像
\begin{equation}
 p:Y\longrightarrow\mathcal S_N,\qquad
 (F,A)\longmapsto(F,\textstyle\bigwedge^3A)
 \label{eq:pg-intrinsic-gerbe-sl3-torsor}
\end{equation}
はZariski局所自明な$\mathrm{SL}_3$-torsorである。
その作用は行列$A$の右からの作用である。
従って$\mathcal S_N$上の可逆層$\mathcal L$の引戻しは
$\mathrm{SL}_3$-線形化を持つ自明束になる。
この線形化を自明化のもとで表すと、
$\mathrm{SL}_3\times Y$上の単元$u$で、$u(e,y)=1$を満たすものを得る。

ここで$\mathrm{SL}_3\times Y$の単元を直接計算する。
多様体としての同型
\begin{equation}
 \mathrm{SL}_3\times\mathbf G_m\xrightarrow{\ \sim\ }\mathrm{GL}_3,
 \qquad (g,t)\longmapsto g\operatorname{diag}(t,1,1)
 \label{eq:pg-intrinsic-gerbe-gl3-variety}
\end{equation}
を用いる。逆写像は
$B\mapsto(B\operatorname{diag}((\det B)^{-1},1,1),\det B)$である。
一方、$\mathrm{GL}_3\times Y$もfactorialなアフィンスキームから
余次元二以上の閉集合を除いた開であり、その単元は
$v(F)(\det B)^j$、$v(F)\in R_N^\times$、$j\in\mathbb Z$だけである。
$\mathrm{SL}_3\times Y$の単元を$\mathrm{SL}_3\times\mathbf G_m\times Y$へ
引き戻すと$t$に依存しないので、上の表示では$j=0$である。
従って
\begin{equation}
 \Gamma(\mathrm{SL}_3\times Y,\OO)^\times
      =\operatorname{pr}_Y^*R_N^\times.
 \label{eq:pg-intrinsic-gerbe-sl3-units}
\end{equation}
$u(e,y)=1$から$u=1$となる。
すなわち$p^*\mathcal L$の自明化は$\mathrm{SL}_3$-不変であり、
torsorに沿って降下して$\mathcal L$を自明化する。
これで$\operatorname{Pic}(\mathcal S_N)=0$を得る。

\emph{同変可逆層と中心重み。}
$\mathfrak G_N$上の可逆層は、$\mathcal S_N$上の
$H_N$-同変可逆層と同じである。
いま証明したPicard群の消滅により、その下にある可逆層を自明化できる。
線形化は$H_N\times\mathcal S_N$上の単元$\varepsilon$で表され、
$\varepsilon(e,s)=1$を満たす。
これを$\operatorname{id}_{H_N}\times p$で$H_N\times Y$へ引き戻す。
この操作には、$H_N$の作用を$Y$へ持ち上げる必要はない。

$H_N\times Y$の関数は
$R_N[A_{i\beta},L_{ij},(\det L)^{-1},\lambda,\lambda^{-1}]$であり、
その単元は$R_N$の単元、$\det L$および$\lambda$の積である。
従って引戻した線形化は
\begin{equation}
 \varepsilon((L,\lambda),p(F,A))
       =v(F)(\det L)^a\lambda^b,
 \qquad v(F)\in R_N^\times,\quad a,b\in\mathbb Z
 \label{eq:pg-intrinsic-gerbe-linearization-character}
\end{equation}
と書ける。恒等元での条件により$v(F)=1$である。
$p$は忠実平坦なので、元の線形化も指標
$(\det L)^a\lambda^b$で与えられる。
自明化を変えると指標の表示が変わり得るが、
$K_{\mathrm{sc}}$は底に自明に作用するため中心重みは変わらない。
逆に任意の指標は自明束の線形化を定め、その中心重みは
\begin{equation}
 (\det(cI_V))^a(c^{-3})^b=c^{Na-3b}
 \label{eq:pg-intrinsic-gerbe-character-weight}
\end{equation}
である。従って中心重みの像はちょうど$N\mathbb Z+3\mathbb Z$となる。

\emph{gerbeの位数。}
一般に代数スタック$M$上の$\mathbf G_m$-gerbe
$\mathfrak G\to M$の類を$\alpha$とすると、
\begin{equation}
 \mathfrak G\text{上に重み$r$の可逆層が存在する}
 \quad\Longleftrightarrow\quad
 r\alpha=0\text{ in }H^2_{\mathrm{fppf}}(M,\mathbf G_m).
 \label{eq:pg-intrinsic-gerbe-weight-period-criterion}
\end{equation}
gerbeのコホモロジー類と重みの一般論について
\cite[Theorem~3.11, Corollary~4.11]{PBerghSchnurerGerbes}を参照できるが、
可逆層についての同値は代数空間を底とする場合に
\cite[Lemma~3.1.1.8, p.~7]{PLieblichPeriodIndex}にも明記されている。
ここでは次の\v{C}ech表示により、代数スタックを底とする場合まで確認する。
必要ならfppf hypercoverを細分し、gerbeの局所対象とそれらの間の同型を取る。
三重の重なりでの合成のずれを$c_{ijk}\in\mathbf G_m$とすると、
これは$\alpha$を表す二次cocycleである。
重み$r$の可逆層を局所的に自明化したとき、その遷移関数$b_{ij}$の条件は
\begin{equation}
 b_{ij}b_{jk}b_{ik}^{-1}=c_{ijk}^{\,r}
 \label{eq:pg-intrinsic-gerbe-weight-cocycle}
\end{equation}
となる（cocycleの向きはこの式に合わせて固定する）。
従ってその存在は$c^r$がcoboundaryであることと同値である。
逆向きにも、この式を満たす$b_{ij}$で局所的な重み$r$の直線を貼り合わせればよい。
この議論は基底がスキームであることを要求しない。

\eqref{eq:pg-intrinsic-gerbe-weight-period-criterion}と
中心重みの像の計算により、$\alpha_N$の位数は$\gcd(N,3)$である。
特に$3\mid N$では重み一の可逆層が存在しないので$\alpha_N\ne0$であり、
指標$\lambda^{-1}$が重み三の可逆層を与えるので$3\alpha_N=0$である。

最後に零切断近傍を考える。
$\mathfrak G_{N,\mathcal U}$上に重み一または負一の可逆層があれば、
零切断への制限が$\mathfrak G_N$上に同じ重みの可逆層を与え、矛盾する。
引戻し類は三で消え、その零切断への制限は非零なので、位数も三のままである。
元の周囲代数の次数一部分には$k(c)$が$c$で作用する。
従って$\mathcal A_d\otimes\mathcal M^{\otimes d}$の形の補正で
中心スカラーを消すには$\operatorname{wt}(\mathcal M)=-1$が必要であり、
そのような補正は存在しない。
\end{proof}

\begin{remark}[障害の射程]
\label{rem:pg-intrinsic-gerbe-obstruction-scope}
この命題は、全階数二の被約底と、その全零切断を含む変形近傍に関する主張である。
前節の節点・尖点の各局所閉層では重み一の可逆層によってgerbeが中立になるが、
その可逆層は全階数二底へ同じ重みの可逆層としては延長できない。
三次Veroneseの各次数で必要となる重み$-3d$は$\lambda^d$で消せるため、
その降下は本命題の障害を受けない。
ここで否定したのはgerbeの中立化と、既存の次数一の族を可逆層のtensorで補正する降下である。
別の量子化された周囲代数、Morita同値な模型、またはカテゴリーの降下の不存在は結論しない。
\end{remark}

\begin{corollary}[全次元での三次Veroneseの降下]
\label{cor:pg-intrinsic-third-veronese-descent}
$N\ge5$とし、定理\ref{thm:pg-intrinsic-rank-two-center}の族を取る。
$\gcd(N,3)$によらず、全被約階数二層の零切断近傍
$\mathcal U_N=[\mathcal W_N/\overline H_N]$上に、
次数別有限局所自由な次数付き代数
$\mathcal A^{[3]}$、$\mathcal B^{[3]}$と、
$\mathcal A^{[3]}$の次数一の大域的な左右正則中心元$\omega$が存在し、
gerbe上への引戻しは
\begin{equation}
 \begin{aligned}
 \mathcal A^{[3]}_d&=\mathcal A_{3d}\otimes\lambda^d,&
 \mathcal B^{[3]}_d&=\mathcal B_{3d}\otimes\lambda^d
 \qquad(d\ge0),\\
 \omega&=\Omega\otimes1,&
 \mathcal B^{[3]}&\simeq\mathcal A^{[3]}/(\omega)
 \end{aligned}
 \label{eq:pg-intrinsic-third-veronese}
\end{equation}
で与えられる。
その古典極限の相対$\operatorname{Proj}$は、
それぞれ普遍の周囲射影空間族と普遍三次超曲面族であり、
新しい次数一の偏極は各幾何学的ファイバー上で$\OO(3)$となる。

$\gcd(N,3)=1$で系\ref{cor:pg-intrinsic-rank-two-descent}の
指標$\theta$を選んだ場合、$\chi=\lambda\theta^3$は
$\overline H_N$の指標となる。その対応する底上の可逆層を
$\mathcal L_\chi$とすると、
\begin{equation}
 \mathcal A^{[3]}_d
   \simeq(\mathcal A^\theta)_{3d}\otimes\mathcal L_\chi^{\otimes d},
 \qquad
 \mathcal B^{[3]}_d
   \simeq(\mathcal B^\theta)_{3d}\otimes\mathcal L_\chi^{\otimes d}.
 \label{eq:pg-intrinsic-veronese-character-comparison}
\end{equation}
従って二つの降下は、三次Veroneseを取り、底の線束の次数冪で捩る操作によって関係づけられる。
\end{corollary}
\begin{proof}
$k(c)=(cI_V,c^{-3})$の$\mathcal A_{3d}$上の作用は$c^{3d}$であり、
$\lambda^d$上の作用は$c^{-3d}$である。
従って両者の積は重み零である。
積の写像では次数が加わるので、これらのtwistは単位と乗法を保つ。
gerbe上の重み零の準連接層の降下
\cite[Proposition~4.10]{PBerghSchnurerGerbes}を次数別に適用すると、
$\mathcal A^{[3]}$と$\mathcal B^{[3]}$を得る。

共変性\eqref{eq:pg-intrinsic-rank-two-covariance}から
$\lambda\Phi(\Omega')=\Omega$であり、$\omega$は実際の大域切断として降りる。
元の中心元の正則性により、各$d\ge1$において
\begin{equation}
 0\longrightarrow\mathcal A^{[3]}_{d-1}
 \xrightarrow{\ \omega\ }\mathcal A^{[3]}_d
 \longrightarrow\mathcal B^{[3]}_d\longrightarrow0
 \label{eq:pg-intrinsic-veronese-quotient-sequence}
\end{equation}
は完全である。これは同変な列
$0\to\mathcal J\otimes\mathcal A_{3d-3}\to\mathcal A_{3d}
\to\mathcal B_{3d}\to0$に$\lambda^d$をかけ、
$\mathcal J\otimes\lambda\simeq\OO$を用いた列である。
従って左端は$\mathcal A_{3d-3}\otimes\lambda^{d-1}=\mathcal A^{[3]}_{d-1}$となる。
従って\eqref{eq:pg-intrinsic-third-veronese}の商の同定と局所自由性を得る。

古典極限ではatlas上で
\begin{equation}
 \bigl(\operatorname{Sym}(V^*)^{(3)}\bigr)_d
     =\operatorname{Sym}^{3d}(V^*),\qquad
 \bigl((\operatorname{Sym}(V^*)/(F))^{(3)}\bigr)_d
     =(\operatorname{Sym}(V^*)/(F))_{3d}
 \label{eq:pg-intrinsic-veronese-classical}
\end{equation}
である。Veroneseは相対$\operatorname{Proj}$を変えず、
底の可逆層の次数冪によるtwistも相対$\operatorname{Proj}$を変えない。
このことは線束を自明化した開上の同定が、重なりでの次数$d$の
スカラー倍によって保たれることから従う。
従ってこれらの同定は商上へ貼り合い、所要の普遍族と偏極を与える。
これは周囲の射影空間族が全体としてベクトル束の射影化になることを仮定しない。

最後に$\theta(k(c))=c$なので
$\chi(k(c))=c^{-3}c^3=1$である。
$\theta^{-3d}\chi^d=\lambda^d$という等式が
\eqref{eq:pg-intrinsic-veronese-character-comparison}を与える。
\end{proof}

\begin{remark}[Veroneseと射影的な圏の範囲]
\label{rem:pg-intrinsic-veronese-qgr}
三次Veroneseの降下は、元の次数一の生成元を降ろすこととは区別する。
特に$3\mid N$の場合、命題\ref{prop:pg-intrinsic-gerbe-period}の障害は残る。
ここではVeronese代数のAS-regular性を主張しない。
AS条件とHilbert級数については、元の代数の結論とその三次ごとの部分を区別する。

各幾何学的ファイバーで、元の代数$R$が左右noetherianである場合には、
既知のVeronese同値により
\begin{equation}
 qgr\,R\simeq qgr\,R^{(3)},\qquad
 \pi M\longmapsto\pi(M^{(3)})
 \label{eq:pg-intrinsic-veronese-qgr}
\end{equation}
を得る。これは$R=A$にも$R=B$にも、それぞれのnoetherian性の下で適用する。
用いるのは次数一生成の場合の標準結果であり
\cite[Theorem~4.1, Proposition~4.8, pp.~2934, 2940--2941]{PSmithVeronese}、
一般の次数付き環の射についての議論には同文献の訂正文の条件が必要である。

この場合の同値と必要な完全性は直接にも確認できる。
$T=R^{(3)}$も左右noetherianである。例えば$T$の斉次右イデアル列$J_i$について、
$J_iR$の安定と$(J_iR)^{(3)}=J_i$から安定性が従い、左側も同様である。
右加群について
$V(M)=M^{(3)}$、$I(N)=N\otimes_T R$とする。
$I$では$N_d$を元の次数$3d$に置く。
すると$V$は完全で$VI\simeq\mathrm{id}$であり、
$IV(M)\to M$の核と余核は$V$によって零になる。
$R$は次数一生成なので、$V(M)=0$なら各斉次元は$R_{\ge3}$に殺される。
従ってこれらの核と余核はtorsionである。
また$T$上の自由分解に$I$を施した複体の$V$像は元の分解なので
\begin{equation}
 V\bigl(\operatorname{Tor}^T_1(N,R)\bigr)=0.
 \label{eq:pg-intrinsic-veronese-tor}
\end{equation}
従って$\pi I$はtorsionを除いた圏で完全となり、訂正文で要求される条件も満たす。
$R$は左右とも$T$上で次数$0,1,2$部分により有限生成されるため、
両関手は有限生成加群を保ち、torsion部分を対応させる。
これが\eqref{eq:pg-intrinsic-veronese-qgr}を与える。

底の線束$\mathcal L$による$R_d\mapsto R_d\otimes\mathcal L^{\otimes d}$のtwistは、
$M_d\mapsto M_d\otimes\mathcal L^{\otimes d}$による次数付き加群の同値を持つ。
これは既知のsheaf $\mathbb Z$-algebraのtwistの特別な場合である
\cite[arXiv v3, \S4.2, p.~9]{PVdBNoncommP1Bundles}。
本節の一般原理への追加は、全support上の中心元に対し、
全階数二層でこの三次偏極の代数族を同時に構成した点である。
この記述だけから、全閉パラメータの$qgr$の滑らかさや、
異なる中心直線から得られる圏の同値は結論しない。
\end{remark}

\subsection{零制限の標準中心直線との整合性}

\begin{proposition}[零制限の中心直線への大域的一致の障害]
\label{prop:pg-intrinsic-zero-matching-obstruction}
$N\ge6$、$m=N-3$とし、
定理\ref{thm:pg-intrinsic-rank-two-center}の底
$\mathcal S_N=\mathcal T_N\times C_3^\times(V)$上で
\begin{equation}
 \mathcal Z_N=
 \{(F,P)\in\mathcal S_N:F|_{W_P}=0\},\qquad
 \mathcal C=\operatorname{Sym}^3(V/W_P)^*
 \label{eq:pg-intrinsic-zero-incidence}
\end{equation}
と書く。$\mathcal Z_N$は非空の滑らかな余次元十の閉部分である。
この上では定理\ref{thm:pg-equivariant-zero-center}の中心三次元
$\Omega^{\mathrm z}=\sym(F)$が全$\hbar$について定義される。

$\mathcal S_N$の全$\hbar=0$切断の近傍で定義された中心三次元
$\Omega$が、Weyl表示で
\begin{equation}
 \Omega=\sym\bigl(F+\hbar^2\omega_2+O(\hbar^4)\bigr)
 \label{eq:pg-intrinsic-zero-comparison-jet}
\end{equation}
を満たすとする。
このとき、$\mathcal Z_N\times\mathbb A^1_\hbar$の零切断に沿う形式完備化上で、
$\OO\Omega$を$\OO\sym(F)$と一致させることはできない。
従って、この零切断全体を含む近傍での一致も不可能である。
この非存在は、全support上の形式的中心元にも成立し、
$\Omega$の群同変性を仮定しない。
従って中心生成元そのものの一致も不可能である。
\end{proposition}

\begin{proof}
\emph{閉部分と非零のHessian切断。}
$\mathbf G=\operatorname{Gr}(3,V)$上の普遍部分束を$\mathcal W$とする。
全三次式の係数空間と$\mathbf G$の積上の制限写像
\begin{equation}
 \operatorname{Sym}^3V^*\otimes\OO_{\mathbf G}
 \longrightarrow\operatorname{Sym}^3\mathcal W^*
 \label{eq:pg-intrinsic-zero-restriction-surjection}
\end{equation}
は階数十の束への全射である。
その核の全空間に滑らかさの開条件を課し、
$C_3^\times(V)\to\mathbf G$へ引き戻したものが
$\mathcal Z_N$である。
従ってこれは$\mathcal S_N$内の滑らかな余次元十の閉部分となる。
また$P\mapsto sP$、$s\in\mathbf G_m$で安定である。

内在的Hessian三次式$\mathscr H(P,F)$を
\eqref{eq:pg-ad-intrinsic-hessian}の縮約で定める。
その$\mathcal Z_N$への制限を$\mathscr H_{\mathrm z}$と書く。
supportに適合する座標では
\begin{equation}
 \begin{aligned}
 P&=\eta\,\partial_{x_1}\wedge\partial_{x_2}\wedge\partial_{x_3},\\
 F&=\tfrac12 x^{\mathsf t}H(u)x+b(u)^{\mathsf t}x+c(u),\\
 \mathscr H_{\mathrm z}&=\eta^2\det H(u)
       \in\mathcal C|_{\mathcal Z_N}.
 \end{aligned}
 \label{eq:pg-intrinsic-zero-hessian-pure-central}
\end{equation}
従って$\mathscr H_{\mathrm z}$は純中心三次式の束の大域切断であり、
$P$のスカラー倍について重み二を持つ。

この切断は零ではない。実際、$m\ge3$に対して
\begin{equation}
 F_0=\sum_{i=1}^3(u_ix_i^2+u_i^3)
          +\sum_{\alpha=4}^m u_\alpha^3,\qquad
 P_0=\partial_{x_1}\wedge\partial_{x_2}\wedge\partial_{x_3}
 \label{eq:pg-intrinsic-zero-obstruction-example}
\end{equation}
と取る。
全偏微分の共通零点では
$2u_ix_i=x_i^2+3u_i^2=0$から$x_i=u_i=0$となり、
$\alpha>3$では$3u_\alpha^2=0$から$u_\alpha=0$となる。
従って$F_0$の射影超曲面は滑らかである。
また$F_0|_{W_{P_0}}=0$かつ
\begin{equation}
 \mathscr H(P_0,F_0)=8u_1u_2u_3\ne0.
 \label{eq:pg-intrinsic-zero-nonzero-hessian}
\end{equation}
この例は$\mathcal Z_N$の非空性も示す。
なお$N=5$では残りの二つの偏微分は$\PP(W_P)=\PP^2$上の二次式であり、
共通零点を持つため、滑らかさと両立しない。
従って本節の$N\ge5$の範囲では$\mathcal Z_N$が非空であることと
$N\ge6$は同値である。

\emph{純中心切断の重み二成分の消滅。}
$\mathcal Q=(V\otimes\OO_{\mathbf G})/\mathcal W$とおく。
$\mathcal C$は$\operatorname{Sym}^3\mathcal Q^*$の引戻しである。
ここで$P$のスカラー作用の線形化は、$W_P$を変えない作用として
$\mathcal C$の各ファイバーを恒等的に同定するものを用いる。
重み$j$とは、切断の係数が$P\mapsto sP$で$s^j$倍となることを意味する。
\eqref{eq:pg-rank-two-cone-direct-images}と$\mathcal T_N$のアフィン性から
\begin{equation}
 \Gamma(\mathcal S_N,\mathcal C)_j
 =\Gamma(\mathcal T_N,\OO)\otimes_\CC
 H^0\left(\mathbf G,
     \operatorname{Sym}^3\mathcal Q^*\otimes\OO_{\mathbf G}(j)\right).
 \label{eq:pg-intrinsic-pure-center-weight-sections}
\end{equation}
補題\ref{lem:pg-rank-two-bott-vanishing}の証明で用いた
$j=2$のBott列は
\begin{equation}
 (N+1,N,N-1,N-1,N-5,N-6,\ldots,0)
 \label{eq:pg-intrinsic-pure-center-weight-two-bott}
\end{equation}
であり、$N-1$が重複する。
そのため$H^1$だけでなく全コホモロジーが消え、特に
\begin{equation}
 \Gamma(\mathcal S_N,\mathcal C)_2=0.
 \label{eq:pg-intrinsic-pure-center-weight-two-zero}
\end{equation}
この結論は$F$の全係数への正則な依存を許したものである。

\emph{任意の中心元の二次係数。}
全零切断近傍の正則なWeyl係数は、$\hbar$進展開の各次数で
$\mathcal S_N$上の大域切断を与える。
命題\ref{prop:pg-universal-casimir-jet}をsupportに適合する局所座標で適用し、
$t=\eta\hbar$と置くと
\begin{equation}
 D_0\left(\omega_2+\frac1{12}\mathscr H(P,F)\right)=0
 \label{eq:pg-intrinsic-any-center-second-coefficient}
\end{equation}
を得る。$D_0$は古典Jacobian括弧による作用素である。
Hessian縮約と$D_0$は座標変更と両立するため、これは大域の等式である。
さらに古典核の\emph{束としての}分裂
\eqref{eq:pg-intrinsic-rank-two-classical-splitting}から、一意な大域切断
\begin{equation}
 a\in\Gamma(\mathcal S_N,\OO),\qquad
 c\in\Gamma(\mathcal S_N,\mathcal C),\qquad
 \omega_2=-\frac1{12}\mathscr H(P,F)+aF+c
 \label{eq:pg-intrinsic-any-center-second-decomposition}
\end{equation}
を得る。点ごとの比例係数を選んでいるのではない。
$F$は滑らかさにより純中心式ではなく、
$\OO F\oplus\mathcal C$という大域直和の二つの射影を用いている。

仮に$\mathcal Z_N$に沿って中心直線が一致すれば、形式的には
$\Omega|_{\mathcal Z_N}=b(\hbar)\sym(F)$、$b(0)=1$となる単元が存在する。
その$\hbar^2$係数を$b_2$とすれば
\begin{equation}
 \omega_2|_{\mathcal Z_N}=b_2F.
 \label{eq:pg-intrinsic-line-matching-second-coefficient}
\end{equation}
式\eqref{eq:pg-intrinsic-any-center-second-decomposition}を制限し、
$\OO F\oplus\mathcal C$の直和で比較すると
\begin{equation}
 c|_{\mathcal Z_N}=\frac1{12}\mathscr H_{\mathrm z}.
 \label{eq:pg-intrinsic-zero-impossible-extension}
\end{equation}
制限写像は$P$のスカラー作用と可換し、
従って各重みへの射影とも可換する。
右辺は非零の重み二切断であるが、左辺の重み二成分は
\eqref{eq:pg-intrinsic-pure-center-weight-two-zero}により零である。
これは矛盾である。

以上は形式的二次jetだけを用いるため、定義開を$\hbar$方向に縮めても変わらない。
同変性を仮定せず、任意の大域切断$c$の重み二成分を取った点にも注意されたい。
\end{proof}

\begin{remark}[相対コホモロジーと中心直線の選択]
\label{rem:pg-intrinsic-zero-relative-obstruction}
$\mathcal Z_N$のイデアル層を$\mathcal I_{\mathrm z}$とすると、完全列
\begin{equation}
 0\longrightarrow\mathcal I_{\mathrm z}\mathcal C
 \longrightarrow\mathcal C
 \longrightarrow i_*\bigl(\mathcal C|_{\mathcal Z_N}\bigr)
 \longrightarrow0
 \label{eq:pg-intrinsic-zero-relative-sequence}
\end{equation}
の接続写像$\partial$について
\begin{equation}
 0\ne
 \partial(\mathscr H_{\mathrm z}/12)
 \in H^1(\mathcal S_N,\mathcal I_{\mathrm z}\mathcal C)_2
 \subset H^1(\mathcal S_N,\mathcal I_{\mathrm z}\mathcal K_0)_2
 \label{eq:pg-intrinsic-zero-relative-class}
\end{equation}
である。
実際、左端の非零性は
$\Gamma(\mathcal S_N,\mathcal C)_2=0$と
$\mathscr H_{\mathrm z}\ne0$から完全性によって従う。
最後の包含は$\mathcal K_0=\OO F\oplus\mathcal C$による直和因子である。
従って、定理\ref{thm:pg-intrinsic-rank-two-center}に用いた
$H^1(\mathcal S_N,\mathcal K_0)=0$を、
零制限上の値を指定する相対消滅へ置き換えることはできない。

一つの大域中心元で
$\mathcal K=\mathcal C\oplus\OO\Omega$と分裂させても、
$\mathbf P_{\mathrm{lin}}(\mathcal K)\setminus
\mathbf P_{\mathrm{lin}}(\mathcal C)$上の所定の閉部分の切断が
大域切断へ延長するとは限らない。
この開部分は$\mathcal Hom(\mathcal K/\mathcal C,\mathcal C)$のaffine torsorであり、
$\Omega$による自明化では中心直線は$\OO(\Omega+\gamma)$、
$\gamma\in\mathcal C$と書ける。
零制限で$\Omega=a_{\mathrm z}\sym(F)+c_{\mathrm z}$と書くと、
指定直線に合わせる条件は$\gamma|_{\mathcal Z_N}=-c_{\mathrm z}$である。
古典極限と偶性からこの差は$\hbar^2$で割れるが、
割った差の最初の係数の接続類が
\eqref{eq:pg-intrinsic-zero-relative-class}となる。
$\tau=\hbar^2$に関するdilatationは、この補正係数を追加データとして記録する。
それだけで所定の切断の延長を与えるものではない。

本命題が排除するのは、指定された周囲代数の同一視の下での中心直線の一致である。
次数三の部分で比較すれば、その生成する斉次中心イデアルの一致も排除される。
一方、中心商の抽象的な非同型や$qgr$圏の非同値は結論していない。
固定supportでの係数空間上の選択や、supportの適切な開部分での選択とも矛盾しない。
\end{remark}

\subsection{中心調整による方法の限界}

\begin{proposition}[全係数空間における中心調整の障害]
\label{prop:pg-equivariant-global-centering-obstruction}
$m\ge1$、$N_V=m+3$、$V=W\oplus Z$、$\dim W=3$、$\dim Z=m$とし、
$\mathcal T_m\subset\operatorname{Sym}^3V^*$を滑らかな全三次式の係数空間とする。
群$G$と$f$への作用は\eqref{eq:pg-equivariant-groups}のものとし、
\begin{equation}
 \mathcal S_m=\mathcal T_m\times\mathbf G_{m,\eta},\qquad
 \eta'=\lambda\det(M)\eta,\qquad \hbar'=\hbar
 \label{eq:pg-global-centering-parameter-action}
\end{equation}
とする。次が成り立つ。
\begin{enumerate}
\item $\mathcal S_m\times\{0\}$全体を含む$G$-不変開
$\mathcal U\subset\mathcal S_m\times\mathbb A^1_\hbar$上に、
正則写像$B:\mathcal U\to\operatorname{Hom}(Z,W)$で
\begin{equation}
 B(f',\eta',\hbar)=(M B(f,\eta,\hbar)-C)N^{-1}
 \label{eq:pg-global-centering-desired-covariance}
\end{equation}
を満たすものは存在しない。
従って、混合係数を含む$f$全体に依存させても、
$\mathcal T_m$全体の中心調整$B_f$はこの変換則を満たすように選べない。
同じ非存在は、$t'=\lambda\det(M)t$と変換する$t$に依存し、
$t=0$全体の近傍で正則である選択にも適用される。
\item 中心平行移動の加法群
\begin{equation}
 U=\operatorname{Hom}(Z,W),\qquad
 u_C\cdot f=f(x-Cu,u)
 \label{eq:pg-global-centering-unipotent-action}
\end{equation}
は$\mathcal T_m$に自由に作用する。
商$Y=[\mathcal T_m/U]$は代数空間で、
$\mathcal T_m\to Y$は非自明な$U$-torsorであり、$Y$はアフィンではない。
\end{enumerate}
本命題は、全ブロック群に共変な中心直線や中心イデアルの不存在を主張しない。
定理\ref{thm:pg-intrinsic-rank-two-center}は中心核を直接扱うことで、この障害を回避する。
\end{proposition}
\begin{proof}
第一の主張について、仮に$B$が存在すれば、$\hbar=0$に制限して
$\mathcal S_m$全体で\eqref{eq:pg-global-centering-desired-covariance}を満たす写像を得る。
滑らかな三次式$F_0\in\operatorname{Sym}^3V^*$を一つ固定し、
Stiefel開部分
\begin{equation}
 \operatorname{St}(3,V)
 =\{A\in\operatorname{Mat}_{N_V\times3}:\operatorname{rank}A=3\}
 \label{eq:pg-global-centering-stiefel-open}
\end{equation}
を考える。
その上で局所的に$A$の列を基底へ延ばし、
$T_A=(A,D)\in\mathrm{GL}(V)$、$D\in\operatorname{Mat}_{N_V\times m}$と書く。
例えば非零の三次小行列を固定した開では、残りの列に所定の座標ベクトルを選べる。
\begin{equation}
 f_A(x,u)=F_0(T_A(x,u)),\qquad
 B_A=B(f_A,1,0),\qquad
 R_A=(I_3\ \ B_A)T_A^{-1}
 \label{eq:pg-global-centering-stiefel-retraction}
\end{equation}
と置けば、$R_A$は局所的に正則な$3\times N_V$行列であり、$R_AA=I_3$である。

基底への延長を変えると、重なり上で
\begin{equation}
 T'_A=T_Ap,\qquad
 p=\begin{pmatrix}I_3&C_0\\0&N_0\end{pmatrix},\qquad
 N_0\in\mathrm{GL}(Z)
 \label{eq:pg-global-centering-frame-transition}
\end{equation}
となる。
$f'_A=f_A\circ p$は、\eqref{eq:pg-equivariant-groups}における
群要素$(p^{-1},1)$の作用である。
この要素のactive部分は$I_3$なので、$\eta=1$と$\hbar=0$は変わらない。
従って\eqref{eq:pg-global-centering-desired-covariance}から
\begin{equation}
 B'_A=B_AN_0+C_0,\qquad
 (I_3\ \ B'_A)p^{-1}=(I_3\ \ B_A),\qquad R'_A=R_A
 \label{eq:pg-global-centering-left-inverse-gluing}
\end{equation}
を得る。
こうして$R_A$は$\operatorname{St}(3,V)$全体の正則行列$R(A)$へ貼り合い、
\begin{equation}
 R(A)A=I_3
 \label{eq:pg-global-centering-impossible-left-inverse}
\end{equation}
を満たす。
これは、非零三重ベクトル$\bigwedge^3A$に依存する補空間選択を、
その枠$A$に引き戻した左逆である。

しかし、階数二以下の行列の閉集合の余次元は
\begin{equation}
 \operatorname{codim}_{\operatorname{Mat}_{N_V\times3}}
       \{\operatorname{rank}A\le2\}=N_V-2=m+1\ge2
 \label{eq:pg-global-centering-rank-boundary-codimension}
\end{equation}
である。
実際、二本の独立な列を固定するchartでは、その二本と第三の列の二つの係数が
自由に動くので、階数二の部分の次元は$2N_V+2$となる。
階数一以下の部分はこれより小さい。
多項式環はUFDであり、余次元二以上の閉集合を除いた開の正則関数は全空間へ延長する。
具体的には、その有理式の分母に残る既約因子があれば、対応する余次元一の因子は
この開と交わるため、そこに極が生じて正則性に反する。
従って$R(A)$の全成分が多項式へ延長する。
\eqref{eq:pg-global-centering-impossible-left-inverse}は稠密開で成立するので
全行列空間上の多項式恒等式となるが、$A=0$への代入で$0=I_3$となり矛盾する。
この議論は$\eta$への依存を許したままであり、$\hbar=0$での制限だけを用いた。
変数$t$を使う場合は$t=\eta\hbar$によってこの場合へ帰着する。

次に$U$作用の自由性を示す。
滑らかな$f$の全偏微分$f_{v_1},\ldots,f_{v_{N_V}}$は二次式で、
その共通零点は錐の原点だけである。
従ってこれらは多項式環の正則列をなし、Koszul完全性により
一次係数のsyzygyを持たない。すなわち
\begin{equation}
 \sum_{i=1}^{N_V} a_i(v)f_{v_i}=0,\quad \deg a_i=1
 \quad\Longrightarrow\quad a_1=\cdots=a_{N_V}=0
 \label{eq:pg-global-centering-linear-syzygy-vanishing}
\end{equation}
である。
この写像は各閉点で単射な有限階数の行列なので、係数空間上で局所的に分裂し、
任意の底変更の後にも単射である。
$D_C=\sum_{i=1}^3(Cu)_i\partial_{x_i}$と書くと、
三次式上で$D_C^4=0$であり、$u_C$の作用は$\exp(-D_C)$である。
$u_C\cdot f=f$なら
\begin{equation}
 0=(\exp(-D_C)-1)f
   =-D_C\left(1-\frac{D_C}{2}+\frac{D_C^2}{6}\right)f
 \label{eq:pg-global-centering-nilpotent-stabilizer}
\end{equation}
となる。括弧内は可逆で$D_C$と可換するため$D_Cf=0$を得る。
\eqref{eq:pg-global-centering-linear-syzygy-vanishing}より$C=0$である。
任意底変更後にも同じ議論が成立するので、安定化群スキームも自明である。
従って自由作用の商スタックは代数空間$Y$であり、
$\mathcal T_m\to Y$は$U$-torsorとなる
\cite[Lemma~80.11.7, Tag~\texttt{06PH}]{PStacks}。

このtorsorが自明であると仮定する。
すると正則写像$\beta:\mathcal T_m\to U$を
\begin{equation}
 \beta(u_C\cdot f)=\beta(f)-C
 \label{eq:pg-global-centering-hypothetical-slice}
\end{equation}
となるように取れる。
還元的部分群$G_0$は$U$に$(M,N,\lambda)\cdot C=MCN^{-1}$で作用する。
$\mathcal T_m$は判別式を可逆にしたアフィン空間なので、
その正則関数と$U$のテンソル積は有理$G_0$-加群である。
従って補題\ref{lem:pg-equivariant-affine-lifting}の証明と同じく、
$\beta$の軌道の有限次元線形包でReynolds平均を取れる。
最大compact部分群$K_0\subset G_0$による表示では
\begin{equation}
 \overline\beta(f)
 =\int_{K_0} h^{-1}\!\cdot\beta(h\cdot f)\,dh
 \label{eq:pg-global-centering-reynolds-slice}
\end{equation}
であり、これは正則な$G_0$-同変写像である。
$h u_C h^{-1}=u_{h\cdot C}$を使えば、被積分関数は全て
$u_C$の作用で$-C$だけ変わる。
従って\eqref{eq:pg-global-centering-hypothetical-slice}は平均後にも成立する。
$G=U\rtimes G_0$なので、$\overline\beta$は
\eqref{eq:pg-global-centering-desired-covariance}を満たす全$G$-共変な中心調整となり、
第一の主張に反する。
よってこの$U$-torsorは非自明である。

もし$Y$がアフィンなら、加法群$U\simeq\mathbf G_a^{3m}$のtorsorは
$H^1(Y,\OO_Y)^{3m}=0$により自明になる
\cite[Theorem~59.22.4, Tag~\texttt{03OY}]{PStacks}。
従って$Y$はアフィンではない。
この結論は、自由な中心平行移動を商した後にアフィン底上の平均化を
全係数空間へ適用する方法にも障害を与える。
\end{proof}

% !TeX root = ../../cubic-quantization.tex
% 2026-11-14 / A-v61.
\section{満supportの同時量子化と全階数のスカラーgerbe}
\label{sec:pg-full-support-quantization}

本節では、正規形に現れるブロックの順序を忘れた係数空間を扱う。
まず全ての非零階数層について、中心スカラーgerbeの位数を求める。
次に$N=3k$、すなわち中心変数が残らない満support層では、
係数空間がアフィンであることを用いて、指定したHessian二次jetを持つ
共変な代数的中心元を同時に選ぶ。
その三次Veroneseは実際のPoisson対の被約モジュライへ降下する。
ブロックの内在性と相対階数層の記述は前節までの結果を使い、
アフィン商・還元群の平均・重み零の降下を既知の入力として区別する。

\subsection{内在的なブロック被覆と全階数層のgerbe}

$N\ge5$、$1\le k\le\lfloor N/3\rfloor$、$m=N-3k$とし、
$V=\CC^N$を固定する。基準となる分解と三重ベクトルを
\begin{equation}
 V=W_1\oplus\cdots\oplus W_k\oplus U,\qquad
 \dim W_s=3,\quad\dim U=m,\qquad
 P_k=\sum_{s=1}^kP_s,\quad
 0\ne P_s\in\bigwedge^3W_s
 \label{eq:pg-intrinsic-multiblock-reference}
\end{equation}
とおく。$U$は以下の表示のために選ぶ補空間であり、
$P_k$から内在的に選ばれているとは主張しない。

\begin{lemma}[安定化群と順序付きブロック被覆]
\label{lem:pg-intrinsic-block-cover}
$W=\operatorname{supp}P_k=\bigoplus_sW_s$とおく。
順序を忘れた集合$\{W_1,\ldots,W_k\}$と各成分$P_s$は、
\cite[縮約二形式による復元補題]{PoissonCubic}により$P_k$から復元される。
$J_{N,k}=\operatorname{Stab}_{\mathrm{GL}(V)}(P_k)$は、
選んだ補空間$U$を用いると
\begin{equation}
 \begin{aligned}
 J_{N,k}
  &=\operatorname{Hom}(U,W)\rtimes
     \bigl((\mathrm{SL}_3^k\rtimes\mathfrak S_k)
                   \times\mathrm{GL}(U)\bigr),\\
 J_{N,k}^{\circ}
  &=\operatorname{Hom}(U,W)\rtimes
     \bigl(\mathrm{SL}_3^k\times\mathrm{GL}(U)\bigr),
 \qquad J_{N,k}/J_{N,k}^{\circ}\simeq\mathfrak S_k
 \end{aligned}
 \label{eq:pg-intrinsic-multiblock-stabilizer}
\end{equation}
となる。ここで各$\mathrm{SL}_3$は$P_s$を保つ群であり、
$\mathfrak S_k$はこの体積形式付きブロックを置換する。
$\mathcal O_{N,k}=\mathrm{GL}(V)/J_{N,k}$上には、
ブロックの集合を表す内在的な次数$k$の有限\'etale被覆がある。
その全順序を選ぶ被覆は
\begin{equation}
 \mathcal O_{N,k}^{\mathrm{ord}}
   =\mathrm{GL}(V)/J_{N,k}^{\circ}
   \longrightarrow\mathcal O_{N,k}
 \label{eq:pg-intrinsic-ordered-block-cover}
\end{equation}
という$\mathfrak S_k$-torsorであり、次数は$k!$である。
これらの構成は底変更と可換である。
\end{lemma}
\begin{proof}
$W$は縮約写像$\bigwedge^2V^*\to V$の像であり、$P_k$から決まる。
$W^*=\bigoplus_sW_s^*$の座標を用いると、
\begin{equation}
 D_{P_k}
  =\{\alpha\in W^*:(\iota_\alpha P_k)\wedge
                    (\iota_\alpha P_k)=0\}
  =\bigcup_{s=1}^kW_s^*
 \label{eq:pg-intrinsic-block-contraction-locus}
\end{equation}
である。ここで右辺は被約なスキームとしての等式である。
実際、左辺の定義イデアルは、異なるブロックの座標の積で生成される。
これは右辺の座標線形空間のイデアルの共通部分である。
これは\cite[縮約二形式による復元補題]{PoissonCubic}の証明そのものであり、
$k=1$では零イデアル、$k=2$でも同じ議論が成立する。
従って$P_k$を保つ変換は$W$を保ち、その制限は$W_s$を置換する。

この変換は$W\oplus U$に関して
\begin{equation}
 L=\begin{pmatrix}A&C\\0&D\end{pmatrix},\qquad
 C\in\operatorname{Hom}(U,W),\quad D\in\mathrm{GL}(U)
 \label{eq:pg-intrinsic-block-stabilizer-matrix}
\end{equation}
と書ける。$A(W_s)=W_{\sigma(s)}$なら、
$\bigwedge^3A(P_s)=P_{\sigma(s)}$が必要十分である。
従って対応する三次元の体積形式に関する行列式は一である。
$C,D$への追加条件はなく、これで
\eqref{eq:pg-intrinsic-multiblock-stabilizer}を得る。
標数零では安定化群は滑らかであり、この表示は閉点だけでなく
代数群の表示である。特に$k=2$にも追加の連続対称性はない。

軌道上では縮約写像の階数が一定なので、その像は階数$3k$の
普遍部分束$\mathcal W\subset V\otimes\OO_{\mathcal O_{N,k}}$となる。
\eqref{eq:pg-intrinsic-block-contraction-locus}を相対化し、
\begin{equation}
 p:\mathbf P(D_P)\longrightarrow\mathcal O_{N,k},\qquad
 \mathcal O_{N,k}^{\mathrm{blk}}
   =\operatorname{Spec}_{\mathcal O_{N,k}}p_*\OO_{\mathbf P(D_P)}
 \label{eq:pg-intrinsic-unordered-block-cover}
\end{equation}
とおく。$\mathrm{GL}(V)\to\mathcal O_{N,k}$で引き戻すと
$\mathbf P(D_P)$は$k$個の$\mathbf P^2$の非交和となる。
従って$\mathcal O_{N,k}^{\mathrm{blk}}$は次数$k$の有限\'etale被覆であり、
その順序付けの$\mathfrak S_k$-torsorが
\eqref{eq:pg-intrinsic-ordered-block-cover}である。
各双対成分から、他の成分の和の零化空間として$W_s$を復元できる。
また$P_s$は$P$の$\bigwedge^3W_s$成分である。
以上は忠実平坦被覆上での構成と降下なので、任意底変更と可換である。
\end{proof}

三次式が$P$と適合するとは、異なる支持ブロックの方向に関する
混合二階偏微分が零であることとする。
これは\cite[HessianによるJacobi条件の線形化]{PoissonCubic}の条件である。
基準点$P_k$における適合三次式の空間を$\mathcal F_{N,k}$と書くと、
\cite[相対Poisson階数層の記述]{PoissonCubic}の証明と同じく
\begin{equation}
 \begin{aligned}
 F&=c(u)+\sum_{s=1}^k\sum_{j=1}^3 f_{s,j}(x_s,u),\\
 \deg_{x_s}f_{s,j}&=j,\qquad\deg_u f_{s,j}=3-j,\\
 \ell_{N,k}:=\dim\mathcal F_{N,k}
   &=\binom{m+2}{3}
     +k\left(10+6m+3\binom{m+1}{2}\right)
 \end{aligned}
 \label{eq:pg-intrinsic-multiblock-cubic-fiber}
\end{equation}
である。その滑らかな三次式の部分を$\mathcal F_{N,k}^{\mathrm{sm}}$とし、
\begin{equation}
 \mathcal S_{N,k}
  =\mathrm{GL}(V)\times^{J_{N,k}}\mathcal F_{N,k}^{\mathrm{sm}}
 \label{eq:pg-intrinsic-multiblock-incidence}
\end{equation}
とおく。これは$\mathcal O_{N,k}$上の階数$\ell_{N,k}$のベクトル束の
判別式非消滅開集合であり、滑らかかつ既約である。
以下ではこの被約なスキーム構造を用いる。
Fermat三次式により非空であることも分かる。
この記述には、純中心三次式$c(u)$をブロックに配分する選択は含まれない。
補題\ref{lem:pg-intrinsic-block-cover}の被覆を引き戻すと、
$\mathcal S_{N,k}$上にも内在的なブロック被覆と順序付け被覆を得る。

\begin{proposition}[全非零階数層の商とgerbeの正確な位数]
\label{prop:pg-intrinsic-all-rank-gerbe}
$\mathfrak P_{n,k}^{\mathrm{red}}$を、$n=N-2$次元の滑らかな三次超曲面と
一般階数$2k$のPoisson構造の相対スタックの被約層とする。
群と作用を
\begin{equation}
 \begin{aligned}
 H_N&=\mathrm{GL}(V)\times\mathbf G_{m,\lambda},\\
 (L,\lambda)\cdot(F,P)
   &=\bigl(\lambda^{-1}F\circ L^{-1},
            \lambda(\bigwedge^3L)P\bigr),\\
 K_{\mathrm{sc}}&=\{(cI_V,c^{-3}):c\in\mathbf G_m\},
 \qquad\overline H_N=H_N/K_{\mathrm{sc}}
 \end{aligned}
 \label{eq:pg-intrinsic-all-rank-action}
\end{equation}
とおく。このとき
\begin{equation}
 [\mathcal S_{N,k}/\overline H_N]
    \simeq\mathfrak P_{n,k}^{\mathrm{red}},\qquad
 \mathfrak G_{N,k}:=[\mathcal S_{N,k}/H_N]
    \longrightarrow\mathfrak P_{n,k}^{\mathrm{red}}
 \label{eq:pg-intrinsic-all-rank-quotient}
\end{equation}
の第二の射は$\mathbf G_m$-gerbeである。
その類を$\alpha_{N,k}$とすると
\begin{equation}
 \begin{aligned}
 \operatorname{im}\bigl(\operatorname{wt}:\operatorname{Pic}(\mathfrak G_{N,k})
                                  \longrightarrow\mathbb Z\bigr)
     &=\gcd(N,3)\mathbb Z,\\
 \operatorname{ord}(\alpha_{N,k})
     &=\gcd(N,3)=\gcd(m,3)
 \end{aligned}
 \label{eq:pg-intrinsic-all-rank-period}
\end{equation}
である。特に$3\mid N$ならこのgerbeは非中立であり、
スカラー重み一の可逆層は存在しない。
$\gcd(N,3)=1$なら中立である。
ここで$\operatorname{wt}$は$K_{\mathrm{sc}}\simeq\mathbf G_m$の作用の重みであり、
$\gcd(0,3)=3$と約束する。
同じ位数は、全零切断を含む
$\mathfrak P_{n,k}^{\mathrm{red}}\times\mathbb A^1_\hbar$の任意の開近傍への
このgerbeの引戻しでも保たれる。
位数の主張は階数層全体、またはその全零切断を含む開近傍についてのものであり、
任意の稠密開集合や一般点における位数を主張しない。
特に満supportの場合$N=3k$には位数は常に三である。
\end{proposition}
\begin{proof}
正規形定理\ref{thm:pg-normal}および\cite[相対Poisson階数層の記述]{PoissonCubic}の証明により、
三次式の代表を選んだ階数$2k$の被約層は
\eqref{eq:pg-intrinsic-multiblock-incidence}である。
ここで被約構造を指定しており、Jacobi方程式で定まるPoissonスキーム全体の
被約性を仮定しているのではない。
方程式の代表を変える作用$(F,P)\mapsto(\lambda^{-1}F,\lambda P)$は自由であり、
\begin{equation}
 1\longrightarrow\mathbf G_{m,\lambda}
 \longrightarrow\overline H_N
 \longrightarrow\mathrm{PGL}(V)\longrightarrow1
 \label{eq:pg-intrinsic-all-rank-effective-group}
\end{equation}
は完全である。従って代表のスカラーを先に除き、次に$\mathrm{PGL}(V)$で割ると、
\eqref{eq:pg-intrinsic-all-rank-quotient}の第一の同型を得る。
$K_{\mathrm{sc}}$は底に自明に作用する中心部分群なので、第二の射はgerbeである。
ブロック被覆は$P$の非零スカラー倍で不変であり、これらの商にも降下する。

指標$(\det L)^r\lambda^s$の$K_{\mathrm{sc}}$上の重みは$Nr-3s$である。
特に$\lambda^{-1}$は重み三を持つ。
$\gcd(N,3)=1$なら$Nr-3s=1$を選べるので、
$\theta=(\det L)^r\lambda^s$は重み一の指標となる。
この場合は
\begin{equation}
 H_N\simeq K_{\mathrm{sc}}\times\ker\theta,
 \qquad\ker\theta\simeq\overline H_N
 \label{eq:pg-intrinsic-all-rank-neutralization}
\end{equation}
と分裂し、gerbeは中立である。

$3\mid N$の場合の逆向きの制約を示す。
$\zeta$を一の原始三乗根とし、$V$の基底を$N/3$個の三つ組に分ける。
最初の$k$個をactiveブロックとし、
\begin{equation}
 F_{\mathrm F}=\sum_{i=1}^N x_i^3,\qquad
 P_k=\sum_{s=1}^k
 e_{3s-2}\wedge e_{3s-1}\wedge e_{3s}
 \label{eq:pg-intrinsic-period-fermat-point}
\end{equation}
を考える。inactive部分の次元$m=N-3k$も三の倍数なので、
全ての三つ組に同じ変換を施せる。
各三つ組の基底を$e_0,e_1,e_2$と書き、添字を三で読むとき
\begin{equation}
 S(e_j)=e_{j+1},\qquad T(e_j)=\zeta^j e_j,
 \qquad STS^{-1}T^{-1}=\zeta^{-1}I_V
 \label{eq:pg-intrinsic-period-heisenberg-commutator}
\end{equation}
とおく。各三つ組で$\det S=\det T=1$であり、
$S,T$はいずれもFermat三次式と各active体積形式を保つ。
従って$(S,1),(T,1)$は$(F_{\mathrm F},P_k)$の$H_N$-安定化群に属し、
その交換子は
\begin{equation}
 (\zeta^{-1}I_V,1)
   =(cI_V,c^{-3})\big|_{c=\zeta^{-1}}\in K_{\mathrm{sc}}
 \label{eq:pg-intrinsic-period-scalar-commutator}
\end{equation}
である。

$\mathfrak G_{N,k}$上の任意の可逆層は、
$\mathcal S_{N,k}$上の$H_N$-同変可逆層である。
その$(F_{\mathrm F},P_k)$におけるファイバーは安定化群の一次元表現であり、
交換子には自明に作用する。
従ってスカラー重みを$a$とすると$\zeta^{-a}=1$、すなわち$3\mid a$である。
$\mathcal S_{N,k}$は既約なので、この重みは底全体で同じ整数である。
重み三の指標が既にあることから、重みの像はちょうど$3\mathbb Z$となる。
この議論は可逆層の下部の自明性やPicard群の計算を必要としない。

式\eqref{eq:pg-intrinsic-gerbe-weight-period-criterion}とそのcocycle証明により、
代数スタック上の$\mathbf G_m$-gerbeの類$\alpha$に対しても、
重み$a$の可逆層の存在は$a\alpha=0$と同値である。
代数空間の場合の既知の記述は\cite[Lemma~3.1.1.8]{PLieblichPeriodIndex}を参照する。
以上により\eqref{eq:pg-intrinsic-all-rank-period}を得る。
全零切断を含む開近傍では、重み三の可逆層を引き戻せるので位数は三を割り、
零切断への制限により元の位数を下回らない。
$\gcd(N,3)=1$の場合の中立化も引き戻せるので、最後の主張が従う。
\end{proof}

\begin{remark}[満supportの場合のアフィン性]
\label{rem:pg-intrinsic-full-support-affinity}
$m=0$の場合、$J_{N,k}=\mathrm{SL}_3^k\rtimes\mathfrak S_k$は還元的である。
還元部分群による斉次空間のアフィン性\cite[Proposition~2.2, p.~2]{PArzhantsevMatsushima}により、
$\mathcal O_{N,k}$はアフィンである。
従ってその上のベクトル束の全空間もアフィンであり、
判別式非消滅の主開集合$\mathcal S_{N,k}$もアフィンである。

$m>0$では$J_{N,k}$のユニポテント根基
$\operatorname{Hom}(U,W)$は非自明であり、軌道$\mathcal O_{N,k}$自体は
Matsushimaの判定法\cite[Theorem~2.1, p.~2]{PArzhantsevMatsushima}によりアフィンではない。
ただし、この事実だけから判別式を除いた$\mathcal S_{N,k}$の非アフィン性や
中心元の大域選択の不存在を結論してはならない。
\end{remark}

\begin{remark}[既存結果との境界と中心元の降下]
\label{rem:pg-intrinsic-all-rank-descent-scope}
ブロック復元そのものは\cite[縮約二形式による復元補題]{PoissonCubic}に、
適合三次式の束と相対階数層の表示は既存の相対成分定理の証明に含まれる。
ここではそれらを有限\'etale被覆として整理し、
全非零階数層のgerbeの位数を同じ交換子で決めた。
ブロックの順序を選んでも、中心方向の補空間や純中心三次式の配分は選ばれない。

中心元の存在を別に確立した$H_N$-不変な零切断近傍では、
次数$d$のスカラー重みが$d$であるという既存の作用則をそのまま使える。
$\gcd(N,3)=1$なら$\theta^{-d}$によるtwistで各次数を降下できる。
全ての$N$について第三Veroneseの各次数は
\begin{equation}
 \mathcal A_{3d}\otimes\lambda^d
 \label{eq:pg-intrinsic-all-rank-third-veronese-weight}
\end{equation}
とtwistすればスカラー重み零となる。
従って共変な中心三次元とその平坦な商が構成されている場合には、
系\ref{cor:pg-intrinsic-third-veronese-descent}と同じ降下が適用できる。
この重み計算だけで、全ての$m,k$に対する同時中心元の存在や、
別の中心直線を用いた商の同型を主張するものではない。
\end{remark}

\subsection{満support層における同変中心元}
以下では$N=3k$、$k\ge2$とする。
補題\ref{lem:pg-intrinsic-block-cover}の記法
$\mathcal O_{N,k}$、$\mathcal S_{N,k}$を用いる。
この場合、$P\in\mathcal O_{N,k}$のsupportは$V$全体であり、
その三次元ブロックは順序を除いて$P$から定まる。
順序を付けた有限\'{e}tale被覆上では
\begin{equation}
 V=\bigoplus_{s=1}^k W_s,\qquad
 P=\sum_{s=1}^kP_s,\qquad
 F=\sum_{s=1}^kg_s,\qquad
 g_s\in\operatorname{Sym}^3W_s^*.
 \label{eq:pg-full-support-canonical-blocks}
\end{equation}
最後の式では$W_s^*$を他のブロックを消す$V^*$の部分空間と同一視する。
従って$g_s=F\circ\operatorname{pr}_{W_s}$であり、純中心項の配分は現れない。
$F$が滑らかであることと全$g_s$が滑らかであることは同値である。
実際、ある$g_s$の非零のgradient零点を他のブロックの零ベクトルと合わせると
$F$の非零のgradient零点となり、逆向きはgradientをブロックごとに見るだけでよい。

注意\ref{rem:pg-intrinsic-full-support-affinity}により、
$\mathcal S_{N,k}$は滑らかな既約アフィンスキームである。
具体的には、還元的安定化群
$J_{N,k}=(\mathrm{SL}_3)^k\rtimes\mathfrak S_k$を用いた
アフィン軌道$\mathrm{GL}(V)/J_{N,k}$上で、
階数$10k$の三次式のベクトル束の判別式非零部分として得られる。
以下の持上げでは、このアフィン性を用いる。

\begin{theorem}[満support層の全係数にわたる同変中心元]
\label{thm:pg-intrinsic-full-support-center}
$N=3k$、$k\ge2$とする。
$H_N=\mathrm{GL}(V)\times\mathbf G_{m,\lambda}$は
\begin{equation}
 F'(X)=\lambda^{-1}F(L^{-1}X),\qquad
 P'=\lambda(\Lambda^3L)P,\qquad
 \hbar'=\hbar
 \label{eq:pg-full-support-group-action}
\end{equation}
によって$\mathcal S_{N,k}\times\mathbb A^1_\hbar$に作用する。
普遍代数を
\begin{equation}
 \mathcal A_{F,P,\hbar}
 =\CC\langle v_1,\ldots,v_N\rangle\Big/
 \left([v_i,v_j]-\hbar\sum_aP^{ija}\sym(F_a)\right)_{i,j}
 \label{eq:pg-full-support-ambient-algebra}
\end{equation}
とする。$H_N$と$\hbar\mapsto-\hbar$で不変なZariski開
\begin{equation}
 \mathcal W_{N,k}\subset\mathcal S_{N,k}\times\mathbb A^1_\hbar,
 \qquad
 \mathcal S_{N,k}\times\{0\}\subset\mathcal W_{N,k}
 \label{eq:pg-full-support-good-open}
\end{equation}
および、その上でWeyl係数が正則な中心三次元$\Omega$が存在し、
次を満たす。
\begin{enumerate}
\item $X=Lv$による代数同型を$\Phi$とすると、
\begin{equation}
 \begin{aligned}
 \Phi(\Omega_{F',P',\hbar})&=\lambda^{-1}\Omega_{F,P,\hbar},&
 \Omega_{F,P,-\hbar}&=\Omega_{F,P,\hbar},\\
 \Omega_{F,P,\hbar}
 &=\sym\left(F-\frac{\hbar^2}{12}\mathscr H(P,F)
                  +O(\hbar^4)\right).
 \end{aligned}
 \label{eq:pg-full-support-equivariant-jet}
\end{equation}
ここで$\mathscr H(P,F)$は\eqref{eq:pg-ad-intrinsic-hessian}であり、
偶性とjetはWeyl係数についての等式である。
\item 割った交換子
\begin{equation}
 \mathcal D(p)=
 \left(\hbar^{-1}[\sym(p),v_i]\right)_{i=1}^N,
 \qquad\mathcal K=\ker\mathcal D
 \label{eq:pg-full-support-divided-kernel}
\end{equation}
を三次多項式の束から$\mathcal A_4^{\oplus N}$への写像と見る。
その核は階数$k$のベクトル束であり、核・像・余核の形成は任意の底変更と両立する。
ブロックに順序を付けた有限\'{e}tale被覆上で、$\mathcal K$は
各ブロックの中心三次直線$\mathcal K_s$の直和となる。
$\Omega$の各$\mathcal K_s$成分は零でなく、
$\mathcal J=\OO\Omega\subset\mathcal K$は$H_N$-同変な直線部分束となる。
\item $\mathcal A$と$\mathcal B=\mathcal A/(\mathcal J)$は次数別有限局所自由である。
各閉ファイバーの$\mathcal A$は次元$N$、AS parameter $N$の
Koszul AS-regular代数であり、$\Omega$は左右正則である。
各閉ファイバーの$\mathcal B$は左右injective dimension $N-1$、
AS parameter $N-3$のAS-Gorenstein代数で、
\begin{equation}
 \begin{aligned}
 \operatorname{Hilb}_{\mathcal A}(z)&=(1-z)^{-N},&
 \operatorname{Hilb}_{\mathcal B}(z)&=(1-z^3)(1-z)^{-N},\\
 \mathcal B/(\hbar)&\simeq
 \operatorname{Sym}(V^*)\otimes\OO_{\mathcal S_{N,k}}/(F).
 \end{aligned}
 \label{eq:pg-full-support-flat-quotient}
\end{equation}
特殊ファイバーの一次Poisson括弧は$\iota_{dF}P$である。
\end{enumerate}
最終的な族の底には、ブロックの順序、各ブロックの基底、
中心三次直線の相対規格化を追加しない。
中心元の一意性や全$\hbar$直線上での存在は主張しない。
\end{theorem}
\begin{proof}
\emph{古典核と局所的な量子核。}
ブロックを順序付け、その体積を$P_s$とする基底を局所的に選べば、
\eqref{eq:pg-full-support-ambient-algebra}は
\begin{equation}
 \mathcal A\simeq\bigotimes_{s=1}^k\mathcal A_{g_s,\hbar},
 \qquad
 (\mathcal K_0)_{(F,P)}=\bigoplus_{s=1}^k\CC g_s
 \label{eq:pg-full-support-tensor-and-classical-kernel}
\end{equation}
となる。テンソル積は係数底上で取る。
古典核の等号を確認する。
三次式$p$がCasimirなら各$s$について
$d_{W_s}p\wedge d g_s=0$である。
$g_s$の三つの偏微分は互いに素なので、
全変数の多項式環の分数体で得る比例係数は多項式である。
斉次次数を比較するとその係数は定数$a_s$であり、
$p-\sum_sa_sg_s$は全変数で微分が零となる。
三次斉次性からこの差は零である。
$g_s$の単項式の支持は互いに異なるので、上の和は直和となる。
この議論は補題\ref{lem:pg-det-classical-kernel}と同じUFDの議論である。
順序を替えると直和の因子が置換されるだけであり、
$\mathcal K_0$とその切断$F=\sum_sg_s$は順序を忘れた底で定義される。

各ブロックに命題\ref{prop:pg-PBW}の係数を動かした開
\eqref{eq:pg-det-universal-open}と、補題\ref{lem:pg-det-divided-kernel}を適用する。
全零切断の近傍で、各ブロック代数は次数別局所自由、
次数三・四のWeyl写像は同型、各割った交換子核$\mathcal K_s$は階数一となる。
異なるブロックの生成元は可換なので、
$\bigoplus_s\mathcal K_s$は全体の$\mathcal D$の核に入る。
テンソル積のブロック多重次数と各$\mathcal K_s$の局所分裂により、
これは$\mathcal A_3$の階数$k$の局所直和因子である。
従って$\mathcal D$は階数$r=\binom{N+2}{3}-k$の商束を経由し、
$(r+1)$次小行列式はすべて恒等的に零となる。
一方、\eqref{eq:pg-full-support-tensor-and-classical-kernel}により、
零切断の全点で階数は$r$である。
階数$r$の開に制限し、小行列式による基本変形を行うと、
$\mathcal D$は局所的に$\operatorname{diag}(I_r,0)$となる。
これで$\mathcal K$、像、余核の局所自由性と任意の底変更との両立を得る。
両辺の階数と各ファイバーでの独立性から
\begin{equation}
 \mathcal K=\bigoplus_{s=1}^k\mathcal K_s
 \label{eq:pg-full-support-quantum-kernel-sum}
\end{equation}
も成り立つ。
ここで混合単項式から追加の三次中心元が現れないことは、
この全交換子行列の階数条件によって保証されている。

これらの条件は代数と割った交換子から定まり、基底変更とブロック置換で保たれる。
従って条件の成立する開は順序を忘れた底へ降り、
全零切断を含む$H_N$-不変開$\mathcal U_{\mathrm{loc}}$を与える。
語の順序反転による反対代数との同型を用い、
$\hbar\mapsto-\hbar$に関しても不変とする。

\emph{アフィン底上の指定jetの直接持上げ。}
$\mathcal S_{N,k}$はアフィンなので、
補題\ref{lem:pg-equivariant-affine-lifting}(1)により
\begin{equation}
 D(q)\subset\mathcal U_{\mathrm{loc}},\qquad
 q\in\Gamma(\mathcal S_{N,k},\OO)[\hbar]^{H_N\times\boldsymbol\mu_2},
 \qquad q\bmod\hbar=1
 \label{eq:pg-full-support-affine-good-open}
\end{equation}
とできる。ここで$\boldsymbol\mu_2$は$\hbar$の符号と語の順序を反転する。
$D(q)$はアフィンであり、その上の$\mathcal K$は有限射影加群に対応する。
特に古典核についても$H^1(\mathcal S_{N,k},\mathcal K_0)=0$であるが、
以下では有限の厚みから直接持ち上げれば十分である。

\begin{equation}
 \Theta=F-\frac{\hbar^2}{12}\mathscr H(P,F)
 \label{eq:pg-full-support-prescribed-theta}
\end{equation}
とおく。\eqref{eq:pg-ad-intrinsic-blocks}と
命題\ref{prop:pg-universal-casimir-jet}を各ブロックに適用すると、
$\mathcal D\Theta\in\hbar^4\mathcal A_4^{\oplus N}$である。
異なるブロックは可換なので、交叉交換子は現れない。
底変更との両立から$\Theta\bmod\hbar^3$は
$\mathcal K/\hbar^3\mathcal K$の大域切断になる。
アフィン性により
\begin{equation}
 \Gamma(D(q),\mathcal K)\longrightarrow
 \Gamma\bigl(V(\hbar^3),\mathcal K/\hbar^3\mathcal K\bigr)
 \quad\text{は全射である。}
 \label{eq:pg-full-support-jet-lift-surjection}
\end{equation}
$\mathscr H(P,F)$は$F$と同じ$\lambda^{-1}$の共変性を持ち、
$\Theta$は偶である。
従って引戻しdatumを$\lambda\Phi$でtwistすると、この剰余切断は
$H_N\times\boldsymbol\mu_2$-不変となる。
補題\ref{lem:pg-equivariant-affine-lifting}(2)を閉部分$V(\hbar^3)$へ適用し、
任意の持上げをReynolds作用素で平均すれば、
$\mathcal D\Omega=0$、$\Omega\equiv\Theta\pmod{\hbar^3}$を満たす
同変かつ偶な代数的切断を得る。
偶性により差の三次係数も零なので、
\eqref{eq:pg-full-support-equivariant-jet}が成立する。
この議論は形式級数の全次数の持上げや完備化を必要としない。

順序を付けた有限\'{e}tale被覆上で、$\Omega$の各
$\mathcal K_s$成分が零になる閉集合を除く。
その和集合はブロック置換で不変なので、除いた開は元の底へ降りる。
古典極限の各成分は非零の$g_s$であるから、全零切断は残る。
この開を$\mathcal W_{N,k}$とすれば、$\Omega$は各ファイバーで非零であり、
所要の同変な直線部分束を生成する。

\emph{正則性と平坦な中心商。}
各閉ファイバーでは、$\mathcal A$は三次元Koszul AS-regular代数の
テンソル積である。
各因子のKoszul分解とAS-Gorenstein条件をテンソルする既存の議論
（定理\ref{thm:pg-ad-all-poisson}）により、
全体の次元とAS parameterはともに$3k=N$となる。

$\Omega$を\eqref{eq:pg-full-support-quantum-kernel-sum}に従って
$\Omega=\sum_s\gamma_s$と書く。
一つのブロック$s$の次数だけでfilterすると、その最高記号は
非零の$\gamma_s\in(\mathcal A_s)_3$である。
三次元因子$\mathcal A_s$はdomainなので、$\gamma_s$の左右乗法は単射である。
他の因子とのテンソル積は体上で単射性を保つ。
最高次数を比較すれば、元の$\Omega$の左右乗法も単射となる。
ここではテンソル積代数自体がdomainであるとは仮定していない。
従って各内部次数$d$について
\begin{equation}
 0\longrightarrow\mathcal J\otimes\mathcal A_{d-3}
 \longrightarrow\mathcal A_d\longrightarrow\mathcal B_d
 \longrightarrow0
 \label{eq:pg-full-support-central-quotient-sequence}
\end{equation}
は、全閉ファイバーで左の射が単射なベクトル束の列である。
最大小行列式によって局所的に分裂し、$\mathcal B_d$は有限局所自由となる。
Hilbert級数と古典極限はこの完全列から従い、
AS-Gorenstein条件は補題\ref{lem:pg-central-cubic}による。
一次Poisson括弧は\eqref{eq:pg-full-support-ambient-algebra}から直ちに従う。
\end{proof}

\begin{remark}[核と選択の範囲]
\label{rem:pg-full-support-center-scope}
この定理の$\mathcal K$は$\hbar$で割った交換子の核である。
$\hbar=0$で通常の中心は次数三全体となるため、通常の中心の形成が
任意の底変更と両立するという主張ではない。
また、各ブロックの中心直線だけから総中心直線が一意に定まるわけではない。
ここで示したのは、古典切断$F$とその二次jetを指定した一つの同変な選択の存在である。
既知の一ブロック中心元、還元群のアフィン商、アフィン底上の同変持上げを用いて、
満support層で分解の順序を保持しない同時の族を構成した。
中心変数がある場合の全supportの同変持上げは、このアフィン性の議論からは従わない。
\end{remark}

\begin{corollary}[満supportの底の次元]
\label{cor:pg-full-support-moduli-dimension}
$N=3k$、$k\ge2$の満support層について
\begin{equation}
 \dim\mathcal S_{N,k}=N^2+2k,\qquad
 \dim[\mathcal S_{N,k}/\overline H_N]=2k.
 \label{eq:pg-full-support-moduli-dimension}
\end{equation}
従って$N=6$では、このPoisson対の被約モジュライ層の次元は四である。
これは底の次元であり、固定$\hbar\ne0$での非可換代数の同型類の像の次元を主張しない。
\end{corollary}
\begin{proof}
$\dim J_{N,k}=8k$なので、$\dim\mathcal O_{N,k}=N^2-8k$である。
許される三次式の束の階数は$10k$であり、滑らかさは非空開条件である。
これより第一の等式を得る。
$\dim\overline H_N=N^2$を引けば第二の等式となる。
\end{proof}

\subsection{満support層への三次偏極での降下}
\begin{corollary}[満support層の三次Veronese族]
\label{cor:pg-full-support-veronese-descent}
$N=3k$、$k\ge2$とし、定理\ref{thm:pg-intrinsic-full-support-center}の族を取る。
全零切断を含む開部分スタック
\begin{equation}
 \mathfrak W_{N,k}=[\mathcal W_{N,k}/\overline H_N]
 \subset\mathfrak P_{n,k}^{\mathrm{red}}\times\mathbb A^1_\hbar
 \label{eq:pg-full-support-effective-open}
\end{equation}
上に、次数別有限局所自由な代数族$\mathcal A^{[3]},\mathcal B^{[3]}$が存在する。
対応するgerbeへ引き戻すと
\begin{equation}
 \mathcal A^{[3]}_d=\mathcal A_{3d}\otimes\lambda^d,
 \qquad
 \mathcal B^{[3]}_d=\mathcal B_{3d}\otimes\lambda^d
 \quad(d\ge0)
 \label{eq:pg-full-support-third-veronese}
\end{equation}
であり、次数一の大域的な左右正則中心元$\omega$について
\begin{equation}
 \mathcal B^{[3]}=\mathcal A^{[3]}/(\omega)
 \label{eq:pg-full-support-veronese-quotient}
\end{equation}
となる。$\hbar=0$での相対$\operatorname{Proj}$は普遍の周囲射影空間族と
三次超曲面族であり、新しい次数一の偏極は各幾何学的ファイバーで$\OO(3)$である。
他方、全層のスカラーgerbeは位数三を持ち、元の次数一の線形化を
一つの可逆層の次数冪で補正して降下させることはできない。
\end{corollary}
\begin{proof}
$K_{\mathrm{sc}}$は$\mathcal A_d,\mathcal B_d$に重み$d$、
指標$\lambda$に重み$-3$で作用する。
従って\eqref{eq:pg-full-support-third-veronese}の各成分の重みは零であり、
乗法と単位もgerbeの重み零の降下
\cite[Proposition~4.10]{PBerghSchnurerGerbes}によって降りる。
中心元の共変性\eqref{eq:pg-full-support-equivariant-jet}は
$\mathcal J\otimes\lambda\simeq\OO$という同変な自明化を与えるので、
$\Omega$は新しい次数一の大域切断$\omega$を定める。
式\eqref{eq:pg-full-support-central-quotient-sequence}の$d$を$3d$に替え、
$\lambda^d$で捻ると、左端は
\begin{equation}
 (\mathcal J\otimes\mathcal A_{3d-3})\otimes\lambda^d
   \simeq\mathcal A_{3d-3}\otimes\lambda^{d-1}
   =\mathcal A^{[3]}_{d-1}
 \label{eq:pg-full-support-veronese-exact-sequence}
\end{equation}
となる。よって商の同定、正則性、次数別局所自由性が従う。
古典極限では通常の三次Veroneseと底の可逆層の次数冪による捻りを取っている。
いずれも相対$\operatorname{Proj}$を保ち、偏極は各ファイバー上で$\OO(3)$となる。
この同定は自明化した底上の通常の同定から降下するので、
周囲族が大域的にベクトル束の射影化であるとは仮定しない。
最後の非存在は命題\ref{prop:pg-intrinsic-all-rank-gerbe}から従う。
\end{proof}

\begin{remark}[三次四様体と比較の範囲]
\label{rem:pg-full-support-fourfold-scope}
$N=6$、$k=2$では、正規形定理により三次四様体の一般階数四のPoisson対を
全て扱っている。底の次元は系\ref{cor:pg-full-support-moduli-dimension}の通り四であり、
その全体で二つの三生成元量子代数と総中心三次商を同時に構成し、
三次偏極によってPoisson対のモジュライへ降ろしたことになる。
この主張は、固定非零$\hbar$での同型類像の次元や一意性を含まない。

三次Veronese自体のAS-regular性は主張しない。
左右noetherianな幾何学的ファイバーに限れば、
注意\ref{rem:pg-intrinsic-veronese-qgr}の直接証明と
Smithの既知のVeronese同値\cite{PSmithVeronese}をそのまま適用できる。
これは相対的な$qgr$の降下や、全ファイバーの滑らかさの証明ではない。
中心変数を持つ複数ブロックの全係数空間への同変中心元の延長と、
異なる中心直線による商の抽象的同型・Morita・$qgr$比較は別に残る。
\end{remark}

% !TeX root = ../../cubic-quantization.tex
\section{Fermat型の明示量子化}\label{sec:pg-special}
\begin{lemma}\label{lem:pg-fermat-center}
結合代数において
\begin{equation}
[x,y]=\hbar z^2,\qquad [y,z]=\hbar x^2,\qquad [z,x]=\hbar y^2
\end{equation}
が成り立つとする。$S=x^3+y^3+z^3$、$T=\sym(xyz)$ とおくと $S-2\hbar^2T$ は中心元である。
ここで $\sym(xyz)$ は六個の語の平均である。
\end{lemma}
\begin{proof}
関係式と積の交換子を展開すると
\begin{equation}
[S,x]=\hbar(zy^2z-yz^2y)=\hbar[zy,yz]=\hbar^2[zy,x^2]
\end{equation}
を得る。一方、$6T$ の六個の語を関係式で並べ替えると
$2T=xzy+zyx+\hbar x^3$ となる。従って $2[T,x]=[zy,x^2]$ であり、
$[S-2\hbar^2T,x]=0$ を得る。循環置換により $y,z$ との交換子も零である。
\end{proof}
特に
\begin{equation}\label{eq:pg-fermat}
F=\frac13\left(\sum_{\alpha=1}^m u_\alpha^3+
\sum_{s=1}^k\sum_{i=1}^3x_{si}^3\right)
\end{equation}
に対する量子化では、各ブロックで $[x_{s1},x_{s2}]=\hbar x_{s3}^2$ とその循環置換を課し、
異なるブロックと中心変数は可換とする。このとき
\begin{equation}\label{eq:pg-fermat-center}
\Omega=\frac13\left[\sum_{\alpha=1}^m u_\alpha^3+
\sum_{s=1}^k\left\{\sum_{i=1}^3x_{si}^3-2\hbar^2\sym(x_{s1}x_{s2}x_{s3})\right\}\right]
\end{equation}
は全パラメータで中心元であり、$\Omega|_{\hbar=0}=F$ を満たす。
定理\ref{thm:pg-quantization}の適切な開集合上で中心正則超曲面の族を与える。


\subsection{Poisson量子化の明示性}
\begin{proposition}\label{prop:poisson-quantization-explicit}
古典的入力である正規形\eqref{eq:pg-normal}
$F=\sum_{s=1}^k f_s(x_s,u)$に対し、
量子化を定める二次関係式は明示的に
\begin{equation}
 [x_{s2},x_{s3}]
 =\hbar\,\operatorname{sym}_R(\partial_{x_{s1}}f_s),
 \quad
 [x_{s3},x_{s1}]
 =\hbar\,\operatorname{sym}_R(\partial_{x_{s2}}f_s),
 \quad
 [x_{s1},x_{s2}]
 =\hbar\,\operatorname{sym}_R(\partial_{x_{s3}}f_s)
 \label{eq:ci-quantization-relations}
\end{equation}
と書ける。異なるブロックと中心変数の間の交換子は零である。
三次中心元$\Omega_\hbar$は零のZariski近傍で存在し、
有限自由加群間の中心性方程式
$z\mapsto([z,x_{si}])_{s,i}$の核を解いて選べる。
しかし任意の入力に一様な閉公式が得られたとは主張しない。
\end{proposition}
\begin{proof}
二次関係式は各ブロックのcyclic potential
$W_s=x_{s1}x_{s2}x_{s3}-x_{s2}x_{s1}x_{s3}
-\hbar\,\operatorname{sym}_R(f_s)$のcyclic derivativeである。
中心元の存在と、核の計算を用いた代数化は
定理\ref{thm:pg-universal-determinant}および定理\ref{thm:pg-quantization}による。
\end{proof}
\begin{example}[Fermat型]
$[x,y]=\hbar z^2$、$[y,z]=\hbar x^2$、
$[z,x]=\hbar y^2$なら
\begin{equation}
 \Omega=x^3+y^3+z^3
      -2\hbar^2\operatorname{sym}(xyz)
 \label{eq:ci-fermat-center}
\end{equation}
は実際に中心元となる。直接の交換子計算は
補題\ref{lem:pg-fermat-center}にある。
\end{example}

% !TeX root = ../../cubic-quantization.tex
\section{非楕円ブロックの古典入力と行列式曲線}\label{cq:nonelliptic-data}
節点・尖点の制限を持つ次の二例を用いる。
座標$(x,y,z,a,b,c,u)$に対応するベクトル基底を$(e_0,\ldots,e_6)$とする。
\begin{equation}\label{cq:nonelliptic-F}
\begin{aligned}
 f_{\mathrm N}&=xyz+y^3+z^3,& f_{\mathrm C}&=xz^2+y^3,\\
 F_\epsilon&=f_\epsilon+a^3+b^3+c^3+u(x^2+a^2+b^2+c^2)+u^3,
 &X_\epsilon&=V(F_\epsilon)\subset\PP^6
 \quad(\epsilon\in\{\mathrm N,\mathrm C\}).
\end{aligned}
\end{equation}
$P_1=e_0\wedge e_1\wedge e_2$、$P_2=e_3\wedge e_4\wedge e_5$、
$L=\CC P_1\oplus\CC P_2$とおく。


$X_{\mathrm N},X_{\mathrm C}$の滑らかさ、および$L$が固定した
Poissonスキームの被約既約成分となることは、
\cite[非楕円型を一般点に持つ既約成分の定理]{PoissonCubic}による。
その証明は古典論文に置き、本稿ではこの明示入力の量子化とその検証を扱う。

\subsection{同次量子ブロックの行列式曲線}
$u=0$の三変数ブロックを、$[y,z]=h\operatorname{sym}(f_x)$とその巡回置換で量子化する。
多重線形化した二次関係式の行列は
\begin{equation}
 M_{\mathrm N}(h)=\begin{pmatrix}
 0&-(1+h/2)z&(1-h/2)y\\
 (1-h/2)z&-3hy&-(1+h/2)x\\
 -(1+h/2)y&(1-h/2)x&-3hz
 \end{pmatrix},
\end{equation}
\begin{equation}
 M_{\mathrm C}(h)=\begin{pmatrix}0&-z&y-hz\\z&-3hy&-x\\-y-hz&x&-hx\end{pmatrix}.
\end{equation}
\begin{proposition}\label{ne:quantum}
\begin{equation}\begin{aligned}
 \det M_{\mathrm N}(h)&=-\frac h4\bigl((h^2+12)xyz+(12-3h^2)(y^3+z^3)\bigr),\\
 \det M_{\mathrm C}(h)&=-3h\bigl(xz^2+y^3-h^2yz^2\bigr).
\end{aligned}\end{equation}
前者は$h\ne0$、$h^2\notin\{4,-12\}$で節点三次曲線、
後者は全ての$h\ne0$で尖点三次曲線である。
\end{proposition}
\begin{proof}
行列式を展開すると二公式を得る。
節点型は二係数が非零のとき$x$の定数倍で$f_{\mathrm N}$に移り、
尖点型は$x'=x-h^2y$により$x'z^2+y^3$となる。
\end{proof}
ここでの主張は多重線形化した関係式の\emph{行列式曲線}についてである。
全point moduleのmoduliとの一致、AS-regular性、全量子化の相対smoothnessは
この計算だけからは主張しない。
$h=0$で行列式は恒等的に零となり、$h$を除いた極限曲線と
可換三変数多項式環のpoint scheme $\PP^2$を区別する。

\subsection{全成分復元における役割}
これらの成分の存在は、\cite{NCCubicMain}の一般点可換復元に反しない。
特異ブロックの点$[1:0:0]$では全体の$u$方向微分$F_u=x^2$が非零であり、
中心変数との混合が全体の滑らかさを保つ。
従って、各ブロックの楕円性を必要条件とする方法ではこれらの成分を扱えない。
全ポテンシャルの対角、中心方向の微分、その量子化との両立が課題となる。

特異曲線の一般化Jacobianにも注意が必要である。
節点三次曲線では$\Pic^0\simeq\mathbb G_m$、尖点型では$\Pic^0\simeq\mathbb G_a$。
前者は正規化$\PP^1$上の二点のfiberを非零スカラーで接着することで、
後者は$\CC[[t^2,t^3]]\subset\CC[[t]]$の単元商の一次係数$t$により得られる。
従って、滑らかな楕円Picardに対して得た有理写像の障害を、これらにそのまま適用してはならない。
一般化Jacobianの導入だけで圏比較が成立するとも主張しない。


% !TeX root = ../../cubic-quantization.tex
\section{非楕円成分の一様な量子化と周囲の相対対角}
\label{nq:sec:quantization}
本節では第\ref{cq:nonelliptic-data}節の固定した二つのPoisson成分を、実在する二パラメータの代数族として量子化する。
形式級数の存在や有限次数のHilbert関数検査だけにとどめず、全次数のflatnessと
\emph{周囲の代数}の有限相対対角分解まで構成する。
残余dg圏の対角は別の対象であり、本節の結論をそのsaturationと読み替えない。

\subsection{一次資料との関係と主定理}
三変数の同次Jacobian代数のAS-regular性はMori--Smith
\cite[Theorem 4.2, Lemma 4.1]{PMoriSmith}の分類を用いる。
中央一次元による拡大にはLe Bruyn--Smith--Van den Bergh
\cite[Theorem 2.6, Corollary 2.7, Theorem 3.1.3]{LBSVNE12}、
低次ポテンシャルを加えるPBW変形とその双加群複体にはBerger--Taillefer
\cite[Theorems 3.1 and 3.6]{PBergerTaillefer}を用いる。
これらの一般定理を新規な成果とは主張しない。
本節では、第\ref{cq:nonelliptic-data}節の特定の節点・尖点成分に対して中心混合項を保った具体的な式を与え、
係数環、例外集合、底変更および対角の種類を明記する。
中心三次元の表示の初出は未確定である。

$r,s$を次数零、$x,y,z,a,b,c,u$を次数一とし、
\begin{equation}
 R_{\mathrm C}=\CC[r,s],\qquad
 R_{\mathrm N}=\CC[r,s,(4-r^2)^{-1}]
\end{equation}
と置く。第一ブロックの代数$D_\epsilon$は$u$を中心として
\begin{align}
 \epsilon=\mathrm C:\quad
 [y,z]&=r(z^2+2ux),&[z,x]&=3ry^2,&[x,y]&=r(xz+zx),
 \label{nq:eq:cusp}\\
 \epsilon=\mathrm N:\quad
 [y,z]&=\tfrac r2(yz+zy)+2rux,&
 [z,x]&=\tfrac r2(zx+xz)+3ry^2,&
 [x,y]&=\tfrac r2(xy+yx)+3rz^2
 \label{nq:eq:node}
\end{align}
で定める。第二ブロック$E_s$も$u$を中心とし、
\begin{equation}\label{nq:eq:fermat}
 [a,b]=s(3c^2+2uc),\quad
 [b,c]=s(3a^2+2ua),\quad
 [c,a]=s(3b^2+2ub)
\end{equation}
で定める。異なるブロックの生成元は可換とする。したがって
\begin{equation}
 \mathcal A_\epsilon=D_\epsilon\otimes_{R_\epsilon[u]}E_s
\end{equation}
は七生成元の次数付き代数である。

\begin{theorem}[二つの非楕円成分の代数的量子化]
\label{nq:thm:family}
以下の式で定める$\Omega_\epsilon\in\mathcal A_{\epsilon,3}$は中心的であり、
$\mathcal B_\epsilon=\mathcal A_\epsilon/(\Omega_\epsilon)$は次を満たす。
\begin{enumerate}
\item $\mathcal A_\epsilon$は次数ごとに有限局所自由で、その各幾何学的ファイバーは
noetherianなKoszul AS-regular domain、global dimension七、AS parameter七である。
\item $\Omega_\epsilon$は各幾何学的ファイバーで非零の中心正則元であり、
$\mathcal B_\epsilon$は次数ごとに有限局所自由で、形成は任意の底変更と両立する。
各ファイバーのHilbert級数とAS-Gorenstein型は
\begin{equation}
 H_{\mathcal B}(t)=\frac{1-t^3}{(1-t)^7},\qquad
 \operatorname{injdim}\mathcal B=6,\qquad \text{AS parameter}=4.
\end{equation}
\item $(r,s)=(0,0)$で$\mathcal B_\epsilon$は第\ref{cq:nonelliptic-data}節の$X_\epsilon$の斉次座標環となる。
任意の$(\lambda,\mu)$に沿う$r=h\lambda,s=h\mu$への引戻しは、
$\lambda\pi_{P_1}+\mu\pi_{P_2}$を半古典極限に持つ。
\item $R_\epsilon[u]$上の$\mathcal A_\epsilon$の対角は、階数列
$1,6,15,20,15,6,1$の有限自由双加群分解を持つ。
この分解は任意の可換$R_\epsilon[u]$-代数への底変更後も完全である。
$R_\epsilon$上でも周囲の対角はperfectである。
\end{enumerate}
この定理は$\operatorname{qgr}\mathcal B_\epsilon$やその残余圏のsmoothnessを含まない。
\end{theorem}

\subsection{第一ブロックのmonic PBW表示}
この小節では$r=h$と書き、係数環を
$Q_{\mathrm C}=\CC[h]$、$Q_{\mathrm N}=\CC[h,(4-h^2)^{-1}]$とする。
$u$を係数に入れ、active語の長さ、次いで辞書式$x>y>z$の順序を用いる。
$\epsilon=\mathrm C$の関係式は
\begin{align}
 yz&\longmapsto zy+hz^2+2hux,\notag\\
 xz&\longmapsto zx-3hy^2,\notag\\
 xy&\longmapsto yx+2hzx-3h^2y^2
 \label{nq:eq:cusp-rewrite}
\end{align}
と等価である。節点型では$q=(2+h)/(2-h)$として
\begin{align}
 yz&\longmapsto qzy+\frac{4h}{2-h}ux,\notag\\
 xz&\longmapsto q^{-1}zx-\frac{6h}{2+h}y^2,\notag\\
 xy&\longmapsto qyx+\frac{6h}{2-h}z^2.
 \label{nq:eq:node-rewrite}
\end{align}
全ての右辺のactive語は左辺より小さい。重なりの曖昧さは$xyz$の一つだけである。
\begin{lemma}[全次数と任意の底変更に対するPBW性]
\label{nq:lem:pbw}
二つの還元系はいずれも合流的であり、
$\{z^iy^jx^k:i,j,k\ge0\}$は$D_\epsilon$の$Q_\epsilon[u]$-基底である。
この主張は任意の可換$Q_\epsilon[u]$-代数への底変更後も成立する。
\end{lemma}
\begin{proof}
節点型で二つの還元経路から得る共通の正規形は
\begin{equation}
 qzyx+\frac{4h}{2-h}ux^2+\frac{6h}{2-h}(z^3-y^3).
\end{equation}
尖点型では両経路が
\begin{equation}\begin{aligned}
 zyx+3hz^2x+2hux^2-3hy^3-9h^2zy^2-6h^3z^2y
 -6h^4z^3\\
 {}-12h^3uyx-24h^4uzx+36h^5uy^2
\end{aligned}\end{equation}
となる。よってDiamond lemma \cite[Theorem 1.2]{BergNE12}を適用できる。
係数は表示した環に属し、還元の先頭係数は全て一である。
同じ有限恒等式が任意の底変更後にも成り立つため、基底の主張も保存される。
\end{proof}

\begin{lemma}[幾何学的ファイバーの正則性]\label{nq:lem:regular4}
$D_\epsilon$の各幾何学的ファイバーは四次元noetherian Koszul AS-regular domainである。
\end{lemma}
\begin{proof}
$u=0$の三変数代数は、尖点型では全$h$、節点型では$h\ne\pm2$で
三次元Calabi--Yau AS-regularである\cite[Lemma 4.1, Theorem 4.2]{PMoriSmith}。
域性とnoetherian性は、次のOre表示からも確認できる。
節点型では量子平面$yz=qzy$上に
\begin{equation}
 \sigma(y)=qy,\quad\sigma(z)=q^{-1}z,\quad
 \delta(y)=\frac{6h}{2-h}z^2,\quad
 \delta(z)=-\frac{6h}{2+h}y^2
\end{equation}
による$x$のOre拡大を取る。
尖点型ではJordan平面$yz-zy=hz^2$上に
\begin{equation}
 \sigma(y)=y+2hz,\quad\sigma(z)=z,\quad
 \delta(y)=-3h^2y^2,\quad\delta(z)=-3hy^2
\end{equation}
を使う。いずれも$\sigma$は可逆、$\delta$は表示した関係を保つ$\sigma$-導分である。
補題\ref{nq:lem:pbw}より$u$は中心正則元である。
中央次数一拡大の定理\cite[Theorem 2.6, Corollary 2.7 and Remark 2.8]{LBSVNE12}
により四次元の結論を得る。ここで特異なのは三次曲線であって、周囲のAS代数ではない。
\end{proof}

\subsection{混合項を含む中心三次元の閉公式}
全ての積を、表示した順序の非可換積として読む。
\begin{proposition}[第一ブロックの中心元]\label{nq:prop:center}
次の元は$D_\epsilon$の中心に属し、$h=0$で$f_\epsilon+ux^2$となる。
\begin{align}
 \Gamma_{\mathrm C}(h)={}&
 z^2x+y^3-3hzy^2+ux^2
 -6h^2uyx-6h^3uzx+12h^4uy^2,
 \label{nq:eq:centerC}\\
 \Gamma_{\mathrm N}(h)={}&
 z^3+\left(\frac{2-h}{2+h}\right)^3y^3
 +\frac{h^2+12}{3(2+h)^2}zyx
 +\frac{h^2+12}{6(2+h)}ux^2\notag\\
 &-\frac{24h^2}{(2+h)^2}uzy
 -\frac{24h^3}{(2+h)^2}u^2x.
 \label{nq:eq:centerN}
\end{align}
\end{proposition}
\begin{proof}
補題\ref{nq:lem:pbw}の還元を両辺の積$\Gamma_\epsilon a$、$a\Gamma_\epsilon$
（$a=x,y,z$）へ適用すると、全てで同じ正規形となる。
付属の\texttt{nonelliptic\_quantization\_certificate.json}に、六つの共通正規形を全係数で収録した。
その係数は尖点で$\mathbb Q[h,u]$、節点で$\mathbb Q[h,u,(4-h^2)^{-1}]$に属し、
コードは差の分子が恒等的に零であることを検査する。
従ってこれは有限個の特殊化による観察ではなく、合流する還元による普遍恒等式の証明である。
$u$との可換性は定義による。$h=0$での式は直接代入で分かる。
\end{proof}
\begin{remark}[量子的な混合補正の必要性]
$\Gamma_\epsilon(h)|_{u=0}+ux^2$という素朴な持上げは一般には中心でない。
付属コードはその非零な交換子の正規形も記録する。
中心混合項を古典式のまま加えるだけでは、量子化の中心三次方程式にならない。
\end{remark}

第二ブロックでは
$\alpha=a+u/3$、$\beta=b+u/3$、$\gamma=c+u/3$、$t=3s$と置く。
関係式\eqref{nq:eq:fermat}は
$[\alpha,\beta]=t(\gamma^2-u^2/9)$とその巡回置換になる。
\begin{lemma}[Fermat側の混合中心元]\label{nq:lem:fermat-center}
\begin{equation}
 \Gamma_{\mathrm F}(s)=
 \alpha^3+\beta^3+\gamma^3
 -3s^2\!\sum_{w\in\mathfrak S_3(\alpha\beta\gamma)}w
 -\frac{u^2}{3}(\alpha+\beta+\gamma)+\frac{2u^3}{9}
 \label{nq:eq:centerF}
\end{equation}
は$E_s$の中心元であり、$s=0$で$a^3+b^3+c^3+u(a^2+b^2+c^2)$となる。
\end{lemma}
\begin{proof}
独立な中心係数$\eta$を用い、
$r_1=xy-yx-t(z^2-\eta)$、$r_2=yz-zy-t(x^2-\eta)$、
$r_3=zx-xz-t(y^2-\eta)$とする。
$D=x^3+y^3+z^3$、$T=xyz+xzy+yxz+yzx+zxy+zyx$に対し、
自由代数$\mathbb Z[t,\eta]\langle x,y,z\rangle$で
\begin{equation}\begin{aligned}
 [3D-t^2T-9\eta(x+y+z),x]={}&
 -3r_1y^2+t^2r_1zx-t^2xr_1z-3yr_1y-t^2zr_1x\\
 &+t^2xzr_1-3y^2r_1+2t^2r_2x^2-3tr_2yz\\
 &-t^2xr_2x-t^2x^2r_2+3tyzr_2
 +3r_3z^2+3zr_3z+3z^2r_3\\
 &+9\eta r_1-9\eta r_3
\end{aligned}\end{equation}
が成立する。語ごとの展開と巡回置換により三交換子を処理できる。
$t=3s,\eta=u^2/9$を代入し、$x,y,z$を$\alpha,\beta,\gamma$に置き換える。
中心スカラー$2u^3/9$を加えると表示式を得る。
\end{proof}

\subsection{相対ポテンシャル複体の実際の完全性}
$Q=Q_\epsilon$、$S=Q[u]$とし、$D=D_\epsilon$を$S$-代数とみる。
自由$S$-代数のcyclic potentialを
\begin{equation}
 W=xyz-yxz-h\operatorname{sym}(f_\epsilon+ux^2),
 \qquad r_i=\partial^{\mathrm{cyc}}_{x_i}W,\quad(x_1,x_2,x_3)=(x,y,z)
\end{equation}
とする。定義関係式は$r_1=r_2=r_3=0$である。

$\partial^{\mathrm{nc}}_j$は語の各$x_j$を取り除き、その左右の部分を分ける演算とする：
\begin{equation}
 \partial^{\mathrm{nc}}_j(a_1\cdots a_d)
 =\sum_{a_\ell=x_j} a_1\cdots a_{\ell-1}\otimes a_{\ell+1}\cdots a_d.
\end{equation}
係数環$S$は微分しない。
\begin{proposition}[底変更と両立する第一ブロックの相対対角]
\label{nq:prop:relative-diagonal}
$D^e_S=D\otimes_S D^{\mathrm{op}}$とする。有限自由双加群の複体
\begin{equation}
 0\longrightarrow D^e_S
 \xrightarrow{d_3}(D^e_S)^3
 \xrightarrow{d_2}(D^e_S)^3
 \xrightarrow{d_1}D^e_S
 \xrightarrow{m}D\longrightarrow0
 \label{nq:eq:potential-resolution}
\end{equation}
を、基底$e_i$（生成元）、$\varepsilon_i$（関係式）について
\begin{equation}\begin{aligned}
 d_1(e_i)&=x_i\otimes1-1\otimes x_i,\\
 d_2(\varepsilon_i)&=\sum_j\sum_w\ \sum_{w=p x_j q} c_{iw}p e_j q,\\
 d_3(1)&=\sum_i(x_i\varepsilon_i-\varepsilon_i x_i)
\end{aligned}\end{equation}
で定める。ここで$r_i=\sum_w c_{iw}w$と展開している。
この複体は完全で、任意の可換$S$-代数への底変更後も完全である。
\end{proposition}
\begin{proof}
$d_1d_2=0$は$r_i=0$から、$d_2d_3=0$はcyclic potentialの恒等式
$\sum_i[x_i,r_i]=0$を非可換微分することから従う。
active次数に関して、複体の自由基底のfiltration次数を$0,1,2,3$と取る。
$ux^2$はactive次数二なので、随伴次数付き複体は$u=0$の同次ポテンシャルの
Koszul双加群分解を$Q[u]$へ拡大したものである。

この同次複体が$Q$上で完全であることを確認する。
補題\ref{nq:lem:pbw}の$u=0$で、各次数成分は有限自由$Q$-加群である。
各複素点に特殊化すると、Mori--Smithの三次元Calabi--Yau代数に対する
同次ポテンシャル分解となり完全である
\cite[Theorem 4.2]{PMoriSmith}\cite[\S3.2]{PBergerTaillefer}。
各次数における複体は有限個の有限自由$Q$-加群からなり、
全閉点のfiberで完全なので$Q$上で完全である。
さらに最後の項も自由であるから、各次数の完全列は$Q$-加群として分裂する。
従って任意の底変更後にも完全性が保たれる。

filtrationは下に有界で、各元のactive次数は有限である。
随伴複体でcycleの最高次数を消し、次数を下げる帰納法により
元の複体の完全性を得る。底変更後も同じmonic PBW表示と同じ随伴分解を持つため、
同じ証明が適用できる。これは既知のPBWポテンシャル分解
\cite[Theorem 3.6の証明]{PBergerTaillefer}を、今回の係数環と底変更について明示したものである。
\end{proof}

\begin{remark}[中心方向の通常微分は使えない]\label{nq:rem:no-derivation}
$h\ne0$のファイバーでは、$\partial_u(u)=1$を満たす次数$-1$の
代数導分$\partial_u:D_\epsilon\to D_\epsilon$は存在しない。
実際、次数から$\partial_u(x),\partial_u(y),\partial_u(z)$はスカラーとなる。
尖点型の$[y,z]-h(z^2+2ux)=0$を微分して$u=0$に落とすと、
$-2h(x+\partial_u(z)z)=0$という一次の関係が必要になるが、PBW性に反する。
節点型でも$-h(2x+\partial_u(y)z+\partial_u(z)y)=0$となって同様である。
相対対角\eqref{nq:eq:potential-resolution}は$u$を係数としており、
この不存在とは矛盾しない。中心三次商の残余圏を調べる際に、
可換多項式の$\partial_u$をそのまま作用させることはできない。
\end{remark}

\subsection{七生成元への組立てと中心三次商}
まず$E_s$を$\CC[s]$上で考える。$u=0$のFermat型三次元Jacobian代数は
全ての$s$で三次元Calabi--Yau AS-regularである
\cite[Theorem 4.2(1)]{PMoriSmith}。
$u$を加えた関係式を
$g_j=f_j+u l_j+\alpha_j u^2$と書くと、
\cite[Theorem 3.1.3]{LBSVNE12}の記法では
$Q=I$、$(x_1,x_2,x_3)=(a,b,c)$、$l_j=-2s x_j$、
$\alpha_j=0$である。
従って$\sum_j(x_jl_j-l_jx_j)=0$であり、
同定理の三つの条件を全て$\gamma_j=0$で満たす。
従って各$s$で$E_s$は中心正則$u$を持つ四次元Koszul AS-regular domainである。
各次数の$\CC[s]$-加群は有限生成で全fiberの次元が一定なので自由となる。
$u$による完全列を次数ごとに分裂させると、$E_s$は$\CC[s,u]$上自由で、
active filtrationの随伴次数付き環は同次Fermat代数の係数拡大となる。
そのポテンシャル複体も命題\ref{nq:prop:relative-diagonal}と同じ議論で、
任意の底変更と両立する長さ三の相対自由対角分解を与える。

\begin{proof}[定理\ref{nq:thm:family}の証明]
中心三次元を
\begin{equation}
 \Omega_\epsilon=\Gamma_\epsilon(r)+\Gamma_{\mathrm F}(s)+u^3
 \label{nq:eq:total-center}
\end{equation}
と定める。中心性と原点への特殊化は既に示した。
両ブロックが$R_\epsilon[u]$上自由なので、そのテンソル積も自由である。
全体の$u=0$での商は、二つの三次元Koszul AS-regular代数のテンソル積であり、
Hilbert級数$(1-t)^{-6}$、AS型$(6,6)$を持つ。
このテンソル積がnoetherian domainであることには注意が要る。
一般のnoetherian代数のテンソル積についての無条件な主張は用いず、
第一ブロックの上記Ore表示を第二のFermat代数上へ係数拡大する。
これにより二変数量子平面またはJordan平面と、もう一つの可逆Ore拡大の反復として得られる。

$u$は全体でも中心正則元である。
中心正則次数一拡大に対するKoszul性の持上げ
\cite[Theorem 2.6, Remark 2.8]{LBSVNE12}と、
$H_{\mathcal A}(t)=(1-t)^{-7}$から、Koszul双対の級数は$(1+t)^7$となりglobal dimensionは七である。
AS-Gorenstein型は中心正則元に対するchange-of-ringsの
\begin{equation}
 \operatorname{Ext}_{\mathcal A}^{i}(k,\mathcal A)
 \simeq\operatorname{Ext}_{\mathcal A/(u)}^{i-1}
             (k,\mathcal A/(u))(1)
\end{equation}
および左側の同じ式から$(7,7)$となる。
noetherian性と域性は$u$-進次数に関する随伴環
$(\mathcal A/(u))[u]$から従う。正次数の$u$を用いるので、固定した全次数における
最高項の消去は有限回で停止する。

各ファイバーで$\Omega_\epsilon\bmod u$の第一ブロック次数$(3,0)$の項は
$\Gamma_\epsilon(r)|_{u=0}$であり、そのPBW正規形は非零である。
第二ブロック次数$(0,3)$の項とは相殺できないため、$\Omega_\epsilon\ne0$である。
周囲がdomainなので左右正則である。
各次数で$\Omega_\epsilon:\mathcal A_{d-3}\to\mathcal A_d$は
全ての幾何学的fiberで単射だから、局所的に分裂した単射であり、
余核は有限局所自由で任意の底変更と両立する。
中心正則次数三商のchange-of-ringsからAS型$(6,4)$を得る。
半古典極限は関係式を$h$で割って$r=h\lambda,s=h\mu$の原点へ制限すればよい。

最後に、$R_\epsilon[u]$上で二つの長さ三の相対ポテンシャル分解をテンソルしtotalizeする。
その自由階数列は$(1+t)^3(1+t)^3=(1+t)^6$であり、完全性と任意の底変更が保たれる。
$R_\epsilon$上では包絡代数から相対包絡代数への商は中心元
$u\otimes1-1\otimes u$による商であり、この元は正則である。
その長さ一のKoszul分解と相対分解から、周囲の対角は絶対的にもperfect（射影次元高々七）となる。
\end{proof}

\subsection{相対残余圏について残る仕事}
本節は二つの固定Poisson成分の稠密開を同時に扱う量子化族を与える。
一つの方向だけを固定した形式弧ではなく、$(r,s)$の両方向を保持する。
周囲の代数の有限対角分解は実在し、その底変更も証明した。
しかし、中心三次商のaffine coneは可換点でも原点で特異である。
\emph{周囲の対角を中心三次商でテンソルするだけ}では、その商や残余圏の対角にはならない。
次に必要なのは、中心三次商の相対局所cohomology・射影化とOrlov windowを含むproperなdg模型、
およびその模型の対角のperfect性である。
この過程で中心方向の混合を除かず、注意\ref{nq:rem:no-derivation}の障害を扱わなければならない。
特に、可換特殊点を含む二変数底全体での相対残余圏の比較は、
本節では証明していない。
一変数一般点での可換模型は、\cite[一般次元の中心三次非可換変形の可換復元定理]{NCCubicMain}で扱う。

\appendix
% !TeX root = ../../cubic-quantization.tex
\section{Hesse型Poisson量子化の明示的な三次中心元}
\label{sec:explicit-hesse-center}
本節では量子化の定理が与える三次中心元を、滑らかなHesse型
三変数ブロックについて閉公式で指定する。
この公式はSklyanin代数の既知の中心三次式
\cite[Theorem 1およびSection 2]{DeLaetLeBruyn}
（原式は\cite[p.~35, Remark (1)]{SmithTate}）
のPoisson量子化パラメータへの代入である。
従って中心元の存在という既知の結果自体を新規な主張とはしない。
本研究で有用なのは、$0$を含む明示的なパラメータ開集合と
$\hbar=0$での留数を正規化した公式を同時に得ることである。

\subsection{公式と適用開集合}
$t\in\CC$、$t^3\ne-1$とし、
\begin{equation}
 f_t(x,y,z)=\frac{x^3+y^3+z^3}{3}+txyz
\end{equation}
と置く。$\hbar$を次数零のパラメータとして、
\begin{equation}
 [y,z]=\hbar\left(x^2+\frac{t}{2}(yz+zy)\right),
 \quad
 [z,x]=\hbar\left(y^2+\frac{t}{2}(zx+xz)\right),
 \quad
 [x,y]=\hbar\left(z^2+\frac{t}{2}(xy+yx)\right)
 \label{eq:hesse-commutators}
\end{equation}
による次数付き代数を$S_{\hbar,t}$と書く。
ここで
\begin{equation}
 S=x^3+y^3+z^3,\quad
 C=xyz+yzx+zxy,\quad
 D=yxz+xzy+zyx,\quad T=(C+D)/6.
\end{equation}
次の$\hbar$の多項式を定める：
\begin{equation}
 U(\hbar,t)=1+\frac{\hbar^2t^2}{2}
  +\frac{\hbar^4t}{6}-\frac{\hbar^4t^4}{48},
 \qquad
 W(\hbar,t)=
 \frac{(4-\hbar^2t^2)
 (12t+\hbar^2(t^3-8))}{48}.
 \label{eq:hesse-UW}
\end{equation}
また
\begin{equation}
 \Delta_{\hbar,t}=
 \left(t+\frac{\hbar^2}{3}
       \left(1+\frac{t^3}{4}\right)\right)^3
  +\left(1-\frac{\hbar^2t^2}{4}\right)^3.
 \label{eq:hesse-discriminant}
\end{equation}
と置く。$\Delta_{0,t}=1+t^3\ne0$である。

\begin{theorem}[Hesse中心三次式の閉公式]
\label{thm:explicit-hesse-center}
$t^3\ne-1$に対し、
\begin{equation}
 U_t=\{\hbar\in\mathbb A^1_{\CC}:
   U(\hbar,t)(1-\hbar^2t^2/4)\Delta_{\hbar,t}\ne0\}
\end{equation}
とする。$0\in U_t$であり、
各$\hbar\in U_t$に対して
$S_{\hbar,t}$は三生成元のKoszul AS-regular domainで、
Hilbert級数$(1-z)^{-3}$を持つ。
元
\begin{equation}
 \boxed{\quad
 \Omega_{\hbar,t}=
   \frac{U(\hbar,t)}3(x^3+y^3+z^3)
   +W(\hbar,t)\operatorname{sym}(xyz)
 \quad}
 \label{eq:hesse-central-element}
\end{equation}
は$S_{\hbar,t}$の次数三の中心正則元であり、
$\Omega_{0,t}=f_t$である。
特に$t=0$では
$\Delta_{\hbar,0}=1+\hbar^6/27$で、
\begin{equation}
 \Omega_{\hbar,0}
 =\frac{x^3+y^3+z^3}{3}
   -\frac{2\hbar^2}{3}\operatorname{sym}(xyz)
\end{equation}
を得る。
\end{theorem}
\begin{proof}
$q=\hbar t/2$、
$a=1-q$、$b=-(1+q)$、$c=-\hbar$と置く。
\eqref{eq:hesse-commutators}は
\begin{equation}
 ayz+bzy+cx^2=0,\quad
 azx+bxz+cy^2=0,\quad
 axy+byx+cz^2=0
\end{equation}
となる。$\hbar\ne0$かつ
$(1-\hbar^2t^2/4)\Delta_{\hbar,t}\ne0$なら、
$abc\ne0$で、標準Sklyaninの
\begin{equation}
 (a^3+b^3+c^3)^3\ne(3abc)^3
\end{equation}
が成り立つ。
実際、左辺の差は
$-27\hbar^3\Delta_{\hbar,t}$である。
従ってこの開部分では点曲線が滑らかな楕円曲線であり、
$S_{\hbar,t}$は三次元Sklyanin AS-regular domainである
\cite{SmithTate,DeLaetLeBruyn}。
$\hbar=0$の代数は可換多項式環である。

既知の中心元の公式は
\begin{equation}
 G=a(b^3-c^3)C+
   b(c^3-a^3)D+
   c(a^3-b^3)S
 \label{eq:hesse-sklyanin-G}
\end{equation}
である\cite[Section 2]{DeLaetLeBruyn}。
一方、三つの定義関係式にそれぞれ$x,y,z$を左から掛けて足すと
\begin{equation}
 C-D=\hbar S+\frac{\hbar t}{2}(C+D)
 \label{eq:hesse-CD}
\end{equation}
を得る。この等式で$G$を$S,C+D$で表し、
$a,b,c$を代入して係数を整理すると
\begin{equation}
 G=-9\hbar\left(
       \frac{U(\hbar,t)}3S+
       \frac{W(\hbar,t)}6(C+D)\right)
   =-9\hbar\Omega_{\hbar,t}.
 \label{eq:hesse-G-identity}
\end{equation}
これは$\hbar$で割る操作をせずに成立する多項式恒等式である。
$\hbar\ne0$の滑らかな楕円Sklyanin開部分では
$G$が中心なので$\Omega_{\hbar,t}$も中心である。
各次数成分が$\hbar$について有限自由であることは
三次元SklyaninのHilbert級数と零ファイバーにおける
同じ級数、および元論文のPBW構成から従う。
従って各交換子$[\Omega_{\hbar,t},x]$等は
有限自由加群の切断として稠密開集合上で零であり、
$U_t$の全体で零となる。
$\Omega_{0,t}=f_t$は$T_0=xyz$と
$U(0,t)=1$、$W(0,t)=t$から分かる。
既知のSklyanin中心元$G$は楕円開部分で非零であり、
代数はdomainだから正則である。
零ファイバーでも$f_t$は非零因子である。
\end{proof}


\begin{corollary}[任意の滑らかな三変数三次式]
\label{cor:arbitrary-smooth-ternary-cubic}
$g(x_1,x_2,x_3)$が滑らかな平面三次曲線を定める
斉次三次式であるとする。$g$のJacobian Poisson括弧
\begin{equation}
 [x_2,x_3]=\hbar\operatorname{sym}(\partial_{x_1}g),
 \quad [x_3,x_1]=\hbar\operatorname{sym}(\partial_{x_2}g),
 \quad [x_1,x_2]=\hbar\operatorname{sym}(\partial_{x_3}g)
\end{equation}
から得られる量子化は、$0$を含むZariski開集合上で
三次中心元の閉公式を持つ。
実際、Hesse標準形\cite[Lemma 1]{ArtebaniDolgachev}により
ある$M\in\mathrm{GL}_3(\CC)$、$\lambda\in\CC^\times$、
$t^3\ne-1$が存在して
\begin{equation}
 g(x)=\lambda f_t(Mx)
\end{equation}
となる。$y=Mx$、$\hbar'=\lambda\det(M)\hbar$とおけば、
\begin{equation}
 \boxed{\ \Omega^g_\hbar(x)=
  \lambda\Omega_{\hbar',t}(y)\big|_{y=Mx}\ }
 \label{eq:arbitrary-ternary-central}
\end{equation}
は三次中心元で、$\Omega^g_0=g$である。
\end{corollary}
\begin{proof}
Poissonの定数三重ベクトル
$\partial_{x_1}\wedge\partial_{x_2}\wedge\partial_{x_3}$
は座標変換$y=Mx$の下で
$\det(M)\,
 \partial_{y_1}\wedge\partial_{y_2}\wedge\partial_{y_3}$
に移る。従って$y$の交換子は
\begin{equation}
 [y_2,y_3]
 =\lambda\det(M)\hbar\,
   \operatorname{sym}(\partial_{y_1}f_t),
\end{equation}
およびその巡回置換となる。
対称化は線形座標変換と可換なので、量子化代数は
$S_{\hbar',t}$への次数付き同型を持つ。
その同型によって定理\ref{thm:explicit-hesse-center}の
中心元を戻し、留数が$g$になるよう$\lambda$倍する。
\end{proof}

\begin{corollary}[中心変数と独立ブロックの閉公式]
\label{cor:hesse-block-center}
相異なる三変数ブロック$x_s=(x_s,y_s,z_s)$と
中心変数$u$について
\begin{equation}
 F=\sum_{s=1}^k
 \left(\frac{x_s^3+y_s^3+z_s^3}{3}
       +t_sx_sy_sz_s\right)+P(u),
 \qquad t_s^3\ne-1,
\end{equation}
$P$は任意の斉次三次式とする。
各ブロックに\eqref{eq:hesse-commutators}を課し、
異なるブロック間と$u$を可換に取る。
各$U_{t_s}$の共通開部分で
\begin{equation}
 \Omega_\hbar
   =\sum_s\Omega_{\hbar,t_s}+P(u)
 \label{eq:hesse-multiblock}
\end{equation}
は次数三の中心元で、$\Omega_0=F$である。
各$u$の一次形式$\ell_{si}(u)$を選んで
$x_{si}$を$x_{si}+\ell_{si}(u)$に置き換えた
非自明な混合項を含むPoissonブロックでも、
同じ置換による明示公式が成立する。
\end{corollary}
\begin{proof}
各$S_{\hbar,t_s}$の中心元は他のブロックと可換であり、
$P(u)$も中心なので和は中心である。
中心線形シフトは生成元の次数付き同型を定め、
同型によって中心元が中心元へ移る。
\end{proof}


\subsection{可換復元の候補となる点曲線}
\begin{proposition}[楕円点曲線のHesseパラメータ]
\label{prop:hesse-point-cubic}
定理\ref{thm:explicit-hesse-center}の条件で$\hbar\ne0$とする。
$S_{\hbar,t}$のpoint schemeを定める滑らかな平面三次曲線は
\begin{equation}
 E_{\hbar,t}:\quad
 x^3+y^3+z^3+3\,\vartheta(\hbar,t)xyz=0,
\quad
 \boxed{\quad
 \vartheta(\hbar,t)=
 \frac{t+\frac{\hbar^2}{3}(1+t^3/4)}
 {1-\hbar^2t^2/4}
 \quad}
 \label{eq:hesse-point-parameter}
\end{equation}
である。$\vartheta(0,t)=t$と形式的に延長され、
\begin{equation}
 \Delta_{\hbar,t}
 =\left(1-\frac{\hbar^2t^2}{4}\right)^3
  \left(1+\vartheta(\hbar,t)^3\right).
 \label{eq:hesse-point-delta}
\end{equation}
特に$\hbar\in U_t$なら、右辺のHesse三次曲線は滑らかである。
\end{proposition}
\begin{proof}
滑らかな楕円型三次元Sklyanin代数のpoint schemeは
\begin{equation}
 abc(x^3+y^3+z^3)-(a^3+b^3+c^3)xyz=0
\end{equation}
である\cite{SmithTate,DeLaetLeBruyn}。
定理の証明で与えた
$a=1-\hbar t/2$、$b=-(1+\hbar t/2)$、$c=-\hbar$
について
\begin{equation}
 abc=\hbar(1-\hbar^2t^2/4),\qquad
 a^3+b^3+c^3
 =-3\hbar t-\hbar^3(1+t^3/4)
\end{equation}
を代入し、共通係数$\hbar$で割ると表示を得る。
最後の恒等式は\eqref{eq:hesse-discriminant}の定義から従う。
\end{proof}
\begin{remark}[幾何学的比較との関係]
この楕円点曲線と中心三次元から得られる非可換cubicの
rotation付き可換復元は\cite[一般次元の中心三次非可換変形の可換復元定理]{NCCubicMain}で扱う。
具体的な一楕円ブロックの比較については
同論文の「一ブロックの明示的可換復元」を参照されたい。
ここでは中心三次量子化の存在・公式と圏同値を区別する。
\end{remark}

\begin{remark}[適用範囲と新規性]
\label{rem:hesse-scope}
一般の$f_s(x_s,u)$には$u$とactive変数の本質的な
混合項があり、必ずしも上の中心シフト形にはならない。
三変数制限が滑らかな場合、その全混合項を含む中心元は
系\ref{cor:pg-smooth-ternary-explicit}の有限多項式公式で得られる。
三変数制限が特異または零でも、全変数の三次超曲面が滑らかなら、
定理\ref{thm:pg-universal-determinant}が
原点で正則な有理式のCramer中心元を与える。
いずれの構成も形式的中心元の存在を入力としない。
また、\eqref{eq:hesse-central-element}は
周囲の量子化と中心元の明示式であり、
残余dg圏の相対saturationや
rotation付き可換復元を結論するものではない。
中心元の基本公式は古典的Sklyanin理論に由来し、
本節のHesseパラメータへの翻訳の新規性は未確認である。
\end{remark}

% !TeX root = ../../cubic-quantization.tex
\section{中心三次量子化からのambientの内在的復元}
\label{sec:pg-ambient-recovery}
本節の二次包絡の主張は三次超曲面に特有でなく、
任意次元・任意の関係次数$d\ge3$について成立する。
本文と同じく、周囲代数を$A$、その中心正則商を$B$と書く。
本節では、系\ref{cor:pg-equivariant-character-descent}の量子化族を含め、
中心商の同型類の判定に用いる一般原理を証明する。
この復元原理だけから、固定した非零変形パラメータでの
非可換代数の同型類の像の次元は結論しない。

\begin{theorem}[任意次元のquadratic envelope]\label{thm:pg-ambient-recovery-general}
$k$を体、$V$を有限次元$k$-ベクトル空間とし、
次数一生成の二次連結次数付き代数
\begin{equation}\label{eq:pg-ambient-B}
 A=T(V)/(R_2)
\end{equation}
を考える。$G\in A_d$を非零斉次元、$d\ge3$とし、
$G$の生成する両側イデアルが$(G)=AG=GA$であると仮定する。
$B=A/(G)$と置くと、
\begin{equation}\label{eq:pg-ambient-envelope}
 A_{\mathrm{can}}(B):=
 \frac{T(B_1)}{\left\langle
 \ker(B_1\otimes B_1\longrightarrow B_2)\right\rangle}
 \simeq A
\end{equation}
であり、また
\begin{equation}\label{eq:pg-ambient-line}
 \ker\bigl((A_{\mathrm{can}}(B))_d
 \longrightarrow B_d\bigr)=kG
\end{equation}
となる。この復元は次数付き代数同型に関して関手的である。
$G$の正則性、中心性、$A$のAS正則性は
二つの復元式自体には必要ない。
\end{theorem}
\begin{proof}
$(G)$は$d$次以上にしか成分を持たず、特に
\begin{equation}\label{eq:pg-ambient-low}
 B_1=A_1=V,\qquad B_2=A_2
\end{equation}
である。よって$B_1\otimes B_1\to B_2$の核は$R_2$に等しい。
$A$は二次関係のみで表示されるので最初の同型が従う。
$(G)_d=A_0G=GA_0=kG$より第二式も従う。
両構成は積と商写像のみによるから同型に関して自然である。
\end{proof}

\begin{theorem}[偏極付き圏と$\mathbb Z$-代数からの復元]
\label{thm:pg-ambient-recovery-polarized}
上の定理の仮定に加え、$B$は左右noetherianな
連結次数付きAS Gorenstein代数でinjective dimensionが
$2$以上であるとする。このとき
\begin{equation}\label{eq:pg-ambient-section}
 B\xrightarrow{\sim}
 \bigoplus_{n\ge0}
 \operatorname{Hom}_{\operatorname{qgr}B}
   (\mathcal O_B,\mathcal O_B(n))
\end{equation}
は積を保つ同型である。
従って偏極付き三つ組
\begin{equation}\label{eq:pg-ambient-triple}
 (\operatorname{qgr}B,\mathcal O_B,(1))
\end{equation}
から$A$と直線$kG$を復元できる。
具体的には$E_i=\mathcal O_B(-i)$として
\begin{equation}\label{eq:pg-ambient-Zalg}
 \mathcal Z_{ij}=
 \operatorname{Hom}_{\operatorname{qgr}B}(E_j,E_i)
 \qquad (i\le j)
\end{equation}
の次数一合成の二次核のみを課した
$\mathbb Z$-代数が$A$の付随$\mathbb Z$-代数となる。
さらに$A$が左右noetherianなら
$\operatorname{qgr}B\to\operatorname{qgr}A$は
完全忠実なAbel関手である。
\end{theorem}
\begin{proof}
AS Gorenstein性とinjective dimensionの仮定により
\begin{equation}\label{eq:pg-ambient-extvan}
 \operatorname{Ext}_B^i(k,B)=0
 \qquad (i=0,1)
\end{equation}
となる。有限次元次数付き加群$L$は$k$のシフトによる
有限filtrationを持つので、
$\operatorname{Hom}_{\operatorname{gr}B}(L,B(n))$と
$\operatorname{Ext}^1_{\operatorname{gr}B}(L,B(n))$は消滅する。
$0\to B_{\ge r}\to B\to B/B_{\ge r}\to0$にこれを適用し、
商圏Homの直極限表示を用いれば
式\eqref{eq:pg-ambient-section}を得る。
積の対応は次数シフトを介した合成の定義による。
このsection algebraの原理はArtin--Zhangの理論
\cite{PAZAmbient}の本仮定下での直接的帰結である。

従って$\mathcal Z_{ij}=B_{j-i}$であり、
各$i$での長さ二の合成の核は
$\ker(B_1\otimes B_1\to B_2)$と一致する。
定理\ref{thm:pg-ambient-recovery-general}により
二次包絡は$A$の付随$\mathbb Z$-代数となる。
最後に$B$-加群を$A$-加群とみなす制限関手は完全忠実で、
有限次元torsionの条件も一致するため商Abel圏に降りる。
この最後の完全忠実性は導来圏での完全忠実性を意味しない。
\end{proof}

\begin{corollary}[一般次元の中心三次量子化への適用]
\label{cor:pg-ambient-recovery-quantization}
第\ref{sec:pg-all-dimensions}節の量子化で得られる、
$N$生成元の二次Koszul AS正則ambient $A_\eta$と
中心正則三次元$\Omega_\eta$の商
$B_\eta=A_\eta/(\Omega_\eta)$について、
定理\ref{thm:pg-ambient-recovery-general}が成立する。
さらに各ファイバー$B_\eta$が上のnoetherian
AS Gorenstein仮定を満たす範囲では、
その偏極付き三つ組からambientと中心三次直線が復元される。
この結論は各ファイバーの内在的同型類についての主張であり、
全Poisson moduli stackからnc cubic moduli stackへの
大域的量子化射を構成したことを意味しない。
\end{corollary}
\begin{proof}
量子化定理のambientは二次で、
中心三次元は正則であるから最初の定理が適用できる。
第\ref{sec:pg-all-dimensions}節で示されたAS Gorenstein性と、
所定のnoetherian条件から第二の定理を適用する。
\end{proof}

\subsection{Serre関手による復元と一般次元での限界}
\label{sec:pg-ambient-serre}
構造対象だけを指定して次数シフトを忘れる復元は、
上記の一般定理より強い問題である。
例えば$N$生成元の三次中心商について
滑らかさとCalabi--Yau型Serre双対性の仮定の下では、
\begin{equation}\label{eq:pg-ambient-serre-r}
 \mathsf S\simeq(3-N)[N-2],\qquad
 \mathsf S^{-1}[N-2]\simeq(N-3)
\end{equation}
となる。この式だけで内在的に得られるのは
$(N-3)$-Veroneseであり、次数一シフトそのものとは限らない。
一般$N$についてVeroneseから元の代数を復元できるとは
ここでは主張しない。またNakayama自己同型が非自明なら
Serre関手にその捻りが加わるため、式
\eqref{eq:pg-ambient-serre-r}の使用には追加仮定が必要である。

\begin{proposition}[二変数中心による偶数Veronese復元]
\label{prop:pg-ambient-even}
$k=\mathbb C$とする。$B$を連結次数付き代数とし、
$S=k[u,v]\subset Z(B)$は次数一の中心多項式環、
$B$は$S$上自由、$Z(B)_2=S_2$と仮定する。
このとき偶数Veronese $R=B^{(2)}$の次数付き同型類から
$B$の次数付き同型類を復元できる。
$S$の平方根に関する二つの持上げは
奇数次数への符号作用で異なる。
\end{proposition}
\begin{proof}
$z\in Z(R)$とし、$a\in B$を奇数次数の元とする。
$ua$は偶数次数であるため、
\begin{equation}\label{eq:pg-ambient-even-center-comm}
 u[z,a]=[z,ua]=0
\end{equation}
である。$B$は$S$上自由なので$u$は正則であり、
$z$は全ての$a$と可換である。従って
\begin{equation}\label{eq:pg-ambient-even-center}
 Z(R)=Z(B)^{(2)},\qquad
 Z(R)_1=S_2
\end{equation}
を得る。$Z(R)_1$の生成する中心部分代数は
\begin{equation}\label{eq:pg-ambient-conic}
 D=k[u^2,uv,v^2]
 \simeq k[a,b,c]/(ac-b^2)
\end{equation}
である。具体的には$D_1$の基底$a,b,c$を選び、唯一の二次関係を
$ac-b^2=0$とする。その二次Veronese表示
$a=u^2$、$b=uv$、$c=v^2$により$S=k[u,v]$を取る。
この選択は復元した代数の同型類に影響しない。
実際、$D$の次数付き自己同型は滑らかなconic
$\operatorname{Proj}D\simeq\mathbb P^1$の自己同型を誘導し、
$\mathcal O_{\mathbb P^1}(2)$を保つ。
従って$k=\mathbb C$上では$\mathrm{GL}_2$の元へ持ち上がり、
持上げの核は共通符号$\{\pm1\}$である。
以下では$R,D$の次数を二倍し、$S$の元の次数と合わせる。
$D$の$S$への中心的拡大を使い、
\begin{equation}\label{eq:pg-ambient-even-localizations}
 T_u=R[(u^2)^{-1}]
       \otimes_{D[(u^2)^{-1}]}S[u^{-1}],
 \quad
 T_v=R[(v^2)^{-1}]
       \otimes_{D[(v^2)^{-1}]}S[v^{-1}]
\end{equation}
と置く。$S[u^{-1}]=D[(u^2)^{-1}]\oplus uD[(u^2)^{-1}]$であり、
$B[u^{-1}]$の奇数部分は$u$とその偶数部分の積に等しい。
従って偶奇の次数を分けた積写像により
$T_u\simeq B[u^{-1}]$、$T_v\simeq B[v^{-1}]$となる。
$B$は$S$上自由であり、$k[u,v]$の高さ二の原点を
除いた正則関数は多項式なので、
\begin{equation}\label{eq:pg-ambient-even-intersection}
 B[u^{-1}]\cap B[v^{-1}]=B
\end{equation}
が共通局所化内で成り立つ。従って$B$を復元する。
二つの平方根の差は$u,v\mapsto-u,-v$であり、
次数の偶奇に関する符号に一致する。
\end{proof}


\begin{lemma}[五生成元Jacobian族の次数二中心]
\label{lem:pg-ambient-center2-Jacobian}
$F\in\mathbb C[x,y,z,u,v]_3$とし、
$u,v$を中心変数として
\begin{equation}\label{eq:pg-ambient-Jacobian-rel}
 [x,y]=\operatorname{sym}(F_z),\quad
 [y,z]=\operatorname{sym}(F_x),\quad
 [z,x]=\operatorname{sym}(F_y)
\end{equation}
で定める二次代数を$A_F$とする。
$f_3=F(x,y,z,0,0)$は滑らかな平面三次を定め、
$x,y,z$-次数のPBW濾過について
\begin{equation}\label{eq:pg-ambient-Jacobian-gr}
 \operatorname{gr}A_F\simeq S_{f_3}\otimes\mathbb C[u,v],
 \qquad Z(S_{f_3})_1=Z(S_{f_3})_2=0
\end{equation}
と仮定する。さらに$G_F\in Z(A_F)_3$は
最高$x,y,z$-次数記号が非零な中心正則元とする。
$B_F=A_F/(G_F)$について
\begin{equation}\label{eq:pg-ambient-Jacobian-center2}
 Z(A_F)_2=Z(B_F)_2=\mathbb C[u,v]_2
\end{equation}
が成立する。
\end{lemma}
\begin{proof}
PBW濾過で中心元の最高記号を取れば、
$S_{f_3}$の一次、二次中心の消滅から
$Z(A_F)_2=\mathbb C[u,v]_2$を得る。
$B_F$の次数二中心元を$Q\in(A_F)_2$へ持ち上げる。
中心性と三次の商核が$\mathbb CG_F$であることから
\begin{equation}\label{eq:pg-ambient-Jacobian-Qcomm}
 [Q,x]=c_xG_F,\quad[Q,y]=c_yG_F,\quad
 [Q,z]=c_zG_F
\end{equation}
と書ける。$G_F$は中心的なのでJacobi恒等式は
$[Q,[x,y]]=[Q,[y,z]]=[Q,[z,x]]=0$を与える。
二次関係へ代入し、中心性と正則性によって
\begin{equation}\label{eq:pg-ambient-Jacobian-Hessian}
 \sum_{i=x,y,z}c_i\,\partial_i\partial_jF=0,
 \qquad j=x,y,z
\end{equation}
を得る。$u=v=0$で制限すれば
$\partial_j(\sum_i c_i\partial_i f_3)=0$となる。
括弧内は斉次二次式ゆえ零である。
非零の定数ベクトル$c$についてこの等式が成立すれば、
$f_3$は一次方向に依存しない錐であり、
滑らかな平面三次であることに反する。
従って$c_x=c_y=c_z=0$。
$Q$は$A_F$の中心元であるから
$Q\in\mathbb C[u,v]_2$となる。
\end{proof}

\begin{corollary}[三次三様体の構造対象付き導来復元]
\label{cor:pg-ambient-derived-threefold}
$k=\mathbb C$とし、$A$を左右noetherianな五次元Koszul
AS正則Calabi--Yau代数とする。
$G\in Z(A)_3$を正則中心元、$B=A/(G)$とする。
$\operatorname{qgr}B$は有限大域次元を持ち、そのSerre関手が
$\mathsf S\simeq(-2)[3]$であると仮定する。
さらに$S=k[u,v]\subset Z(B)$が次数一中心多項式環で、
$B$は$S$上自由、$Z(B)_2=S_2$であるとする。
すると
\begin{equation}\label{eq:pg-ambient-derived-input}
 (D^b(\operatorname{qgr}B),\mathcal O_B)
\end{equation}
だけから$B$の次数付き同型類、ambient $A$、
中心三次直線$kG$を復元できる。
次数シフト$(1)$は入力として指定しなくてよい。
\end{corollary}
\begin{proof}
Serre関手は内在的なので、
$\mathsf S^{-1}[3]\simeq(2)$から
\begin{equation}\label{eq:pg-ambient-derived-Veronese}
 \bigoplus_{n\ge0}
 \operatorname{Hom}(\mathcal O_B,\mathsf S^{-n}[3n]\mathcal O_B)
 \simeq B^{(2)}
\end{equation}
が得られる。ここでは、$B$が左右noetherianなAS Gorenstein代数で
injective dimension $4$を持つことから、
定理\ref{thm:pg-ambient-recovery-polarized}のsection algebra同型を用いた。
命題\ref{prop:pg-ambient-even}で$B$を復元し、
定理\ref{thm:pg-ambient-recovery-general}を適用する。
\end{proof}

\begin{remark}[既存の三次三様体族と未解決部分]
Jacobian型五生成元族については、
純三変数の滑らかな三次式の低次数中心の消滅と
Hessianの非退化性を用い、一般点で
$Z(B)_2=k[u,v]_2$を示す補題がある。
$B$の$k[u,v]$上の自由性には、ambientの相対PBW性に加えて
商の最高active記号を確認する必要がある。
例えば$\operatorname{gr}A=S_{f_3}\otimes k[u,v]$で、
$G$の最高active記号が$S_{f_3}$の正則な三次元$g_3$ならば、
\begin{equation}\label{eq:pg-ambient-quotient-symbol}
 \operatorname{gr}(A/(G))\simeq(S_{f_3}/(g_3))\otimes k[u,v]
\end{equation}
となる。右辺の各active次数の$k$-基底を持ち上げ、
最高記号について帰納することで$B$の$k[u,v]$-自由基底を得る。
ambientが自由であるという条件だけでは商の自由性は従わず、
例えば中心元$u^3$による商は$k[u,v]$上torsionを持つ。
ただしこの事実だけで$\operatorname{qgr}B$の
有限大域次元が全閉点で保証されるとはいわない。
これらの追加条件なしに一般次元でSerre復元を主張しない。
また構造対象を忘れた任意の導来同値、
Kuznetsov成分のみからのambient復元は結論しない。
\end{remark}

% !TeX root = cubic-quantization.tex
% 2026-11-21 / Q-v1. Bibliography for cubic-quantization.tex.
\begin{thebibliography}{99}
\small
\bibitem{PoissonCubic}
\emph{滑らかな三次超曲面のPoisson構造：正規形・退化軌跡・モジュライ},
未公表研究稿、2026年11月21日、A-v62。
\href{https://github.com/uedakazushi/poisson-cubic/blob/main/papers/poisson-cubic.tex}{主原稿}。
古典入力はPoisson正規形定理（\texttt{thm:pg-normal}）、
HessianによるJacobi条件の線形化（\texttt{lem:pg-linearize}）、
縮約二形式によるブロック復元（\texttt{ne:blocks}）、
相対Poisson階数層の記述（\texttt{thm:pg-fermat-relative-count}の第三項）、
および非楕円成分の具体例（\texttt{ne:main}）である。

\bibitem{NCCubicMain}
\emph{非可換三次$ n$-foldの可換復元とモジュライ：
Kuznetsov成分・構造対象付きモジュライ・Poisson幾何},
未公表研究稿、2026年11月21日。
\href{https://github.com/uedakazushi/nc-cubic/blob/main/papers/nc-cubic.tex}{主原稿}。
本稿からは一般次元の中心三次非可換変形の可換復元定理
（\texttt{gdr:thm:generic}、完備局所比較の入力は\texttt{cft:thm:complete}）を参照する。
その比較を本稿の代数的量子化の証明入力には用いない。

\bibitem{PArzhantsevMatsushima}
I.~V.~Arzhantsev,
\emph{Invariant Ideals and Matsushima's Criterion},
Communications in Algebra \textbf{36} (2008), no.~12, 4368--4374.
\href{https://doi.org/10.1080/00927870802179347}{doi:10.1080/00927870802179347};
\href{https://arxiv.org/abs/math/0506430v1}{arXiv:math/0506430v1}.
本文で確認した参照箇所は、この公開原稿の\S2、
Theorem~2.1とProposition~2.2, p.~2である。

\bibitem{PEtingofLie}
P.~Etingof,
\emph{Lie Groups and Lie Algebras}, MIT lecture notes, 2024 course version.
\href{https://live.ocw.mit.edu/courses/18-755-lie-groups-and-lie-algebras-ii-spring-2024/mit18_755_s24_lec_full.pdf}{MIT OpenCourseWare 公開原稿}.
本文ではCorollaries~13.7--13.8、Proposition~16.9、\S19.4、Definition~21.13、Lemma~21.14、
Theorem~25.17、\S32.2を参照する。
特にPBW対称化と随伴作用の比較はp.~75、
Proposition~16.9, p.~88、Definition~21.13とLemma~21.14, p.~114、
Theorem~25.17, p.~137。

\bibitem{PMoriSmith}
I.~Mori and S.~P.~Smith,
The classification of 3-Calabi--Yau algebras with 3 generators and 3 quadratic relations,
\emph{Mathematische Zeitschrift} \textbf{287} (2017), 215--241.
\url{https://doi.org/10.1007/s00209-016-1824-5};
\href{https://arxiv.org/abs/1502.07403v2}{arXiv:1502.07403v2}.
定理番号はarXiv v2による：Theorems~1.3, 1.4, 4.2、
Lemma~4.1、Corollaries~1.6, 2.7, 2.9, 4.3。
Theorem~1.4, p.~5はArtin--Tate--Van den Berghの判定の再掲、
Theorem~4.2はpp.~16--17。
点曲線のHessian公式はLemma~7.1, Theorem~7.2, pp.~23--24。

\bibitem{PBergerTaillefer}
R.~Berger and R.~Taillefer,
Poincar\'e--Birkhoff--Witt deformations of Calabi--Yau algebras,
\emph{Journal of Noncommutative Geometry} \textbf{1} (2007), 241--270.
\href{https://arxiv.org/abs/math/0610112}{arXiv:math/0610112}.
\url{https://doi.org/10.4171/JNCG/6}.
刊行版Theorem~1.1(i), p.~242のpotentialによるPBW変形、
Theorems~3.1, 3.6, pp.~246, 251を用いる。

\bibitem{PKontsevich}
M.~Kontsevich,
Deformation quantization of Poisson manifolds,
\emph{Letters in Mathematical Physics} \textbf{66} (2003), 157--216.
\href{https://doi.org/10.1023/B:MATH.0000027508.00421.bf}{doi:10.1023/B:MATH.0000027508.00421.bf};
原稿名は \emph{Deformation quantization of Poisson manifolds, I},
\href{https://arxiv.org/abs/q-alg/9709040}{arXiv:q-alg/9709040} (1997).
参照箇所は原稿のSections~4.4--4.6（Maurer--Cartan変形functorとformality）、
6.4（$L_\infty$関係）、8.2および8.3.3.1（中心のtangent mapとwheel）。

\bibitem{PDolgushev}
V.~Dolgushev,
The Van den Bergh duality and the modular symmetry of a Poisson variety,
\emph{Selecta Mathematica, New Series} \textbf{14} (2009), 199--228.
\url{https://doi.org/10.1007/s00029-008-0062-z};
\href{https://arxiv.org/abs/math/0612288v3}{arXiv:math/0612288v3} (2007).
原稿版Proposition~1, p.~9、およびTheorem~3, pp.~15--16。

\bibitem{EG}
P. Etingof and V. Ginzburg,
\emph{Noncommutative del Pezzo surfaces and Calabi--Yau algebras},
Journal of the European Mathematical Society \textbf{12} (2010), 1371--1416.
刊行版Theorem~3.3.2(ii), p.~1385（低次potential変形の非スカラー中心元）、
その証明\S8.3, pp.~1410--1411、
Proposition~8.1.1(i)とその証明, pp.~1407--1408
（指定した古典剰余を持つ形式中心元、相対係数環と還元群への対応）、
\S9, p.~1414（Rainsによる明示中心公式）を参照。
Sections~3.5--3.6、Theorem~3.6.2も関連する。
\url{https://doi.org/10.4171/JEMS/235}.
\href{https://arxiv.org/abs/0709.3593v4}{arXiv:0709.3593v4};
\href{https://ems.press/content/serial-article-files/31746}{刊行版本文}。

\bibitem{PCMR}
L.~Chekhov, M.~Mazzocco and V.~Rubtsov,
\emph{Quantised Painlev\'e monodromy manifolds, Sklyanin and Calabi--Yau algebras},
Adv. Math. \textbf{376} (2021), article 107442.
\url{https://doi.org/10.1016/j.aim.2020.107442}.
本文確認は\href{https://arxiv.org/abs/1905.02772v2}{arXiv:1905.02772v2}
（accepted version）。式(4.38), p.~21、Lemma~4.6・式(4.39), p.~22、
Table~5, p.~41を参照。

\bibitem{LBSVNE12} L. Le Bruyn, S. P. Smith and M. Van den Bergh,
\emph{Central extensions of three dimensional Artin--Schelter regular algebras},
Math. Z. \textbf{222} (1996), 171--212.
Theorem 2.6, Corollary 2.7, Remark 2.8, Theorem 3.1.3.
\url{https://sites.math.washington.edu/~smith/Research/Central-extns.pdf}.

\bibitem{BergNE12} G. M. Bergman,
\emph{The diamond lemma for ring theory}, Adv. Math. \textbf{29} (1978), 178--218, Theorem 1.2.
\url{https://doi.org/10.1016/0001-8708(78)90010-5}.

\bibitem{SmithTate}
S. P. Smith and J. Tate,
\emph{The center of the 3-dimensional and 4-dimensional Sklyanin algebras},
K-Theory \textbf{8} (1994), 19--63.
Corollary 2.3 (p.~26); Theorem 3.7(c) (p.~38).
\url{https://doi.org/10.1007/BF00962089};
\url{https://sites.math.washington.edu/~smith/Research/tate.pdf}.

\bibitem{DeLaetLeBruyn}
K. De Laet and L. Le Bruyn,
\emph{The geometry of representations of 3-dimensional Sklyanin algebras},
Algebras and Representation Theory \textbf{18} (2015), 761--776.
arXiv:1405.1158v1, Section 2, p.~3 (canonical central cubic formula);
the result is a reformulation of the Artin--Schelter and Smith--Tate formula.
\url{https://arxiv.org/abs/1405.1158}.

\bibitem{ArtebaniDolgachev}
M. Artebani and I. Dolgachev,
\emph{The Hesse pencil of plane cubic curves},
L'Enseignement Math\'ematique (2) \textbf{55} (2009), no.~3--4, 235--273,
Lemma 1 (every smooth plane cubic admits a Hesse canonical form).
\url{https://doi.org/10.4171/LEM/55-3-3};
\url{https://arxiv.org/abs/math/0611590}.

\bibitem{PStacks}
The Stacks Project Authors,
\emph{The Stacks Project},
Criteria for Representability, Theorem~97.17.2
(Tag~\texttt{06FI}, quotient stacks),
More on Algebra, Lemma~15.119.3 (Tag~\texttt{0BCH}, UFDのPicard群),
Algebra, Lemmas~10.97.1--10.97.2 (Tags~\texttt{00MA}, \texttt{00MB}, 有限加群の完備化),
Bootstrap, Lemma~80.11.7 (Tag~\texttt{06PH}, 自由作用の商),
\'Etale Cohomology, Theorem~59.22.4
(Section tag~\texttt{03OY}, 準連接層のコホモロジー比較),
Groupoids in Algebraic Spaces, Lemma~78.27.1
(Tag~\texttt{06PD}, 商のgerbe判定),
Sheaves on Algebraic Stacks, Proposition~96.14.3
(Tag~\texttt{06WT}, 同変準連接層の降下).
\url{https://stacks.math.columbia.edu/tag/06FI};
\url{https://stacks.math.columbia.edu/tag/0BCH};
\url{https://stacks.math.columbia.edu/tag/00MA};
\url{https://stacks.math.columbia.edu/tag/00MB};
\url{https://stacks.math.columbia.edu/tag/06PH};
\url{https://stacks.math.columbia.edu/tag/03OY};
\url{https://stacks.math.columbia.edu/tag/06PD};
\url{https://stacks.math.columbia.edu/tag/06WT}.

\bibitem{PLurieBWB}
J.~Lurie,
\emph{A proof of the Borel--Weil--Bott theorem}, 著者公開ノート。
Theorem~5とその証明, p.~3を参照。PDF内の刊行年は記載されていない。
\url{https://www.math.harvard.edu/~lurie/papers/bwb.pdf}.

\bibitem{PGuoLuSharpe}
J.~Guo, Z.~Lu and E.~Sharpe,
\emph{Classical sheaf cohomology rings on Grassmannians},
\href{https://arxiv.org/abs/1605.01410v2}{arXiv:1605.01410v2} (2017).
\S2.2, Theorem~2.3, pp.~5--6のGrassmann多様体上のBorel--Weil--Bott規約、
およびTheorem~2.4の証明における負の重みの扱いを参照。

\bibitem{PDerksenInvariants}
H.~Derksen,
\emph{Computation of Invariants for Reductive Groups},
Advances in Mathematics \textbf{141} (1999), 366--384.
Theorem~2.1, pp.~368--369のHilbert有限生成定理と、
\S2, p.~369のReynolds作用素および不変部分加群に関する性質を参照。
\url{https://sites.lsa.umich.edu/hderksen/wp-content/uploads/sites/614/2018/05/A.I.a.4.pdf}.

\bibitem{PMilneGroups}
J.~S.~Milne,
\emph{Algebraic Groups: The Theory of Group Schemes of Finite Type over a Field},
Cambridge University Press, 著者公開2022年版。
Propositions~7.12, 7.17, pp.~142--143（軌道と安定化群による商）を参照。
\url{https://www.jmilne.org/math/Books/iAG2022.pdf}.

\bibitem{PBerghSchnurerGerbes}
D.~Bergh and O.~M.~Schn\"urer,
\emph{Decompositions of derived categories of gerbes and of families of Brauer--Severi varieties},
Documenta Mathematica \textbf{26} (2021), 1465--1500.
Theorem~3.11（gerbeの類）、Proposition~4.10, p.~1487（重み零の準連接層の降下）、
Corollary~4.11, p.~1488（指標による類の変更と重み）を参照。
\url{https://doi.org/10.4171/DM/846};
\url{https://ems.press/content/serial-article-files/26626?nt=1}.

\bibitem{PLieblichPeriodIndex}
M.~Lieblich,
\emph{Twisted sheaves and the period-index problem},
Compositio Mathematica \textbf{144} (2008), no.~1, 1--31.
Lemma~3.1.1.8, p.~7（重み$n$の可逆層とgerbe類の$n$倍の消滅）を参照。
同節の底は代数空間であり、本稿では代数スタックの場合のcocycle証明も記す。
\url{https://doi.org/10.1112/S0010437X07003144}.

\bibitem{PSmithVeronese}
S.~P.~Smith,
\emph{Maps between non-commutative spaces},
Transactions of the American Mathematical Society \textbf{356} (2004), no.~7, 2927--2944.
Theorem~4.1, p.~2934、Proposition~4.8, pp.~2940--2941と、
p.~2941の次数一生成の場合の段落を参照。
\url{https://doi.org/10.1090/S0002-9947-03-03411-1};
\url{https://sites.math.washington.edu/~smith/Research/Maps.pdf}.
訂正文：同著者、\emph{Corrigendum to ``Maps between non-commutative spaces''},
同誌\textbf{368} (2016), no.~11, 8295--8302,
\url{https://doi.org/10.1090/tran/6908}.
訂正文の実読箇所はarXiv:1507.02497v2、Lemma~1.1, pp.~1--4である。
\url{https://arxiv.org/abs/1507.02497v2}.

\bibitem{PVdBNoncommP1Bundles}
M.~Van den Bergh,
\emph{Non-commutative $\mathbb P^1$-bundles over commutative schemes},
Transactions of the American Mathematical Society \textbf{364} (2012), no.~12, 6279--6313.
\url{https://doi.org/10.1090/S0002-9947-2012-05469-9}.
実読箇所はarXiv:math/0102005v3、\S4.2, p.~9のLemma~4.2.3直後の段落であり、
可逆sheaf-bimoduleによるsheaf $\mathbb Z$-algebraの捩りと次数付き加群の同値を用いる。
\url{https://arxiv.org/abs/math/0102005v3}.

\bibitem{PAZAmbient}
M.~Artin and J.~J.~Zhang,
\emph{Noncommutative projective schemes},
Adv. Math. \textbf{109} (1994), 228--287,
\href{https://doi.org/10.1006/aima.1994.1087}{doi:10.1006/aima.1994.1087}.
\end{thebibliography}

\end{document}
